% AA Classe Del Documento
\documentclass[11pt, a4paper]{article}

% AA Tipografia E Lingua
% BB Lingua E Codifica Font
\usepackage[italian]{babel}
\usepackage[T1]{fontenc}
\usepackage{lmodern}

% BB Geometria Pagina E Microtipografia
\usepackage[margin=2cm, top=1.5cm, bottom=1.5cm]{geometry}
\usepackage{parskip}
\usepackage{microtype}

% AA Gestione Dei Colori
% BB Pacchetto Colori
\usepackage[table]{xcolor}

% BB Palette Colori Personalizzata
\definecolor{darkgreen}{RGB}{0,100,0}
\definecolor{theoryblue}{RGB}{0,84,147}
\definecolor{backcolour}{RGB}{245,245,245}
\definecolor{codegray}{rgb}{0.5,0.5,0.5}

\definecolor{codebg}{HTML}{F8F9FA}
\definecolor{codeframe}{HTML}{E2E8F0}
\definecolor{codekeywords}{HTML}{7C3AED}
\definecolor{codedocker}{HTML}{0284C7}
\definecolor{codecomments}{HTML}{64748B}
\definecolor{codestrings}{HTML}{0D9488}
\definecolor{codenumbers}{HTML}{94A3B8}
\definecolor{codetext}{HTML}{1E293B}

% AA Matematica E Notazione Scientifica
\usepackage{amsmath, amssymb, amsfonts}
\usepackage{mathtools}
\usepackage{siunitx}

% BB OTTIMIZZAZIONE MATRICI LUNGHE PER LA STAMPA
\setlength{\arraycolsep}{10pt}
\usepackage{bm}

% AA Tabelle Ed Allineamenti
\usepackage{booktabs}
\usepackage{array}
\usepackage{multirow}
\usepackage{tabularx}
\usepackage{arydshln}

% AA Grafica E Disegno Vettoriale
% BB Gestione Figure E Posizionamento
\usepackage{graphicx}
\graphicspath{{./Images/CloudComputing/}}
\usepackage{subcaption}
\usepackage{float}

% BB TikZ E Pgfplots
\usepackage{tikz}
\usetikzlibrary{arrows.meta, patterns, angles, quotes, calc}
\usetikzlibrary{positioning}
\usepackage{pgfplots}
\pgfplotsset{compat=1.18}

% AA Listati Di Codice (Listings)
% BB Configurazione Generale
\usepackage{listings}
\usepackage{upquote}

% BB Configurazione Generale
\lstset{
    basicstyle=\ttfamily\footnotesize\color{codetext},
    backgroundcolor=\color{codebg},
    frame=single,
    rulecolor=\color{codeframe},
    framesep=6pt,
    xleftmargin=8pt,
    xrightmargin=8pt,
    breaklines=true,
    breakatwhitespace=true,
    showstringspaces=false,
    keepspaces=true,
    columns=flexible,
    upquote=true,
    tabsize=2,
    captionpos=b,
    % Stili e Colori Predefiniti Globali
    keywordstyle=\bfseries\color{codekeywords},
    commentstyle=\itshape\color{codecomments},
    stringstyle=\color{codestrings}
}

% BB Definizione Dei Linguaggi

% CC YAML
\lstdefinelanguage{yaml}{
    keywords={true,false,null,yes,no,on,off},
    keywordstyle=\bfseries\color{codekeywords},
    comment=[l]{\#},
    commentstyle=\itshape\color{codecomments},
    stringstyle=\color{codestrings},
    morestring=[b]',
    morestring=[b]",
    sensitive=false
}

% CC Dockerfile
\lstdefinelanguage{dockerfile}{
    keywords={FROM, RUN, CMD, LABEL, EXPOSE, ENV, ADD, COPY, ENTRYPOINT,
              VOLUME, USER, WORKDIR, ARG, ONBUILD, STOPSIGNAL, HEALTHCHECK,
              SHELL, AS},
    keywordstyle=\bfseries\color{codedocker},
    comment=[l]{\#},
    commentstyle=\itshape\color{codecomments},
    stringstyle=\color{codestrings},
    morestring=[b]",
    morestring=[b]',
    sensitive=true
}

% CC INI
\lstdefinelanguage{ini}{
    comment=[l]{\#},
    morecomment=[l]{;},
    commentstyle=\itshape\color{codecomments},
    stringstyle=\color{codestrings},
    morestring=[b]",
    morestring=[b]',
    keywordstyle=\bfseries\color{codekeywords},
    sensitive=false
}

% CC JSON
\lstdefinelanguage{json}{
    keywords={true,false,null},
    keywordstyle=\bfseries\color{codekeywords},
    comment=[l]{//},
    morecomment=[s]{/*}{*/},
    commentstyle=\itshape\color{codecomments},
    stringstyle=\color{codestrings},
    morestring=[b]",
    sensitive=true
}

% CC PYTHON (Dialetto nativo pulito)
\lstset{defaultdialect=[3]Python}
\lstdefinelanguage{python}[]{Python}{
    morekeywords={as, async, await, True, False, None},
    sensitive=true
}

% CC BASH / SHELL
\lstdefinelanguage{bash}[]{sh}{
    morekeywords={docker, cd, run, build, stop, rm, rmi, push, pull, network, compose, openstack, rq, worker, keystone-manage, time, sleep, echo, exit, sudo, export, systemctl},
    sensitive=true
}

% CC HTML (Autonomo, previene il loop infinito del parser XML/HTML)
\lstdefinelanguage{html}{
    keywords={img, src, alt, class, id, div, span, table, tr, td, th, href, a, p, h1, h2, h3, form, input, button},
    keywordstyle=\bfseries\color{codekeywords},
    comment=[s]{<!--}{-->},
    commentstyle=\itshape\color{codecomments},
    stringstyle=\color{codestrings},
    morestring=[b]",
    morestring=[b]',
    sensitive=false
}

% CC SQL / CYPHER (Autonomo con supporto per // di Cypher)
\lstdefinelanguage{sql}{
    keywords={MATCH, RETURN, DETACH, DELETE, WHERE, CREATE, SET, WITH, ORDER, BY, ASC, DESC, LIMIT, SELECT, INSERT, INTO, VALUES, FROM, UPDATE, DROP, TABLE, JOIN, ON, AND, OR, NOT},
    keywordstyle=\bfseries\color{codekeywords},
    comment=[l]{//},
    morecomment=[l]{--},
    commentstyle=\itshape\color{codecomments},
    stringstyle=\color{codestrings},
    morestring=[b]',
    morestring=[b]",
    sensitive=false
}

% AA Algoritmi E Pseudocodice
\usepackage[ruled,vlined,linesnumbered]{algorithm2e}
\SetKwData{Dati}{Input}
\SetKwData{Risultato}{Output}
\SetKw{Inizializza}{Inizializza}
\SetKwFor{Per}{per}{fare}{fine per}
\SetKwComment{Commento}{// }{}

% BB Costrutti Di Controllo Aggiuntivi
\SetKwIF{Se}{AltrSe}{Altrimenti}{Se}{Allora}{Altrimenti Se}{Altrimenti}{Fine Se}
\SetKwFor{Mentre}{mentre}{Fare}{Fine Mentre}

% AA Struttura, Titoli Ed Ambienti
% BB Personalizzazione Titoli Sezioni
\usepackage{titlesec}
\titleformat{\section}{\normalfont\LARGE\bfseries}{\thesection}{1em}{}
\titleformat{\subsection}{\normalfont\Large\bfseries}{\thesubsection}{1em}{}
\titleformat{\subsubsection}{\normalfont\large\bfseries}{\thesubsubsection}{1em}{}

% BB Teoremi Ed Ambienti Ausiliari
\usepackage{comment}
\newtheorem{theorem}{Definizione}

% BB Liste Ed Elenchi Puntati
\usepackage{enumitem}
\setlist{noitemsep}
\usepackage{pifont}

% AA Riferimenti, Link Ipertestuali E Segnalibri
\usepackage[unicode]{hyperref}
\hypersetup{
    colorlinks=true,
    linkcolor=blue,
    urlcolor=blue,
    citecolor=blue,
    pdftitle={2) Cloud Computing},
    pdfauthor={Francesco Lo Verde},
    pdfpagemode=FullScreen,
}
\usepackage{bookmark}

% AA Metadati Del Documento
\title{Cloud Computing}

\date{\today}

\author{Francesco Lo Verde}
\begin{document}
\maketitle

\section{Fondamenti Del Cloud Computing: Dalla Concretezza Fisica Al Modello Di Servizio}

Nel contesto dell'ingegneria informatica moderna, il paradigma del \textbf{Cloud Computing} viene frequentemente descritto attraverso metafore astratte che rischiano di occultarne la reale natura infrastrutturale. Lungi dall'essere un'entità eterea o puramente concettuale, il cloud è uno spazio fisico rigidamente delimitato, caratterizzato da vincoli termici, requisiti energetici, estensioni geografiche e confini di sicurezza materiale.

Alla base di qualsiasi servizio erogato «on the cloud» vi sono imponenti strutture in calcestruzzo e acciaio: i \textbf{datacenter iperscalari}. All'interno di questi complessi, distribuiti strategicamente a livello globale, alloggiano decine di migliaia di server blade organizzati in file serrate di rack standardizzati da 19 pollici. Ciascun rack ospita risorse eterogenee di calcolo (\emph{compute}), memorizzazione (\emph{storage}) e commutazione (\emph{networking}), alimentate da linee ridondate e raffreddate da sofisticati impianti di condizionamento. Quando gli utenti eseguono un'applicazione o archiviano dati sul cloud, tali operazioni si traducono materialmente nell'allocazione di transistor di memoria RAM, scritture su array di dischi a stato solido o magnetici, e instradamento di pacchetti attraverso switch di aggregazione \emph{Top-of-Rack} (ToR) e dorsali in fibra ottica.

Un esempio paradigmatico di tale substrato è rappresentato dalla cosiddetta \textbf{Azure Fabric} di Microsoft. Con tale termine si identifica la totalità delle risorse fisiche (chassis, server, alimentatori, apparati di rete) coordinate da un'infrastruttura di controllo distribuita (il \emph{Fabric Controller}). Quest'ultimo funge da vero e proprio sistema operativo distribuito dell'infrastruttura, occupandosi del monitoraggio dello stato di salute dell'hardware, del provisioning automatico dei nodi, della gestione dei cicli di vita delle macchine virtuali e della riconfigurazione dinamica della topologia in presenza di guasti hardware. I servizi cloud ad alto livello poggiano interamente sull'orchestrazione deterministica di queste risorse fisiche tangibili.

% [IMMAGINE DEL FILE DI SLIDE CON NOME "CloudComputing.pdf"] Fotografie aeree di un datacenter iperscalare e corridoio interno con file di rack server illuminati — Slide 2 / Pagina 2 del PDF

\subsection{Caratteristiche Essenziali Del Modello Cloud}

La definizione standardizzata del cloud computing (formalizzata originariamente dal National Institute of Standards and Technology, NIST, e recepita nella didattica ingegneristica) si articola su cinque caratteristiche operative fondamentali:

\begin{enumerate}
    \item \textbf{Accesso Adeguato e Universale Alla Rete (\emph{Broad Network Access}):} le funzionalità computazionali sono disponibili attraverso la rete e accessibili tramite meccanismi standard, tali da promuovere l'uso da parte di piattaforme client eterogenee (quali workstation, laptop, dispositivi mobili e nodi IoT). L'infrastruttura richiede collegamenti ad altissima banda sia all'interno dei singoli datacenter sia lungo le dorsali geografiche private (\emph{inter-datacenter links}), consentendo la replica sincrona o asincrona dei dati a bassa latenza.

    \item \textbf{Servizio Misurato (\emph{Measured Service}) e Modello \emph{Pay-Per-Use}:} i sistemi cloud controllano e ottimizzano automaticamente l'allocazione delle risorse mediante capacità di misurazione a un livello di astrazione appropriato al tipo di servizio (ad esempio tempo di utilizzo della CPU, storage occupato, volume di traffico instradato o transazioni attive). Il cliente non sostiene investimenti iniziali per sovrastimare la capacità di picco (\emph{Capital Expenditure}, CapEx), ma remunera unicamente il consumo effettivo rendicontato come costo operativo (\emph{Operational Expenditure}, OpEx).

    \item \textbf{Condivisione Delle Risorse e Multi-Tenancy (\emph{Resource Pooling}):} le risorse computazionali del provider sono aggregate in un pool condiviso per servire una molteplicità di clienti (\emph{tenants}) attraverso un modello multi-tenant. Risorse fisiche e virtuali eterogenee vengono assegnate e riassegnate dinamicamente in base alla domanda, inducendo l'indipendenza dalla localizzazione fisica (\emph{location independence}): il consumatore raramente controlla o conosce l'esatta allocazione geografica dell'hardware, potendo tipicamente specificare solo la regione geografica o la \emph{availability zone}.

    \item \textbf{Rapida Elasticità (\emph{Rapid Elasticity}) e Scalabilità:} le risorse possono essere approvvigionate e rilasciate elasticamente, talvolta in modo del tutto autonomo, per ridimensionarsi rapidamente in funzione del carico. Dal punto di vista del committente, la capacità computazionale disponibile per il provisioning appare virtualmente illimitata e acquisibile in qualsiasi quantità e istante.

    \item \textbf{Self-Service Su Richiesta (\emph{On-Demand Self-Service}):} il consumatore può richiedere e configurare autonomamente capacità computazionale (ad esempio potenza di calcolo, storage e connettività di rete) senza richiedere l'interazione umana diretta con il personale del fornitore di servizi.
\end{enumerate}

\subsection{L'Opportunità Del Cloud: Analisi Costi-Benefici E Limiti Applicativi}

La transizione verso ambienti cloud pubblici è governata da rigorosi bilanciamenti architetturali ed economici. I benefici primari che spingono le organizzazioni ad adottare il paradigma cloud includono:

\begin{itemize}
    \item \textbf{Contrazione dei Costi Strutturali:} l'azzeramento dell'onere di acquisto di server fisici ad alta gamma e la drastica semplificazione nella gestione di complesse topologie di rete locali. Si riduce contestualmente la dimensione del reparto operativo sistemistico interno, delegando al provider la manutenzione di base delle macchine e dei cavidotti.
    \item \textbf{Elasticità di Fronte a Sovraccarichi Improvvisi:} in caso di picchi inattesi di traffico, l'infrastruttura elastica assorbe la domanda senza saturazione, allocando istanze aggiuntive in pochi istanti. Ottenere il medesimo risultato in un'infrastruttura proprietaria (\emph{on-premises}) richiederebbe l'acquisto permanente e il mantenimento inattivo di una notevole sovraccapacità hardware, con pesanti inefficienze economiche.
    \item \textbf{Resilienza e Disaster Recovery Geografico:} i dati e i servizi possono essere replicati su nodi dislocati in regioni remote del pianeta, garantendo la sopravvivenza dei carichi applicativi di fronte a disastri naturali di vasta scala, incendi, interruzioni prolungate della rete elettrica o persino conflitti bellici.
\end{itemize}

Nonostante tali vantaggi, la migrazione al cloud non costituisce una soluzione universale priva di controindicazioni. Nello specifico, si riscontrano impedimenti strutturali quali:

\begin{itemize}
    \item \textbf{Sistemi Monolitici a Carico Costante:} un'applicazione basata su un'architettura legacy tradizionale di tipo monolitico, caratterizzata da un profilo di carico stabile e privo di variabilità, non trae alcun beneficio tangibile dai meccanismi di elasticità dinamica offerti dal cloud.
    \item \textbf{Complessità di Re-Ingegnerizzazione (\emph{Refactoring}):} trasporre un software non strutturato per l'orchestrazione distribuita comporta una riscrittura profonda per disaccoppiare lo stato (\emph{stateless design}), segmentare la logica in microservizi e gestire la persistenza su database distribuiti. Tali attività introducono costi di ingegnerizzazione notevoli e lunghi tempi di transizione, spesso superiori ai risparmi immediati previsti.
\end{itemize}

\subsection{Profili Di Sicurezza: Vantaggi Difensivi E Criticità Di Governance}

La sicurezza rappresenta il fattore più dibattuto nell'adozione del cloud. Da un lato, il cloud provider gode di economie di scala tali da giustificare investimenti ingenti nell'applicazione tempestiva di patch di sicurezza sui livelli bassi dello stack, nello sviluppo di procedure avanzate di hardening dei sistemi operativi e nella selezione rigorosa di personale sistemistico sottoposto a controlli sul background. Inoltre, i datacenter del cloud pubblico garantiscono standard elevatissimi di sicurezza perimetrale fisica: varchi con autenticazione biometrica, sorveglianza continuativa e rigidi protocolli di accesso alle sale macchine, difficilmente replicabili all'interno di un'azienda di medie dimensioni.

D'altro canto, la trasposizione di dati e carichi computazionali su piattaforme terze solleva problematiche di governo dei dati:

\begin{itemize}
    \item \textbf{Cessione Del Controllo Diretto:} l'onboarding di software sul cloud comporta la delega del controllo operativo primario al fornitore del servizio. La responsabilità della corretta segregazione dei dati e dell'integrità dei meccanismi di cifratura non ricade più esclusivamente sul team di sicurezza interno.
    \item \textbf{Sovranità Del Dato e Giurisdizione Territoriale:} le multinazionali del cloud distribuiscono i dati su nodi mondiali. Qualora tali datacenter risiedano in giurisdizioni estere con quadri legislativi differenti o geopoliticamente ostili, possono sorgere conflitti di conformità legale. Un tema critico è determinare la responsabilità giuridica dell'azienda proprietaria del dato in caso di accesso indebito o perquisizione informatica autorizzata da governi esteri.
    \item \textbf{Diritto All'Oblio ed Effettiva Cancellazione:} in ottemperanza ai regolamenti sulla privacy (come il GDPR), un cittadino può richiedere la distruzione integrale e verificabile delle proprie informazioni personali. In ambienti distribuiti, dove i blocchi di storage vengono replicati e distribuiti tramite architetture di caching, copie shadow e snapshot asincrone, garantire e certificare tecnicamente che ogni traccia del record sia stata fisicamente rimossa rappresenta una sfida formidabile.
    \item \textbf{Politiche Di Notifica Dei Data Breach:} le clausole contrattuali devono specificare rigorosamente entro quali finestre temporali il provider sia vincolato a notificare una violazione dei sistemi perimetrali (argomento dettagliato nella letteratura di settore da Winkler nel 2012). Parallelamente, si articola la responsabilità e la responsabilità civile dell'azienda cliente qualora essa identifichi una falla nei propri applicativi residenti sul cloud senza notificarla tempestivamente al provider.
\end{itemize}

\newpage

\section{Tassonomia Dei Modelli Di Servizio E Matrice Di Responsabilità}

Il panorama del Cloud Computing articola le proprie offerte lungo uno spettro continuo di astrazione e automazione. Ad un estremo dello spettro si collocano soluzioni che demandano all'utente la totale configurazione dell'ambiente operativo virtuale; all'estremo opposto, l'infrastruttura sottostante scompare del tutto dalla vista dello sviluppatore, gestita interamente in modalità autonoma da piattaforme automatizzate.

\subsection{Lo Spettro Dei Servizi: IaaS, PaaS, SaaS E FaaS}

I quattro modelli di servizio canonici si differenziano in base alla linea di demarcazione che separa le responsabilità dell'utente da quelle del cloud provider:

\begin{figure}[H]
\centering
\begin{tikzpicture}[
    scale=0.95, every node/.style={font=\footnotesize},
    box/.style={rectangle, draw=black, minimum width=3.4cm, minimum height=0.6cm, align=center, inner sep=2pt},
    layerusr/.style={box, fill=theoryblue!25},
    layerprv/.style={box, fill=backcolour},
    title/.style={font=\bfseries\small, align=center, minimum height=1cm}
]
    % Colonne: Traditional IT, IaaS, PaaS, SaaS
    \node[title] at (0, 6.2) {On-Premises\\(Traditional IT)};
    \node[title] at (4, 6.2) {Infrastructure\\as a Service (IaaS)};
    \node[title] at (8, 6.2) {Platform\\as a Service (PaaS)};
    \node[title] at (12, 6.2) {Software\\as a Service (SaaS)};

    % Livelli:
    % 1: Applications, 2: Data, 3: Runtime, 4: Middleware, 5: OS, 6: Virtualization, 7: Servers, 8: Storage, 9: Networking

    % Traditional IT (Tutto User)
    \node[layerusr] at (0, 5.0) {Applicazioni};
    \node[layerusr] at (0, 4.3) {Dati};
    \node[layerusr] at (0, 3.6) {Runtime};
    \node[layerusr] at (0, 2.9) {Middleware};
    \node[layerusr] at (0, 2.2) {Sistema Operativo};
    \node[layerusr] at (0, 1.5) {Virtualizzazione};
    \node[layerusr] at (0, 0.8) {Server Fisici};
    \node[layerusr] at (0, 0.1) {Storage};
    \node[layerusr] at (0, -0.6) {Networking};

    % IaaS (User: App, Data, Runtime, Middleware, OS; Provider: Rest)
    \node[layerusr] at (4, 5.0) {Applicazioni};
    \node[layerusr] at (4, 4.3) {Dati};
    \node[layerusr] at (4, 3.6) {Runtime};
    \node[layerusr] at (4, 2.9) {Middleware};
    \node[layerusr] at (4, 2.2) {Sistema Operativo};
    \node[layerprv] at (4, 1.5) {Virtualizzazione};
    \node[layerprv] at (4, 0.8) {Server Fisici};
    \node[layerprv] at (4, 0.1) {Storage};
    \node[layerprv] at (4, -0.6) {Networking};

    % PaaS (User: App, Data; Provider: Rest)
    \node[layerusr] at (8, 5.0) {Applicazioni};
    \node[layerusr] at (8, 4.3) {Dati};
    \node[layerprv] at (8, 3.6) {Runtime};
    \node[layerprv] at (8, 2.9) {Middleware};
    \node[layerprv] at (8, 2.2) {Sistema Operativo};
    \node[layerprv] at (8, 1.5) {Virtualizzazione};
    \node[layerprv] at (8, 0.8) {Server Fisici};
    \node[layerprv] at (8, 0.1) {Storage};
    \node[layerprv] at (8, -0.6) {Networking};

    % SaaS (Provider: Tutto)
    \node[layerprv] at (12, 5.0) {Applicazioni};
    \node[layerprv] at (12, 4.3) {Dati};
    \node[layerprv] at (12, 3.6) {Runtime};
    \node[layerprv] at (12, 2.9) {Middleware};
    \node[layerprv] at (12, 2.2) {Sistema Operativo};
    \node[layerprv] at (12, 1.5) {Virtualizzazione};
    \node[layerprv] at (12, 0.8) {Server Fisici};
    \node[layerprv] at (12, 0.1) {Storage};
    \node[layerprv] at (12, -0.6) {Networking};

    % Legenda
    \node[layerusr, minimum width=0.6cm, minimum height=0.4cm] at (1.75, -1.5) {};
    \node[right, font=\footnotesize] at (2.15, -1.55) {Gestito dall'utente (\emph{Customer})};
    \node[layerprv, minimum width=0.6cm, minimum height=0.4cm] at (9.0, -1.5) {};
    \node[right, font=\footnotesize] at (9.4, -1.55) {Gestito dal fornitore (\emph{Cloud Provider})};
\end{tikzpicture}
\caption{Modello di responsabilità condivisa (\emph{Shared Responsibility Model}) attraverso i differenti paradigmi cloud.}
\label{fig:shared_responsibility}
\end{figure}

% [IMMAGINE DEL FILE DI SLIDE CON NOME "CloudComputing.pdf"] Diagramma a colonne orizzontali sulla ripartizione delle responsabilità operative fra cliente e fornitore nei modelli On-Premises, IaaS, PaaS e SaaS — Slide 11 / Pagina 11 del PDF

\subsubsection{Infrastructure as a Service (IaaS)}
In questo modello, il provider si limita a noleggiare le risorse infrastrutturali di base: cicli di calcolo su processori virtualizzati, capacità di archiviazione a blocchi o a oggetti, e componenti di interconnessione di rete (quali interfacce virtuali, sottoreti e routing di base).

L'ingegnere o il sistemista ha la piena facoltà di configurare la macchina virtuale, installare il sistema operativo desiderato (distribuzioni Linux, Windows Server), impostare i demoni di supporto (\emph{middleware}) e definire politiche avanzate di filtraggio del traffico (firewall host-based, regole di instradamento e bilanciatori di carico). IaaS offre la massima flessibilità architetturale e portabilità verso ambienti on-premises, a fronte della completa responsabilità sull'amministrazione, aggiornamento e messa in sicurezza dei sistemi operativi e dei layer sovrastanti.
\emph{Esempio di riferimento:} \textbf{Amazon Elastic Compute Cloud (EC2)}.

\subsubsection{Platform as a Service (PaaS)}
Il modello PaaS eleva il piano di astrazione: l'infrastruttura fisica, il livello di ipervisione, il sistema operativo, i componenti middleware e i relativi ambienti di esecuzione (\emph{runtime}) sono interamente gestiti dal provider. All'utente viene fornito un ambiente di hosting pronto all'uso, corredato di compilatori, interpreti, librerie di supporto e strumenti di build per specifici linguaggi di programmazione (ad esempio Java, Python, .NET, Node.js).

Lo sviluppatore non deve gestire indirizzi IP privati, allocazione della RAM a livello di kernel o patch del sistema operativo: è sufficiente specificare la dimensione o il numero di istanze applicative richieste e caricare il codice sorgente o il pacchetto compilato. La piattaforma scala e instrada automaticamente le richieste.
\emph{Esempi tipici:} \textbf{AWS Elastic Beanstalk}, \textbf{Google App Engine}.

\subsubsection{Software as a Service (SaaS)}
Nel modello SaaS, l'intero stack applicativo è realizzato, gestito, manutenuto e ospitato dal fornitore di servizi. Il consumatore finale vi accede generalmente tramite un'interfaccia client leggera (\emph{thin client}), tipicamente un browser web o un'app mobile dedicata.

L'utente non ha alcun controllo sull'infrastruttura, sulla rete, sui server, sui sistemi operativi, né tantomeno sulla struttura interna dell'applicazione, potendo esclusivamente interagire con le logiche di business fornite ed impostare un ventaglio ristretto di configurazioni specifiche dell'utenza. Le applicazioni SaaS operano diffusamente secondo logiche di abbonamento periodico (\emph{subscription}) o in modalità freemium.
\emph{Esempi standard:} \textbf{Microsoft Office 365}, \textbf{Google Docs}, \textbf{Dropbox}, \textbf{Citrix GoToMeeting}.

\subsubsection{Function as a Service (FaaS) E Serverless Computing}
Il paradigma FaaS costituisce l'estremo contemporaneo dello spettro cloud: esso azzera completamente il concetto stesso di server o di istanza permanente visibile all'utente (\textbf{Serverless Architecture}). Lo sviluppatore decompone l'applicazione in blocchi logici atomici ed effimeri (le funzioni), che vengono registrati sulla piattaforma cloud associandoli a specifici eventi scatenanti (\emph{event triggers}, ad esempio il caricamento di un file su uno storage a oggetti, una chiamata API HTTP o un messaggio inserito in una coda).

Le istanze di esecuzione non rimangono attive in attesa di richieste: la piattaforma istanzia dinamicamente un container lightweight al verificarsi dell'evento, esegue il codice della funzione e distrugge l'ambiente non appena l'elaborazione è completata. Il modello di fatturazione risulta ultra-granulare, remunerando unicamente il tempo di esecuzione effettivo del codice (spesso misurato con risoluzione al millisecondo) e la memoria consumata, riducendo a zero i costi nei periodi di inattività. Questo pattern rappresenta l'infrastruttura abilitante ideale per le moderne architetture a microservizi disaccoppiati.

\subsection{Approfondimento Tecnico: Il Ruolo Del «Runtime»}

Per comprendere appieno la cesura architetturale introdotta dal modello PaaS rispetto al modello IaaS, è opportuno chiarire formalmente il concetto sistemistico di \textbf{Runtime}.

Nel dominio dei sistemi di calcolo, con il termine \emph{runtime} si identificano l'insieme delle istruzioni, dei demoni e delle strutture dati che vengono mandate in esecuzione durante la vita attiva del programma, operando in modo trasparente rispetto al codice sorgente redatto esplicitamente dallo sviluppatore, ma risultando indispensabili per il suo funzionamento.

\begin{itemize}
    \item \textbf{Distinzione tra Codice di Libreria e Codice di Runtime:} sebbene entrambi risiedano nello spazio di memoria dell'applicazione, il codice di libreria implementa funzionalità algoritmiche accessorie invocate esplicitamente (ad esempio librerie grafiche, matematiche o crittografiche). Il codice di runtime, al contrario, implementa la semantica fondamentale del linguaggio di programmazione stesso. Esso gestisce il layout della memoria (allocazione dinamica dello heap e dello stack), l'accesso alle variabili globali, i meccanismi di passaggio dei parametri nei frame di stack delle funzioni, la gestione delle eccezioni e l'interfacciamento a basso livello con le chiamate di sistema del kernel (\emph{system calls}).
    \item \textbf{Impatto a seconda del paradigma di programmazione:} linguaggi a basso livello compilati direttamente in codice macchina (come il linguaggio C) possiedono un runtime estremamente snello (tipicamente circoscritto al modulo d'avvio `crt0`, responsabile della preparazione dello stack e del passaggio di `argc` e `argv` al `main`). Linguaggi dinamici, a oggetti o gestiti (quali Objective-C con il dynamic message passing, oppure Java e Python con Virtual Machine dotate di \emph{Garbage Collector}) possiedono runtime estesi e complessi. Il compilatore genera codice macchina formulando precise assunzioni operative sul runtime sottostante.
\end{itemize}

In un'infrastruttura IaaS, l'amministratore deve occuparsi dell'installazione, configurazione e aggiornamento di questi ambienti (ad esempio la JVM per Java, l'interprete Python e i moduli binari CPython correlati). Nel modello PaaS, tale onere scompare: il runtime è offerto come servizio integrato e garantito dal fornitore di piattaforma.

\subsection{Composizione Ed Ecosistema Dei Servizi Cloud}

Le categorie di servizio non vivono in isolamento, ma si sovrappongono e si compongono dinamicamente nell'ecosistema industriale:

\begin{itemize}
    \item \textbf{SaaS Costruiti Sopra IaaS/PaaS Terzi:} è consuetudine per i fornitori di software distribuire le proprie soluzioni SaaS allocando le risorse sottostanti presso fornitori IaaS o PaaS di livello mondiale. Un caso storico celebre è rappresentato da \textbf{Dropbox}, che originariamente eseguiva e archiviava la totalità del proprio carico di storage e calcolo sulle istanze EC2 e i bucket S3 di Amazon Web Services, prima di strutturare un'infrastruttura proprietaria dedicata (documentato da Metz nel 2016).
    \item \textbf{Ibridazione Dei Servizi Nei Grandi Provider:} piattaforme quali \textbf{Microsoft Azure} non operano rigidamente secondo un solo modello, ma espongono contestualmente sia servizi IaaS puri (macchine virtuali configurabili dall'utente) sia servizi PaaS (Azure App Services, Azure SQL Database), consentendo architetture ibride.
    \item \textbf{Esposizione Di Modelli PaaS da Sistemi SaaS:} la piattaforma \textbf{Salesforce.com}, originariamente concepita come puro SaaS per la gestione delle relazioni con i clienti (CRM), ha integrato e commercializzato la piattaforma \textbf{Force.com}. Quest'ultima costituisce un'offerta PaaS a tutti gli effetti, che permette a sviluppatori terzi di scrivere codice (tramite il linguaggio proprietario Apex) per creare estensioni ed applicazioni verticali perfettamente integrate con i database dell'ecosistema Salesforce.
\end{itemize}

% [IMMAGINE DEL FILE DI SLIDE CON NOME "CloudComputing.pdf"] Schema a strati verticali e orizzontali dell'ecosistema cloud con marchi di mercato ripartiti nei livelli SaaS, PaaS e IaaS — Slide 24 / Pagina 24 del PDF

\newpage

\section{Modelli Di Maturità Architetturale Nel SaaS: Da Single-Tenant A Multi-Tenant}

Nella transizione di un'applicazione monolitica classica verso il modello Software as a Service, i progettisti software si trovano ad affrontare la sfida di servire una molteplicità di clienti (\emph{tenants}) garantendo l'isolamento dei dati, la sicurezza, la manutenibilità del codice e la sostenibilità economica dei costi infrastrutturali.

A tale scopo, Microsoft ha codificato un modello di riferimento evolutivo in quattro stadi, noto come \textbf{SaaS Maturity Levels}.

% [IMMAGINE DEL FILE DI SLIDE CON NOME "CloudComputing.pdf"] Diagramma a 4 quadranti illustrante i 4 livelli di maturità SaaS con omini (tenants), istanze applicative e bilanciatori di carico — Slide 13 e 16 / Pagine 13 e 16 del PDF

\subsection{I Quattro Livelli Di Maturità}

\begin{table}[H]
\centering
\small
\begin{tabularx}{\textwidth}{l p{3.5cm} X p{3cm}}
\toprule
\textbf{Livello} & \textbf{Nome Modello} & \textbf{Caratteristiche Architetturali} & \textbf{Meccanismo di Scalabilità} \\
\midrule
\textbf{Level 1} & \emph{Ad-Hoc / Custom} & Single-Tenant puro. Ciascun cliente riceve un'istanza software separata e un codice sorgente personalizzato (\emph{fork}). Nessuna condivisione di codice, database o hardware. & \emph{Scale-Up} per singolo server dedicato. \\
\addlinespace
\textbf{Level 2} & \emph{Configurable} & Single-Tenant su base comune. Istanza separata per ogni cliente, ma il codice sorgente è unificato. Le differenziazioni funzionali avvengono solo tramite metadati di configurazione. & \emph{Scale-Up} per singola istanza dedicata. \\
\addlinespace
\textbf{Level 3} & \emph{Configurable, Multi-Tenant-Efficient} & Multi-Tenant nativo. Un'unica istanza applicativa condivisa serve simultaneamente molteplici clienti. La separazione dei dati è demandata a logiche logiche interne. & \emph{Scale-Up} hardware sul server centrale condiviso. \\
\addlinespace
\textbf{Level 4} & \emph{Scalable, Configurable, Multi-Tenant-Efficient} & Multi-Tenant multi-tier. Farm di istanze identiche eseguite su un cluster di macchine eterogenee, precedute da bilanciatori di carico (\emph{Tenant Load Balancer}). & \emph{Scale-Out} orizzontale elastico e dinamico. \\
\bottomrule
\end{tabularx}
\caption{Matrice riassuntiva dei quattro livelli di maturità per architetture Software as a Service (SaaS).}
\label{tab:saas_maturity}
\end{table}

\subsection{Analisi Dettagliata Dei Paradigmi: Single-Tenant Vs Multi-Tenant}

\subsubsection{Paradigma Single-Tenant (Maturity Levels 1 e 2)}
Nel modello single-tenant (noto commercialmente anche come modello \textbf{Hosted}), ad ogni cliente viene assegnata un'infrastruttura dedicata. Non vi è alcuna condivisione né a livello logico né a livello di memorizzazione. All'acquisizione di un nuovo cliente, il fornitore provvede al provisioning di nuove risorse hardware (o partizioni logiche distinte all'interno del medesimo server Web, come differenti virtual host o siti dedicati su Microsoft IIS).

\begin{itemize}
    \item \textbf{Vantaggi Ingegneristici:}
    \begin{itemize}[noitemsep]
        \item \emph{Rapidità di roll-out iniziale:} l'applicazione aziendale preesistente non deve subire modifiche logiche profonde per riconoscere il concetto di tenant; è sufficiente replicarla in una nuova macchina virtuale.
        \item \emph{Bassa complessità architetturale iniziale:} non richiede competenze avanzate di ingegneria SaaS o isolamento dati dinamico; le problematiche hardware sono interamente delegate al fornitore IaaS.
        \item \emph{Supporto per applicazioni legacy non web-native:} rende agevole erogare via cloud software desktop o client/server tradizionali (es. applicazioni sviluppate in Visual Basic), distribuendole tramite soluzioni di virtualizzazione dello schermo come Citrix XenApp/Terminal Services.
        \item \emph{Conformità a requisiti di sicurezza stringenti:} clienti operanti in settori ad alta regolamentazione (es. istituti bancari, enti sanitari, difesa) rifiutano categoricamente la coesistenza dei propri dati negli stessi database o spazi di indirizzamento di altre entità concorrenti.
        \item \emph{Personalizzazione estrema (solo Livello 1):} permette di creare versioni altamente ad-hoc per singoli grandi clienti commerciali modificando direttamente il codice sorgente.
    \end{itemize}

    \item \textbf{Svantaggi e Limiti Operativi:}
    \begin{itemize}[noitemsep]
        \item \emph{Complessità di manutenzione esponenziale:} qualora sia necessario distribuire una correzione critica di sicurezza (\emph{bug fix} o patch), l'operazione deve essere replicata manualmente o orchestrata $n$ volte, dove $n$ indica il numero di clienti serviti. Se nel Livello 1 ciascun cliente possiede un branch di codice differente, il porting della patch richiede sforzi titanici di reimplementazione.
        \item \emph{Costi operativi (OpEx) insostenibili su larga scala:} l'allocazione di risorse hardware dedicate e la frammentazione gestionale riducono drasticamente i margini di profitto all'aumentare della clientela.
    \end{itemize}
\end{itemize}

\subsubsection{Paradigma Multi-Tenant (Maturity Levels 3 e 4)}
La multi-tenancy rappresenta una proprietà architetturale fondamentale: un unico sistema software è progettato strutturalmente per riconoscere il contesto dell'utente (singolo operatore, gruppo o organizzazione) ed erogare viste, regole aziendali, schemi e flussi operativi specifici in virtù dei metadati associati a quel particolare tenant.

Ogni singola funzionalità deve essere riprogettata per operare costantemente filtrata dal \textbf{Tenant Context}. La sfida più complessa risiede nella \textbf{segregazione dei dati}: le query al database non devono mai permettere la fuoriuscita accidentale di record appartenenti ad altri inquilini, ricorrendo a colonne discriminatorie (\emph{Tenant ID}), schemi di database separati, o partizionamento a livello di motore di archiviazione.

\begin{itemize}
    \item \textbf{Vantaggi Ingegneristici:}
    \begin{itemize}[noitemsep]
        \item \emph{Efficienza economica ed economie di scala:} i costi di esercizio dell'infrastruttura sono ripartiti sull'intera base di utenza, permettendo al fornitore software indipendente (ISV) di proporre prezzi di abbonamento competitivi che stimolano un circolo virtuoso di adozione di mercato (\emph{snowball effect}).
        \item \emph{Manutenibilità unificata e qualità del software:} il team di ingegneria sviluppa, testa e aggiorna una sola ed unica versione del codice sorgente. Il rilascio di un aggiornamento ha effetto istantaneo per l'intera base utenti.
        \item \emph{Estendibilità guidata da metadati:} se un cliente richiede attributi informativi aggiuntivi, questi non vengono codificati alterando le tabelle fisiche di base ma iniettati tramite moduli di custom-fields gestiti dal layer logico dell'applicazione.
        \item \emph{Scalabilità per volumi di massa:} architetture come Salesforce.com gestiscono centinaia di migliaia di tenant all'interno del medesimo sistema distribuito; una scala operativa impraticabile da sostenere tramite istanze single-tenant isolate.
    \end{itemize}

    \item \textbf{Svantaggi e Vincoli di Sviluppo:}
    \begin{itemize}[noitemsep]
        \item \emph{Elevato investimento iniziale:} riprogettare una base di codice esistente per supportare nativamente la multi-tenancy richiede profonde modifiche a livello di architettura software e del dato.
        \item \emph{Rischio d'impresa e competenze richieste:} la realizzazione di un'architettura multi-tenant efficiente impone competenze specialistiche di alto profilo per la gestione di caching distribuito, isolamento della memoria e failover; qualora assenti all'interno dell'azienda, tali competenze impongono il ricorso a partnership esterne costose.
    \end{itemize}
\end{itemize}

\subsubsection{Discriminante Tra Livello 3 E Livello 4}
Sebbene entrambi i livelli implementino la multi-tenancy, la discriminante fondamentale risiede nella topologia infrastrutturale e nella strategia di dimensionamento:

\begin{itemize}
    \item \textbf{Livello 3 (\emph{Scale-Up}):} l'applicazione multi-tenant risiede su un'unica istanza o su un singolo server logico. All'aumentare dei tenant, l'unica strategia per prevenire il degrado prestazionale consiste nell'ampliare la capacità hardware della macchina (incremento verticale di CPU, memoria RAM o throughput del sottosistema I/O). Tale approccio incontra presto limiti fisici ed economici invalicabili.
    \item \textbf{Livello 4 (\emph{Scale-Out}):} il sistema è segmentato in una \textbf{farm di istanze identiche multi-tenant stateless}, attestate dietro a bilanciatori di carico intelligenti (\emph{Tenant Load Balancer}). La scalabilità si ottiene orizzontalmente aggiungendo macchine virtuali a basso costo nel pool computazionale. L'infrastruttura garantisce l'\textbf{elasticità dinamica}, allocando o dismettendo nodi applicativi in tempo reale al variare del carico aggregato generato dai tenant.
\end{itemize}

\newpage

\section{Modelli Di Rilascio E Piattaforme Infrastrutturali}

La fruibilità e l'accessibilità operativa dei servizi cloud dipendono dal modello di deployment prescelto, il quale definisce il perimetro di condivisione e la governance dell'infrastruttura fisica.

\subsection{Tassonomia Dei Modelli Di Rilascio (\emph{Deployment Models})}

La standardizzazione identifica quattro modelli primari:

\begin{enumerate}
    \item \textbf{Public Cloud (Cloud Pubblico):} l'infrastruttura computazionale è di proprietà di una grande azienda (\emph{Cloud Provider}) che alloca e noleggia le risorse a condizioni contrattuali definite verso l'utenza generale o grandi settori industriali.
    \begin{itemize}[noitemsep]
        \item \emph{Caratteristiche distintive:} architettura fortemente omogenea, policy di accesso comuni, condivisione totale delle risorse fisiche tramite virtualizzazione multi-tenant, locazione a canone variabile e massime economie di scala.
    \end{itemize}

    \item \textbf{Private Cloud (Cloud Privato):} l'infrastruttura è operata a beneficio esclusivo di una singola organizzazione aziendale. Può risiedere fisicamente all'interno dell'azienda (\emph{on-premise}) oppure essere ospitata all'esterno presso un fornitore terzo (\emph{off-premise}), ma l'accesso è rigidamente limitato agli utenti facenti capo alla medesima entità giuridica.
    \begin{itemize}[noitemsep]
        \item \emph{Caratteristiche distintive:} infrastruttura tipicamente eterogenea, politiche di sicurezza su misura (\emph{tailored policies}), risorse interamente dedicate e controllo end-to-end garantito su stack hardware e software.
    \end{itemize}

    \item \textbf{Community Cloud (Cloud Comunitario):} l'infrastruttura è condivisa da una cerchia ristretta di organizzazioni appartenenti a filiere affini o contraddistinte da requisiti operativi omologhi (quali missioni strategiche congiunte, rigide norme di sicurezza governativa o standard stringenti di conformità normativa).
    \begin{itemize}[noitemsep]
        \item \emph{Dinamica e Benefici:} più imprese con target compatibili condividono il medesimo hardware per ripartire le spese di investimento e di gestione. L'architettura garantisce la segregazione e riservatezza dei dati di ciascun membro, consentendo al contempo l'apertura selettiva di varchi di cooperazione in tempo reale. Il successo del modello fa perno sulla stretta convergenza di buone pratiche e sulla reciproca fiducia contrattuale.
        \item \emph{Esempio concreto:} l'infrastruttura \textbf{IBM SoftLayer} dedicata alle agenzie federali degli Stati Uniti d'America, concepita per garantire requisiti stringenti di conformità governativa all'interno di una piattaforma comune riservata.
    \end{itemize}

    \item \textbf{Hybrid Cloud (Cloud Ibrido):} composizione organica di due o più tipologie eterogenee di cloud (pubblico, privato o comunitario) che mantengono la loro individualità strutturale, ma risultano integrate e coordinate da tecnologie standard o proprietarie atte ad assicurare la portabilità di dati ed applicazioni.
    \begin{itemize}[noitemsep]
        \item \emph{Meccanismo del Cloud Bursting:} un'applicazione basata su un cloud privato soddisfa il carico ordinario attraverso l'infrastruttura interna; nel momento in cui si registra un picco di richieste computazionali che saturerebbe la capacità on-premises, il sistema devia automaticamente il carico eccedente (\emph{bursting}) verso istanze dinamiche allocate su un cloud pubblico, rientrando nei confini interni al termine dell'evento.
    \end{itemize}
\end{enumerate}

% [IMMAGINE DEL FILE DI SLIDE CON NOME "CloudComputing.pdf"] Schemi architetturali delle connessioni Enterprise-to-Cloud per Public, Private, Community e Hybrid Cloud — Slide 20, 21, 22, 23 / Pagine 20, 21, 22, 23 del PDF

\subsection{Piattaforme Open Source Di Virtualizzazione IaaS}

L'implementazione concreta di ambienti cloud privati e comunitari ha trovato ampio sviluppo grazie alla maturazione di progetti open source per l'automazione e il provisioning delle risorse fisiche sottostanti:

\begin{table}[H]
\centering
\small
\begin{tabularx}{\textwidth}{l p{4cm} l p{2.8cm} l}
\toprule
\textbf{Piattaforma} & \textbf{Origine e Storia} & \textbf{Licenza} & \textbf{Linguaggi} & \textbf{Difficoltà Install.} \\
\midrule
\textbf{OpenStack} & Fondata congiuntamente da Rackspace e dalla NASA nel 2010. Divenuta lo standard de facto industriale. & Apache v2.0 & Python & Elevata (\emph{difficult}) \\
\addlinespace
\textbf{CloudStack} & Rilasciata inizialmente da Citrix nella Apache Software Foundation nel 2011. & Apache v2.0 & Java, C & Media (\emph{medium}) \\
\addlinespace
\textbf{Eucalyptus} & Nata da ricerche accademiche presso la Rice University e la UCSB; acquisita da HP. Compatibilità nativa con AWS API. & GPL v3.0 & Java, C & Media (\emph{medium}) \\
\addlinespace
\textbf{OpenNebula} & Piattaforma snella promossa e sponsorizzata da OpenNebula Systems. & Apache v2.0 & C++, C, Ruby, Java, Shell, lex, yacc & Media (\emph{medium}) \\
\bottomrule
\end{tabularx}
\caption{Quadro comparativo delle principali piattaforme software open source per la gestione di infrastrutture IaaS.}
\label{tab:open_source_iaas}
\end{table}

\newpage

\section{Vulnerabilità E Continuità Operativa: L'Amplificazione Delle Minacce Nel Cloud}

Le vulnerabilità informatiche che affliggono i sistemi distribuiti rappresentano minacce storiche della disciplina sistemistica (\emph{evergreen}). Tuttavia, il paradigma del Cloud Computing introduce un fattore di amplificazione delle conseguenze a causa dell'altissima densità di dati aggregati e della dipendenza strutturale degli utenti dai canali di comunicazione del provider.

Le principali classi di minaccia amplificate dal contesto cloud comprendono:
\begin{itemize}
    \item \textbf{Inattività e Fuori Servizio (\emph{Downtimes}):} l'interruzione dei servizi di un grande fornitore paralizza simultaneamente migliaia di imprese clienti che non possiedono meccanismi locali di fallback.
    \item \textbf{Phishing e Compromissione Credenziali (\emph{Password Cracking}):} nei sistemi multi-tenant, la violazione dell'account amministrativo o delle chiavi API di un'organizzazione espone l'intera infrastruttura virtuale del tenant (macchine, database, reti software-defined) a cancellazione o manipolazione malevola.
    \item \textbf{Botnet e Vettori Di Diffusione Malware:} l'immensa capacità computazionale e la larghezza di banda offerta dai nodi IaaS compromessi possono essere sfruttate da attori malevoli come piattaforme ideali per l'esecuzione di attacchi Distributed Denial of Service (DDoS) globali.
    \item \textbf{Perdita Irreversibile Dei Dati (\emph{Data Loss}):} sebbene la replica sia un pilastro del cloud, errori architetturali nei sistemi di memorizzazione o falle nei processi di snapshot possono tradursi nella distruzione definitiva di moli colossali di informazioni.
\end{itemize}

\subsection{Caso Di Studio: Il Disastro Microsoft/Danger Sidekick (2009)}

Uno degli eventi storici più significativi nell'evoluzione della percezione del rischio nei servizi cloud è rappresentato dall'incidente occorso nell'autunno del 2009 ai danni dell'infrastruttura dei dispositivi mobili \textbf{T-Mobile Sidekick}, gestita da \textbf{Danger} (sussidiaria acquisita da Microsoft nel 2008).

Il dispositivo basava la propria esperienza d'uso su un'architettura marcatamente cloud-centrica: i dati personali dell'utente (rubrica contatti, calendari, liste delle attività e gallerie fotografiche) non risiedevano in via primaria sulla memoria fisica del terminale, ma venivano costantemente sincronizzati e archiviati sui server centrali del provider.

A seguito di un intervento di riconfigurazione e manutenzione dei database di produzione, accompagnato da un contestuale guasto hardware sistemico e dal fallimento imprevisto delle procedure di ripristino dai supporti di backup di emergenza (\emph{backup failure}), l'infrastruttura subì una corruzione irreparabile dei dati di archiviazione. La comunicazione ufficiale rilasciata agli utenti evidenziò in modo drammatico l'impatto della perdita di sovranità materiale sui propri dati:

\begin{quote}
\emph{«Regrettably, based on Microsoft/Danger's latest recovery assessment of their systems, we must now inform you that personal information stored on your device—such as contacts, calendar entries, to-do lists or photos—that is no longer on your Sidekick almost certainly has been lost as a result of a server failure at Microsoft/Danger.»}
\end{quote}

L'episodio mise in luce i seguenti insegnamenti ingegneristici per l'architettura dei sistemi cloud:
\begin{enumerate}
    \item \textbf{Il Rischio Del Single Point of Failure Operativo:} demandare l'integrità del patrimonio informativo a un singolo ecosistema logico senza un piano di disaster recovery indipendente espone l'utente al disastro sistemico in caso di errore procedurale del provider.
    \item \textbf{La Fallacia Del Backup Non Testato:} la disponibilità teorica di meccanismi di replica non equivale a una garanzia di ripristino; i backup devono essere sottoposti a continui test empirici di consistenza e ripristinabilità (\emph{recovery testing}).
    \item \textbf{Architetture Ibride e Caching Locale Resiliente:} la progettazione di sistemi client/cloud non deve mai privare il nodo periferico di una copia locale consistente dei dati operativi primari, garantendo la possibilità di operare anche a fronte dell'azzeramento o dell'inaccessibilità prolungata dei nodi centrali.
\end{enumerate}

% AA -------------------------------------------------------------------------------

\section{Sicurezza Nella Virtualizzazione E Nel Cloud Computing}

La transizione dai data center proprietari alle infrastrutture di calcolo virtualizzate e ai fornitori di servizi in cloud pubblico ha trasformato radicalmente il paradigma architetturale dei sistemi distribuiti. Se da un lato il consolidamento dell'hardware e la gestione dinamica delle risorse offrono flessibilità e abbattimento dei costi di investimento in conto capitale (\emph{Capital Expenditure}, CapEx), dall'altro introducono vettori d'attacco inediti e complessità sistemiche non presenti nelle infrastrutture tradizionali fisicamente segregate.

\subsection{Vulnerabilità, Incidenti E Rischi Reali Nel Cloud}

La comprensione delle problematiche di sicurezza del cloud non può prescindere dall'analisi dei riscontri operativi sul campo. L'interdipendenza tra i componenti software di orchestrazione e le risorse condivise espone il modello operativo ad almeno quattro categorie critiche di eventi avversi: interruzioni di servizio non pianificate (\emph{downtime}), attacchi basati su ingegneria sociale (\emph{phishing}), abuso computazionale per attacchi crittografici e infiltrazione di software malevolo (\emph{malware} e \emph{botnet}).

\subsubsection{Interruzioni Di Servizio E Vulnerabilità Operative (Downtimes)}
La concentrazione di migliaia di carichi di lavoro su un'infrastruttura condivisa rende i guasti infrastrutturali particolarmente devastanti. A differenza dei sistemi \emph{on-premises}, dove un guasto hardware circoscritto impatta una singola applicazione aziendale, un malfunzionamento a livello di \emph{control plane}, di connettività o di alimentazione elettrica in un data center cloud provoca effetti a cascata su scala multinazionale.

% [IMMAGINE DEL FILE DI SLIDE CON NOME "CloudComputing"] Screenshot di rassegna stampa su storiche interruzioni di servizio: Microsoft Windows Azure, Rackspace e Amazon S3 Availability Event — Slide 22 / Pagina 1 del PDF

I casi storici documentati evidenziano le principali criticità operative:
\begin{itemize}
    \item \textbf{Guasti infrastrutturali fisici}: l'incidente che ha coinvolto il data center di Dallas del provider Rackspace ha dimostrato come la rottura di una singola connessione elettrica e il successivo fallimento dei generatori di emergenza possano causare la perdita prolungata di connettività, costringendo il fornitore a rimborsare milioni di dollari in crediti contrattuali di livello di servizio (\emph{Service Level Agreement}, SLA).
    \item \textbf{Propagazione di errori nel piano di controllo}: l'evento di disservizio che ha colpito Amazon S3 nel luglio 2008 ha evidenziato la fragilità dei sistemi distribuiti basati su algoritmi di consenso interno. Un accumulo di messaggi corotti all'interno del protocollo di propagazione dello stato ha saturato la capacità di elaborazione dei server di storage, innescando un collasso generale delle richieste API globali.
    \item \textbf{Interruzioni del software di orchestrazione}: i primi disservizi occorsi a piattaforme PaaS nascenti quali Microsoft Windows Azure hanno sottolineato come difetti logici nelle routine di aggiornamento automatico o nella gestione delle tabelle di routing virtuali possano rendere inaccessibili decine di migliaia di istanze contemporaneamente.
\end{itemize}

\subsubsection{Abuso Di Piattaforma: Phishing E Furto Di Credenziali}
Il modello di fruizione centralizzato via web rende il cloud sia un bersaglio sia un potenziale moltiplicatore di attacchi basati su ingegneria sociale. La centralizzazione delle identità digitali fa sì che la compromissione di un singolo set di credenziali garantisca l'accesso non solo alle caselle email, ma a interi panieri di servizi applicativi federati.

% [IMMAGINE DEL FILE DI SLIDE CON NOME "CloudComputing"] Esempi di attacchi phishing mirati: furto credenziali su Twitter e massiccia esfiltrazione di password Hotmail — Slide 23 / Pagina 2 del PDF

Le vulnerabilità evidenziate dai casi storici (come gli attacchi di massa tramite repliche fraudolente delle pagine di login di Twitter o la pubblicazione online di elenchi contenenti decine di migliaia di password del servizio Microsoft Hotmail) dimostrano che l'accesso alle risorse virtualizzate è vincolato alla robustezza del solo livello applicativo di autenticazione. Nei sistemi multi-tenant, un utente compromesso all'interno del medesimo tenant può fungere da cavallo di Troia per inoculare messaggi di phishing verso gli altri nodi della rete fidata interna.

\subsubsection{La Potenza Computazionale Come Arma: Password Cracking Su Risorse Elastiche}
La disponibilità su richiesta (\emph{on-demand}) di risorse computazionali massicce a costi minimi ha radicalmente modificato l'asimmetria tra difensore e attaccante. Prima dell'avvento dei cluster elastici su cloud pubblico, sferrare attacchi a forza bruta (\emph{brute-force}) contro algoritmi di hashing crittografico richiedeva l'acquisto e la manutenzione di costosi cluster fisici dedicati.

% [IMMAGINE DEL FILE DI SLIDE CON NOME "CloudComputing"] Articoli di settore (PCWorld, stacksmashing.net) sull'uso delle istanze GPU EC2 per il cracking di chiavi crittografiche WPA e password hash — Slide 24 / Pagina 3 del PDF

L'introduzione di istanze computazionali avanzate, dotate di processori grafici dedicati (\emph{Graphics Processing Units}, GPU) come le macchine \texttt{cg1.4xlarge} di Amazon EC2 (equipaggiate storicamente con schede NVIDIA Tesla), ha reso accessibile il calcolo parallelo massivo per pochi dollari all'ora. Un attaccante può istanziare temporaneamente decine di istanze GPU ed eseguire strumenti ottimizzati (come CUDA-Multiforce):
\begin{itemize}
    \item \textbf{Violazione di protocolli wireless}: algoritmi basati su chiavi pre-condivise come il WPA-PSK possono essere attaccati analizzando centinaia di migliaia di combinazioni al secondo per un costo di frazioni di dollaro, invalidando la sicurezza perimetrale degli standard legacy.
    \item \textbf{Inadeguatezza degli hash non rallentati}: funzioni di hash non progettate specificamente per resistere alla parallelizzazione massiva (come SHA-1 o MD5) cedono in poche decine di minuti a fronte di dizionari complessi. Questa realtà impone l'adozione stringente di funzioni di derivazione di chiave computazionalmente onerose e parametrizzabili nel tempo di calcolo e nell'uso di memoria, quali \textbf{bcrypt}, \textbf{scrypt} o \textbf{PBKDF2}.
\end{itemize}

\subsubsection{Infiltrazione Di Infrastrutture: Botnet E Command And Control Nel Cloud}
La neutralità e l'alta affidabilità fornite dai principali cloud provider sono state storicamente sfruttate anche da organizzazioni criminali per ospitare infrastrutture malevole ad alta resilienza.

% [IMMAGINE DEL FILE DI SLIDE CON NOME "CloudComputing"] Notizia della violazione del Treasury Department e schema operativo del trojan Zeus attestato su infrastruttura Amazon EC2 — Slide 25 / Pagina 4 del PDF

L'analisi dei vettori d'infezione storici mostra due pattern fondamentali:
\begin{enumerate}
    \item \textbf{Compromissione di carichi di lavoro governativi/aziendali}: nel caso dell'intrusione ai danni dei server del \emph{Bureau of Engraving and Printing} (BEP) del Dipartimento del Tesoro statunitense ospitati su fornitori terzi, gli attaccanti non hanno violato l'hardware sottostante ma hanno sfruttato vulnerabilità applicative per inserire snippet HTML malevoli all'interno dei siti legittimi. Il traffico dei cittadini veniva reindirizzato verso kit d'exploit remoti (come \emph{Eleonore Exploit Pack}), trasformando un portale istituzionale in un distributore di malware.
    \item \textbf{Ospitalità di server Command and Control (C\&C)}: varianti del famigerato trojan bancario \textbf{Zeus} hanno impiegato direttamente macchine virtuali su Amazon EC2 come nodi di ricezione dei dati esfiltrati e distribuzione di comandi via script PHP (\texttt{gate.php}). La reputazione positiva associata ai blocchi di indirizzi IP dei provider cloud rende estremamente difficile per i sistemi di Intrusion Detection perimetrali distinguere il traffico di una botnet da normali comunicazioni aziendali.
\end{enumerate}

\subsection{Sicurezza Della Virtualizzazione: Proprietà, Hypervisor E Snapshot}

La virtualizzazione rappresenta il fondamento abilitante dell'intero paradigma cloud. Essa consente l'esecuzione contemporanea e trasparente di molteplici sistemi operativi ospiti (\textbf{Guest OS}) su un'unica piattaforma hardware fisica sottostante, mediante l'interposizione di uno strato software o firmware denominato \textbf{Hypervisor} o \textbf{Virtual Machine Monitor} (VMM).

% [IMMAGINE DEL FILE DI SLIDE CON NOME "CloudComputing"] Architettura di virtualizzazione a strati: Hardware Layer, Host OS, Hypervisor, Host Applications e macchine virtuali isolate — Slide 27 / Pagina 6 del PDF

\subsubsection{Isolamento E Il Compromesso Architetturale}
Prima di valutare la robustezza della virtualizzazione, è opportuno rinfrescare brevemente la distinzione tra le principali architetture di esecuzione:
\begin{itemize}
    \item \textbf{Condivisione del sistema operativo (Sistemi Monolitici o Container)}: le applicazioni condividono il medesimo kernel del sistema operativo ospitante. L'isolamento logico è demandato a meccanismi propri del kernel (come \emph{namespaces} e \emph{control groups} in Linux). Qualsiasi violazione o \emph{kernel panic} compromette l'intero sistema.
    \item \textbf{Macchine Virtuali (Hypervisor-based)}: ogni carico di lavoro opera all'interno di un'istanza virtuale completa, provvista di un proprio kernel indipendente e di dispositivi hardware emulati o paravirtualizzati.
    \item \textbf{Bare-Metal segregato}: ogni applicazione risiede su un server fisico dedicato, disconnesso o fisicamente separato a livello di bus, memoria e canali I/O.
\end{itemize}

Dal punto di vista della sicurezza, l'isolamento offerto da una macchina virtuale si colloca esattamente a metà strada:
\begin{equation*}
    \text{Isolamento Multi-Processo (Stesso OS)} < \mathbf{\text{Isolamento VM (Hypervisor)}} < \text{Isolamento Fisico Bare-Metal}
\end{equation*}

Una VM offre un confinamento notevolmente superiore rispetto all'esecuzione di due processi distinti sullo stesso sistema operativo, poiché una violazione dei privilegi a livello di spazio utente (\emph{user space}) o all'interno del kernel guest rimane, in linea teorica, confinata entro la sandbox virtuale gestita dall'hypervisor.

Tuttavia, l'isolamento della VM è \textbf{inferiore} a quello offerto da due server fisicamente distinti. Le istanze virtuali continuano a condividere:
\begin{enumerate}
    \item Il silicio del processore fisico (\emph{core}, pipeline, branch predictor);
    \item I livelli di memoria cache condivisa (L2, L3/LLC);
    \item Il controller di memoria, i bus di sistema e i canali del controller di rete (NIC).
\end{enumerate}
Questa superficie fisica comune espone le macchine a sofisticati \textbf{canali laterali} (\emph{side-channel attacks}). Esempi storici e fondamentali sono rappresentati dalle vulnerabilità \textbf{Meltdown} e \textbf{Spectre}: manipolando l'esecuzione speculativa e l'accesso transitorio alle cache della CPU, un processo guest malevolo può misurare con precisione nanometrica i tempi di risposta della cache per ricostruire byte per byte segreti appartenenti all'hypervisor o ad altre macchine virtuali collocate sul medesimo socket fisico.

\subsubsection{Le Problematiche Dello Snapshot E Dello State Restore}
Una delle funzionalità più apprezzate della virtualizzazione è la possibilità di catturare lo stato istantaneo (\textbf{Snapshot}) della memoria volatile (RAM), dei registri della CPU e dell'immagine di disco di una macchina virtuale, per poi ripristinarlo (\textbf{State Restore}) in caso di errore o necessità di ripristino configurativo.

Tuttavia, tale meccanismo introduce un paradosso critico di sicurezza che ribalta l'efficacia delle tradizionali strategie di difesa:

% [IMMAGINE DEL FILE DI SLIDE CON NOME "CloudComputing"] Diagramma temporale del ciclo Snapshot e State Restore: reintroduzione nel sistema di virus eliminati e annullamento del ciclo di patching — Slide 28 / Pagina 7 del PDF

Il problema architetturale risiede nella natura non lineare del tempo nei sistemi virtualizzati. Si consideri la sequenza temporale schematizzata nelle lezioni:
\begin{enumerate}
    \item All'istante $t_0$, la macchina virtuale viene infettata da un artefatto malevolo (oppure presenta una vulnerabilità nel software di sistema). L'infezione resta silente.
    \item All'istante $t_1$, l'amministratore o l'utente salva lo stato della VM (creazione di uno Snapshot) per conservare uno stato applicativo coerente, ad esempio a seguito della modifica di un file di configurazione importante (File A).
    \item All'istante $t_2$, a causa di un'operazione accidentale, il File A viene cancellato.
    \item All'istante $t_3$, il sistema di sicurezza rileva l'infezione silente originata in $t_0$.
    \item All'istante $t_4$, il personale di sicurezza interviene: il malware viene rimosso, i file compromessi vengono ripuliti e il sistema viene aggiornato con le patch correttive del fornitore.
    \item Successivamente, all'istante $t_5$, un operatore necessita di recuperare il File A cancellato in $t_2$ e decide di eseguire un'operazione di \textbf{State Restore} tornando allo Snapshot catturato a $t_1$.
\end{enumerate}

Le conseguenze di questa operazione sono disastrose:
\begin{itemize}
    \item \textbf{Annullamento delle patch (\emph{Patch Undoing})}: il ripristino sovrascrive interamente l'immagine del disco e la memoria, eliminando le correzioni applicate all'istante $t_4$ e riportando il sistema operativo a uno stato vulnerabile.
    \item \textbf{Rigenerazione delle minacce}: il virus, fisicamente rimosso, viene re-iniettato direttamente in memoria o su disco dall'immagine di snapshot, tornando a essere attivo ed eludendo i controlli che ne consideravano conclusa la bonifica.
    \item \textbf{Persistenza a basso livello di worm endemici}: la proliferazione di snapshot obsoleti archiviati nei sistemi di storage fa sì che, nel corso dei mesi, l'avvio accidentale di vecchie immagini ri-immetta costantemente nell'infrastruttura di rete worm e botnet ritenuti debellati.
\end{itemize}

\subsubsection{Vulnerabilità Dell'Hypervisor E Fuga Dalla Sandbox (VM Escape)}
Poiché l'hypervisor controlla l'accesso alle risorse hardware reali e funge da arbitro assoluto tra le diverse macchine virtuali, qualsiasi falla logica al suo interno compromette irreversibilmente l'intero confine di sicurezza.

Le vulnerabilità a livello di VMM si classificano principalmente in:
\begin{enumerate}
    \item \textbf{Esecuzione di codice nel guest con privilegi elevati}: l'attaccante sfrutta errori nella paravirtualizzazione dei registri della CPU o nell'emulazione dell'architettura hardware per ottenere privilegi di amministratore (\emph{root} o \emph{ring 0}) all'interno del solo ambiente virtuale.
    \item \textbf{Fuga dalla macchina virtuale (\textbf{VM Escape})}: rappresenta l'evento più critico in assoluto nell'ambito della sicurezza del cloud. Un'applicazione malevola o un utente ostile, operando all'interno della macchina guest isolata, riesce a superare la barriera dell'emulazione e ad eseguire codice arbitrario direttamente nel contesto dell'hypervisor o del sistema operativo ospitante (\textbf{Host OS}). In questo scenario, l'attaccante ottiene il controllo simultaneo di tutte le altre macchine virtuali co-allocate sullo stesso server fisico.
\end{enumerate}

% [IMMAGINE DEL FILE DI SLIDE CON NOME "CloudComputing"] Bollettino di sicurezza VMware VMSA-2009-0006 ed elenchi di vulnerabilità CVE relative a hypervisor Xen e VirtualBox — Slide 29 / Pagina 8 del PDF

L'analisi dei bollettini di sicurezza e degli identificatori \textbf{CVE} (\emph{Common Vulnerabilities and Exposures}) fornisce evidenza empirica di questo fenomeno:
\begin{itemize}
    \item \textbf{VMware VMSA-2009-0006 (CVE-2009-1244)}: una falla critica nella funzione di visualizzazione a schermo (\emph{virtual machine display function}) dei prodotti VMware Workstation, Player ed ESX consentiva a un codice in esecuzione con privilegi all'interno del Guest OS di scatenare un buffer overflow ed eseguire codice arbitrario all'interno del contesto dell'Host OS.
    \item \textbf{Xen CVE-2008-1943}: difetti nell'interprete di emulazione dei backend di blocco o di rete che permettevano violazioni di memoria con conseguente blocco o compromissione dell'intero dominio di controllo privilegiato (\emph{Dom0}).
    \item \textbf{VirtualBox CVE-2010-3583}: errori nella gestione dei driver di condivisione e integrazione (\emph{Guest Additions}) che aprivano varchi diretti alla corruzione della memoria dell'applicazione host.
\end{itemize}

\subsection{Le Nuove Sfide Di Sicurezza Nell'Infrastruttura Cloud}
Il passaggio a un'architettura completamente virtualizzata introduce un ventaglio di problematiche gestionali e conformative che non possono essere risolte mediante i vecchi paradigmi di protezione perimetrale:

\begin{enumerate}
    \item \textbf{Imputabilità (\emph{Accountability})}: l'astrazione multilivello e l'alta frequenza di allocazione e deallocazione delle risorse rendono complessa l'attribuzione univoca delle azioni fraudolente o accidentali ai singoli soggetti giuridici o utenti applicativi.
    \item \textbf{Assenza di perimetro fisico (\emph{No Security Perimeter})}: nei data center tradizionali, la sicurezza si fondava su un perimetro fisico ben definito (firewall di frontiera, DMZ, cablaggi strutturati isolati). Nel cloud, le macchine virtuali condividono interfacce di rete fisiche e gli indirizzi IP possono essere allocati dinamicamente da pool globali; il controllo d'accesso deve pertanto essere ripensato e applicato rigorosamente su base \textbf{host-by-host}.

    % [IMMAGINE DEL FILE DI SLIDE CON NOME "CloudComputing"] Illustrazione concettuale dell'assenza di confini fisici tangibili nel paradigma Cloud Computing — Slide 31 / Pagina 10 del PDF

    \item \textbf{Superficie d'attacco estesa (\emph{Larger Attack Surface})}: oltre allo stack del sistema operativo guest e all'applicazione, la superficie di esposizione include ora l'intero strato dell'hypervisor, i driver di virtualizzazione e le API di orchestrazione esposte su rete pubblica per la gestione dell'infrastruttura.
    \item \textbf{Nuovi canali laterali (\emph{New Side Channels})}: la co-locazione di tenant concorrenti sulla stessa piattaforma hardware rende possibili attacchi basati sull'analisi temporale delle cache, sul consumo di potenza energetica o sul jitter di trasmissione di pacchetti su bus condivisi.
    \item \textbf{Carenza di auditabilità (\emph{Lack of Auditability})}: l'impossibilità, per il cliente, di ispezionare fisicamente i dispositivi di memorizzazione o analizzare il traffico di rete grezzo a livello di switch fisici limita profondamente la trasparenza forense in caso di incidenti.
    \item \textbf{Conformità normativa (\emph{Regulatory Compliance})}: legislazioni stringenti sulla protezione dei dati personali impongono vincoli rigidi sulla sovranità geografica dei dati e sull'accesso non autorizzato da parte di fornitori esteri.
\end{enumerate}

\newpage

\section{Piattaforme E Servizi Dei Principali Cloud Provider: AWS E Azure}

L'implementazione industriale dei concetti di virtualizzazione ha portato all'affermazione dei grandi fornitori pubblici di risorse computazionali. Tra questi, Amazon Web Services (AWS) e Microsoft Azure rappresentano le incarnazioni di riferimento, rispettivamente, dell'approccio infrastrutturale altamente granulare (\textbf{IaaS}) e dell'evoluzione incentrata sui servizi applicativi e sulle piattaforme gestite (\textbf{PaaS}).

% [IMMAGINE DEL FILE DI SLIDE CON NOME "CloudComputing"] Panoramica dei principali attori e servizi del Cloud Computing: EC2, S3, RDS, Azure e OpenStack — Slide 32 / Pagina 11 del PDF

\subsection{Amazon Elastic Compute Cloud (EC2): Calcolo Elastico E Risorse Virtualizzate}

Amazon Elastic Compute Cloud (Amazon EC2) è un servizio web che eroga capacità computazionale ridimensionabile (\emph{resizable compute capacity}) direttamente all'interno dei data center globali di Amazon. Il modello operativo elimina l'esigenza di pianificare investimenti preventivi per l'acquisizione di hardware, consentendo l'allocazione e il de-provisioning di server virtuali in pochi minuti, pagando esclusivamente per le risorse effettivamente consumate (\emph{pay-as-you-go}).

\subsubsection{Architettura Fondamentale: Amazon Machine Images (AMI) E Istanze}
L'astrazione fondamentale di EC2 poggia su due concetti distinti ma complementari: il modello statico dell'immagine e la sua esecuzione temporale dinamica.

\begin{theorem}[Amazon Machine Image (AMI)]
Un'Amazon Machine Image è un modello preconfigurato (\emph{template}) in sola lettura che incapsula la configurazione software completa necessaria per l'avvio di un sistema: il sistema operativo di base, il server applicativo, eventuali librerie intermedie e il codice dell'applicazione specifica.
\end{theorem}

% [IMMAGINE DEL FILE DI SLIDE CON NOME "CloudComputing"] Schema di derivazione: lancio di molteplici istanze indipendenti a partire da una singola AMI — Slide 34 / Pagina 13 del PDF

Da una singola AMI è possibile lanciare una o molteplici \textbf{istanze} (\emph{Instances}), le quali rappresentano le copie operative (\emph{running copies}) e attive dell'immagine template. Il ciclo di vita di un'istanza è delimitato dai comandi di arresto (\texttt{stop}), terminazione definitiva (\texttt{terminate}) o da un guasto fatale dell'hardware sottostante. In caso di avaria, l'infrastruttura garantisce il principio di riproducibilità: una nuova istanza identica può essere ricreata istantaneamente a partire dall'AMI originaria.

% [IMMAGINE DEL FILE DI SLIDE CON NOME "CloudComputing"] Differenziazione dell'archetipo hardware: istanziazione di Large, High-CPU o High-Memory da un template comune — Slide 35 / Pagina 14 del PDF

Un aspetto cruciale del disaccoppiamento logico risiede nella scelta dell'\textbf{Instance Type} (Tipo di Istanza). L'instance type rappresenta l'archetipo hardware virtuale assegnato all'istanza:
\begin{itemize}
    \item Definisce l'allocazione proporzionale di risorse fisiche: quantità di memoria RAM, numero di core virtuali della CPU (vCPU), tipologia di storage di transito e prestazioni dell'interfaccia di rete.
    \item A partire dalla medesima AMI è possibile avviare istanze di classi prestazionali radicalmente eterogenee (ad esempio nodi generici \emph{Large}, nodi ottimizzati per il calcolo parallelo \emph{High-CPU Extra Large}, oppure nodi per basi di dati in memoria \emph{High-Memory Double Extra Large}).
\end{itemize}

L'avvio di un'istanza può essere completamente automatizzato: oltre alle configurazioni statiche incluse nell'AMI pubblica o proprietaria, è possibile iniettare stringhe di configurazione o script di bootstrap (denominati storicamente \emph{User Data}). In fase di accensione, l'istanza esegue lo script, installa le versioni software aggiornate, scarica i sorgenti dinamici e si registra autonomamente nei bilanciatori di carico.

\subsubsection{Topologia Di Distribuzione: Regioni E Availability Zones}
Per ovviare ai limiti dei singoli data center e soddisfare requisiti stringenti di affidabilità e bassa latenza, l'infrastruttura globale di AWS è strutturata gerarchicamente:

% [IMMAGINE DEL FILE DI SLIDE CON NOME "CloudComputing"] Mappa architetturale della distribuzione geografica: Regioni AWS contenenti molteplici Availability Zones isolate — Slide 36 / Pagina 15 del PDF

\begin{itemize}
    \item \textbf{Regioni (\emph{Regions})}: aree geografiche completamente separate nel mondo (ad esempio Nord America, Europa o Asia). La scelta di allocazione in una determinata regione risponde a vincoli di latenza verso l'utenza finale, differenziali di costo e conformità legislativa sul trattamento del dato (\emph{data residency}). Le regioni sono interamente disaccoppiate tra loro: nessun dato o risorsa viene replicata al di fuori di essa senza un'esplicita istruzione dell'utente.
    \item \textbf{Zone di Disponibilità (\textbf{Availability Zones} - AZ)}: ogni regione racchiude molteplici zone distinte, isolate fisicamente e progettate per costituire domini di guasto indipendenti. Ciascuna AZ è allocata in data center fisicamente separati da distanze di sicurezza rispetto a disastri naturali o blackout, ma interconnessi mediante reti in fibra ottica dedicate a bassissima latenza e alto throughput.
\end{itemize}
La prassi ingegneristica per ottenere la tolleranza ai guasti consiste nel distribuire le istanze di una medesima applicazione su più Availability Zones entro la stessa regione: la caduta catastrofica di un intero data center non pregiudica la disponibilità operativa del servizio complessivo.

\subsection{Strategie Di Archiviazione: Instance Store, EBS E S3}

La gestione della memorizzazione nel cloud si articola su tre tecnologie primarie, differenziate per livello di astrazione, persistenza temporale e modello di consistenza.

% [IMMAGINE DEL FILE DI SLIDE CON NOME "CloudComputing"] Schema di fissaggio di un volume persistente di tipo Amazon EBS a un'istanza EC2 — Slide 38 / Pagina 17 del PDF

\subsubsection{Instance Store: Storage A Blocchi Effimero}
L'\textbf{Instance Store} è uno storage a livello di blocco fornito da dischi fisicamente integrati nel server fisico (host) su cui è allocata la macchina virtuale.
\begin{itemize}
    \item \textbf{Caratteristiche}: garantisce latenze di accesso ridotte e costi computazionali nulli o minimi.
    \item \textbf{Comportamento alla persistenza}: è uno storage \textbf{non persistente} (effimero). I dati scritti permangono esclusivamente finché l'istanza è in esecuzione; in caso di arresto forzato dell'istanza (\texttt{stop}) o di terminazione (\texttt{terminate}), o a seguito di un guasto hardware dell'host fisico, i dati residenti sull'instance store vengono cancellati in modo irreversibile. È indicato per memorie cache temporanee, code di calcolo volatili o carichi di lavoro replicati a livello applicativo.
\end{itemize}

\subsubsection{Amazon Elastic Block Store (EBS): Dischi Virtuali Persistenti}
Amazon EBS fornisce volumi di archiviazione a livello di blocco indipendenti dal ciclo di vita della singola macchina virtuale. Concettualmente assimilabili a un disco rigido virtuale o a una LUN di una Storage Area Network (SAN), i volumi EBS presentano caratteristiche ingegneristiche fondamentali:
\begin{itemize}
    \item \textbf{Persistenza}: i dati sopravvivono all'arresto e al riavvio dell'istanza.
    \item \textbf{Flessibilità di montaggio}: un volume può essere formattato con un file system tradizionale, utilizzato per database relazionali e disconnesso (\texttt{detach}) da un'istanza sorgente per essere immediatamente collegato (\texttt{attach}) a un'altra istanza all'interno della medesima AZ.
\end{itemize}

% [IMMAGINE DEL FILE DI SLIDE CON NOME "CloudComputing"] Diagramma di creazione di Snapshot EBS su bucket S3 e successivo provisioning verso un nuovo volume per altra istanza — Slide 39 / Pagina 18 del PDF

% [IMMAGINE DEL FILE DI SLIDE CON NOME "CloudComputing"] Meccanica operativa di distacco di un volume EBS da Instance A e successivo aggancio a Instance B — Slide 40 / Pagina 19 del PDF

Per garantire la sicurezza e il backup dei dati, è possibile creare \textbf{EBS Snapshot}. Si tratta di copie punto-nel-tempo (\emph{point-in-time}) a salvataggio incrementale archiviate automaticamente in Amazon S3. A partire da uno snapshot, è possibile generare nuovi volumi EBS identici per replicare dati di produzione in ambienti di test o per migrare dischi tra zone di disponibilità differenti.

\subsubsection{Amazon Simple Storage Service (S3): Storage Degli Oggetti}
Amazon S3 rappresenta la realizzazione industriale dell'\textbf{Object Storage}. A differenza dello storage a blocchi (EBS) o dello storage a file (NFS/POSIX), S3 memorizza il dato come un'entità opaca (\textbf{oggetto}) accompagnata da metadati e identificata univocamente da una chiave stringa arbitraria all'interno di un contenitore denominato \textbf{Bucket}.

Le caratteristiche operative e architetturali di S3 includono:
\begin{enumerate}
    \item \textbf{Granularità degli oggetti}: la dimensione di un singolo oggetto può spaziare da 1 byte fino a 5 terabyte.
    \item \textbf{Scalabilità infinita}: il numero di oggetti stoccabili all'interno di un bucket è illimitato; non vi è necessità di dimensionare file system o tabelle di allocazione.
    \item \textbf{Spazio dei nomi e accesso}: ogni oggetto è reperibile univocamente mediante un percorso globale strutturato nella forma:
    \begin{equation*}
        \texttt{http(s)://nome-bucket.s3.amazonaws.com/chiave-oggetto}
    \end{equation*}
    \item \textbf{Sovranità geografica del dato}: un bucket viene istanziato in una specifica regione. Gli oggetti al suo interno \textbf{non lasciano mai} la regione selezionata a meno che il proprietario non ne richieda esplicitamente il trasferimento programmato (garantendo il rispetto delle normative sulla residenza dei dati).
    \item \textbf{Sicurezza e controllo degli accessi}: S3 fornisce granularità d'accesso mediante Access Control List (ACL), Bucket Policies e integrazione con meccanismi di cifratura sia in transito (HTTPS/TLS) sia a riposo (\emph{Encryption at Rest} lato client o lato server tramite AWS KMS).
    \item \textbf{Interfacce standard}: il consumo delle API avviene secondo i paradigmi standard REST e SOAP, consentendo l'integrazione immediata tramite qualsiasi client HTTP o SDK applicativo.
    \item \textbf{Ottimizzazione dei costi di trasferimento dati}: l'interscambio di enormi dataset tra Amazon S3 e le istanze di calcolo EC2 allocate all'interno della medesima regione geografica avviene senza incorrere in alcun costo di trasferimento dati (\emph{zero data transfer fees}), rendendo S3 il repository ideale per pipeline di analisi massiva (\emph{Big Data Analytics}).
\end{enumerate}

\subsection{Gestione Dei Database Nel Cloud: Self-Hosted Vs Amazon RDS}

Nel momento in cui un'applicazione distribuita su EC2 richiede la persistenza su database relazionale, l'architetto del cloud può adottare due strategie alternative:

\begin{enumerate}
    \item \textbf{Database Self-Hosted su EC2}: si acquisisce un'istanza di calcolo generica (avviata tramite un'AMI opportuna), si monta un volume EBS persistente e si installa manualmente il motore di database (es. MySQL, PostgreSQL o Oracle). Questa soluzione garantisce il pieno controllo sui file di configurazione e sui parametri del kernel del sistema operativo, ma scarica integralmente sull'utente l'onere delle attività di amministrazione (\emph{DBA tasks}): applicazione delle patch del software, gestione della frammentazione, rotazione dei log, predisposizione della replica e monitoraggio dei dischi.
    \item \textbf{Amazon Relational Database Service (RDS)}: rappresenta un servizio completamente gestito (\emph{Managed Service}). Il provider astrae il sistema operativo sottostante e si fa carico dell'automazione delle routine manutentive: esecuzione di backup continui con ripristino point-in-time, aggiornamento automatico delle patch di sicurezza del motore, allocazione dello storage elastico e configurazione in alta affidabilità multi-AZ con failover trasparente.
\end{enumerate}

\subsection{Networking E Sicurezza Perimetrale In AWS}

Come esposto nell'analisi delle criticità della virtualizzazione, l'assenza di confini fisici tangibili impone che l'isolamento e la segregazione del traffico siano demandati a costrutti logici di rete implementati a livello di hypervisor.

\subsubsection{Indirizzamento IP Pubblico Ed Elastic IP}
Al momento del provisioning all'interno della rete EC2, a ciascuna istanza viene assegnato per impostazione predefinita un indirizzo IP pubblico dinamico. Se l'istanza viene riavviata o subisce un guasto e viene rigenerata, il suo indirizzo IP pubblico cambia inevitabilmente, rompendo le correlazioni DNS o le configurazioni hardcoded dei client esterni.

Per risolvere questa criticità, AWS adotta il meccanismo degli \textbf{Elastic IP (EIP)}. Un Elastic IP è un indirizzo IPv4 pubblico statico riservato per l'account utente che può essere riassociato istantaneamente e a caldo da un'istanza all'altra (ad esempio reindirizzando il traffico da un'istanza guasta a un nodo di riserva passivo), mascherando i fallimenti infrastrutturali al client esterno.

\subsubsection{I Security Group Come Firewall Stateful}
Il principale strumento di segregazione a livello di istanza è rappresentato dai \textbf{Security Group} (Gruppi di Sicurezza). Un Security Group funge da firewall virtuale operante a livello di interfaccia di rete virtuale (vNIC).

Le proprietà cardine di un Security Group sono:
\begin{itemize}
    \item \textbf{Controllo del traffico in ingresso (\emph{Inbound}) e in uscita (\emph{Outbound})}: per impostazione predefinita, tutto il traffico in ingresso è negato e deve essere esplicitamente autorizzato tramite regole dichiarative (whitelist).
    \item \textbf{Natura Stateful}: i Security Group sono firewall \textbf{con stato} (\emph{stateful}). Se viene aperta una regola che ammette una connessione TCP su una specifica porta in ingresso, il pacchetto di risposta generato dall'istanza per quella medesima sessione viene autorizzato a uscire automaticamente, a prescindere dalle regole configurate per il traffico outbound.
    \item \textbf{Riferimento ad altri Security Group}: la caratteristica architetturale più potente consiste nella facoltà di specificare, come sorgente di una regola di traffico, non solo una notazione di sottorete CIDR (es. \texttt{192.168.1.0/24}), ma l'\textbf{ID di un altro Security Group}. In questo modo, l'accesso viene autorizzato unicamente alle istanze dinamicamente registrate in quel gruppo, a prescindere dal loro effettivo indirizzo IP.
\end{itemize}

\subsubsection{Implementazione Dell'Architettura A Tre Livelli (Three-Tier)}
La composizione dinamica dei Security Group trova la sua applicazione classica nell'implementazione sicura dell'architettura web a tre livelli (\textbf{Three-Tier Architecture}), composta dal livello web, dal livello applicativo e dal livello di memorizzazione/database:

% [IMMAGINE DEL FILE DI SLIDE CON NOME "CloudComputing"] Diagramma architetturale di una Three-Tier Web Application protetta da regole concatenate di Security Group — Slide 24 / Pagina 24 del PDF

La matrice di isolamento mostrata nelle lezioni si articola secondo le seguenti direttive di flusso:

\begin{enumerate}
    \item \textbf{Livello Web (Web Servers Tier)}:
    \begin{itemize}
        \item \emph{Ingresso Pubblico}: ammette traffico proveniente dall'intero Internet pubblico (\texttt{0.0.0.0/0}) unicamente sulle porte \textbf{TCP 80 (HTTP)} e \textbf{TCP 443 (HTTPS)}.
        \item \emph{Gestione e Manutenzione}: la porta di amministrazione \textbf{TCP 22 (SSH)} non è aperta verso l'esterno, ma consente l'accesso unicamente alle connessioni provenienti dal Security Group degli Application Server (oppure da un host bastion autorizzato).
    \end{itemize}
    \item \textbf{Livello Applicativo (App Servers Tier)}:
    \begin{itemize}
        \item \emph{Ingresso Funzionale}: non è accessibile direttamente dall'Internet pubblico. Le connessioni applicative sono autorizzate \textbf{esclusivamente} se la sorgente appartiene al \emph{Security Group dei Web Servers}.
        \item \emph{Gestione e Manutenzione}: ammette connessioni sulla porta \textbf{TCP 22 (SSH)} provenienti rigorosamente ed esclusivamente dal blocco di indirizzi IP associato alla \emph{subnet aziendale corporate}. Gli ingegneri e gli operatori effettuano l'accesso SSH primario su questo livello, e da qui possono attestarsi a cascata sulle altre macchine virtuali della farm.
    \end{itemize}
    \item \textbf{Livello Dati (DB Servers Tier)}:
    \begin{itemize}
        \item \emph{Ingresso Funzionale}: la regola di sicurezza del database accetta comunicazioni (es. su porta MySQL 3306 o Oracle 1521) \textbf{unicamente ed esclusivamente} se generate da istanze appartenenti al \emph{Security Group degli App Servers}.
        \item Qualsiasi tentativo di contatto diretto proveniente dai web server o dalla rete Internet pubblica viene scartato silenziosamente a monte dall'hypervisor.
    \end{itemize}
\end{enumerate}

Questo approccio a catena realizza il principio del \textbf{privilegio minimo} a livello di rete, impedendo che l'eventuale compromissione di un server web esposto su Internet possa tradursi in una violazione diretta e immediata del database dei dati di produzione.

\subsection{La Piattaforma Microsoft Azure}

In parallelo all'affermazione del modello IaaS di Amazon, l'offerta cloud di Microsoft ha debuttato storicamente come \textbf{Azure Services Platform} (successivamente evoluta in Microsoft Azure), orientandosi in origine principalmente verso il paradigma della piattaforma come servizio (\textbf{Platform-as-a-Service} - PaaS).

% [IMMAGINE DEL FILE DI SLIDE CON NOME "CloudComputing"] Prospetto del modello architetturale iniziale dell'Azure Services Platform e integrazione con lo stack applicativo Microsoft — Slide 26 / Pagina 26 del PDF

\subsubsection{Principi E Benefici Operativi}
L'approccio iniziale di Azure poggiava sulla fornitura di un ambiente d'esecuzione per applicazioni personalizzate compilate su framework proprietari (in particolare l'ecosistema \textbf{.NET}), ospitate su macchine virtuali preconfigurate e manutenute interamente da Microsoft. I capisaldi di questo modello includono:
\begin{itemize}
    \item \textbf{Virtualizzazione trasparente}: la gestione del ciclo di vita dell'hypervisor e del sistema operativo è totalmente opaca per lo sviluppatore, che deve occuparsi esclusivamente della logica applicativa.
    \item \textbf{Scalabilità automatica}: allocazione dinamica dei processi e delle istanze in base all'evoluzione del carico di richieste.
    \item \textbf{Utility Computing}: imputazione dei costi su consumo effettivo di risorse computazionali e storage a supporto del software.
\end{itemize}

\subsubsection{Componenti Chiave Dell'Architettura Azure Services Platform}
L'ecosistema iniziale dell'offerta Microsoft si articolava in moduli infrastrutturali e servizi applicativi federati:

\begin{itemize}
    \item \textbf{Windows Azure (Cloud Services OS)}: rappresenta lo strato del sistema operativo cloud di base. Fornisce l'ambiente di hosting per i servizi, il file system distribuito, lo storage scalabile e il coordinamento dell'infrastruttura. Il coordinamento delle risorse è demandato storicamente al \textbf{Fabric Controller}, un controller software distribuito responsabile di monitorare lo stato di salute di tutti i server fisici, ridispiegare automaticamente i nodi applicativi in caso di guasto hardware e coordinare gli aggiornamenti di versione del software senza interruzione del servizio.
    \item \textbf{Microsoft .NET Services}: una suite di servizi di alto livello per accelerare lo sviluppo distribuito, comprendente l'\emph{Access Control} (gestione dell'autenticazione basata su claim federati), strumenti di gestione dei flussi (\emph{Workflow}) e un componente fondamentale: il \textbf{Hosted Service Bus}. Il Service Bus consente la comunicazione fluida e sicura (\emph{loosely-coupled}) tra applicazioni in esecuzione \emph{on-premises} e servizi migrati nel cloud, superando agevolmente problematiche di NAT e configurazioni restrittive dei firewall perimetrali aziendali.
    \item \textbf{Microsoft SQL Services / SQL Data Services (SDS)}: l'estensione cloud del motore relazionale Microsoft SQL Server. Concepito originariamente per offrire database accessibili direttamente via internet tramite protocolli web e query REST/SOAP, si è poi strutturato come motore relazionale cloud in grado di memorizzare e analizzare dati strutturati, semi-strutturati e non strutturati, gestendo autonomamente sincronizzazione, backup e reporting.
    \item \textbf{Microsoft SharePoint Services \& CRM Online}: l'offerta di componenti applicativi aziendali fruibili come servizio, garantendo la federazione delle identità tra i server fisici locali del cliente e il cloud Microsoft lungo l'intero spettro \emph{server-service continuum}.
    \item \textbf{Live Services e Live Framework}: un'interfaccia basata su API per connettere dispositivi eterogenei e applicazioni consumer (come l'ecosistema federato Windows Live e Office Live), orchestrando la condivisione dello stato utente e la sincronizzazione di dati personali attraverso una rete di centinaia di milioni di profili.
\end{itemize}

Questa impostazione ha segnato la convergenza tra infrastruttura elastica di calcolo, orchestrazione dei componenti software e convergenza dei servizi per l'impresa e il consumatore finale, ponendo le basi per i moderni modelli architetturali del cloud ibrido.

% AA -------------------------------------------------------------------------------

\section{Infrastrutture Cloud Enterprise Ed Open-Source: Panoramica Di Microsoft Azure Ed OpenStack}

L'evoluzione dei paradigmi di computazione distribuita ha portato alla convergenza verso modelli di erogazione delle risorse basati su virtualizzazione massiva, orchestrazione automatizzata e astrazione dell'hardware fisico sottostante. In questo contesto, le organizzazioni si trovano a valutare soluzioni eterogenee che spaziano dai servizi di cloud pubblico proprietario a piattaforme di orchestrazione infrastrutturale open-source per cloud privato ed ibrido.

\subsection{L'Architettura Di Microsoft Azure}

Microsoft Azure rappresenta una delle principali piattaforme di \emph{public cloud hyperscale}. La sua architettura di base è concepita per operare su vasta scala all'interno di \emph{data center} proprietari distribuiti globalmente e gestiti direttamente da Microsoft. Il principio ingegneristico cardine di tali infrastrutture risiede nell'impiego estensivo di hardware convenzionale (\emph{commodity hardware}): anziché affidarsi a mainframe verticali ad altissimo costo, l'infrastruttura aggrega cluster sterminati di server commerciali standard x86-64, delegando la resilienza, la tolleranza ai guasti e il bilanciamento del carico al livello software dell'ipervisore e del sistema operativo di controllo distribuito (\emph{Fabric Controller}).

Dal punto di vista dello sviluppo applicativo, Azure supera la tradizionale impronta incentrata esclusivamente sull'ecosistema Microsoft. Sebbene offra un supporto nativo ottimizzato per il framework .NET e linguaggi quali C\# e Visual Basic, la piattaforma garantisce interoperabilità completa verso ambienti a codice nativo non gestito (\emph{unmanaged code}, come C e C++) e runtime open-source largamente diffusi, tra cui Java, PHP, Node.js e Python, unitamente a framework di comunicazione enterprise quali Windows Communication Foundation (WCF) e ASP.NET.

Un pilastro fondamentale dell'infrastruttura Azure è costituito dalla diversificazione delle opzioni di memorizzazione persistente, progettate per soddisfare requisiti contrastanti di consistenza, latenza e complessità del modello dati:
\begin{itemize}[noitemsep]
    \item Windows Azure Storage (Blob e strutture semplici): orientato al salvataggio non strutturato ad altissima scalabilità. Il servizio per oggetti binari di grandi dimensioni (BLOB -- \emph{Binary Large Objects}) è ottimizzato per ospitare flussi video, immagini ad alta risoluzione, backup e file di log. L'accesso avviene tramite primitive stateless conformi all'architettura REST (\emph{Representational State Transfer}) su protocollo HTTP/HTTPS, garantendo un disaccoppiamento netto tra il client consumatore e il sottosistema di storage.
    \item SQL Azure Database: implementa il paradigma relazionale tradizionale transazionale, offrendo la piena compatibilità con il linguaggio SQL e le garanzie ACID (\emph{Atomicity, Consistency, Isolation, Durability}). Il servizio astrae completamente l'onere di configurazione del motore di base di dati, gestendo la replica geografica e l'alta disponibilità in modalità Platform as a Service (PaaS).
    \item Connettività distribuita: la piattaforma integra componenti di rete e middleware dedicati (quali Service Bus e reti ibride) per abilitare la comunicazione sicura e performante tra applicazioni distribuite residenti sia all'interno del perimetro del cloud, sia \emph{on-premises} in centri di calcolo aziendali remoti.
\end{itemize}

\subsection{La Piattaforma OpenStack Per L'Infrastruttura Come Servizio}

Parallelamente ai fornitori di cloud pubblico, l'esigenza di realizzare architetture Infrastructure as a Service (IaaS) all'interno di data center privati ha portato alla standardizzazione industriale di OpenStack.

OpenStack si configura come un vero e proprio sistema operativo distribuito per il cloud (\emph{cloud operating system}) il cui scopo primario consiste nell'aggregare, governare e orchestrare grandi pool di risorse fisiche eterogenee di calcolo (\emph{compute}), archiviazione (\emph{storage}) e connettività di rete (\emph{networking}) distribuite all'interno dell'intero data center.

% [IMMAGINE DEL FILE DI SLIDE CON NOME "CloudComputing"] Schema architetturale a blocchi di OpenStack: interazione tra Dashboard, Compute, Networking, Storage, Shared Services e Standard Hardware — Slide 47 / Pagina 3 del PDF

L'architettura logica di OpenStack è rigorosamente modulare e a microservizi:
\begin{enumerate}[noitemsep]
    \item Interfaccia Utente e Dashboard (Horizon): fornisce una console web centralizzata che consente agli amministratori di monitorare lo stato di salute dei nodi fisici e definire quote di utilizzo, permettendo al contempo agli utenti finali di allocare e gestire autonomamente macchine virtuali, indirizzi IP pubblici e volumi disco. Tutte le funzionalità esposte dalla GUI si appoggiano direttamente a un insieme completo di API RESTful.
    \item Modulo di Calcolo (Nova): si occupa del ciclo di vita delle istanze virtuali, interfacciandosi con differenti ipervisori (principalmente KVM/QEMU) per pianificare e istanziare le macchine virtuali sui nodi fisici tramite algoritmi di \emph{scheduling} avanzati.
    \item Modulo di Rete (Neutron): implementa il paradigma Software-Defined Networking (SDN), astraendo la topologia di rete fisica per fornire reti virtuali private multi-tenant isolate, switch virtuali, router software e funzionalità di incapsulamento avanzate (come VXLAN e GRE).
    \item Moduli di Memorizzazione (Cinder e Swift): forniscono rispettivamente storage a blocchi persistente e rimovibile da collegare alle istanze virtuali (Cinder) e storage a oggetti scalabile e ridondato su scala geografica (Swift).
    \item Servizi Condivisi e Astrazione Hardware: al di sotto dei motori operativi operano servizi trasversali condivisi, tra cui il componente di autenticazione e controllo degli accessi basato su ruoli (Keystone) e il catalogo delle immagini disco dei sistemi operativi (Glance). L'intero stack opera al di sopra di hardware standard convenzionale, garantendo l'assoluta indipendenza dai fornitori di apparati proprietari (\emph{vendor lock-in}).
\end{enumerate}

\newpage

\section{Progettazione Di Applicazioni Cloud-Native: Principi E Tassonomia}

L'approccio ingegneristico volto alla realizzazione di sistemi per il cloud non può limitarsi alla mera trasposizione (\emph{lift-and-shift}) di applicativi tradizionali su macchine virtuali remote. Progettare in modalità Cloud-Native richiede la piena consapevolezza delle proprietà intrinseche delle piattaforme distribuite, massimizzando i vantaggi di elasticità ed isolamento e affrontando in modo proattivo le complessità relative alla latenza di rete, ai fallimenti parziali e alla concorrenza distribuita.

\subsection{Transizione Dal Monolite Ai Microservizi Distribuiti}

Nelle architetture monolitiche convenzionali, l'intera logica di business, lo strato di accesso ai dati e l'interfaccia utente risiedono all'interno del medesimo processo operativo e condividono un unico schema di persistenza. Tale conformazione presenta evidenti limiti di manutenibilità: una modifica introdotta in un singolo modulo richiede la ricompilazione e la ridistribuzione dell'intero applicativo, mentre la scalabilità può avvenire unicamente in senso verticale oppure replicando l'intero monolite su vasta scala, con un massiccio spreco di risorse computazionali.

Nelle applicazioni cloud-native, il sistema viene sistematicamente decomposto in servizi autonomi, decentralizzati e di dimensioni ridotte (\emph{microservices}). Ciascun servizio incapsula un dominio applicativo specifico ed è dotato di un ciclo di vita indipendente. I microservizi cooperano scambiandosi messaggi attraverso due canali architetturali principali:
\begin{itemize}[noitemsep]
    \item Interfacce Applicative Sincrone (API): basate prevalentemente sul paradigma RESTful sopra protocollo HTTP/HTTPS o su chiamate procedurali binarie ad alte prestazioni come gRPC, utilizzate per interazioni di tipo richiesta-risposta a bassa latenza.
    \item Messaggistica Asincrona e Gestione Degli Eventi: mediata da broker di messaggi distribuiti (es. Apache Kafka, RabbitMQ). Tale meccanismo disaccoppia temporalmente e spazialmente il produttore dal consumatore dell'informazione; i servizi reagiscono al manifestarsi di eventi di dominio senza richiedere il blocco dei flussi di esecuzione.
\end{itemize}

Questo disaccoppiamento consente la scalabilità orizzontale granulare: il sistema scala verso l'esterno (\emph{scale-out}) allocando dinamicamente nuove repliche e istanze dei soli microservizi che registrano un incremento temporaneo di carico, riducendole (\emph{scale-in}) quando la domanda decresce.

\subsection{Modelli Di Servizio E Tassonomia Dello Storage}

L'ambiente cloud categorizza la potenza di calcolo secondo livelli crescenti di astrazione:
\begin{itemize}[noitemsep]
    \item Infrastructure as a Service (IaaS): il fornitore mette a disposizione risorse fisiche virtualizzate (macchine virtuali, interfacce di rete, spazio disco grezzo). L'utente ha il pieno controllo del sistema operativo, del runtime e delle configurazioni software.
    \item Platform as a Service (PaaS): il livello del sistema operativo, delle patch di sicurezza e del framework di runtime è interamente gestito dalla piattaforma cloud. Lo sviluppatore vi distribuisce unicamente il codice sorgente o il contenitore applicativo.
    \item Software as a Service (SaaS): l'intero stack applicativo è fornito come servizio pronto all'uso, accessibile via browser o client dedicato, liberando completamente l'utente da oneri infrastrutturali.
    \item Function as a Service (FaaS / Serverless): il paradigma di massima astrazione computazionale. Il codice è incapsulato in singole funzioni atomiche e \emph{stateless}, la cui esecuzione è scatenata da specifici eventi esterni (invocazioni HTTP, modifiche sul database, messaggi in coda). L'allocazione della memoria e della CPU scala istantaneamente da zero al fabbisogno reale, fatturando la frazione esatta di millisecondi di elaborazione attiva.
\end{itemize}

In maniera speculare, la memorizzazione dei dati deve riflettere la natura eterogenea del dominio applicativo. Non esiste un singolo motore di memorizzazione universale; la scelta progettuale deve ponderare fattori quali il throughput, il modello relazionale, la latenza di lettura, la cardinalità e i vincoli di consistenza (\emph{Polyglot Persistence}):
\begin{itemize}[noitemsep]
    \item Basi di dati relazionali (RDBMS): basate su schemi rigidi e transazioni ACID, ideali per dati finanziari e transazionali complessi.
    \item Archivi Chiave-Valore (Key-Value Stores): strutture dati associative ad altissimo throughput e minima latenza, perfette per memorizzare sessioni e profili utente.
    \item Basi di dati documentali (Document Databases): memorizzano i dati sotto forma di gerarchie semistrutturate (es. JSON o BSON), ideali per cataloghi prodotti ed e-commerce con schemi evolutivi.
    \item Basi di dati a grafo (Graph Databases): focalizzate sulla persistenza efficiente di nodi e relazioni, impiegate per reti sociali e motori di raccomandazione.
    \item Basi di dati colonnari (Wide-Column Stores): orientate all'elaborazione analitica su vasta scala (\emph{Big Data}) e all'interrogazione aggregata su enormi volumi.
    \item Motori di ricerca testuale (Search Engine Databases): architetture basate su indici invertiti, progettate per la ricerca full-text ad alta velocità.
    \item Basi di dati per serie temporali (Time Series Databases): ottimizzate per flussi di inserimenti continui con timestamp, indispensabili per telemetria, IoT e monitoraggio infrastrutturale.
    \item File System Condivisi (Shared Files): repository di storage di rete conformi a protocolli quali NFS o SMB per l'accesso concorrente a file non strutturati.
\end{itemize}

\subsection{I Quattro Principi Architetturali Delle Applicazioni Cloud-Native}

Affinché un sistema distribuito possa definirsi resiliente, governabile e scalabile, la progettazione deve aderire rigorosamente a quattro direttrici fondamentali:

\begin{enumerate}
    \item Progettazione orientata all'autoriparazione (Design for Self-Healing):
    In un'infrastruttura distribuita su larga scala, il fallimento di singoli componenti (nodi di calcolo, link di rete, unità di memoria) non rappresenta un'anomalia, bensì un evento statisticamente ordinario. Il sistema deve integrare meccanismi attivi in grado di rilevare i guasti in tempo reale (mediante sonde di \emph{liveness} e \emph{readiness}), isolare i nodi degradati, reindirizzare il traffico di rete e provvedere automaticamente alla reinizializzazione o sostituzione delle istanze danneggiate, ristabilendo lo stato nominale di funzionamento senza alcun intervento manuale.

    \item Ridondanza globale (Make Everything Redundant):
    Ogni elemento architetturale che costituisce il percorso critico del dato o dell'elaborazione deve essere duplicato fisicamente o logicamente al fine di eliminare categoricamente ogni singolo punto di fallimento (SPOF -- \emph{Single Point of Failure}). Ciò implica la replicazione geografica dei nodi di calcolo attraverso differenti \emph{Availability Zones}, la configurazione di data store in replica master-replica o cluster multi-master, e la ridondanza degli apparati di bilanciamento del carico.

    \item Minimizzazione del coordinamento tra servizi (Minimize Coordination):
    In un contesto distribuito, l'esecuzione di protocolli di sincronizzazione bloccanti o di transazioni distribuite tradizionali (quali il \emph{Two-Phase Commit}, 2PC) introduce un collo di bottiglia critico che degrada le prestazioni e vanifica i tentativi di scalabilità. È necessario ridurre al minimo l'accoppiamento temporale, privilegiando comunicazioni asincrone mediante code ed eventi, accettando modelli di consistenza eventuale (\emph{Eventual Consistency}) che consentono ai singoli sottosistemi di progredire nella loro computazione in modo autonomo e concorrente.

    \item Progettazione per la scalabilità orizzontale (Design for Horizontal Scaling):
    I microservizi devono essere concepiti secondo il paradigma \emph{stateless}: le istanze operative non devono conservare lo stato delle sessioni client nella propria memoria locale. La rimozione o l'aggiunta di nuove repliche a runtime deve risultare totalmente trasparente agli utenti. In questo modo l'infrastruttura può contrarre o espandere dinamicamente il numero di istanze in esecuzione al variare istantaneo del carico di lavoro (\emph{horizontal auto-scaling}).
\end{enumerate}

\newpage

\section{Pattern Architetturali Per La Gestione Dei Dati Nel Cloud}

La persistenza e la manipolazione dei dati costituiscono la sfida più complessa all'interno dei sistemi distribuiti nel cloud, condizionando direttamente parametri quali latenza, throughput, consistenza ed affidabilità complessiva.

\subsection{Le Sfide Nella Gestione Distribuita Del Dato}

Nelle architetture moderne, i dati risiedono fisicamente dislocati su server multipli e in regioni geograficamente distanti per ragioni di vicinanza all'utente finale, conformità normativa, tolleranza ai disastri naturali e scalabilità. Tale frammentazione impone sfide severe:
\begin{itemize}[noitemsep]
    \item Mantenimento della consistenza: garantire che viste differenti del dato convergano al medesimo stato informativo a fronte di scritture concorrenti distribuite, gestendo i compromessi formulati dal teorema CAP (Consistenza, Disponibilità, Tolleranza alle partizioni).
    \item Sincronizzazione distribuita: propagare gli aggiornamenti attraverso canali di comunicazione soggetti a latenze variabili, ritrasmissioni e partizioni di rete.
    \item Sicurezza e conformità (Triade CIA): i dati devono essere protetti mediante crittografia sia quando residenti nei supporti di memorizzazione (at rest), sia durante l'attraversamento delle reti fisiche e virtuali (in transit), regolamentando l'accesso tramite meccanismi di autenticazione e autorizzazione capillari al fine di garantire riservatezza (\emph{Confidentiality}), integrità (\emph{Integrity}) e disponibilità (\emph{Availability}).
\end{itemize}

\subsection{Tassonomia Dei Pattern Di Data Management}

Per affrontare in modo sistematico le problematiche sopra descritte, l'ingegneria del software per il cloud ha formalizzato una serie di pattern architetturali standardizzati, riassunti nella Tabella~\ref{tab:data_patterns}.

\begin{table}[H]
\centering
\small
\renewcommand{\arraystretch}{1.3}
\begin{tabularx}{\textwidth}{l p{3.8cm} X}
\toprule
\textbf{Pattern} & \textbf{Obiettivo Fondamentale} & \textbf{Descrizione Operativa} \\
\midrule
\textbf{Cache-Aside} & Ottimizzazione delle letture e scarico del database principale & Carica i dati nella memoria cache su richiesta (\emph{on-demand}) a partire dal database sottostante al primo accesso utile. \\
\textbf{CQRS} & Segregazione funzionale tra lettura e scrittura & Separa le operazioni di lettura dei dati da quelle di modifica/aggiornamento strutturando interfacce, modelli e database dedicati. \\
\textbf{Event Sourcing} & Tracciabilità assoluta e consistenza storica & Memorizza l'evoluzione dello stato del dominio come sequenza immutabile di eventi di sola aggiunta (\emph{append-only}). \\
\textbf{Index Table} & Accelerazione delle query secondarie & Definisce tabelle di indici ausiliarie su campi frequentemente interrogati in data store che non supportano indici secondari nativi. \\
\textbf{Materialized View} & Efficienza su query complesse e aggregazioni & Genera e mantiene viste prepopolate del dato ottimizzate per specifici pattern di interrogazione quando lo schema originale non è adeguato. \\
\textbf{Sharding} & Superamento dei limiti fisici di storage e calcolo & Suddivide orizzontalmente un unico data store logico in una pluralità di partizioni fisiche autonome denominate \emph{shards}. \\
\textbf{Static Content Hosting} & Riduzione del carico di calcolo e abbattimento latenza & Distribuisce risorse statiche (HTML, immagini, script) direttamente su storage a oggetti cloud e reti CDN dedicate. \\
\textbf{Valet Key} & Scarico del traffico di I/O dal server applicativo & Fornisce al client un token di accesso limitato e a tempo che consente il trasferimento diretto e sicuro del file con lo storage cloud. \\
\bottomrule
\end{tabularx}
\caption{Quadro sinottico dei Cloud Design Pattern per la gestione dei dati.}
\label{tab:data_patterns}
\end{table}

\newpage

\section{Il Pattern Cache-Aside (Lazy Loading)}

Il pattern Cache-Aside (noto anche come caricamento pigro o \emph{lazy loading}) rappresenta la soluzione di caching più diffusa per mitigare le problematiche di latenza di lettura ed evitare il sovraccarico computazionale sul database primario.

\subsection{Principio Di Funzionamento E Flusso Operativo}

Nel modello Cache-Aside, la cache è posizionata concettualmente "a fianco" (\emph{kept aside}) rispetto al flusso di elaborazione: il sottosistema di cache non interagisce mai direttamente con il database. È onere esclusivo dell'applicazione gestire in modo esplicito la lettura e la scrittura verso entrambi i componenti.

% [IMMAGINE DEL FILE DI SLIDE CON NOME "CloudComputing"] Diagramma architetturale del pattern Cache-Aside e sequence diagram delle interazioni tra User, App, Cache e Database — Slide 54-55 / Pagine 10-12 del PDF

Il ciclo operativo di gestione di una richiesta di lettura si articola nelle seguenti fasi sequenziali:
\begin{enumerate}
    \item Richiesta del dato: l'applicazione riceve una query di lettura (ad esempio i dettagli di un utente o un catalogo prodotti).
    \item Verifica in Cache: l'applicazione interroga prioritariamente la memoria cache ad alta velocità (es. Redis o Memcached).
    \item Condizione di Hit (Cache Hit): qualora la chiave richiesta sia presente in memoria, il dato viene immediatamente restituito al chiamante, evitando l'accesso al database e completando l'operazione in tempi dell'ordine del sub-millisecondo.
    \item Condizione di Miss (Cache Miss): se il dato non è presente, l'applicazione procede al recupero dell'informazione interrogando il database persistente primario.
    \item Aggiornamento della Cache: l'applicazione provvede a scrivere il valore prelevato dal database all'interno della cache, associandovi una specifica politica temporale. Le successive interrogazioni per la medesima chiave risulteranno pertanto in un \emph{Cache Hit}.
    \item Politica di Scadenza (\emph{Time To Live} - TTL): per scongiurare il rischio che informazioni obsolete rimangano indefinitamente all'interno della cache a fronte di modifiche concorrenti sul database, ogni record viene marcato con un intervallo di scadenza temporale (TTL). Alla scadenza del timer, la chiave viene invalidata ed eliminata.
\end{enumerate}

\subsection{Confronto Con Le Strategie Read-Through, Write-Through E Write-Behind}

Per comprendere a fondo i compromessi del Cache-Aside, è opportuno formalizzare le differenze architetturali che intercorrono rispetto ai sistemi di cache trasparenti (spesso integrati in soluzioni enterprise come Oracle Coherence):

\begin{itemize}
    \item Read-Through: l'applicazione considera la cache come il proprio unico punto di contatto per le letture. In caso di assenza del dato (\emph{miss}), non è il codice applicativo a rivolgersi al database, bensì il gestore interno della cache (\emph{CacheStore}) che delega il caricamento trasparente dal database sottostante, inserisce la voce in memoria e la recapita all'applicazione.
    \item Write-Through: quando l'applicazione aggiorna o inserisce un dato, scrive unicamente sulla cache. È quest'ultima che si fa carico di aggiornare in modo sincrono il database primario. L'operazione di salvataggio (es. `put(...)`) non ritorna l'esito positivo finché la scrittura su disco nel database sottostante non è andata a buon fine. Questo approccio garantisce forte coerenza ma non apporta alcun beneficio prestazionale alle operazioni di scrittura, le quali scontano comunque la latenza della persistenza fisica.
    \item Write-Behind (o Write-Back): l'applicazione esegue la scrittura unicamente nella cache, la quale rilascia immediatamente il controllo al thread chiamante. La persistenza verso il database sottostante avviene in modalità completamente asincrona: le voci modificate vengono inserite in una coda gestita internamente (\emph{write-behind queue}) e riversate verso la base di dati dopo un ritardo configurabile (ad esempio da pochi secondi fino a diversi minuti). Tale architettura accelera drasticamente il throughput di scrittura e consente l'aggregazione di modifiche ripetute sul medesimo record, ma espone al rischio di perdita di dati in caso di crash improvviso del nodo di cache prima del flush su disco.
\end{itemize}

Nel pattern Cache-Aside standard, qualora un'applicazione modifichi un dato, può applicare una variazione del Write-Through modificando direttamente il database relazionale e invalidando contestualmente (cancellando) la voce corrispondente dalla cache. La prima richiesta di lettura successiva subirà un fisiologico miss, riallineando la cache con il valore appena aggiornato.

\subsection{Contesti Di Inapplicabilità: La Gestione Dello Stato Di Sessione Nelle Web Farm}

Sebbene il pattern Cache-Aside si riveli eccellente per richieste a carico imprevedibile, esso presenta specifici vincoli di applicabilità:
\begin{itemize}[noitemsep]
    \item Dati statici invarianti: se l'insieme dei dati è limitato, immutabile e allocabile interamente nello spazio di memoria disponibile, la strategia ideale consiste nel precaricare la cache all'avvio (\emph{cache priming}) disattivando la scadenza temporale.
    \item Stato di sessione in cluster di server Web (Web Farm): memorizzare le informazioni di sessione applicativa dei client nelle memorie locali dei singoli nodi applicativi mediante Cache-Aside introduce anomalie sistemiche severe.
\end{itemize}

In un'architettura web bilanciata orizzontalmente, affidare lo stato di sessione alla cache locale di ciascun web server comporta le seguenti criticità:
\begin{itemize}[noitemsep]
    \item Inconsistenza del dato di sessione: se le richieste successive del medesimo utente sono instradate a nodi differenti dal bilanciatore, il server di destinazione risulterà privo della sessione memorizzata localmente dal predecessore, provocando errori di autenticazione o svuotamenti del carrello elettronico.
    \item Complessità di scalabilità ed invalidazione: sincronizzare la cache di sessione su un cluster di centinaia di server web genera un traffico broadcast ingestibile che satura la banda e degrada la memoria di ciascuna macchina.
\end{itemize}

Per aggirare tale problema, storicamente si è fatto ricorso alla Client-Server Affinity (nota comunemente come Sticky Sessions), configurando il bilanciatore di carico per forzare tutte le richieste originate da un determinato utente sempre sul medesimo nodo fisico di backend. Tuttavia, questa dipendenza costituisce un anti-pattern critico nelle architetture cloud-native:
\begin{enumerate}[noitemsep]
    \item Single Point of Failure (SPOF): se il server su cui è confinato l'utente sperimenta un guasto, la sessione viene irrimediabilmente perduta, anche qualora tutti gli altri nodi della farm siano operativi.
    \item Bilanciamento del carico inefficace: gruppi eterogenei di utenti particolarmente attivi legati al medesimo server possono saturarne le risorse, lasciando gli altri server del cluster sostanzialmente inattivi.
    \item Complessità di manutenzione: l'esecuzione di procedure di rolling update o il riavvio controllato di un nodo impone la brusca disconnessione di tutti gli utenti vincolati a quella specifica istanza.
\end{enumerate}

Nelle moderne architetture cloud si impone pertanto l'approccio puramente stateless per lo strato di elaborazione web: le sessioni vengono delocalizzate su cluster di database in-memory centralizzati (es. Redis Cluster) oppure codificate direttamente all'interno di token crittografici auto-consistenti gestiti dal client, quali i JSON Web Tokens (JWT).

\subsection{Implementazione Dettagliata In Python}

Il seguente listato illustra l'implementazione del pattern Cache-Aside mediante il client Redis in Python, evidenziando la gestione delle condizioni di hit e miss e la configurazione del parametro TTL.

\begin{lstlisting}[language=python]
def get_user_data(user_id):
    # 1. Verifica preliminare della presenza del dato in cache
    user = redis_cache.get(user_id)
    if user:
        print("Cache Hit! Dato recuperato istantaneamente dalla memoria.")
        return user

    # 2. Condizione di Cache Miss: accesso alla persistenza relazionale
    print("Cache Miss! Interrogazione sul database relazionale...")
    user = database.query(f"SELECT * FROM users WHERE id={user_id}")

    # 3. Scrittura manuale del dato in cache per le richieste future
    # Impostazione di un Time To Live (TTL) pari a 3600 secondi (1 ora)
    redis_cache.set(user_id, user, expire=3600)

    return user
\end{lstlisting}

\newpage

\section{Il Pattern CQRS (Command And Query Responsibility Segregation)}

Nel panorama della progettazione software enterprise, il pattern CQRS (Command and Query Responsibility Segregation) segrega esplicitamente le operazioni di modifica da quelle di lettura, superando i vincoli architetturali tipici dei modelli dati unificati.

\subsection{Limiti Del Modello Unificato E Motivazioni Architetturali}

Nelle architetture tradizionali, il medesimo modello concettuale del dato e la stessa rappresentazione tabellare su database vengono impiegati indistintamente per effettuare tutte le operazioni di tipo CRUD (\emph{Create, Read, Update, Delete}). Nelle applicazioni a complessità crescente, tuttavia, i requisiti posti a carico delle letture e delle scritture divergono in modo inconciliabile:
\begin{itemize}[noitemsep]
    \item La scrittura esige rigidi controlli transazionali, validazione di regole di business complesse, normalizzazione relazionale per scongiurare anomalie e garanzie di consistenza immediata.
    \item La lettura richiede invece aggregazioni eterogenee, velocità di risposta istantanea, proiezioni parziali e formati denormalizzati strutturati su misura per i componenti dell'interfaccia utente.
\end{itemize}

L'impiego di un modello unico costringe a compromessi penalizzanti, introducendo costose operazioni di \emph{JOIN} su schemi altamente normalizzati, conflitti di lock sul database e la manipolazione di oggetti applicativi eccessivamente pesanti rispetto alle reali necessità di visualizzazione.

Il pattern CQRS risolve tale problematica dividendo il sistema in due rami computazionali indipendenti:
\begin{enumerate}
    \item Asimmetria e Scalabilità Indipendente (\emph{Independent Scalability}):
    In pressoché tutti i moderni sistemi distribuiti e web-based, il volume delle operazioni di lettura supera quello delle scritture di diversi ordini di grandezza (frequente rapporto di 1000 a 1). Con CQRS è possibile scalare orizzontalmente lo stack di lettura allocando repliche del servizio e istanze distribuite di sola lettura (\emph{Read Replicas}), mantenendo al contempo il ramo di scrittura confinato a un numero ridotto di istanze altamente performanti.

    \item Persistenza Poliglotta Ottimizzata (\emph{Polyglot Persistence}):
    Ciascun ramo architetturale può impiegare il motore di persistenza più confacente alle proprie necessità. Il lato dei comandi può avvalersi di un database relazionale conforme allo standard SQL per garantire rigore transazionale e consistenza, mentre il lato delle query può sincronizzarsi verso archivi NoSQL a documenti, cache in-memory o motori di indicizzazione dedicati (quali Elasticsearch) per velocizzare ricerche complesse e full-text.

    \item Ottimizzazione Verso La UI Tramite DTO:
    Il ramo di lettura estrae direttamente strutture dati denormalizzate pronte all'uso, note come Data Transfer Objects (DTO), eliminando le trasformazioni dati a runtime e le serializzazioni onerose sul server.
\end{enumerate}

\subsection{I Due Pilastri: Comandi E Query}

% [IMMAGINE DEL FILE DI SLIDE CON NOME "CloudComputing"] Schema logico della segregazione tra ramo di comando (Validation, Commands, Domain Logic, Write Data Store) e ramo di interrogazione (Read Data Store, Read Model, Queries/DTO) — Slide 60 / Pagina 18 del PDF

L'architettura stabilisce una netta differenziazione semantica tra le due tipologie di interazione:
\begin{itemize}
    \item I Comandi (Ramo di Scrittura):
    Un comando è un'istruzione imperativa che esprime inequivocabilmente l'intento dell'utente di alterare lo stato del sistema (es. \texttt{CreateOrder}, \texttt{CancelReservation}). I comandi incorporano la logica di business, le regole di validazione formale e applicativa e agiscono su un'entità di consistenza atomica definita in Domain-Driven Design (DDD) come aggregato. Essi completano la transazione e non ritornano dati di dominio, ma restituiscono esclusivamente una conferma di successo o il codice dell'errore riscontrato. Il modello di scrittura riflette fedelmente i processi aziendali.

    \item Le Query (Ramo di Lettura):
    Una query è una richiesta esplicita di estrazione di informazioni. Essa interroga il sistema per ottenere uno stato senza apportare alcuna alterazione allo stato dei dati (\emph{operazione priva di effetti collaterali}). Le query scavalcano completamente la complessa logica di dominio e gli strati di validazione, interrogando direttamente il modello di lettura ed estraendo proiezioni sintetiche per il client. I modelli di lettura possono essere distrutti e ricreati arbitrariamente a partire dallo storico delle scritture senza rischio di perdita di informazione.
\end{itemize}

Uno svantaggio architetturale da considerare risiede nell'impossibilità di impiegare strumenti automatici di generazione del codice basati su ponte relazionale-oggetto (ORM -- \emph{Object-Relational Mapping}) mediante \emph{scaffolding}: non essendovi più una corrispondenza biunivoca 1:1 tra tabelle normalizzate del database e classi del dominio, la logica di estrazione e persistenza va sviluppata esplicitamente.

\subsection{Ruolo Ed Efficienza Dei Data Transfer Objects (DTO)}

All'interno di un'architettura distribuita a microservizi, la comunicazione remota tra processi (tramite invocazioni REST, gRPC o middleware di rete) comporta un significativo costo computazionale dovuto al tempo di andata e ritorno dei pacchetti (RTT -- \emph{Round-Trip Time}).

Un Data Transfer Object (DTO) è un pattern strutturale che risponde a questa problematica agendo da contenitore dati leggero (\emph{Plain Old Data Structure} o PODS). Le sue proprietà ingegneristiche fondamentali comprendono:
\begin{itemize}[noitemsep]
    \item Assenza di logica di business: il DTO è privo di comportamenti o metodi operativi; ospita unicamente campi, proprietà e attributi serializzabili.
    \item Aggregazione e abbattimento delle chiamate di rete: invece di effettuare chiamate RPC multiple verso diversi endpoint per raccogliere i dati di un cliente, i suoi ordini e il bilancio contabile, il server popola un unico DTO aggregato restituito tramite un'unica trasmissione di rete.
    \item Serializzazione efficiente: i campi sono nativamente formattati per favorire la conversione verso formati standard da trasmettere su rete (JSON, XML o buffer binari Protocol Buffers).
    \item Disaccoppiamento dei contratti di interfaccia: impedisce la fuga di dettagli implementativi dello schema interno del database verso l'esterno; modifiche alle tabelle interne non impattano l'interfaccia esposta al client.
\end{itemize}

\subsection{Sincronizzazione A Eventi E Problematiche Di Consistenza}

% [IMMAGINE DEL FILE DI SLIDE CON NOME "CloudComputing"] Flusso completo CQRS distribuito da AWS Prescriptive Guidance: Command API, Write DB, Event Trigger, Read DB, Query API — Slide 61 / Pagina 21 del PDF

La dislocazione di database fisicamente separati per il ramo di lettura e di scrittura impone una pipeline di allineamento continuo, articolata nelle seguenti sette fasi operative:
\begin{enumerate}[noitemsep]
    \item Il client invia il comando di modifica invocando le API di scrittura.
    \item Lo strato applicativo elabora il comando eseguendo autorizzazione, validazione dei vincoli e computazione del dominio.
    \item Lo stato modificato viene salvato e reso persistente nel database di scrittura.
    \item A fronte del commit su disco, il sistema di scrittura pubblica un evento di notifica (es. inviandolo a un broker di messaggi distribuito).
    \item Il componente di sincronizzazione del ramo di lettura consuma l'evento, calcola la proiezione denormalizzata e persiste la vista materializzata nel database ottimizzato per le query.
    \item Il client o la UI inoltra una richiesta di lettura alle API dedicate del ramo di interrogazione.
    \item Il query engine estrae il dato dal database di lettura, restituendo il DTO al chiamante.
\end{enumerate}

% [IMMAGINE DEL FILE DI SLIDE CON NOME "CloudComputing"] Dettaglio della separazione fisica dei dati con vista materializzata e sincronizzazione guidata dagli eventi — Slide 62 / Pagina 22 del PDF

Poiché il database di scrittura e il message broker di coordinamento non possono essere inglobati in un'unica transazione distribuita atomica senza sacrificare drasticamente le prestazioni, l'architettura abbandona la consistenza forte a favore della consistenza eventuale (\emph{Eventual Consistency}): il ramo di lettura risulterà allineato con un lieve ritardo temporale rispetto al momento della scrittura. Per prevenire la corruzione dei dati a fronte di trasmissioni duplicate di messaggi di rete causate da fallimenti parziali, il sottosistema di ricezione degli eventi deve tassativamente implementare il principio di idempotenza (\emph{idempotent message processing}), garantendo che l'applicazione ripetuta del medesimo evento produca il medesimo stato finale.

\subsection{Criteri Di Applicabilità E Controindicazioni}

L'adozione di CQRS deve derivare da un'analisi comparativa tra la complessità architetturale introdotta e i benefici architetturali conseguiti. I contesti d'impiego raccomandati includono:
\begin{itemize}[noitemsep]
    \item Domini altamente collaborativi: scenari in cui una moltitudine di utenti accede e modifica concorrentemente gli stessi dati. La granularità specifica dei singoli comandi minimizza l'insorgenza di conflitti di unione (\emph{merge conflicts}).
    \item Interfacce utente orientate al task (Task-Based UI): flussi di lavoro complessi strutturati a passi, in cui il caricamento di una vista necessita di aggregazioni avanzate mentre l'input dell'utente mappa esattamente un'azione di dominio.
    \item Asimmetria strutturale del traffico: contesti in cui le interrogazioni superano enormemente le transazioni di modifica e impongono l'indipendenza delle metriche di scalabilità e throughput.
    \item Specializzazione dei team di sviluppo: permette a un team di doministi di concentrarsi sulla complessa logica dell'aggregato di scrittura, delegando a un altro team l'ottimizzazione del front-end e delle viste veloci di lettura.
    \item Integrazione con Event Sourcing: architetture in cui lo stato non è sovrascritto ma calcolato dalla riproduzione sequenziale degli eventi storici.
\end{itemize}

Al contrario, CQRS si configura come un anti-pattern quando:
\begin{itemize}[noitemsep]
    \item Le regole di business dell'applicazione sono semplici o minimali.
    \item L'interfaccia utente è interamente incentrata su modelli CRUD standard con mappatura biunivoca su tabelle relazionali.
\end{itemize}

\subsection{Implementazione Dettagliata In Python Di Un'Architettura CQRS}

Di seguito viene presentata l'implementazione completa dei due rami operativi in linguaggio Python, evidenziando le tecniche di validazione del comando, la prevenzione da attacchi di tipo SQL Injection e la gestione della lettura a doppio stadio con memoria cache e replica di lettura.

\subsubsection{Ramo Dei Comandi (Write Model)}

Nel modello di scrittura, il comando trasporta l'intento operativo mentre l'handler esegue le verifiche di business e l'inserimento parametrizzato nel database relazionale:

\begin{lstlisting}[language=python]
# Modello di Comando (DTO di Scrittura): incapsula l'intento di acquisto
class CreateOrderCommand:
    def __init__(self, order_id: str, product_name: str, quantity: int):
        self.order_id = order_id
        self.product_name = product_name
        self.quantity = quantity

# Gestore del Comando: incapsula le regole di business e la persistenza
class OrderCommandHandler:
    def handle(self, command: CreateOrderCommand):
        # 1. Validazione delle regole di business
        if command.quantity < 1:
            raise ValueError("Errore di validazione: la quantita minima ordinabile e 1.")

        # 2. Persistenza sicura su Database Relazionale Transazionale
        # L'uso di parametri posizionali (%s) previene vulnerabilita di SQL Injection:
        # la query viene compilata preventivamente e i parametri non subiscono concatenazione.
        write_db.execute(
            "INSERT INTO orders (id, product, qty) VALUES (%s, %s, %s)",
            (command.order_id, command.product_name, command.quantity)
        )

        # 3. Notifica asincrona per allineamento del Read Model
        print(f"Comando completato: Ordine {command.order_id} registrato con successo.")
\end{lstlisting}

\subsubsection{Ramo Delle Query (Read Model)}

Il modello di lettura utilizza una query atomica e implementa una strategia ottimizzata a due livelli (prima verifica su cache in-memory, seguita dal fallback su replica di lettura denormalizzata):

\begin{lstlisting}[language=python]
# Modello di Query: definisce i parametri di estrazione
class GetOrderDetailsQuery:
    def __init__(self, order_id: str):
        self.order_id = order_id

# Gestore della Query: esegue l'interrogazione ottimizzata per la UI
class OrderQueryHandler:
    def handle(self, query: GetOrderDetailsQuery):
        # 1. Percorso rapido: verifica della vista precalcolata in cache
        order_data = redis_cache.get(f"view:order:{query.order_id}")
        if order_data:
            return order_data

        # 2. Percorso lento: recupero dalla vista materializzata su Read-Replica DB
        return read_db.fetch_one(
            "SELECT * FROM orders_view WHERE id = %s",
            (query.order_id,)
        )
\end{lstlisting}

\subsubsection{Flusso Di Esecuzione Client}

Il codice client illustra la netta separazione delle pipeline di esecuzione:

\begin{lstlisting}[language=python]
# --- ESECUZIONE DEL COMANDO (Ramo di Scrittura) ---
create_cmd = CreateOrderCommand(
    order_id="ORD-99",
    product_name="Cloud Architecture Book",
    quantity=1
)
command_service = OrderCommandHandler()
command_service.handle(create_cmd)

# --- ESECUZIONE DELLA QUERY (Ramo di Lettura) ---
get_query = GetOrderDetailsQuery(order_id="ORD-99")
query_service = OrderQueryHandler()
result = query_service.handle(get_query)

print(f"Risultato aggregato per la UI: {result}")
\end{lstlisting}

% AA -------------------------------------------------------------------------------

\section{Pattern Architetturali Per La Gestione Dei Dati Nel Cloud}

Nelle moderne architetture distribuite e orientate ai servizi cloud, la persistenza e l'accesso ai dati presentano vincoli e sfide prestazionali radicalmente differenti rispetto ai sistemi monolitici basati su database relazionali tradizionali (\emph{Relational Database Management System}, RDBMS). Mentre in un contesto centralizzato l'integrità referenziale, le transazioni conformi alle proprietà ACID (\emph{Atomicity, Consistency, Isolation, Durability}) e le operazioni dirette di creazione, lettura, aggiornamento e cancellazione (\emph{Create, Read, Update, Delete} -- CRUD) risultano consolidate, nel paradigma cloud-native emergono problematiche legate alla concorrenza massiva, alla scalabilità orizzontale, alla latenza geografica e alla divergenza tra i requisiti di memorizzazione e quelli di consultazione.

Per risolvere tali criticità, l'ingegneria del software per il cloud adotta specifici \textbf{Cloud Design Pattern}, ovvero soluzioni architetturali formalizzate e riusabili volte a garantire scalabilità, resilienza, disaccoppiamento ed efficienza computazionale. I quattro pattern cardine esaminati in questa trattazione sono:
\begin{enumerate}[noitemsep]
    \item \textbf{Event Sourcing Pattern}: che ridefinisce il concetto di persistenza sostituendo il salvataggio dello stato attuale con il tracciamento cronologico e immutabile degli eventi di business;
    \item \textbf{Index Table Pattern}: che sopperisce all'assenza di indici secondari nativi negli archivi distribuiti NoSQL mediante la costruzione di tabelle di indicizzazione dedicate;
    \item \textbf{Materialized View Pattern}: che disaccoppia la memorizzazione dei dati dalla loro interrogazione, generando viste precalcolate e ottimizzate per specifiche query;
    \item \textbf{Sharding Pattern}: che affronta i limiti fisici delle singole macchine frammentando orizzontalmente gli archivi di dati su più nodi indipendenti.
\end{enumerate}

\newpage

% ==============================================================================
% SEZIONE 2: EVENT SOURCING
% ==============================================================================

\section{Pattern Event Sourcing}

\subsection{Limiti Del Modello CRUD Tradizionale E Motivazioni Architetturali}

La maggior parte delle applicazioni tradizionali memorizza le informazioni aggiornando direttamente lo stato corrente delle entità nel database. In un'architettura basata sul modello CRUD, ogni modifica comporta la sovrascrittura distruttiva del record precedente: un'operazione di tipo \texttt{UPDATE} cancella lo stato antecedente per far posto a quello nuovo, fornendo una visione puntuale assimilabile a un singolo fotogramma (\emph{freeze frame}).

Tale paradigma mostra profonde limitazioni architetturali quando applicato a sistemi distribuiti, collaborativi e ad alta concorrenza:
\begin{itemize}
    \item \textbf{Conflitti di aggiornamento concorrente (\emph{Data Update Conflicts})}: nei domini collaborativi in cui più utenti o servizi concorrenti tentano di aggiornare simultaneamente la medesima entità di dati, le operazioni si concentrano su un singolo record fisico. Ciò richiede meccanismi di blocco (\emph{pessimistic locking}) o controlli di versione (\emph{optimistic concurrency control}), i quali degradano drasticamente il throughput, causano frequenti aborti transazionali o generano fenomeni di contesa (\emph{lock contention}).
    \item \textbf{Perdita della storia e dell'intento di business (\emph{Loss of History})}: sovrascrivere lo stato cancella la sequenza causale che ha condotto a tale configurazione. A meno di implementare complessi e separati meccanismi di auditing secondari, il database perde la motivazione semantica del cambiamento: non è possibile determinare a posteriori \emph{perché}, \emph{come} o \emph{attraverso quali stadi} un attributo ha assunto il suo valore finale.
\end{itemize}

\begin{theorem}
L'\textbf{Event Sourcing} è un design pattern architetturale in cui lo stato di un'applicazione non viene persistito memorizzando la sua fotografia corrente, bensì mediante una sequenza ordinata e immutabile di eventi di dominio accaduti nel tempo, registrata all'interno di un archivio ad accodamento esclusivo (\emph{append-only store}).
\end{theorem}

In questo paradigma, lo stato attuale non è un dato primario salvato in una tabella relazionale, ma un \emph{valore derivato}: per conoscere la configurazione corrente di un'entità, l'applicazione preleva l'intera cronistoria degli eventi a essa associati e li riesegue sequenzialmente a partire dall'istante iniziale.

\subsection{Principi Fondamentali E Struttura Dell'Event Store}

L'elemento architetturale centrale è l'\textbf{Event Store}, un database transazionale specializzato e intrinsecamente ad accodamento (\emph{append-only}), all'interno del quale le operazioni di modifica (\texttt{UPDATE}) e cancellazione fisica (\texttt{DELETE}) sono architetturalmente inibite. Tra le tecnologie adoperate per l'implementazione figurano sia database nativamente concepiti per l'event sourcing (es. \emph{EventStoreDB}) sia broker di messaggistica distribuiti ad alta persistenza basati su log distribuito (es. \emph{Apache Kafka}).

L'archivio degli eventi contiene esclusivamente record strutturati secondo tre componenti fondamentali, schematizzati nella Tabella~\ref{tab:event_store_components}.

\begin{table}[H]
\centering
\small
\caption{Componenti strutturali elementari registrati all'interno di un Event Store.}
\label{tab:event_store_components}
\begin{tabularx}{\textwidth}{l p{5cm} X}
\toprule
\textbf{Componente} & \textbf{Descrizione} & \textbf{Esempio Pratico} \\
\midrule
\textbf{Events} & Record immutabili formulati al passato che descrivono azioni di business irreversibilmente concluse nel dominio applicativo. & \texttt{OrderCreated}, \texttt{ItemAddedToOrder}, \texttt{PaymentProcessed} \\
\addlinespace
\textbf{Metadata} & Informazioni contestuali necessarie all'instradamento, all'ordinamento causale e alla gestione distribuita dell'evento. & Aggregate ID (identificativo entità), Sequence Number (progressivo univoco), Timestamp, Event Type \\
\addlinespace
\textbf{Payload} & Dati specifici del dominio trasportati dall'evento, serializzati tipicamente in formato JSON o protocol buffer. & \texttt{\{"orderId": "ORD-99", "product": "Book", "quantity": 1\}} \\
\bottomrule
\end{tabularx}
\end{table}

\subsection{Dinamica Operativa: Ricostruzione Dello Stato E Disaccoppiamento}

Il codice applicativo non manipola direttamente tabelle relazionali mediante istruzioni SQL DML, ma formula \emph{eventi} che descrivono in modo imperativo e inequivocabile ogni transizione di stato avvenuta nel sistema (es. \texttt{AddedItemToOrder}).

% [IMMAGINE DEL FILE DI SLIDE CON NOME "slides_lezione.pdf"] Schema architetturale del pattern Event Sourcing: generazione eventi dal carrello (Presentation -> Cart created -> Item 1 added -> Item 2 added -> Item 1 removed -> Shipping info added), persistenza nell'Event Store, proiezione asincrona verso Materialized View e consumo da parte di sistemi esterni o query — Slide 68 / Pagina 2 del PDF

La persistenza della sequenza di eventi garantisce proprietà cruciali per l'ingegneria del software distribuiti:
\begin{enumerate}[noitemsep]
    \item \textbf{Ricostruzione degli stati passati (\emph{Time Travel})}: poiché l'archivio conserva la totalità delle transizioni, è possibile calcolare l'esatto stato del dominio in qualsiasi istante temporale $t_{\text{past}}$ semplicemente interrompendo il playback degli eventi a tale marcatore temporale.
    \item \textbf{Gestione di variazioni retroattive}: è possibile inserire correzioni o ricalcolare retroattivamente le consistenze (ad esempio per storni contabili o rettifiche fiscali) senza dover alterare i dati storici, registrando eventi compensativi.
    \item \textbf{Fonte autoritativa unica (\emph{System of Record})}: l'Event Store costituisce la sorgente primaria e immutabile della verità aziendale.
\end{enumerate}

L'Event Store agisce inoltre come un componente di pubblicazione: ogni nuovo evento aggiunto al registro viene pubblicato verso l'esterno tramite meccanismi di sottoscrizione. I consumatori (\emph{subscribers}) possono reagire in modo asincrono:
\begin{itemize}
    \item Notificando sistemi esterni o attivando workflow transazionali secondari;
    \item Aggiornando viste di lettura denormalizzate (\emph{Materialized Views}), utilizzate per alimentare le interfacce utente senza gravare sulla catena principale di scrittura.
\end{itemize}
Ciò garantisce un completo disaccoppiamento architetturale: il componente che emette l'evento non possiede alcuna dipendenza o conoscenza dei sistemi che lo elaboreranno a valle. La materializzazione delle entità può avvenire:
\begin{itemize}
    \item \textbf{On-Demand}: quando l'applicazione riceve una richiesta utente, carica l'aggregato ricostruendo dinamicamente lo stato corrente tramite il re-play degli eventi associati;
    \item \textbf{Scheduled / Asynchronous Task}: un processo pianificato elabora costantemente il flusso degli eventi e aggiorna una vista materializzata permanente a supporto del livello di presentazione.
\end{itemize}

\subsection{Caso Di Studio: Sistema Di Gestione Delle Prenotazioni Per Conferenze}

Per chiarire il vantaggio competitivo dell'Event Sourcing rispetto all'approccio CRUD convenzionale, si consideri un sistema di gestione conferenze (\emph{Conference Management System}) incaricato di monitorare le prenotazioni confermate e determinare la disponibilità residua di posti per gli acquirenti.

Come principio architetturale:
\begin{quote}
\emph{L'Event Sourcing non è finalizzato unicamente a ``sapere quanti posti rimangono'', bensì a modellare e guidare le decisioni di business nel tempo sulla base della cronologia degli eventi.}
\end{quote}

Il problema della capienza residua può essere indirizzato secondo due modalità strutturali:
\begin{enumerate}
    \item \textbf{Approccio con entità contatore (CRUD tradizionale)}: il database include un'entità che registra il numero aggregato di prenotazioni. A fronte di ogni acquisto o disdetta, il sistema esegue un incremento o decremento (\texttt{UPDATE bookings SET total = total + 1}). Sebbene concettualmente semplice, questa soluzione presenta un severo collo di bottiglia di scalabilità: a ridosso della chiusura delle iscrizioni, migliaia di utenti concorrenti tentano di accedere e bloccare in scrittura la medesima riga di database, generando massicci fallimenti per timeout o deadlock.
    \item \textbf{Approccio Event Sourcing}: prenotazioni e cancellazioni vengono registrate esclusivamente come eventi immutabili di append (\texttt{SeatReserved}, \texttt{SeatReservationCancelled}). Il calcolo dei posti disponibili è demandato al replay degli eventi. L'operazione di scrittura richiede unicamente l'accodamento dell'evento, massimizzando il throughput della memorizzazione senza conflitti di blocco.
\end{enumerate}

% [IMMAGINE DEL FILE DI SLIDE CON NOME "slides_lezione.pdf"] Diagramma di sequenza e blocchi funzionali per la prenotazione posti nel sistema per conferenze: Command Handler, aggregato SeatAvailability, Event Store con flussi di prenotazione e cancellazione — Slide 69 / Pagina 4 del PDF

Il flusso sequenziale operativo per la prenotazione di due posti si articola nei seguenti passaggi:
\begin{enumerate}[noitemsep]
    \item L'interfaccia utente (UI) emette un comando di prenotazione (\emph{Command}). Il comando viene intercettato da un \textbf{Command Handler}, un modulo software disaccoppiato dall'interfaccia preposto alla validazione preliminare.
    \item Viene istanziato l'aggregato di dominio, denominato \textbf{SeatAvailability}, interrogando l'Event Store per recuperare e riprodurre tutti gli eventi storici di prenotazione e cancellazione inerenti alla conferenza selezionata.
    \item \textbf{Ottimizzazioni prestazionali}: per evitare il replay completo di milioni di eventi dall'inizio dei tempi, si adottano le \textbf{Snapshot} (ovvero punti di ripristino periodici che congelano lo stato aggregato a una determinata versione $N$, consentendo di rieseguire solo gli eventi da $N+1$ a oggi) e il mantenimento in memoria volatile (\emph{in-memory cache}) delle istanze degli aggregati caldi.
    \item Il Command Handler invoca il metodo di business del modello di dominio (\texttt{SeatAvailability.reserve(2)}), verificando che la capienza massima non sia stata violata.
    \item L'aggregato emette e registra internamente un nuovo evento contenente il numero di posti allocati.
    \item Il sistema accoda fisicamente il nuovo evento nell'Event Store.
\end{enumerate}
In caso di disdetta, si esegue la medesima catena: il Command Handler genera un evento di cancellazione (\texttt{SeatCancelled}), che viene accodato in coda all'archivio, garantendo una tracciabilità di audit completa e infallibile.

\subsection{Integrazione Tra Event Sourcing E CQRS}

L'Event Sourcing trova la sua massima efficacia quando implementato congiuntamente al pattern \textbf{CQRS} (\emph{Command Query Responsibility Segregation}). La separazione concettuale prevede:
\begin{itemize}
    \item \textbf{Lato Command (Write Model)}: riceve i comandi operativi, applica le regole del dominio aziendale e persiste i dati esclusivamente sotto forma di eventi atomici all'interno dell'Event Store.
    \item \textbf{Lato Query (Read Model)}: non esegue mai scansioni computazionalmente onerose sull'Event Store; un processo asincrono di fondo intercetta gli eventi emessi e aggiorna una \textbf{proiezione} denormalizzata su un database secondario (SQL piatto o archivio NoSQL). I client in lettura interrogano unicamente la proiezione precalcolata.
\end{itemize}

Lo snippet Python seguente illustra la logica di generazione e accodamento di un evento contabile, formalizzando la disgiunzione tra stato derivato e tracciamento dell'intento:

\begin{lstlisting}[language=python]
def apply_transaction(account_id, amount):
    """
    In Event Sourcing, non aggiorniamo direttamente una colonna 'balance'.
    Accodiamo invece un nuovo evento immutabile alla cronistoria dello stream.
    """
    transaction_event = {
        "event_type": "MoneyDeposited" if amount > 0 else "MoneyWithdrawn",
        "account_id": account_id,
        "amount": abs(amount),
        "occurred_at": datetime.utcnow().isoformat(),
        "metadata": {
            "source": "mobile_app",
            "currency": "EUR"
        }
    }

    # L'Event Store opera rigidamente come archivio 'append-only'
    event_store.append(
        stream_id=account_id,
        event=transaction_event
    )
\end{lstlisting}

\subsection{Scenari D'Uso, Vantaggi E Controindicazioni}

Il pattern Event Sourcing risulta raccomandato nei seguenti scenari:
\begin{itemize}[noitemsep]
    \item Quando è tassativo catturare l'intento applicativo, lo scopo e la motivazione semantica associata a ciascuna variazione di stato;
    \item Quando occorre minimizzare o abbattere integralmente la contesa derivante da scritture concorrenti su singoli record;
    \item In domini in cui è vitale conservare una traccia di controllo certificata (\emph{audit trail}), ripristinare configurazioni storiche pregresse o consentire rollback precisi;
    \item Quando gli eventi rappresentano l'astrazione naturale del business applicativo;
    \item Quando si desidera disaccoppiare radicalmente la fase di acquisizione dei dati da quella della loro elaborazione e propagazione.
\end{itemize}

Al contrario, l'adozione di Event Sourcing è sconsigliata o controproducente in presenza di:
\begin{itemize}[noitemsep]
    \item Domini semplici, con scarsa o nulla logica di business, tipicamente orientati a semplici CRUD dove la complessità aggiunta non è ammortizzabile;
    \item Sistemi che impongono consistenza forte immediata (\emph{strong consistency}) e aggiornamento in tempo reale sincrono delle viste di visualizzazione;
    \item Ambiti applicativi in cui lo storico delle operazioni non apporta valore di business;
    \item Sistemi a bassa contesa di aggiornamento, dominati da inserimenti sequenziali autonomi.
\end{itemize}

\newpage

% ==============================================================================
% SEZIONE 3: INDEX TABLE
% ==============================================================================

\section{Pattern Index Table}

\subsection{Contesto Del Problema: Accesso Ai Dati E Limitazioni Dei DBMS NoSQL}

Nei sistemi di gestione dei dati distribuiti (in particolare nei database cloud di tipo NoSQL quali archivi chiave-valore o documentali), le entità vengono partizionate e indicizzate primariamente attraverso una \textbf{chiave primaria} (\emph{Primary Key} o \emph{Partition Key}). L'applicazione sfrutta tale chiave per localizzare con estrema rapidità ($O(1)$) il record e la partizione fisica di riferimento.

Tuttavia, emergono severe problematiche quando le query devono accedere alle entità sfruttando criteri basati su attributi differenti dalla chiave primaria. Ad esempio, in una tabella clienti in cui la chiave primaria è il \texttt{CustomerID}, se un servizio necessita di recuperare periodicamente tutti i clienti residenti in una data città (\texttt{Town}), in assenza di un indice secondario il motore di storage è costretto a eseguire una scansione globale dell'intero dataset (\emph{full table scan}). In ambiente cloud, tale scansione comporta consumi massicci di unità di lettura I/O, latenze elevate e costi monetari proibitivi.

Mentre i database relazionali implementano nativamente strutture ad albero bilanciato (es. $B^+$-Tree) per supportare indici secondari gestiti automaticamente dal motore DBMS, la stragrande maggioranza dei datastore NoSQL distribuiti cloud non offre tale funzionalità, oppure la supporta imponendo pesanti limitazioni architetturali.

\subsection{Strategie Di Emulazione Degli Indici Secondari}

Qualora il datastore non supporti indici secondari nativi, l'architettura applicativa può emularli manualmente realizzando e manutenendo una o più \textbf{Index Table} dedicate. A seconda della natura delle query e del rapporto tra letture e scritture, si delineano tre strategie architetturali.

\subsubsection{Strategia 1: Denormalizzazione Completa}

Questa strategia prevede la duplicazione totale dell'intero record all'interno di ciascuna tabella di indice, organizzando fisicamente le tabelle secondarie in base alla chiave di ricerca desiderata (es. una tabella indicizzata per \texttt{Town} e una tabella indicizzata per \texttt{LastName}).

% [IMMAGINE DEL FILE DI SLIDE CON NOME "slides_lezione.pdf"] Rappresentazione schematica della Strategia 1: completa denormalizzazione con duplicazione dell'intero dataset nelle tabelle di indice secondario (Town e LastName) — Slide 74 / Pagina 10 del PDF

\begin{itemize}[noitemsep]
    \item \textbf{Vantaggi}: la query secondaria viene risolta con una singola operazione di lettura diretta ($O(1)$), massimizzando la velocità di estrazione senza richiedere alcun join logico applicativo.
    \item \textbf{Svantaggi e Limiti}: richiede una quantità elevata di spazio di archiviazione a causa della ridondanza totale. Inoltre, in presenza di dati altamente dinamici, l'overhead computazionale e di rete necessario per propagare ogni aggiornamento a tutte le tabelle duplicate diventa insostenibile. È indicata solo per dataset prevalentemente statici ad altissimo volume di lettura.
\end{itemize}

\subsubsection{Strategia 2: Indice Normalizzato Con Riferimento Alla Fact Table}

La seconda strategia definisce tabelle indice snelle, strutturate associando la chiave secondaria al mero identificatore primario (\texttt{Primary Key}) dell'entità. I dati anagrafici completi risiedono unicamente in un archivio centrale, formalmente denominato \textbf{Fact Table}.

% [IMMAGINE DEL FILE DI SLIDE CON NOME "slides_lezione.pdf"] Diagramma della Strategia 2: tabelle indice normalizzate con Secondary Key e puntatore Customer Reference (ID) verso la Fact Table centrale — Slide 75 / Pagina 11 del PDF

\begin{itemize}[noitemsep]
    \item \textbf{Vantaggi}: abbattimento radicale dello spazio di memorizzazione occupato e sensibile riduzione dei costi di manutenzione durante le scritture, poiché le modifiche agli attributi informativi non impattano l'indice.
    \item \textbf{Svantaggi}: penalizzazione della latenza. L'applicazione deve obbligatoriamente compiere \textbf{due letture sequenziali} (\emph{two-lookup penalty}): la prima sulla Index Table per ottenere la chiave primaria, e la seconda sulla Fact Table per recuperare l'entità effettiva.
\end{itemize}

\subsubsection{Strategia 3: Indice Parzialmente Normalizzato}

Questa strategia si colloca come compromesso architetturale bilanciato: l'Index Table memorizza la chiave secondaria, il puntatore alla chiave primaria e una replica dei soli \textbf{campi interrogati con elevata frequenza} (es. nome e cognome). Gli attributi a basso accesso rimangono confinati nella Fact Table.

% [IMMAGINE DEL FILE DI SLIDE CON NOME "slides_lezione.pdf"] Schema della Strategia 3: indice parzialmente normalizzato contenente la Secondary Key, il Customer Reference e i dati consultati frequentemente — Slide 76 / Pagina 12 del PDF

In questo modo, la maggioranza delle query standard si conclude con un'unica lettura rapida sull'indice. Solo nei rari casi in cui l'utente richieda il dettaglio completo dell'entità, il sistema effettua il secondo accesso alla Fact Table.

\subsection{Chiavi Composite E Supporto Al Dati Shardati}

Spesso le interrogazioni applicative combinano vincoli su attributi multipli (ad esempio: \emph{``Seleziona tutti i clienti residenti a Redmond aventi cognome Smith''}). In tali circostanze, l'Index Table adotta una \textbf{chiave composita}, ottenuta tramite la concatenazione logica degli attributi coinvolti:
\[
k_{\text{composite}} = \texttt{Town} \mathbin{\Vert} \texttt{LastName}
\]

% [IMMAGINE DEL FILE DI SLIDE CON NOME "slides_lezione.pdf"] Struttura dell'Index Table basata su Composite Key (Town + LastName) ordinata gerarchicamente, collegata alla Fact Table — Slide 77 / Pagina 13 del PDF

L'ordinamento logico dei record all'interno della tabella indice segue gerarchicamente il primo campo e, a parità di valore, il secondo. Tale tecnica risulta essenziale quando applicata a dataset frammentati orizzontalmente (\emph{sharded}), specialmente qualora la chiave di partizionamento originale sia sottoposta a funzioni crittografiche di hashing che disperdono i dati impedendo interrogazioni per intervallo.

\subsection{Considerazioni Progettuali, Cardinalità E Distribuzione Delle Query}

La realizzazione manuale di tabelle di indicizzazione comporta compromessi di consistenza ed efficienza che ne delimitano l'impiego:
\begin{itemize}
    \item \textbf{Assenza di indici speculativi}: ogni indice secondario impone costi di sincronizzazione e storage. Vanno creati esclusivamente indici a supporto di query frequenti ed essenziali, evitando indici ipotetici o saltuari.
    \item \textbf{Consistenza eventuale (\emph{Eventual Consistency})}: in architetture distribuite su vasta scala, l'aggiornamento simultaneo di Fact Table e Index Table richiede transazioni distribuite (\emph{two-phase commit}), le quali minano la disponibilità. Pertanto, la sincronizzazione viene spesso delegata a processi asincroni, esponendo il sistema a disallineamenti temporanei.
    \item \textbf{Campi non discriminanti (\emph{Nondiscriminating Keys})}: un attributo è definito non discriminante quando presenta una cardinalità estremamente limitata rispetto alla dimensione del dataset (ad esempio l'attributo binario sul genere). In tale scenario, l'indice produce bucket enormi di record e non riduce significativamente lo spazio di ricerca, risultando controproducente.
    \item \textbf{Distribuzioni fortemente asimmetriche (\emph{Skewed Data})}: se una chiave assume il medesimo valore per il $90\%$ dei record complessivi, l'uso dell'indice per tale valore comporterebbe un overhead maggiore rispetto alla scansione sequenziale. Tuttavia, l'indice si rivela altamente efficiente qualora l'applicazione interroghi prevalentemente la minoranza residuale ($10\%$).
\end{itemize}

\newpage

% ==============================================================================
% SEZIONE 4: MATERIALIZED VIEW
% ==============================================================================

\section{Pattern Materialized View}

\subsection{Discrepanza Tra Modello Di Scrittura E Modello Di Lettura}

Nei sistemi informativi di livello enterprise, le modalità con cui i dati vengono persistiti non coincidono con i formati richiesti dai client in consultazione. In sede di memorizzazione, architetti e amministratori privilegiano requisiti quali:
\begin{itemize}[noitemsep]
    \item Normalizzazione formale (per azzerare le anomalie di aggiornamento);
    \item Integrità referenziale e confini transazionali rigidi;
    \item Incapsulamento in aggregati documentali completi in ambito NoSQL.
\end{itemize}
Tuttavia, quando un modulo di rendicontazione o un cruscotto grafico richiede un estratto sintetico (ad esempio l'ammontare venduto aggregato per categoria merceologica), la struttura originaria costringe il motore a estrarre milioni di entità dettagliate, decodificarle in memoria ed eseguire computazionalmente costosi join e aggregazioni aggregate.

\begin{theorem}
Il \textbf{Materialized View Pattern} consiste nel generare in anticipo viste prepopolate e denormalizzate sui dati residenti in uno o più archivi eterogenei, formattandole nella struttura attesa dai consumatori finali per ottimizzare le prestazioni di estrazione.
\end{theorem}

\subsection{Caratteristiche E Ciclo Di Vita Delle Viste Materializzate}

Una vista materializzata differisce profondamente da una vista logica relazionale tradizionale: quest'ultima riesegue la query SQL di aggregazione a ogni invocazione, mentre la vista materializzata calcola i risultati a priori e li memorizza fisicamente su disco o memoria volatile.

% [IMMAGINE DEL FILE DI SLIDE CON NOME "slides_lezione.pdf"] Flusso del Pattern Materialized View: tabelle di origine come source of truth (Orders, Accounts, Items), generazione asincrona della vista denormalizzata e accesso read-only dell'applicazione — Slide 81 / Pagina 17 del PDF

Le proprietà distintive di una Materialized View includono:
\begin{itemize}
    \item \textbf{Natura usa-e-getta (\emph{Disposable})}: la vista non costituisce il sistema di record primario. Qualora l'archivio della vista subisca un'avaria o si corrompa, esso può essere rigenerato integralmente a partire dai dati sorgente.
    \item \textbf{Accesso in sola lettura}: l'applicazione non modifica mai direttamente la vista materializzata. Essa funge a tutti gli effetti da \textbf{cache persistente e specializzata}.
    \item \textbf{Precomputazione e trasformazione}: include colonne calcolate, conversioni valutarie, join multidatabase e raggruppamenti statistici preeseguiti.
\end{itemize}

\subsection{Strategie Di Aggiornamento E Consistenza Dei Dati}

Poiché i dati sorgente mutano nel tempo, la vista materializzata deve essere sincronizzata per preservare un grado accettabile di aderenza alla realtà. Le strategie di rigenerazione contemplano tre modalità operative:
\begin{enumerate}
    \item \textbf{Reattiva guidata dagli eventi (\emph{Event-Driven})}: ogni inserimento o modifica nei dati sorgente genera un evento che attiva un processo di ricalcolo o aggiornamento incrementale della vista. Questa soluzione minimizza la latenza di sincronizzazione, ma può generare un sovraccarico sistemico insostenibile qualora i dati di origine siano estremamente volatili.
    \item \textbf{Pianificata (\emph{Scheduled Task})}: un job batch rigenera la vista a intervalli prefissati (es. ogni notte o ogni ora). Introduce una discrepanza transitoria (\emph{stale data}), ma stabilizza il consumo di risorse.
    \item \textbf{Manuale o su trigger esterno}: attivata esplicitamente da procedure amministrative o flussi di orchestrazione.
\end{enumerate}

La disgiunzione temporale tra la modifica dei dati di origine e la loro propagazione impone l'accettazione del modello di \textbf{consistenza eventuale}. Inoltre, nei sistemi ad alta scalabilità, la vista può essere posizionata su tecnologie differenti rispetto al database primario (ad esempio memorizzata all'interno di un'istanza \emph{Redis} in-memory o su un indice \emph{Elasticsearch}), beneficiando del paradigma di \textbf{persistenza poliglotta}.

\subsection{Integrazione Con Microservizi E Sistemi Eterogenei}

Nel contesto dei microservizi, il principio cardine del Domain-Driven Design (DDD) impone che ogni servizio governi in modo esclusivo il proprio archivio di persistenza (\emph{Database-per-Service}). Quando una funzionalità richiede dati distribuiti tra bounded context eterogenei (es. microservizio Catalogo, microservizio Ordini e microservizio Spedizioni), interrogarli tramite chiamate REST sincrone in cascata degrada il tempo di risposta e introduce un accoppiamento stretto.

La Materialized View risolve tale criticità: un servizio dedicato ascolta gli eventi pubblicati dai vari microservizi e consolida i dati necessari all'interno di una vista materializzata locale, consentendo interrogazioni aggregate istantanee senza alcuna invocazione di rete verso i microservizi originari.

\subsection{Confronto Con Cache-Aside E Ruolo In CQRS}

È opportuno formalizzare le differenze concettuali tra pattern affini:
\begin{itemize}
    \item \textbf{Rispetto al Cache-Aside}: in un'architettura Cache-Aside convenzionale, la cache memorizza entità temporanee destinate a scadere (gestite mediante policy di time-to-live o algoritmi LRU), e l'applicazione deve gestire l'eventuale \emph{cache miss}. Al contrario, la Materialized View è \textbf{persistente}, prestrutturata stabilmente per soddisfare una specifica query e non prevede fallimenti di accesso da colmare a runtime sulla Fact Table.
    \item \textbf{Rispetto a CQRS}: la Materialized View rappresenta lo strumento tecnologico e implementativo d'elezione per erigere il lato di lettura (\textbf{Read Model}) di un'architettura CQRS accoppiata a Event Sourcing.
\end{itemize}

\subsection{Esempio Pratico: Calcolo Incrementale Dei Totali}

Si consideri un flusso di transazioni commerciali non strutturate archiviate come stream di verità. Una Materialized View preposta all'alimentazione di un cruscotto aziendale mantiene costantemente aggiornati i totali aggregati per categoria, consentendo letture a complessità costante:

% [IMMAGINE DEL FILE DI SLIDE CON NOME "slides_lezione.pdf"] Schema implementativo della vista materializzata con codice di aggiornamento incrementale: ricezione evento di vendita, scrittura su tabella aggregata (Redis/DynamoDB) ed estrazione ultra-rapida — Slide 85 / Pagina 23 del PDF

\begin{lstlisting}[language=python]
# Struttura di persistenza della vista materializzata (es. Redis Hash o tabella NoSQL)
materialized_view = {}

def update_sales_view(event):
    """
    Funzione reattiva invocata all'accadimento di un nuovo evento di vendita.
    Aggiorna in modo incrementale il totale della categoria associata.
    """
    category = event['category']
    amount = event['amount']

    if category in materialized_view:
        materialized_view[category] += amount
    else:
        materialized_view[category] = amount

def get_sales_dashboard():
    """
    Lettura ultra-rapida: nessun calcolo di aggregazione a runtime.
    Restituisce immediatamente i dati precomputati in anticipo.
    """
    return materialized_view
\end{lstlisting}

\newpage

% ==============================================================================
% SEZIONE 5: SHARDING
% ==============================================================================

\section{Pattern Sharding}

\subsection{Limiti Dello Scalamento Verticale Negli Archivi Di Dati Cloud}

Nelle applicazioni cloud a livello globale, il volume di dati gestito cresce esponenzialmente nel tempo, superando le capacità di memorizzazione e di elaborazione di una singola istanza computazionale. Un database monolitico ospitato su un unico server si scontra inevitabilmente con quattro colli di bottiglia fisici:
\begin{enumerate}
    \item \textbf{Spazio di archiviazione (\emph{Storage Space})}: un server possiede una quantità finita di dischi fisici. Sebbene sia possibile sostituire i dischi con unità più capienti o aggregare volumi SAN/NAS, si raggiunge un limite hardware oltre il quale l'espansione locale è preclusa.
    \item \textbf{Risorse computazionali (\emph{Compute Resources})}: all'aumentare dei client concorrenti, la CPU e la memoria RAM del singolo server si saturano, determinando accodamenti, tempi di risposta degradati e timeout dell'applicazione.
    \item \textbf{Larghezza di banda di rete (\emph{Network Bandwidth})}: la velocità di ingresso e uscita dei dati è limitata dalla capacità della scheda di interfaccia di rete (NIC) del nodo primario.
    \item \textbf{Vincoli geografici e normativi (\emph{Geography \& Compliance})}: requisiti di sovranità del dato (es. GDPR) impongono che i dati generati da utenti residenti in una specifica giurisdizione rimangano fisicamente entro i confini regionali; un database centralizzato non può soddisfare tale conformità preservando al contempo basse latenze.
\end{enumerate}

Lo scalamento verticale (\emph{Vertical Scaling} o \emph{Scale-Up}), consistente nell'incrementare potenza di calcolo, memoria e dischi della singola macchina, costituisce unicamente una misura temporanea e comporta costi marginali asintoticamente divergenti. Le applicazioni cloud necessitano di scalare quasi indefinitamente, imponendo l'adozione dello \textbf{scalamento orizzontale} (\emph{Horizontal Scaling} o \emph{Scale-Out}).

\subsection{Concetto Di Partizionamento Orizzontale E Architettura Sharded}

\begin{theorem}
Lo \textbf{Sharding} è un pattern architetturale basato sulla suddivisione orizzontale di un archivio di dati in un insieme di partizioni indipendenti, denominate \textbf{shard}. Ogni shard possiede il medesimo schema logico ma ospita un sottoinsieme disgiunto dei dati, operando come un database autonomo su un nodo di storage dedicato.
\end{theorem}

L'architettura partizionata conferisce molteplici vantaggi ingegneristici:
\begin{itemize}[noitemsep]
    \item \textbf{Scalabilità orizzontale}: l'infrastruttura può espandersi indefinitamente mediante l'aggiunta di nuovi nodi di storage a basso costo (\emph{commodity/off-the-shelf hardware});
    \item \textbf{Abbattimento della contesa}: le operazioni di I/O si distribuiscono parallelamente sui nodi, eliminando colli di bottiglia centralizzati;
    \item \textbf{Prossimità geografica}: gli shard possono essere dislocati nei datacenter fisicamente più vicini alle comunità di utenti che vi accedono primariamente.
\end{itemize}

\subsection{Selezione Della Shard Key E Logica Di Routing}

L'elemento progettuale determinante nella configurazione dello sharding è la scelta della \textbf{Shard Key} (o \emph{Partition Key}), ovvero l'attributo statico (o insieme di attributi) che governa la distribuzione deterministica dei record tra i vari nodi.
\begin{itemize}
    \item \textbf{Requisito di immutabilità}: la Shard Key deve essere rigidamente statica; variazioni nei valori di chiave imporrebbero la migrazione fisica inter-nodo del record, degradando le prestazioni.
    \item \textbf{Posizionamento della logica di routing}: l'indirizzamento delle richieste verso il corretto shard può risiedere:
    \begin{itemize}
        \item A livello del \textbf{codice di accesso ai dati dell'applicazione} (routing logico gestito tramite librerie client);
        \item Nativamente all'interno del \textbf{motore di database} (sharding trasparente all'applicazione, come in MongoDB o Amazon DynamoDB).
    \end{itemize}
\end{itemize}

L'astrazione della locazione fisica consente la migrazione o il partizionamento ulteriore dei dati in modo trasparente rispetto alla logica di business, a fronte dell'overhead computazionale richiesto per calcolare l'instradamento di ciascun dato richiesto.

\subsection{Bilanciamento Del Carico, Ribilanciamento E Query Non Allineate}

La distribuzione delle entità negli shard mira a evitare la formazione di nodi sovraccarichi (\emph{hotspot}). Tuttavia, l'architettura sharded introduce sfide rilevanti:
\begin{itemize}
    \item \textbf{Ribilanciamento dinamico}: qualora un sottoinsieme di dati si accresca in modo anomalo rispetto agli altri, l'infrastruttura deve supportare meccanismi di migrazione e ripartizionamento a caldo per riequilibrare l'occupazione delle risorse.
    \item \textbf{Query non allineate alla Shard Key}: la strategia di sharding deve rispecchiare fedelmente il pattern delle interrogazioni dominanti. In sistemi multi-tenant, ad esempio, se lo shard key coincide con il \texttt{TenantID}, tutte le interrogazioni basate su tale identificatore raggiungeranno direttamente il singolo nodo di competenza. Viceversa, qualora sia necessario ricercare un tenant per nome o collocazione geografica, il sistema non è in grado di determinare a priori lo shard di destinazione e deve trasmettere la query in broadcast a tutti i nodi (\emph{scatter-gather query}), penalizzando il throughput globale. In tale contingenza, diviene imperativo combinare lo sharding con il pattern \textbf{Index Table} per ricostruire mappe di instradamento ausiliarie.
\end{itemize}

% AA -------------------------------------------------------------------------------------------------------------

% ==============================================================================
% SEZIONE: SHARDING PATTERN
% ==============================================================================

\section{Il Pattern Di Partizionamento Orizzontale: Sharding Pattern}

Nelle architetture distribuite scalabili, la gestione di moli crescenti di dati (\emph{volume}) e di una frequenza elevata di transazioni di lettura e scrittura (\emph{throughput}) rappresenta uno dei colli di bottiglia computazionali più gravosi. Quando una singola istanza di base di dati o un singolo server fisico esaurisce le proprie capacità di scalabilità verticale (\emph{scale-up}), legata a limiti intrinseci di CPU, memoria RAM o I/O del sottosistema di archiviazione, la transizione verso la scalabilità orizzontale (\emph{scale-out}) diventa una necessità architetturale. Il pattern di Sharding (noto anche come partizionamento orizzontale) risponde a questa esigenza suddividendo un archivio logico monolitico in una moltitudine di partizioni indipendenti, denominate shard, ciascuna contenente un sottoinsieme disgiunto dell'intero insieme dei dati.

Un principio strutturale fondamentale risiede nella distinzione tra l'unità logica di partizione e la topologia fisica sottostante: non deve necessariamente sussistere una corrispondenza uno-a-uno tra uno shard e la macchina host che lo ospita. Un singolo server fisico o una singola macchina virtuale può fungere da host per molteplici shard concorrenti. La determinazione della collocazione fisica o logica di ciascun record dipende univocamente dalla scelta della chiave di sharding (\emph{shard key}), un attributo (o una tupla di attributi) associato all'entità di business, in base al quale la logica di partizionamento assegna e localizza il dato.

In ambito architetturale, si identificano tre macro-strategie per l'assegnazione e l'instradamento delle richieste verso gli shard: la strategia Lookup, la strategia Range e la strategia Hash.

% ------------------------------------------------------------------------------
% Sottosezione: Strategia Lookup
% ------------------------------------------------------------------------------
\subsection{La Strategia Lookup E Il Partizionamento Virtuale}

La strategia Lookup basa il proprio funzionamento su una componente logica di mappatura esplicita (\emph{routing map}). In questo paradigma, la logica di instradamento intercetta ciascuna richiesta dati proveniente dall'applicazione, estrae il valore della \emph{shard key} e consulta una tabella o un indice centralizzato che memorizza la mappatura univoca tra la chiave e lo shard di destinazione.

Un tipico contesto applicativo è rappresentato dai sistemi \textbf{multi-tenant}, nei quali un'infrastruttura condivisa serve una molteplicità di clienti aziendali indipendenti (i \emph{tenant}). In questa configurazione, l'identificativo univoco del cliente (\texttt{tenant\_id}) viene impiegato sistematicamente come \emph{shard key}. Tale impostazione garantisce che l'intera totalità dei record e dei dati transazionali relativi a uno specifico tenant sia confinata all'interno di un unico shard: più tenant possono condividere il medesimo shard, ma i dati di un singolo tenant non vengono mai frammentati tra partizioni differenti, preservando l'integrità relazionale locale e semplificando le query contestuali.

% [IMMAGINE DEL FILE DI SLIDE CON NOME "slide_estratto.pdf"] Schema architetturale della Lookup Sharding Strategy in un ambiente multi-tenant con tabella di routing e mapping tenant ID verso Shard A, Shard B, Shard C, Shard N — Slide 88 / Pagina 1 del PDF

La mappatura tra la chiave di sharding e l'architettura di memorizzazione fisica può strutturarsi secondo due paradigmi:
\begin{itemize}[noitemsep]
    \item \textbf{Mappatura fisica diretta:} ciascun valore o raggruppamento di \emph{shard key} punta direttamente all'indirizzo di rete o all'istanza di partizione fisica corrispondente. Tale schema presenta una rigidità intrinseca: qualora sia necessario ribilanciare il carico o trasferire dati a causa della saturazione di un nodo, le tabelle di configurazione dell'applicazione o i riferimenti infrastrutturali devono essere aggiornati esplicitamente.
    \item \textbf{Partizionamento virtuale (\emph{Virtual Partitioning}):} introduce uno strato di indirezione architetturale. La \emph{shard key} viene associata a un insieme prefissato di \textbf{shard virtuali} ($V_1, V_2, \dots, V_K$), i quali a loro volta vengono mappati trasparentemente su un numero inferiore di partizioni fisiche ($P_1, P_2, \dots, P_M$ con $M \le K$). L'applicazione indirizza le query esclusivamente verso la partizione virtuale. Qualora un server fisico necessiti di decongestionamento, il sistema trasferisce un sottoinsieme di shard virtuali verso un nuovo nodo fisico e aggiorna il mapping virtuale-fisico a runtime, senza richiedere alcuna ricompilazione o modifica del codice client né l'alterazione delle chiavi di partizione.
\end{itemize}

% ------------------------------------------------------------------------------
% Sottosezione: Strategia Range
% ------------------------------------------------------------------------------
\subsection{La Strategia Range E Le Query Su Intervalli}

La strategia Range organizza e raggruppa elementi correlati all'interno del medesimo shard ordinandoli sequenzialmente rispetto a una \emph{shard key} monotona o temporale. Questa strategia è concepita specificamente per massimizzare l'efficienza dei carichi di lavoro caratterizzati da query su intervalli (\emph{range queries}), ovvero operazioni in cui il client non richiede un singolo elemento puntuale, bensì un insieme ordinato di record contigui.

Se, ad esempio, un'applicazione gestionale o di e-commerce deve interrogare frequentemente la totalità degli ordini evasi all'interno di un mese specifico, la memorizzazione di tali record all'interno del medesimo shard in ordine strettamente cronologico consente al motore di database di risolvere l'interrogazione tramite una scansione sequenziale di blocchi di memoria contigui. Se gli ordini fossero distribuiti casualmente tra nodi differenti, l'applicazione dovrebbe eseguire una miriade di singole query puntuali (\emph{point queries}) distribuite su rete, raccogliendo e riordinando i risultati a livello applicativo con un disastroso decadimento delle prestazioni.

% [IMMAGINE DEL FILE DI SLIDE CON NOME "slide_estratto.pdf"] Schema architetturale della Range Sharding Strategy con raggruppamento temporale degli ordini su Shard A (ottobre), Shard B (novembre), Shard C (dicembre) — Slide 89 / Pagina 2 del PDF

Per implementare con efficacia tale raggruppamento, si adottano frequentemente chiavi composte (\emph{composite keys}). In una chiave composita a base temporale, l'elemento gerarchicamente più significativo è costituito dal mese dell'ordine, seguito a cascata dal giorno, dall'ora e dal timestamp di transazione:
\[
    K_{\text{composite}} = \langle \text{Mese}, \, \text{Giorno}, \, \text{Ora}, \, \text{ID\_Univoco} \rangle
\]
In molti motori NoSQL o archivi a partizione scalabile, il modello a due componenti si articola formalmente in:
\begin{enumerate}[noitemsep]
    \item \textbf{Partition Key:} determina univocamente lo shard fisico o logico in cui risiede il record;
    \item \textbf{Row Key (o Clustering/Sort Key):} individua univocamente il record all'interno dello shard e governa l'ordinamento fisico dei dati sulle strutture di memorizzazione (es. tabelle LSM o B-Tree).
\end{enumerate}
Le entità che devono essere interrogate congiuntamente condividono il medesimo valore di \emph{partition key}, possedendo al contempo valori distinti e sequenziali di \emph{row key}.

% ------------------------------------------------------------------------------
% Sottosezione: Strategia Hash
% ------------------------------------------------------------------------------
\subsection{La Strategia Hash E La Mitigazione Degli Hotspot}

La strategia Hash mira all'obiettivo diametralmente opposto rispetto alla concentrazione sequenziale: essa si prefigge di garantire la massima omogeneità nella distribuzione dei dati e dei carichi operativi, azzerando o minimizzando l'insorgenza di hotspot (nodi o partizioni che assorbono una quota sproporzionata del volume di traffico o dello spazio disco).

In questo schema, l'assegnazione dello shard avviene tramite l'applicazione di una funzione di hashing deterministica $h(\cdot)$ calcolata su uno o più attributi identificativi dell'entità:
\[
    \text{Shard Destinazione} = h(K) \pmod N
\]
oppure tramite schemi di hashing coerente (\emph{consistent hashing}), dove $N$ rappresenta il numero complessivo delle partizioni. La funzione di hashing deve garantire proprietà di pseudo-casualità ed equa dispersione, garantendo che chiavi generate in sequenza cronologica o con valori numerici contigui vengano proiettate su partizioni fisicamente distinte.

% [IMMAGINE DEL FILE DI SLIDE CON NOME "slide_estratto.pdf"] Schema architetturale della Hash Sharding Strategy con funzione hash che mappa tenant sequenziali (55 e 56) su shard differenti (Shard B e Shard N) — Slide 90 / Pagina 3 del PDF

Per comprendere il netto vantaggio computazionale della strategia Hash rispetto alla strategia Range, si consideri uno scenario applicativo multi-tenant in cui la registrazione di nuovi utenti generi identificativi sequenziali ($1, 2, \dots, n, n+1, \dots, m$). Se venisse adottata una strategia Range, i tenant da $1$ a $n$ confluirebbero nello shard A, mentre i tenant da $n+1$ a $m$ risiederebbero nello shard B. Nei moderni servizi software, i clienti registrati di recente mostrano tipicamente pattern di attività e transazione di ordini di grandezza superiori rispetto ai clienti storici silenti. Ne consegue che l'attività computazionale e le richieste I/O si concentrerebbero quasi interamente sull'ultimo shard creato (fenomeno noto come \emph{write hotspot}), lasciando i server precedenti in uno stato di sottoutilizzo. La strategia Hash, calcolando il valore $h(\texttt{tenant\_id})$, assegna tenant contigui (come i tenant $55$ e $56$) a partizioni completamente indipendenti (ad esempio Shard B e Shard N), ripartendo uniformemente il carico di scrittura e di elaborazione sull'intero cluster.

% ------------------------------------------------------------------------------
% Sottosezione: Confronto e Criteri Architetturali
% ------------------------------------------------------------------------------
\subsection{Analisi Comparativa E Criteri Di Progettazione Basati Sul Dominio}

La scelta della strategia ottimale richiede un bilanciamento analitico tra complessità gestionale, flessibilità di ripartizionamento ed efficienza delle interrogazioni. La Tabella~\ref{tab:sharding_strategies} sintetizza i trade-off operativi di ciascun approccio.

\begin{table}[H]
\centering
\small
\caption{Confronto analitico tra le strategie di sharding.}
\label{tab:sharding_strategies}
\begin{tabularx}{\textwidth}{l p{4.2cm} p{4.2cm} p{4.5cm}}
\toprule
\textbf{Strategia} & \textbf{Vantaggi Principali} & \textbf{Svantaggi e Criticità} & \textbf{Complessità di Rebalancing} \\
\midrule
\textbf{Lookup} &
Massimo controllo sulla configurazione. Supporto eccellente al partizionamento virtuale. Disaccoppia logica client da allocazione fisica. &
Overhead prestazionale dovuto alla doppia interazione (accesso alla tabella di lookup prima dell'interrogazione reale). &
Minima, specialmente con virtual shard: si aggiorna la tabella di lookup senza alterare il codice applicativo. \\
\midrule
\textbf{Range} &
Semplicità implementativa. Ottimizzazione estrema per range query e scansioni temporali contigue. &
Scarsa uniformità di bilanciamento. Elevato rischio di \emph{hotspot} su scritture sequenziali e chiavi adiacenti. &
Elevata e onerosa: suddividere o fondere intervalli contigui richiede la riallocazione massiva di ampie porzioni di dati. \\
\midrule
\textbf{Hash} &
Distribuzione uniforme del volume e del throughput. Routing diretto deterministico senza tabelle di stato intermedie. &
Completamente inefficiente per le range query (richiede interrogazioni broad-cast su tutti gli shard). &
Difficile: la variazione del numero di shard richiede algoritmi complessi come il \emph{consistent hashing} per limitare lo spostamento dati. \\
\bottomrule
\end{tabularx}
\end{table}

Al di là della pura equipartizione statistica, la progettazione dello sharding in ambienti industriali reali deve recepire i requisiti funzionali, economici e legali dell'applicazione:
\begin{itemize}[noitemsep]
    \item \textbf{Sharding basato sul carico di lavoro (\emph{Workload}):} i tenant caratterizzati da un traffico estremo o da profili computazionali altamente volatili possono essere isolati all'interno di shard dedicati ad alte prestazioni. Questa segregazione impedisce che un tenant iperattivo degradi le prestazioni dei vicini (il problema del cosiddetto \emph{noisy neighbor}), stabilizzando la latenza per la restante utenza.
    \item \textbf{Sharding geolocalizzato (\emph{Location}):} la distribuzione dei dati segue la residenza geografica del cliente. Questa strategia abilita la pianificazione delle finestre di manutenzione e backup durante le ore notturne locali di ciascuna regione, preservando la piena operatività e disponibilità del servizio nei nodi delle aree geografiche attive.
    \item \textbf{Segmentazione per Livello di Servizio (\emph{SLA Tiering}):} i clienti strategici ad alto valore (\emph{high-value}) vengono posizionati su shard dedicati, scarsamente popolati e con risorse computazionali riservate. Al contrario, l'utenza standard o \emph{free-tier} viene consolidata in shard multi-tenant densamente popolati per ottimizzare i costi di gestione infrastrutturale.
    \item \textbf{Conformità normativa e isolamento formale:} nei settori regolamentati (sanità, finanza, GDPR), la riservatezza e i vincoli legali impongono che i record di determinati tenant non risiedano su archivi condivisi, bensì su server e database fisicamente disgiunti.
\end{itemize}

\newpage

% ==============================================================================
% SEZIONE: STATIC CONTENT HOSTING PATTERN
% ==============================================================================

\section{Pattern Per L'Offloading Dei Dati: Static Content Hosting}

Nelle architetture web e cloud tradizionali, i server applicativi (istanze compute basate su VM o container) gestiscono sia la logica di business computazionalmente intensiva, sia la consegna delle risorse statiche quali immagini, file di stile CSS, librerie JavaScript, font tipografici e documentazione statica (es. file PDF). Tale commistione introduce una grave inefficienza sistemica: i thread di lavoro e la memoria dell'application server, dimensionati e prezzati per sostenere calcoli complessi, accessi a database e logica transazionale, vengono vincolati per tempi prolungati a banali operazioni di streaming I/O di rete per trasferire file immutabili.

% ------------------------------------------------------------------------------
% Sottosezione: Concetto e Architettura
% ------------------------------------------------------------------------------
\subsection{Disaccoppiamento Delle Risorse E Meccanismo Di Offloading}

Il pattern Static Content Hosting risponde a questa problematica disaccoppiando in modo netto i file statici dal codice computazionale dinamico, delegando l'archiviazione e l'erogazione di tali risorse a servizi di \emph{object storage} gestiti in cloud (quali \textbf{Amazon S3} o \textbf{Azure Blob Storage}) oppure a reti geografiche di distribuzione dei contenuti (\textbf{Content Delivery Network}, o CDN).

La ratio architetturale è duplice:
\begin{enumerate}[noitemsep]
    \item \textbf{Preservazione del Compute Tier:} le istanze di calcolo vengono scaricate dalla gestione del traffico I/O pesante, consentendo di ridurne drasticamente il numero o la taglia orizzontale, minimizzando i costi dei nodi computazionali ad alta intensità di CPU/RAM.
    \item \textbf{Ottimizzazione della Consegna:} i servizi di storage gestito e le CDN sono architettati per garantire throughput aggregato elevato, elevata concorrenza e latenze minime sfruttando una rete distribuita di nodi \emph{Edge} dislocati in prossimità fisica rispetto all'utente finale.
\end{enumerate}

% [IMMAGINE DEL FILE DI SLIDE CON NOME "slide_estratto.pdf"] Confronto concettuale tra erogazione errata di contenuti statici tramite server applicativo e architettura corretta tramite AWS S3 e Amazon CloudFront — Slide 93 / Pagina 6 del PDF

Nel modello errato (\emph{wrong way}), un framework web come Flask o Django espone un endpoint di routing applicativo che apre localmente il file dal file system, alloca buffer di memoria ed esegue lo streaming verso il client, tenendo aperta una connessione di rete persistente sul server di backend:

\begin{lstlisting}[language=python, caption={Modello inefficiente: streaming di contenuti statici dal server applicativo.}]
# Approccio inefficiente: spreco di worker threads e compute resources
@app.route('/images/product.jpg')
def serve_image():
    # 1. Apertura del file da disco locale
    # 2. Lettura del payload binario
    # 3. Trasferimento al client attraverso il socket del server web
    # 4. Mantenimento della connessione TCP aperta sul processo applicativo
    return send_file('product.jpg')
\end{lstlisting}

Nel modello ottimizzato (\emph{correct way}), l'application server non viene minimamente sollecitato per il recupero dei file binari. Il rendering dell'interfaccia si limita a generare codice HTML contenente URI assoluti facenti capo all'infrastruttura di storage o al nome a dominio della CDN:

\begin{lstlisting}[language=python, caption={Configurazione e generazione template con disaccoppiamento su CloudFront/S3.}]
STATIC_FILES_HOSTING = {
    "provider": "Amazon S3",
    "bucket": "my-static-site",
    "url": "https://d3p7j5s9.cloudfront.net",
    "files": {
        "css/style.css": "https://d3p7j5s9.cloudfront.net/css/style.css",
        "js/app.js": "https://d3p7j5s9.cloudfront.net/js/app.js",
        "images/logo.png": "https://d3p7j5s9.cloudfront.net/images/logo.png"
    }
}
\end{lstlisting}

\begin{lstlisting}[language=html, caption={Frammento di codice HTML generato: accesso diretto al server statico edge.}]
<!-- L'application server non viene toccato: download diretto dal nodo edge piu vicino -->
<img src="https://d3p7j5s9.cloudfront.net/images/product.jpg" alt="Prodotto">
\end{lstlisting}

% ------------------------------------------------------------------------------
% Sottosezione: Problematiche e Considerazioni Progettuali
% ------------------------------------------------------------------------------
\subsection{Considerazioni Progettuali, Resilienza E Sicurezza}

L'implementazione concreta del pattern impone un'attenta analisi di molteplici fattori operativi:
\begin{itemize}[noitemsep]
    \item \textbf{Esposizione di Endpoint HTTP/HTTPS:} i bucket di storage devono supportare protocolli web standard su porta 80 e 443, garantendo la compatibilità con certificati SSL/TLS per la trasmissione sicura e la conformità ai vincoli di navigazione dei moderni browser.
    \item \textbf{Integrazione con CDN e Costi:} l'impiego di una CDN posiziona i contenuti nei \emph{Point of Presence} (PoP) globali. Sebbene massimizzi il throughput e riduca i tempi di \emph{Round-Trip Time} (RTT), l'infrastruttura introduce costi di traffico in uscita (\emph{data egress}) che devono essere bilanciati rispetto al risparmio computazionale.
    \item \textbf{Resilienza e Geo-Replicazione:} i container di storage in cloud sono tipicamente geo-replicati per impostazione predefinita a fronte di disastri nei data center. In tali scenari, l'indirizzo IP di erogazione può mutare per motivi di routing globale, ma il puntatore FQDN/URL rimane stabile, garantendo continuità operativa.
    \item \textbf{Complessità di Deployment e Versionamento:} disaccoppiando l'archivio statico dall'eseguibile applicativo, il ciclo di rilascio diventa eterogeneo. Modifiche a file JavaScript o fogli di stile CSS richiedono l'aggiornamento preventivo del bucket di storage e una strategia rigorosa di \emph{cache busting} (ad esempio, inserendo l'hash crittografico del contenuto nel nome del file, come \texttt{app.v2.hash987.js}), onde evitare che i client eseguano codice client-side disallineato rispetto alle API del backend.
    \item \textbf{Supporto a Domini Personalizzati e CORS:} qualora il servizio di storage non supporti nativamente domini personalizzati, i link devono puntare a un dominio differente da quello principale, imponendo un'adeguata configurazione delle policy di \emph{Cross-Origin Resource Sharing} (CORS).
    \item \textbf{Sicurezza degli Accessi (Read vs Write):} i contenitori di storage devono essere configurati per l'accesso pubblico in sola lettura (\emph{public read}). Risulta vitale accertarsi che i permessi di scrittura pubblica (\emph{public write}) siano rigorosamente bloccati, per impedire che attori malevoli carichino payload arbitrari. Qualora i dati statici abbiano natura confidenziale, l'accesso anonimo deve essere precluso, delegando il controllo ad appositi token ad accesso limitato (meccanismo formalizzato dal pattern \textbf{Valet Key}).
\end{itemize}

% ------------------------------------------------------------------------------
% Sottosezione: Ambiti di Applicabilità
% ------------------------------------------------------------------------------
\subsection{Ambiti Di Applicabilità E Limitazioni Operative}

Il pattern risulta particolarmente vantaggioso nei seguenti scenari:
\begin{itemize}[noitemsep]
    \item Riduzione drastica dei costi infrastrutturali per siti e portali web con alta presenza di asset multimediali e documentali;
    \item Condivisione centralizzata di risorse statiche verso client eterogenei o applicazioni ospitate in ambienti multicloud e server \emph{on-premises};
    \item Tracciamento analitico dei consumi: l'impiego di account di archiviazione dedicati consente di disaggregare chiaramente i costi di banda passante e memorizzazione dai costi orari di elaborazione della CPU.
\end{itemize}

Viceversa, l'adozione del pattern è controindicata nelle seguenti condizioni:
\begin{itemize}[noitemsep]
    \item \textbf{Elaborazione dinamica vincolata:} quando il documento richiede modifiche personalizzate o elaborazioni al volo prima dell'invio (ad esempio, l'apposizione di una filigrana digitale, un timestamp crittografico personalizzato o la redazione di campi riservati su un PDF a seconda dell'identità dell'utente);
    \item \textbf{Volumi trascurabili:} qualora la quantità di contenuti statici sia minima, la complessità architetturale derivante dalla gestione di deployment separati e dalla risoluzione di domini terzi supera ampiamente il beneficio economico marginale ottenuto.
\end{itemize}

\newpage

% ==============================================================================
% SEZIONE: VALET KEY PATTERN
% ==============================================================================

\section{Gestione Sicura Degli Accessi Diretti: Valet Key Pattern}

Nei sistemi distribuiti che gestiscono grandi flussi di dati (upload di file multimediali, download di archivi, ingestione di telemetrie IoT), il transito forzato di ogni flusso dati attraverso il server applicativo consuma risorse vitali di elaborazione, allocazione di memoria e banda passante, introducendo un grave collo di bottiglia architetturale. L'alternativa ingenua, consistente nel consentire al client di accedere direttamente al data store remoto, imporrebbe la condivisione delle credenziali primarie di sicurezza dell'archivio (le chiavi master o segreti di connessione amministrativa). Tale concessione distruggerebbe i confini di sicurezza del sistema: un client malevolo o compromesso non potrebbe più essere arginato e l'application server perderebbe il proprio ruolo di guardiano (\emph{gatekeeper}).

% ------------------------------------------------------------------------------
% Sottosezione: Principi Operativi del Valet Key
% ------------------------------------------------------------------------------
\subsection{Architettura Del Token Ad Accesso Limitato (Passepartout)}

Il pattern Valet Key (noto metaforicamente anche come \emph{"Passepartout"}) risolve questo conflitto consentendo al client di interagire direttamente con lo storage per il trasferimento effettivo dei dati, senza passare per il server applicativo, ma preservando un controllo d'accesso granulare e centralizzato. Il server applicativo non trasferisce i byte, bensì genera ed emette un token crittografico con permessi limitati nel tempo e nell'ambito (quale un \emph{Shared Access Signature} - SAS su Azure, o una \emph{Pre-signed URL} su Amazon S3) che il servizio di storage sottostante è in grado di convalidare in modo del tutto autonomo.

% [IMMAGINE DEL FILE DI SLIDE CON NOME "slide_estratto.pdf"] Flusso temporale e architettura del Valet Key Pattern: 1. Richiesta risorsa, 2. Generazione token da parte dell'applicazione, 3. Restituzione token al client, 4. Accesso diretto alla risorsa target tramite token — Slide 98 / Pagina 12 del PDF

Il ciclo operativo si articola nelle seguenti fasi:
\begin{enumerate}[noitemsep]
    \item Il client autenticato invia all'applicazione una richiesta logica di lettura o scrittura relativa a una risorsa;
    \item L'applicazione valida la richiesta, applica le regole di business e, agendo in qualità di autorità di sicurezza, produce un token temporaneo firmato crittograficamente;
    \item L'applicazione restituisce il token (o l'URI pre-firmato che lo ingloba) al client richiedente;
    \item Il client contatta direttamente lo storage provider presentando il token, avviando il download o l'upload del payload senza alcun coinvolgimento computazionale del backend.
\end{enumerate}

Questo approccio offre considerevoli vantaggi di gestione delle identità: elimina la necessità di creare, profilare e distruggere continuamente utenti a livello del sistema IAM (\emph{Identity and Access Management}) dello storage, sostituendo l'intero ciclo di vita amministrativo con la generazione stateless di una firma a runtime limitata per:
\begin{itemize}[noitemsep]
    \item \textbf{Ambito e granularità delle risorse:} il token può autorizzare l'accesso a un intero contenitore/tabella, o restringersi rigorosamente a una specifica partizione, riga, o singolo file binario identificato dal suo percorso logico;
    \item \textbf{Insieme minimale di operazioni consentite:} ad esempio, concessione di sola lettura (\texttt{Read-Only}) per la visualizzazione di un documento confidenziale, o autorizzazione di sola scrittura (\texttt{Write-Only}) per l'upload di un referto, precludendo l'elencazione o la sovrascrittura di altri file;
    \item \textbf{Finestra temporale di validità (\emph{Time-to-Live} o TTL):} la chiave scade automaticamente al termine del periodo prefissato, rendendo inutile qualsiasi tentativo di riutilizzo tardivo.
\end{itemize}

% ------------------------------------------------------------------------------
% Sottosezione: Considerazioni e Mitigazione dei Rischi
% ------------------------------------------------------------------------------
\subsection{Considerazioni Implementative E Mitigazione Dei Rischi}

L'adozione del pattern Valet Key comporta una parziale perdita di controllo immediato durante la fase di trasferimento dei dati. Di conseguenza, l'ingegnerizzazione del sistema deve presidiare rigorosamente i seguenti aspetti di sicurezza e scalabilità:

\begin{itemize}[noitemsep]
    \item \textbf{Gestione del periodo di validità e revoca anticipata:}
    Se il token viene intercettato o esfiltrato, esso garantisce a un eventuale attaccante l'accesso incondizionato alla risorsa target per tutta la durata della sua validità. È essenziale impostare finestre temporali di validità minime (ad esempio pochi minuti). Nel caso di trasferimenti massivi che richiedono ore, l'applicazione deve predisporre meccanismi di rinnovo del token. Inoltre, qualora il client notifichi il completamento dell'operazione, il server può invalidare esplicitamente la chiave intervenendo su policy lato storage o ruotando le chiavi di firma.

    \item \textbf{Controllo del comportamento dell'utente e quote dimensionali:}
    La maggior parte delle implementazioni dei sistemi di object storage non consente di vincolare, all'interno del token stesso, la dimensione massima del payload caricabile o il numero totale di accessi consentiti. Un utente malintenzionato in possesso di una chiave di upload potrebbe saturare lo spazio disco caricando file da terabyte o eseguire scritture a ciclo continuo (comportando gravosi attacchi di tipo \emph{Denial-of-Wallet}). Per mitigare tale limite, si possono imporre vincoli di upload lato client forzati da eventi di notifica (\emph{storage trigger/webhooks}) che avvisano l'applicazione al termine del salvataggio, consentendo di abortire o eliminare file anomali prima dell'ingestione definitiva.

    \item \textbf{Validazione e sanificazione dei dati in transito:}
    Poiché i dati bypassano l'application server durante il caricamento, il backend non può eseguire una scansione in tempo reale del payload. Dati malevoli potrebbero corrompere il sistema o nascondere codice malevolo (\emph{malware}). Risulta pertanto imperativo isolare i file caricati in un'area di \emph{staging} non accessibile ad altri utenti, attivando pipeline asincrone di ispezione, validazione formale e bonifica prima di autorizzarne l'uso nel dominio di business.

    \item \textbf{Auditing sistematico e contabilità:}
    Le API dello storage devono essere configurate per registrare analiticamente ogni transazione (tentativi di accesso, successi, fallimenti di autenticazione per token scaduto, volume di byte trasferiti). Tali log costituiscono la base per la fatturazione dell'uso effettivo della banda e per individuare anomalie o configurazioni errate nelle policy di accesso.

    \item \textbf{Canali di consegna sicuri e protezione in transito:}
    La trasmissione della Valet Key tra application server e client deve avvenire tassativamente su canali cifrati mediante protocollo HTTPS/TLS. Parimenti, l'accesso diretto del client al data store deve essere forzato esclusivamente su sessioni TLS, scongiurando rischi di \emph{eavesdropping} e attacchi di tipo \emph{Man-in-the-Middle}.
\end{itemize}

\newpage

% ==============================================================================
% SEZIONE: PANORAMICA CLOUD DESIGN PATTERNS
% ==============================================================================

\section{Tassonomia Dei Cloud Design Pattern Per Microservizi}

Nello sviluppo di sistemi distribuiti e architetture native per il cloud (\emph{cloud-native}), la progettazione del software deve confrontarsi strutturalmente con la natura inaffidabile, eterogenea e dinamica delle reti. A tal fine, la letteratura architetturale ha catalogato una serie di pattern ricorrenti volti a risolvere problematiche di disaccoppiamento, scalabilità, resilienza ed evoluzione dei sistemi legacy. La Tabella~\ref{tab:cloud_patterns_summary} illustra la classificazione dei principali pattern di progettazione e implementazione comunemente adottati nei moderni ambienti cloud.

\begin{table}[H]
\centering
\small
\caption{Quadro sinottico dei pattern architetturali per il cloud.}
\label{tab:cloud_patterns_summary}
\begin{tabularx}{\textwidth}{l X}
\toprule
\textbf{Pattern} & \textbf{Sintesi Funzionale E Scopo Architetturale} \\
\midrule
\textbf{Ambassador} &
Crea servizi helper locali (\emph{out-of-process}) che intercettano e gestiscono le richieste di rete per conto di un'applicazione client. \\
\addlinespace
\textbf{Anti-Corruption Layer} &
Implementa uno strato di mediazione, facciata o adattatore tra un'applicazione moderna e un sistema legacy per preservare il modello di dominio. \\
\addlinespace
\textbf{Backends for Frontends (BFF)} &
Sviluppa servizi di backend separati e ottimizzati per specifiche tipologie di interfaccia client (es. app mobile, browser web, dispositivi IoT). \\
\addlinespace
\textbf{CQRS} &
Segrega le operazioni di lettura dati (\emph{Query}) da quelle di modifica dello stato (\emph{Command}) tramite interfacce e modelli distinti. \\
\addlinespace
\textbf{Compute Resource Consolidation} &
Consolida molteplici task o carichi di lavoro computazionali su una singola unità di calcolo per ottimizzare l'allocazione delle risorse. \\
\addlinespace
\textbf{External Configuration Store} &
Estrae le informazioni di configurazione dal package applicativo accentrandole in un repository accessibile dinamicamente a runtime. \\
\addlinespace
\textbf{Gateway Aggregation} &
Raggruppa molteplici richieste individuali verso backend differenti all'interno di un'unica richiesta aggregata gestita da un gateway. \\
\addlinespace
\textbf{Gateway Offloading} &
Delega funzionalità condivise o trasversali (es. autenticazione SSL, logging, rate limiting) a un proxy di ingresso gateway. \\
\addlinespace
\textbf{Gateway Routing} &
Instrada le richieste verso servizi di backend distribuiti esponendo un singolo endpoint unificato verso i client esterni. \\
\addlinespace
\textbf{Leader Election} &
Coordina le azioni svolte da un insieme di istanze cooperanti eleggendo un leader responsabile della gestione dei carichi critici. \\
\addlinespace
\textbf{Pipes and Filters} &
Scompone un'elaborazione complessa in una pipeline lineare di stadi di elaborazione discreti e riutilizzabili. \\
\addlinespace
\textbf{Sidecar} &
Installa componenti ausiliari all'interno di un processo o container separato ma co-allocato per garantire isolamento e incapsulamento. \\
\addlinespace
\textbf{Static Content Hosting} &
Offload di asset statici su archivi di storage cloud o CDN per erogarli direttamente al client svincolando i server di calcolo. \\
\addlinespace
\textbf{Strangler Fig} &
Permette la migrazione incrementale di un'architettura monolitica legacy sostituendo gradualmente singoli moduli con microservizi. \\
\bottomrule
\end{tabularx}
\end{table}

\newpage

% ==============================================================================
% SEZIONE: AMBASSADOR PATTERN
% ==============================================================================

\section{Disaccoppiamento Delle Funzionalità Di Rete: Ambassador Pattern}

Nel paradigma a microservizi, un'applicazione o un servizio client (denominato tipicamente \emph{Servizio A}) necessita costantemente di comunicare via rete con servizi remoti, database transazionali e API eterogenee. Delegare direttamente alla logica di business applicativa la gestione delle problematiche di rete di basso livello (quali meccanismi di retry con backoff esponenziale, circuit breaking, terminazione e negoziazione TLS, autenticazione reciproca, tracciamento distribuito e instradamento dinamico) comporta conseguenze deleterie: il codice sorgente diventa eccessivamente complesso, fragile e accoppiato con specifiche librerie esterne.

Questa problematica assume contorni critici nell'integrazione di sistemi legacy o componenti \emph{brownfield}, i cui linguaggi di implementazione o i cicli di manutenzione non consentono l'inclusione di moderni SDK cloud.

% ------------------------------------------------------------------------------
% Sottosezione: Definizione e Ruolo Architetturale
% ------------------------------------------------------------------------------
\subsection{Definizione E Modello Architetturale}

L'Ambassador Pattern risolve radicalmente questa separazione di responsabilità: esso introduce un'istanza proxy \emph{out-of-process} (collocata fisicamente all'interno dello stesso ambiente host del client) che si interpone tra il servizio applicativo e il mondo esterno, assumendosi il carico di orchestrare e mettere in sicurezza le comunicazioni di rete.

L'ambassador può essere concettualizzato come una specializzazione del Sidecar pattern specificamente orientata alla gestione della connettività e delle chiamate di rete in uscita.

% [IMMAGINE DEL FILE DI SLIDE CON NOME "slide_estratto.pdf"] Schema a blocchi dell'Ambassador Pattern: Application code che invia richieste locali all'Ambassador (Retry, Circuit breaking, Monitoring, Security), il quale inoltra verso il Remote Service — Slide 103 / Pagina 19 del PDF

Le strategie di deployment di un'istanza Ambassador dipendono dall'ambiente di esecuzione:
\begin{itemize}[noitemsep]
    \item \textbf{Sidecar Container in ambienti orchestrati (es. Kubernetes):} l'ambassador viene istanziato come container secondario all'interno dello stesso Pod del container applicativo. Condividendo lo stesso \emph{network namespace} (e quindi l'indirizzo di loopback \texttt{localhost}), il client comunica con il proxy a latenze minime tramite socket locali;
    \item \textbf{Demone di sistema o servizio Windows/Linux:} qualora più processi indipendenti risiedano su un host comune non containerizzato, l'ambassador può essere eseguito come processo di background condiviso, fungendo da gestore unificato del traffico in uscita per tutte le applicazioni locali.
\end{itemize}

I benefici architetturali risultano evidenti:
\begin{enumerate}[noitemsep]
    \item \textbf{Indipendenza dal linguaggio:} il proxy opera a livello di protocollo di rete (es. HTTP/1.1, HTTP/2, gRPC); di conseguenza, le funzionalità avanzate di resilienza possono essere impiegate da servizi scritti in Python, Java, C++ o COBOL senza dover riscrivere una sola riga di logica di rete nativa.
    \item \textbf{Standardizzazione dell'osservabilità:} metriche quali latenza di rete, percentuali di errore, jitter e consumo di banda vengono raccolte in modo uniforme dall'ambassador ed esportate verso sistemi di monitoraggio (es. Prometheus), disaccoppiando le metriche applicative da quelle infrastrutturali.
    \item \textbf{Separazione delle responsabilità (\emph{Separation of Concerns}):} il team di sviluppo software può concentrarsi interamente sulle regole di business, delegando ai team di Site Reliability Engineering (SRE) o infrastruttura la configurazione dei criteri di retry, timeout, crittografia e discovery dei servizi.
\end{enumerate}

% ------------------------------------------------------------------------------
% Sottosezione: Idempotenza e Criticità
% ------------------------------------------------------------------------------
\subsection{Criticità Progettuali: Latenza, Contesto E Idempotenza}

L'adozione dell'Ambassador Pattern richiede un'attenta valutazione delle implicazioni prestazionali e logiche del sistema:

\begin{itemize}[noitemsep]
    \item \textbf{Overhead di latenza:} sebbene la comunicazione inter-processo su \texttt{localhost} sia rapida, essa introduce pur sempre un doppio cambio di contesto (\emph{context switch}) e l'attraversamento dello stack TCP/IP locale, risultando controindicata per microservizi caratterizzati da vincoli stringenti di latenza sub-millisecondo (\emph{ultra-low latency}).
    \item \textbf{Idempotenza e sicurezza dei retry automatici:}
    La delega delle politiche di re-inoltro automatico (\emph{retry}) all'ambassador è priva di rischi esclusivamente se le chiamate remote sottostanti sono idempotenti.

    \begin{theorem}[Idempotenza Operativa]
    Un'operazione $f$ si definisce \textbf{idempotente} se e solo se l'applicazione ripetuta dell'operazione produce il medesimo effetto sullo stato del sistema ottenuto tramite una singola invocazione:
    \[
        f(f(x)) = f(x) \quad \forall x \in \mathcal{S}
    \]
    Nel contesto delle reti, la cancellazione di un elemento da un insieme o l'aggiornamento di una risorsa tramite il metodo HTTP \texttt{PUT} o \texttt{DELETE} sono tipicamente idempotenti. Al contrario, un'operazione transazionale non idempotente (quale l'invio di un pagamento con \texttt{POST}) potrebbe causare addebiti duplicati qualora l'ambassador tenti una ri-esecuzione a causa di un banale timeout di risposta.
    \end{theorem}

    \item \textbf{Propagazione del contesto e override del client:} il client deve poter comunicare all'ambassador metadati di contesto (es. tramite header HTTP personalizzati come \texttt{X-No-Retry} o \texttt{X-Max-Retries: 5}), per consentire alla logica di business di scavalcare le configurazioni generiche del proxy quando necessario.
\end{itemize}

% ------------------------------------------------------------------------------
% Sottosezione: Simulazione Algoritmica
% ------------------------------------------------------------------------------
\subsection{Implementazione E Simulazione Di Un Ambassador In Python}

Per chiarire il flusso di interazione runtime tra il client applicativo, l'ambassador locale e il servizio remoto inaffidabile, si riporta una simulazione didattica completa basata sui listati delle slide di lezione.

Il componente remoto (\texttt{RemoteService}) simula un servizio esterno affetto da instabilità di rete, introducendo una probabilità di fallimento del $50\%$ tramite sollevamento di un'eccezione di timeout:

\begin{lstlisting}[language=python, caption={Simulazione di un servizio remoto inaffidabile.}]
import time
import random

# 1. Remote Service (Il sistema esterno inaffidabile)
class RemoteService:
    def do_work(self):
        # Simula un errore di rete transitorio con probabilita del 50%
        if random.choice([True, False]):
            raise Exception("Network Timeout")
        return "Data Success"
\end{lstlisting}

La classe \texttt{ServiceAmbassador} implementa il proxy locale intelligente: essa incapsula la connessione al servizio remoto ed espone al client il metodo \texttt{get\_data()}, all'interno del quale esegue fino a tre tentativi con una pausa di backoff pari a un secondo prima di arrendersi:

\begin{lstlisting}[language=python, caption={Implementazione del proxy Ambassador con politiche di retry e backoff.}]
# 2. Ambassador (Proxy locale intelligente per connettivita, retry e logging)
class ServiceAmbassador:
    def __init__(self):
        # Nel costruttore istanzia il servizio remoto simulandone la connessione
        self.remote_service = RemoteService()

    def get_data(self):
        max_retries = 3
        # Ciclo di tentativi controllato dall'Ambassador
        for attempt in range(max_retries):
            try:
                print(f"Ambassador: Attempt {attempt + 1} to call remote service...")
                # Tentativo di invocazione del servizio remoto
                return self.remote_service.do_work()
            except Exception as e:
                # Intercetta e logga l'errore senza propagarlo al client
                print(f"Ambassador: Error - {e}")
                # Semplice politica di backoff: attesa di 1 secondo prima del nuovo tentativo
                time.sleep(1)

        # Qualora tutti i tentativi falliscano, restituisce una segnalazione di errore gestita
        return "Error: Remote service unavailable"
\end{lstlisting}

Infine, la classe \texttt{Client} rappresenta l'applicazione reale. L'applicazione instanzia localmente l'ambassador e interagisce esclusivamente con esso; il codice del client è completamente svincolato da blocchi \texttt{try-catch} complessi, ritardi o logiche di recupero dagli errori di rete:

\begin{lstlisting}[language=python, caption={Codice del Client applicativo focalizzato esclusivamente sulla logica di business.}]
# 3. Client Application (Focalizzata unicamente sulla logica di business)
class Client:
    def __init__(self):
        # Il client dialoga solo con il local ambassador, non con la rete remota
        self.ambassador = ServiceAmbassador()

    def run(self):
        print("Client: Requesting data...")
        # Riceve direttamente il risultato consolidato (successo o errore)
        result = self.ambassador.get_data()
        print(f"Client: Final result is: {result}")

if __name__ == "__main__":
    client = Client()
    client.run()
\end{lstlisting}

\newpage

% ==============================================================================
% SEZIONE: GATEWAY AGGREGATION PATTERN
% ==============================================================================

\section{Ottimizzazione Della Comunicazione Client-Backend: Gateway Aggregation Pattern}

La transizione verso un'architettura orientata ai microservizi comporta la scomposizione di un sistema monolitico in decine o centinaia di servizi granulari indipendenti. Sebbene tale approccio favorisca l'indipendenza dei team e il deploy continuo, esso introduce una grave inefficienza nella comunicazione verso i client esterni (quali applicazioni mobile, interfacce Single-Page Web - SPA o dispositivi IoT).

% ------------------------------------------------------------------------------
% Sottosezione: Chattiness e Concetto
% ------------------------------------------------------------------------------
\subsection{La Problematica Della Comunicazione Eccessiva (Chattiness)}

Per renderizzare un'unica schermata (ad esempio, la pagina di visualizzazione di un ordine in un e-commerce), il dispositivo client necessita contemporaneamente di informazioni provenienti da molteplici domini: il catalogo prodotti, il servizio prezzi e sconti, lo stato del profilo utente, il tracking delle spedizioni e le recensioni dei clienti.

Qualora il client debba interrogare autonomamente ciascun microservizio, si manifesta il fenomeno della chattiness (comunicazione eccessivamente frammentata e frequente tra client e backend). Tale dinamica presenta gravi svantaggi:
\begin{itemize}[noitemsep]
    \item \textbf{Penalizzazione da latenza di rete:} le reti mobili cellulari (4G/5G) o wireless geografiche sono caratterizzate da latenze intrinseche non trascurabili. L'esecuzione di $N$ round-trip sequenziali moltiplica il tempo complessivo di caricamento dell'interfaccia, degradando la \emph{User Experience};
    \item \textbf{Consumo di risorse sul dispositivo client:} l'apertura simultanea di molteplici connessioni TCP e sessioni crittografiche TLS comporta un consumo elevato di batteria e banda sui dispositivi mobili;
    \item \textbf{Esposizione della topologia interna:} i client esterni devono conoscere gli indirizzi, le porte e i contratti specifici di decine di servizi interni, rendendo complessa qualsiasi operazione di refactoring architetturale del backend.
\end{itemize}

% ------------------------------------------------------------------------------
% Sottosezione: Risoluzione Architetturale
% ------------------------------------------------------------------------------
\subsection{Architettura Di Aggregazione Tramite API Gateway}

Il pattern Gateway Aggregation affronta questa problematica interponendo un'entità centrale (\textbf{API Gateway}) tra i client esterni e i microservizi di backend. Il gateway funge da punto unico di accesso: riceve una singola richiesta logica dal client, la scompone in chiamate concorrenti indirizzate ai rispettivi microservizi, attende le singole risposte, aggrega i dati e restituisce al client una risposta consolidata in un unico payload compatto.

% [IMMAGINE DEL FILE DI SLIDE CON NOME "slide_estratto.pdf"] Confronto architetturale tra chiamate multiple client-servizi (chattiness) e flusso mediato con Gateway Aggregation in 4 passaggi — Slide 109 / Pagina 26 del PDF

La dinamica operativa, illustrata analiticamente dal diagramma architetturale, si sviluppa attraverso le seguenti quattro fasi:
\begin{enumerate}[noitemsep]
    \item \textbf{Invio della richiesta aggregata (Fase 1):} l'applicazione client invia un'unica richiesta HTTP/HTTPS al Gateway API;
    \item \textbf{Decomposizione e fan-out (Fase 2):} il Gateway scompone la richiesta principale ed esegue chiamate concorrenti verso i singoli microservizi (\emph{Service 1}, \emph{Service 2}, \emph{Service 3}) sfruttando la rete interna del data center ad altissima velocità e bassissima latenza;
    \item \textbf{Raccolta delle risposte interne (Fase 3):} ciascun microservizio elabora la richiesta in modo del tutto autonomo e restituisce la propria risposta parziale al Gateway;
    \item \textbf{Composizione e risposta unificata (Fase 4):} il Gateway fonde i payload ricevuti in una struttura dati coerente (tipicamente un unico documento JSON) e lo invia al client, completando la transazione a fronte di una singola interazione su rete geografica.
\end{enumerate}

% ------------------------------------------------------------------------------
% Sottosezione: Considerazioni Progettuali
% ------------------------------------------------------------------------------
\subsection{Considerazioni Progettuali Ed Esecuzione Resiliente}

La centralizzazione del flusso di traffico all'interno del Gateway di aggregazione comporta precise responsabilità di progettazione architetturale:

\begin{itemize}[noitemsep]
    \item \textbf{Preservazione dell'autonomia e disaccoppiamento dei servizi:}
    Il gateway non deve mai diventare un elemento di accoppiamento stretto tra i microservizi sottostanti. Ciascun servizio di backend deve continuare a funzionare in modo totalmente indipendente, senza dipendenze nascoste indotte dalle logiche di aggregazione del gateway.

    \item \textbf{Prossimità topologica al backend:}
    Per trarre effettivo beneficio dal pattern, il Gateway deve essere posizionato all'interno dello stesso cloud provider, regione o rete virtuale (VPC) in cui risiedono i microservizi, minimizzando la latenza delle chiamate interne di \emph{fan-out}.

    \item \textbf{Mitigazione del Single Point of Failure (SPOF):}
    Essendo il gateway l'unico punto di ingresso per il client, un suo disservizio comporterebbe l'indisponibilità dell'intera applicazione. È mandatorio implementare cluster di gateway ad alta disponibilità dislocati su molteplici zone di disponibilità (\emph{Availability Zones}), protetti da bilanciatori di carico elastici di livello 4 o 7.

    \item \textbf{Prevenzione dei colli di bottiglia e dimensionamento elastico:}
    L'aggregazione di molteplici risposte richiede allocazione di CPU per la serializzazione/deserializzazione JSON e memoria per il buffer temporaneo dei payload. Il gateway deve scalare orizzontalmente in modo automatico per sostenere picchi improvvisi di throughput.

    \item \textbf{Resilienza e guasti a cascata (\emph{Cascading Failures}):}
    Il gateway deve essere sottoposto a rigorosi stress-test per accertarsi che il rallentamento di un singolo servizio di backend non provochi l'esaurimento dei thread pool del gateway stesso, trascinando l'intero sistema in un blocco sistemico.

    \item \textbf{Gestione dei timeout e risposte parziali (\emph{Partial Degradation}):}
    Se uno dei servizi invocati (es. il modulo per le raccomandazioni commerciali) subisce un ritardo eccessivo o non risponde entro un intervallo prefissato, il gateway deve interrompere l'attesa tramite timeout e restituire al client un insieme parziale di dati. L'interfaccia client gestirà con grazia l'assenza del dato non critico, garantendo che le informazioni vitali (quali il riepilogo del pagamento o il carrello) vengano comunque visualizzate dall'utente senza ritardi.
\end{itemize}

% AA -----------------------------------------------------------------------------------------

\section{Pattern Architetturali Per Il Gateway: Aggregation E Routing}

Nelle architetture distribuite a microservizi, la superficie espositiva verso il mondo esterno rappresenta uno dei nodi critici per l'intero sistema. Se i client (applicazioni mobili, \emph{Single Page Application} o client di terze parti) dovessero interagire direttamente con i singoli microservizi interni, emergerebbero gravi problemi di accoppiamento, latenza di rete e vulnerabilità di sicurezza. L'introduzione di un componente intermedio specializzato, comunemente identificato come \textbf{API Gateway}, permette di governare il traffico in ingresso, disaccoppiando l'infrastruttura backend dalle interfacce utente. All'interno di questo perimetro si distinguono due pattern architetturali fondamentali: l'\textbf{Aggregation} e il \textbf{Routing}.

\subsection{Gateway Aggregation Pattern: Principi E Anti-Pattern}

Il pattern \textbf{Gateway Aggregation} risponde alla necessità di aggregare molteplici richieste individuali verso differenti microservizi di backend in una singola richiesta logica da parte del client. Nelle reti ad alta latenza o a banda variabile (come le reti radiomobili cellulari), il costo associato al \emph{round-trip time} (RTT) per stabilire connessioni multiple e scambiare messaggi sequenziali o paralleli degrada pesantemente l'esperienza utente (\emph{chattiness} della rete). Il gateway riceve una singola richiesta dal client, effettua il \emph{fan-out} invocando in parallelo i vari servizi backend coinvolti, attende le risposte, le ricompone e le restituisce in un unico payload unificato.

\subsubsection{L'Anti-Pattern Del Gateway Con Accoppiamento Stretto}

Un errore architetturale comune nell'adozione di questo pattern consiste nel delegare al gateway responsabilità di trasformazione dati, manipolazione semantica o calcolo di logica di business.

% [IMMAGINE DEL FILE DI SLIDE CON NOME "CloudComputing"] Diagramma del codice e note esplicative sull'anti-pattern del Gateway che introduce accoppiamento stretto e logica di business — Slide 1 / Pagina 111 del PDF

Nel momento in cui il gateway viene contaminato da compiti che non gli competono, si generano vincoli rigidi e fragilità operative:
\begin{enumerate}[noitemsep]
    \item \textbf{Creazione di dipendenze nascoste}: se il gateway estrae un campo dal payload di un servizio (ad esempio, le preferenze dell'utente fornite da \texttt{user\_service}) per iniettarlo nei parametri di chiamata di un altro servizio (come \texttt{orders\_service}), si instaura un accoppiamento implicito tra due microservizi che dovrebbero rimanere concettualmente e operativamente indipendenti.
    \item \textbf{Assunzione di logica applicativa}: l'esecuzione di conversioni valutarie, calcoli matematici su prezzi o formattazioni complesse trasforma il gateway in un monolito nascosto. Se tale logica è replicata in più gateway (ad esempio in architetture \emph{Backend for Frontend} distinte per Web e Mobile), qualsiasi aggiornamento di tasso di cambio o correzione di bug richiede la ricompilazione e il ridispiegamento coordinato di tutti i gateway coinvolti.
    \item \textbf{Rottura immediata al cambio di schema}: una modifica sintattica introdotta da un team backend causa il fallimento a cascata del gateway.
\end{enumerate}

Si consideri lo scenario concreto illustrato nel listato seguente:

\begin{lstlisting}[language=python, caption={Anti-pattern di Gateway Aggregation con accoppiamento e logica di business.}]
# Anti-pattern: il gateway effettua calcoli e crea dipendenze incrociate
def get_user_dashboard(self, user_id):
    # 1. Chiamata a user_service
    user_data = self.user_service.get_user(user_id)

    # PROBLEMA: il gateway crea una dipendenza!
    # orders_service ora dipende dal formato interno restituito da user_service
    orders_data = self.orders_service.get_orders(
        user_id,
        user_data['preferred_format']  # Il gateway propaga dati da un servizio all'altro
    )

    # PROBLEMA ULTERIORE: combinazione dati hard-coded e calcoli di business
    for order in orders_data:
        # Il gateway assume che orders_service restituisca prezzi in euro,
        # ed esegue una moltiplicazione matematica per la conversione
        order['price_euro'] = order['price'] * user_data['euro_rate']

    return {
        'user': user_data,
        'orders': orders_data
    }
\end{lstlisting}

L'instabilità intrinseca di questo approccio emerge chiaramente nel ciclo di vita dell'applicazione:
\begin{itemize}[noitemsep]
    \item \textbf{Giorno 1}: il payload restituito da \texttt{user\_service} è coerente con le assunzioni del gateway:
    \begin{lstlisting}[language=python]
user_data = {'preferred_format': 'JSON', 'euro_rate': 1.1}
orders_service.get_orders(user_id, 'JSON') # Esito: OK
    \end{lstlisting}
    \item \textbf{Giorno 2}: il team responsabile di \texttt{user\_service} rinomina le chiavi del dizionario applicativo (ad esempio localizzandole o normalizzandole):
    \begin{lstlisting}[language=python]
user_data = {'formato_preferito': 'JSON', 'tasso_euro': 1.1} # Campi rinominati!
    \end{lstlisting}
    \item \textbf{Fallimento}: il gateway tenta di accedere a \texttt{user\_data['preferred\_format']} sollevando immediatamente un'eccezione non gestita:
    \begin{lstlisting}[language=bash]
KeyError: 'preferred_format' no longer exists!
    \end{lstlisting}
\end{itemize}

\subsubsection{Implementazione Corretta E Disaccoppiamento Funzionale}

L'implementazione corretta del pattern impone che l'API Gateway operi esclusivamente come aggregatore puro (\emph{pure aggregator}), agendo a livello di trasporto e composizione di payload, senza alterare la semantica interna dei dati.

% [IMMAGINE DEL FILE DI SLIDE CON NOME "CloudComputing"] Codice di implementazione del pattern di Gateway corretto con aggregazione pura e delega a servizio di arricchimento dedicato — Slide 2 / Pagina 112 del PDF

Ciascun servizio backend deve mantenere autonomamente il proprio formato e rimanere totalmente all'oscuro dell'esistenza degli altri servizi. Qualora siano necessarie operazioni di calcolo, aggregazione semantica o arricchimento del dato (\emph{enrichment}), queste devono essere delegate a uno specifico microservizio dedicato (\emph{enhancement service}), preservando l'isolamento del gateway e garantendo la testabilità unitaria della logica di business.

\begin{lstlisting}[language=python, caption={Implementazione architetturalmente corretta del Gateway Aggregator.}]
class CorrectGateway:
    def get_user_dashboard(self, user_id):
        # 1. Le chiamate sono mutuamente indipendenti
        user_data = self.user_service.get_user(user_id)
        orders_data = self.orders_service.get_orders(user_id)

        # 2. Il gateway aggrega i dati MA NON LI TRASFORMA
        return {
            'user': user_data,        # Esattamente come ricevuto dal servizio
            'orders': orders_data,    # Esattamente come ricevuto dal servizio
            'timestamp': time.now()
        }

    def get_enhanced_dashboard(self, user_id):
        # 3. Se e richiesta una trasformazione, si utilizza un servizio DEDICATO
        user_data = self.user_service.get_user(user_id)

        # Invocazione del servizio dedicato per l'arricchimento degli ordini
        enhanced_orders = self.orders_enhancement_service.enhance(
            user_id,
            user_data['preferences']  # Accettabile solo se e un contratto esplicito di un servizio dedicato
        )

        return {
            'user': user_data,
            'orders': enhanced_orders
        }
\end{lstlisting}

\subsubsection{Condizioni Di Applicabilità E Controindicazioni}

L'adozione del Gateway Aggregation deve essere guidata da precisi requisiti infrastrutturali e prestazionali:
\begin{itemize}[noitemsep]
    \item \textbf{Casi d'uso ideali}:
    \begin{itemize}[noitemsep]
        \item Un client deve interagire con molteplici microservizi dislocati nel backend per completare un'unica operazione logica d'interfaccia.
        \item Il client opera su canali di comunicazione caratterizzati da latenza elevata o jitter significativo (es. reti radiomobili cellulari 4G/5G, connessioni satellitari), dove l'accumulo di RTT multipli comporterebbe tempi di caricamento inaccettabili.
    \end{itemize}
    \item \textbf{Controindicazioni ed esclusioni}:
    \begin{itemize}[noitemsep]
        \item Se l'esigenza architetturale è ridurre le chiamate ripetute tra un client e un \textbf{singolo} servizio backend lungo molteplici operazioni, il pattern non è idoneo; la soluzione corretta consiste nell'implementare operazioni di tipo \emph{batch} direttamente sull'API di quel servizio.
        \item Quando il client risiede nello stesso contesto di rete locale o nella medesima regione cloud dei servizi di backend (bassa latenza, larghezza di banda abbondante), l'introduzione di un gateway aggregatore introduce un salto di rete (\emph{network hop}) superfluo e un potenziale collo di bottiglia elaborativo.
    \end{itemize}
\end{itemize}

---

\subsection{Gateway Routing Pattern}

Se l'aggregazione si occupa di unire dati provenienti da sorgenti multiple in un'unica risposta composita, il pattern \textbf{Gateway Routing} funge da vero e proprio \textbf{vigile urbano} (\emph{traffic policeman}) dell'architettura applicativa. Invece di far comunicare direttamente il client con gli indirizzi IP e le porte fisiche dei singoli microservizi, il Gateway Routing espone un \textbf{singolo punto di ingresso pubblico} (\emph{single point of entry}), ispeziona le richieste in transito e le instrada verso il corretto endpoint backend.

\subsubsection{Differenza Chiave Tra Aggregation E Routing}

La distinzione fondamentale tra i due pattern risiede nella cardinalità della mappatura delle comunicazioni di rete:
\begin{itemize}[noitemsep]
    \item \textbf{Gateway Aggregation}:
    \[
        1 \text{ richiesta client} \longrightarrow N \text{ chiamate backend} \longrightarrow 1 \text{ risposta unificata}
    \]
    \item \textbf{Gateway Routing}:
    \[
        1 \text{ richiesta client} \longrightarrow 1 \text{ chiamata backend specifica} \longrightarrow 1 \text{ risposta diretta}
    \]
\end{itemize}

Il gateway opera concettualmente come un \textbf{Reverse Proxy} a livello applicativo (\textbf{Layer 7} del modello ISO/OSI). Esso ispeziona l'\emph{Uniform Resource Identifier} (URI), i metodi HTTP, le intestazioni (\emph{headers}) o i parametri di query, inoltrando il traffico all'istanza corretta in maniera del tutto trasparente per il consumatore del servizio.

\subsubsection{Scenari Operativi E Valore Ingegneristico}

Senza un layer di instradamento centralizzato, qualsiasi variazione nell'infrastruttura backend ricade direttamente sul client, costringendo a rilasciare aggiornamenti continui dell'applicazione utente. Il Gateway Routing risolve tre problematiche cardine:

\begin{enumerate}
    \item \textbf{Consumo di molteplici servizi disparati}: in una piattaforma di commercio elettronico, il sistema è scomposto in servizi quali ricerca (\emph{search}), recensioni (\emph{reviews}), carrello (\emph{cart}), pagamento (\emph{checkout}) e storico ordini (\emph{order history}). Ognuno possiede endpoint, protocolli o cicli di rilascio differenti. Se un servizio viene rifattorizzato, diviso in due nuovi microservizi (\emph{service decomposition}) o fuso con un altro, il client non subisce alcun impatto: l'URL pubblico rimane immutato, variando unicamente le regole di instradamento interne al gateway.

    % [IMMAGINE DEL FILE DI SLIDE CON NOME "CloudComputing"] Diagramma architetturale di Gateway Routing applicato a molteplici servizi disparati (Search, Checkout, Order history, Cart, Reviews) — Slide 6 / Pagina 116 del PDF

    \item \textbf{Gestione di istanze multiple dello stesso servizio (Elasticità e Bilanciamento)}: per garantire alta disponibilità e sostenere picchi di carico, un microservizio può essere istanziato in decine di repliche distribuite su più nodi o regioni geografiche. La registrazione e deregistrazione dinamica delle istanze (\emph{service discovery}) viene interamente gestita dal gateway, occultando l'elasticità dell'infrastruttura al client.

    % [IMMAGINE DEL FILE DI SLIDE CON NOME "CloudComputing"] Diagramma architetturale di instradamento verso istanze multiple dello stesso servizio allocate in Region 1 e Region 2 — Slide 7 / Pagina 117 del PDF

    \item \textbf{Dispiegamento di versioni multiple (Strategie Di Rilascio)}: l'aggiornamento dei servizi mediante strategie avanzate come i \textbf{Blue-Green Deployments} o i rilasci \emph{Canary} richiede di convogliare frazioni controllate di traffico tra versioni coesistenti (es. versione \texttt{V1} consolidata e versione \texttt{V1.1} di test). Il gateway consente di regolare le percentuali di instradamento a livello di configurazione dinamica; qualora insorgano anomalie nella nuova versione, il ripristino (\emph{rollback}) è immediato, commutando il 100\% del traffico sulla versione precedente senza richiedere alcuna modifica al codice client.

    % [IMMAGINE DEL FILE DI SLIDE CON NOME "CloudComputing"] Diagramma di instradamento Layer 7 verso versioni multiple concorrenti (Search Service V1 e V1.1) per Blue-Green deployment — Slide 8 / Pagina 118 del PDF
\end{enumerate}

\subsubsection{Problematiche, Rischi E Considerazioni Di Progettazione}

L'introduzione del Gateway Routing comporta specifici vincoli architetturali che devono essere mitigati in fase di progetto:
\begin{itemize}[noitemsep]
    \item \textbf{Singolo punto di fallimento} (\emph{Single Point of Failure, SPOF}): centralizzare il traffico su un unico componente espone l'intero sistema a disservizi totali in caso di crash del gateway. È tassativo strutturarlo con meccanismi di clustering ad alta affidabilità e ridondanza attiva.
    \item \textbf{Collo di bottiglia prestazionale} (\emph{Bottleneck}): dovendo ispezionare gli header Layer 7 di tutti i pacchetti in transito, il gateway può saturare CPU e memoria. Sono necessari test di carico approfonditi (\emph{stress testing}) per dimensionare l'autoscaling del gateway e scongiurare fallimenti a cascata.
    \item \textbf{Granularità dell'instradamento}: l'instradamento a Livello 7 può essere regolato in funzione di indirizzi IP sorgente, porte, percorsi URL o header HTTP custom.
    \item \textbf{Gateway Globale vs Regionale}: nelle architetture multi-regione si adotta una gerarchia di routing. Un gateway globale distribuisce le richieste tra le diverse regioni geografiche, mentre un gateway regionale (quale un \emph{Application Gateway}) implementa il bilanciamento fine all'interno del singolo data center (ad esempio tra cluster di macchine virtuali o pod Kubernetes).
    \item \textbf{Isolamento di rete e perimetro di sicurezza}: il gateway funge da unico endpoint pubblicamente accessibile. I servizi backend a valle devono essere rigorosamente confinati all'interno di una rete virtuale privata (\emph{Private Virtual Network} o VPC), impedendone qualsiasi esposizione diretta a Internet.
\end{itemize}

\newpage

\section{Pattern Per L'Affidabilità: Queue-Based Load Leveling}

Nei sistemi cloud, la disponibilità (\emph{availability}) e la resilienza sono minacciate dalla natura altamente variabile dei carichi di lavoro. La tabella seguente sintetizza i pattern di progettazione per l'affidabilità (\emph{Reliability}):

\begin{table}[H]
\centering
\small
\begin{tabularx}{\textwidth}{lX}
\toprule
\textbf{Pattern} & \textbf{Descrizione Sintetica} \\
\midrule
\textbf{Deployment Stamps} & Dispiega copie multiple e indipendenti dei componenti applicativi, inclusi i data store, per isolare guasti e scalare. \\
\textbf{Geodes} & Distribuisce servizi backend su nodi geografici eterogenei, ognuno in grado di servire qualsiasi richiesta da qualunque regione. \\
\textbf{Health Endpoint Monitoring} & Implementa controlli funzionali interni all'applicazione accessibili a intervalli regolari da tool esterni. \\
\textbf{Queue-Based Load Leveling} & Inserisce una coda di messaggi come buffer tra un'attività e il servizio invocato per livellare picchi improvvisi di carico. \\
\textbf{Throttling} & Controlla e limita il consumo di risorse per istanza, utente o intero servizio quando si superano le soglie di tolleranza. \\
\bottomrule
\end{tabularx}
\caption{Quadro sinottico dei Cloud Design Patterns dedicati all'Affidabilità.}
\label{tab:reliability_patterns}
\end{table}

\subsection{Analisi Del Problema: Il Sovraccarico Intermittente}

Nelle architetture tradizionali a invocazione sincrona, se un'applicazione o un insieme di compiti concorrenti genera un volume imprevedibile di richieste verso un servizio downstream (ad esempio un database relazionale o un'API esterna), il carico istantaneo può superare drasticamente la capacità nominale del destinatario.

% [IMMAGINE DEL FILE DI SLIDE CON NOME "CloudComputing"] Architettura di fallimento sincrono diretto verso data store e risoluzione tramite Queue-Based Load Leveling con Service Bus e Function App — Slide 14 / Pagina 124 del PDF

Le conseguenze immediate di un mancato assorbimento del carico sono:
\begin{itemize}[noitemsep]
    \item \textbf{Degrado delle prestazioni}: aumento repentino della latenza a causa dell'accodamento incontrollato nei thread di rete.
    \item \textbf{Timeout e guasti a cascata}: le richieste dei client vanno in timeout esaurendo i pool di connessione, mentre il servizio backend collassa a causa della contesa interna sulle risorse fisiche (CPU, I/O su disco, lock sul database).
    \item \textbf{Costi elevati di provisioning}: per prevenire tali anomalie senza buffer, l'infrastruttura backend dovrebbe essere sovradimensionata per sostenere costantemente il \emph{carico di picco} (\emph{peak load}), pur rimanendo sottoutilizzata per la maggior parte del tempo.
\end{itemize}

\subsection{Principio Di Funzionamento E Architettura Asincrona}

Il pattern \textbf{Queue-Based Load Leveling} disaccoppia temporalmente il generatore del lavoro (\emph{producer} o task) dall'esecutore (\emph{consumer} o servizio), interponendo tra loro una \textbf{coda di messaggi} con semantica asincrona.

% [IMMAGINE DEL FILE DI SLIDE CON NOME "CloudComputing"] Diagramma concettuale del Queue-Based Load Leveling: arrivo richieste a tasso variabile, buffer del message queue ed estrazione a tasso costante da parte del servizio — Slide 13 / Pagina 123 del PDF

Il flusso operativo si articola nelle seguenti fasi:
\begin{enumerate}[noitemsep]
    \item Il task applicativo non attende l'elaborazione; incapsula i parametri della richiesta in un messaggio, lo deposita nella coda e termina immediatamente la propria operazione restituendo una ricevuta di presa in carico (tipicamente un codice HTTP \texttt{202 Accepted}).
    \item La coda funge da polmone d'accumulo (\emph{buffer}), trattenendo in modo persistente i messaggi durante le ondate transitorie di traffico.
    \item Il servizio ricevente (\emph{consumer}, ad esempio una \emph{Function App} \emph{serverless} o un demone dedicato) preleva i messaggi dalla coda a una velocità controllata e sostenibile, elaborandoli sequenzialmente o con un grado prefissato di parallelismo.
\end{enumerate}

\subsubsection{Formalizzazione E Dimensionamento Della Coda}

Una delle domande ingegneristiche fondamentali sollevate dal pattern riguarda il \textbf{dimensionamento della coda} (\emph{"How to size the queue?"}). Dal punto di vista della valutazione delle prestazioni, il sistema può essere modellato attraverso la \textbf{teoria delle code}.

Consideriamo una coda con arrivi a tasso medio $\lambda$ (pacchetti o messaggi al secondo) e un tasso di servizio medio $\mu$ per ciascun'istanza del worker. Quando il traffico subisce un picco temporaneo, si verifica localmente una condizione in cui:
\[
    \lambda_{\text{picco}} > \mu_{\text{totale}}
\]
Se il sistema fosse privo di buffer, la probabilità di blocco o di fallimento $P_b$ saturerebbe rapidamente a 1. Inserendo una coda con capacità massima di memorizzazione $K$ messaggi, l'eccedenza istantanea viene accumulata nel buffer. Applicando la \textbf{Legge di Little}, nello stato di regime la relazione tra il numero medio di messaggi nel sistema $\overline{N}$, il tasso medio di arrivo effettivo $\lambda$ e il tempo medio di permanenza $\overline{W}$ è definita da:
\[
    \boxed{\overline{N} = \lambda \cdot \overline{W}}
\]
Se il picco transitorio ha una durata temporale $\Delta t$, il volume di messaggi accumulati nella coda cresce approssimativamente al tasso:
\[
    \Delta N \approx \int_{0}^{\Delta t} \Big(\lambda(t) - \mu(t)\Big) \, dt
\]
Pertanto, la capacità $K$ del sistema di messaggistica deve essere dimensionata in modo tale da garantire:
\[
    K \ge \max_{\Delta t} \left[ \int_{0}^{\Delta t} \Big(\lambda(t) - \mu(t)\Big) \, dt \right]
\]
evitando la saturazione del disco o della memoria del broker di messaggi, evento che comporterebbe lo scarto (\emph{message drop}) dei messaggi.

\subsection{Metafora Della Ristorazione Ed Esempio Di Implementazione}

Per comprendere l'impatto architetturale del pattern, si consideri la classica analogia della ristorazione:
\begin{itemize}[noitemsep]
    \item \textbf{Senza coda (comunicazione sincrona bloccante)}: il cliente ordina al cameriere; il cameriere si reca in cucina e attende fermo che il cuoco prepari la pietanza per 30 minuti; solo allora torna dal cliente. Nel frattempo gli altri tavoli restano abbandonati, il thread pool dei camerieri è saturo e l'intero ristorante è bloccato.
    \item \textbf{Con coda (comunicazione asincrona non bloccante)}: il cliente ordina; il cameriere appende la comanda a una bacheca (la coda) e torna \emph{immediatamente} a servire gli altri tavoli; la cucina preleva le comande ed evade gli ordini al proprio ritmo costante. Al termine, il sistema emette una notifica: \emph{"Ordine n° 123 pronto!"}.
\end{itemize}

L'implementazione del lato produttore (\emph{producer side}) all'interno di un'API di backend si traduce in una logica non bloccante:

\begin{lstlisting}[language=python, caption={Esempio di accodamento asincrono sul produttore per il livellamento del carico.}]
def web_request_handler(user_data):
    """
    PRODUCER SIDE: Mima un endpoint di API Gateway.
    Invece di attendere il completamento di un task oneroso,
    livella il carico delegando l'elaborazione asincrona a una coda.
    """
    # Inserimento immediato del messaggio nella coda di elaborazione
    # L'esecuzione non si blocca, evitando di esaurire il thread pool dell'API
    job = task_queue.enqueue(intensive_background_job, user_data)

    # Ritorno immediato al client con ricevuta HTTP 202 Accepted
    return {
        "status": "Accepted",
        "job_id": job.get_id(),
        "message": "The system is busy, but your request has been queued."
    }
\end{lstlisting}

\subsection{Vantaggi Operativi E Punti Di Attenzione}

L'applicazione del Queue-Based Load Leveling apporta benefici misurabili:
\begin{enumerate}[noitemsep]
    \item \textbf{Massimizzazione della disponibilità}: guasti transitori o manutenzioni del servizio consumatore non interrompono l'operatività del produttore, che continua a inserire messaggi nella coda.
    \item \textbf{Ottimizzazione dei costi di computazione}: non è più necessario dimensionare l'infrastruttura backend per il picco massimo; è sufficiente dimensionarla per sostenere il \textbf{carico medio} (\emph{average load}), poiché la coda assorbe l'energia d'urto dei picchi.
    \item \textbf{Prevenzione dell'attivazione del throttling}: evita che il backend superi le soglie rigide di traffico imposte da provider o database.
\end{enumerate}

Tuttavia, vi sono criticità architetturali inderogabili evidenziate di seguito:
\begin{itemize}[noitemsep]
    \item \textbf{Incompatibilità con requisiti a bassa latenza}: il pattern è inutilizzabile se il client esige una risposta sincrona immediata.
    \item \textbf{Necessità di Rate Control nei consumatori}: il consumatore (\emph{worker}) deve controllare tassativamente la propria frequenza di elaborazione. Se un'applicazione serverless a valle (come un'istanza di \emph{Azure Functions}) dovesse effettuare un autoscaling selvaggio per consumare istantaneamente tutti i messaggi dalla coda, scaricherebbe nuovamente l'intero picco sul database sottostante, vanificando il livellamento.
    \item \textbf{Comunicazione unidirezionale}: le code sono intrinsecamente monodirezionali. Qualora il task richieda un responso, occorre implementare pattern asincroni complessi basati su \emph{Correlation ID}, code di risposta dedicate (\emph{Reply Queues}) o canali WebSocket/WebHook.
    \item \textbf{Persistenza e affidabilità del broker}: in caso di crash del gestore di code, se i messaggi non sono scritti su supporti persistenti (\emph{durable queues}), si rischia la perdita irreversibile dei dati. Devono essere considerate le politiche di consegna (\emph{at-least-once} delivery) e la gestione dei messaggi non recapitabili (\emph{Dead-Letter Queue}, DLQ).
\end{itemize}

\newpage

\section{Pattern Per L'Affidabilità: Throttling}

Sebbene l'autoscaling costituisca una delle prerogative centrali del cloud computing, esso non rappresenta una soluzione panacea per gestire qualsiasi incremento di carico. Quando le richieste computazionali eccedono la capacità massima allocabile o quando i tempi di reazione dell'infrastruttura non sono sufficientemente rapidi, il pattern \textbf{Throttling} interviene limitando proattivamente il consumo delle risorse per prevenire il collasso sistemico.

\subsection{Il Limite Dell'Autoscaling: Il Provisioning Lag}

Il carico sui sistemi cloud è tipicamente variabile: picchi nelle ore lavorative, elaborazioni analitiche pesanti a fine mese, o scatti improvvisi e imprevedibili (\emph{bursts}). L'autoscaling adatta dinamicamente le risorse stanziate al fabbisogno reale. Tuttavia, il processo di provisioning non è mai istantaneo:
\begin{itemize}[noitemsep]
    \item La misurazione della metrica di carico e il superamento delle soglie richiede finestre di osservazione (generalmente da 1 a 5 minuti).
    \item L'allocazione di nuove macchine virtuali, l'avvio del sistema operativo o il \emph{pull} di immagini di container complessi con successiva inizializzazione del runtime richiede tempi tecnici (\emph{cold start}).
\end{itemize}
Durante questa finestra temporale transitoria (denominata provisioning lag o finestra di deficit delle risorse), la domanda supera drasticamente l'offerta computazionale. Se il sistema non interviene, i nodi attivi esauriscono memoria e CPU, provocando il fallimento totale del servizio. In tale intervallo, il Throttling funge da meccanismo di protezione temporaneo fino al completamento dell'estensione elastica.

\subsection{Strategie Di Attuazione Del Throttling}

Il throttling può essere implementato secondo diverse direttrici architetturali:

\begin{enumerate}
    \item \textbf{Rigetto delle richieste basato su soglia (Rate Limiting)}: il sistema conteggia le chiamate effettuate da uno specifico utente o \emph{tenant}. Se il numero di invocazioni supera $n$ richieste al secondo in un intervallo fissato, le richieste eccedenti vengono respinte con l'errore standard HTTP \texttt{429 Too Many Requests}.

    % [IMMAGINE DEL FILE DI SLIDE CON NOME "CloudComputing"] Diagramma multi-tenant con instradamento e applicazione del rate limiting (Adatum a 5 req/s, Fabrikam a 10 req/s, Contoso respinto a 150 req/s) — Slide 23 / Pagina 133 del PDF

    \item \textbf{Degradazione controllata delle funzionalità (\emph{Feature Degradation})}: in condizioni di stress, invece di bloccare l'intero sistema, si disabilitano selettivamente i servizi secondari a beneficio di quelli essenziali. Ad esempio, una piattaforma di streaming video può ridurre la risoluzione dal formato 4K a definizioni inferiori, oppure un portale e-commerce può sospendere il motore di raccomandazione personalizzato pur garantendo il funzionamento ininterrotto del carrello e dei pagamenti.

    % [IMMAGINE DEL FILE DI SLIDE CON NOME "CloudComputing"] Grafico ad aree dell'utilizzo di risorse nel tempo per le funzionalità A, B e C, che illustra la sospensione temporanea della Feature B tra T1 e T2 — Slide 20 / Pagina 130 del PDF

    Come evidenziato dal grafico ad aree delle risorse, supponiamo che un'applicazione operi con tre funzionalità:
    \begin{itemize}[noitemsep]
        \item \textbf{Feature A}: elaborazione video / algoritmi di machine learning (critica).
        \item \textbf{Feature B}: reportistica e analisi dati batch (differibile).
        \item \textbf{Feature C}: cache in-memory / erogazione di contenuti statici (critica).
    \end{itemize}
    All'istante $T_1$, la somma delle risorse consumate tocca la soglia di guardia (\emph{soft limit}). Per scongiurare il collasso, il gestore disabilita temporaneamente la sola Feature B, rilasciando immediatamente le risorse che essa impegnava. Tra $T_1$ e $T_2$, le funzionalità critiche A e C continuano a operare normalmente. Una volta che l'assorbimento complessivo scende al di sotto del limite di sicurezza (istante $T_2$), la Feature B viene ripristinata in piena operatività.

    \item \textbf{Combinazione sinergica con Autoscaling}: il throttling non è necessariamente alternativo all'autoscaling, ma ne costituisce spesso il completamento temporaneo.

    % [IMMAGINE DEL FILE DI SLIDE CON NOME "CloudComputing"] Grafico temporale che mostra l'azione combinata di autoscaling e throttling come misura tampone tra T1 e T2 — Slide 21 / Pagina 131 del PDF

    All'istante $T_1$ il sistema intercetta il superamento del \emph{soft limit} e avvia la procedura di scalabilità orizzontale (\emph{scale-out}). Nel periodo $[T_1, T_2]$, durante il quale le nuove istanze non sono ancora pronte all'uso, il throttling entra in azione limitando il traffico per mantenere le metriche sotto la capacità massima. All'istante $T_2$, con i nuovi nodi operativi, la soglia di capacità totale aumenta e le restrizioni di throttling vengono rimosse.

    \item \textbf{Differimento delle operazioni e Priority Queue}: nelle architetture multi-tenant con accordi sul livello di servizio differenziati (\emph{Service Level Agreements}, SLA), in caso di sovraccarico le richieste degli utenti \emph{Premium} vengono servite a latenza zero, mentre le elaborazioni dei clienti con piani gratuiti vengono incanalate in code a bassa priorità o differite, notificando l'utente di riprovare più tardi.

    \item \textbf{Protezione dei servizi terzi}: il pattern riduce la concorrenza verso API esterne degradate, evitando che i log si riempiano di eccezioni inutili e scongiurando le tempeste di tentativi a vuoto (\emph{retry storms}). Librerie specializzate (quali ad esempio \emph{NServiceBus}) integrano costrutti di questo tipo.
\end{enumerate}

\subsection{Implementazione Pratica Nel Cloud: Ingress Controller Kubernetes}

Nelle infrastrutture basate su contenitori e orchestratori come Kubernetes, l'applicazione del throttling viene demandata frequentemente al controller di ingresso (\emph{Ingress Controller}), il quale applica politiche di \textbf{Rate Limiting} sfruttando algoritmi classici quali il \emph{Token Bucket} o il \emph{Leaky Bucket}.

Il manifest Kubernetes seguente mostra la configurazione dichiarativa di annotazioni su un Ingress basato su NGINX per limitare il throughput di richieste ammesse:

% [IMMAGINE DEL FILE DI SLIDE CON NOME "CloudComputing"] Manifest YAML per Ingress Kubernetes con configurazione di limit-rps, limit-burst-multiplier e whitelist — Slide 24 / Pagina 134 del PDF

\begin{lstlisting}[language=yaml, caption={Configurazione dichiarativa di Throttling e Rate Limiting tramite NGINX Ingress Controller.}]
apiVersion: networking.k8s.io/v1
kind: Ingress
metadata:
  name: api-gateway-ingress
  annotations:
    # 1. Abilitazione del Rate Limiting: definisce il limite nominale (5 richieste al secondo per IP)
    nginx.ingress.kubernetes.io/limit-rps: "5"

    # 2. Capacita di burst: moltiplicatore che ammette un buffer temporaneo prima del blocco rigido
    # Qui ammette fino a 5 * 2 = 10 richieste istantanee
    nginx.ingress.kubernetes.io/limit-burst-multiplier: "2"

    # 3. Whitelist: indirizzi IP esenti dall'applicazione del throttling (es. rete interna)
    nginx.ingress.kubernetes.io/limit-whitelist: "192.168.1.1"
spec:
  ingressClassName: nginx
  rules:
  - host: api.example.com
    http:
      paths:
      - path: /data
        pathType: Prefix
        backend:
          service:
            name: python-backend-service
            port:
              number: 80
\end{lstlisting}

\subsubsection{Considerazioni Architetturali Sul Throttling}

\begin{itemize}[noitemsep]
    \item \textbf{Scelta precoce di design}: introdurre il throttling in un'applicazione già esistente è estremamente oneroso. La gestione degli errori di rate limit deve essere prevista nativamente nei contratti API fin dalla fase di concezione.
    \item \textbf{Semantica degli errori e trasparenza verso il client}: le risposte di throttling devono trasmettere chiaramente codici univoci (HTTP 429) corredati dall'header standard \texttt{Retry-After: <secondi>}, permettendo al client di programmare la riesecuzione della richiesta senza generare traffico parassita.
    \item \textbf{Reattività del sistema di telemetria}: la capacità di limitare il traffico e rimuovere il blocco una volta rientrato il picco richiede una pipeline di monitoraggio ad altissima frequenza.
\end{itemize}

\newpage

\section{Pattern Per L'Alta Disponibilità: Bulkhead}

Un sistema software distribuito resiliente deve essere in grado di tollerare guasti parziali impedendo che l'indisponibilità di un singolo microservizio comprometta l'intero ecosistema applicativo. La tabella seguente illustra il quadro sinottico dei pattern orientati all'Alta Disponibilità (\emph{High Availability}):

\begin{table}[H]
\centering
\small
\begin{tabularx}{\textwidth}{lX}
\toprule
\textbf{Pattern} & \textbf{Descrizione Sintetica} \\
\midrule
\textbf{Deployment Stamps} & Replica dell'intero stack applicativo (incluso lo strato di memorizzazione) in unità indipendenti. \\
\textbf{Geodes} & Nodi geografici distribuiti e perfettamente omogenei, ciascuno attivo e capace di soddisfare qualunque richiesta. \\
\textbf{Health Endpoint Monitoring} & Verifica programmata e costante dello stato di salute dei componenti mediante endpoint dedicati. \\
\textbf{Bulkhead} & Isolamento dei componenti e delle risorse computazionali in compartimenti stagni per circoscrivere i guasti. \\
\textbf{Circuit Breaker} & Gestione delle anomalie transitorie interrompendo le chiamate verso servizi che richiedono tempo indeterminato per il ripristino. \\
\bottomrule
\end{tabularx}
\caption{Quadro sinottico dei Cloud Design Patterns dedicati all'Alta Disponibilità.}
\label{tab:ha_patterns}
\end{table}

\subsection{L'Analogia Navale E La Propagazione Del Guasto}

Il pattern Bulkhead (in italiano: \emph{paratia stagna}) trae il proprio nome dall'ingegneria navale. Gli scafi delle navi moderne sono suddivisi internamente mediante paratie verticali a tenuta stagna. Se lo scafo subisce uno squarcio, l'acqua penetra all'interno di un unico compartimento: la paratia impedisce all'acqua di riversarsi nelle sezioni adiacenti, mantenendo la nave a galla ed evitandone il naufragio.

Nei sistemi distribuiti, in assenza di isolamento, si assiste frequentemente al fenomeno del \textbf{guasto a cascata} (\emph{cascading failure}):
\begin{itemize}[noitemsep]
    \item Supponiamo che un client chiami due servizi, $S_1$ e $S_2$, condividendo un unico bacino di risorse (ad esempio lo stesso pool di thread o di connessioni TCP).
    \item Se il servizio $S_1$ subisce un rallentamento o smette di rispondere, i thread dedicati a servirne le richieste rimangono bloccati in attesa dei timeout di rete.
    \item A fronte di nuove richieste in ingresso verso $S_1$, l'intero pool di connessioni del client si esaurisce rapidamente.
    \item A questo punto, il client non dispone più di connessioni libere neppure per invocare il servizio $S_2$, sebbene quest'ultimo sia perfettamente sano e operativo. Un'anomalia confinata a un singolo servizio secondario finisce per paralizzare l'intera applicazione.
\end{itemize}

Lo stesso problema si manifesta in ottica \emph{multi-consumer}: se un singolo utente o client invia una raffica spropositata di richieste a un servizio comune, può saturarne interamente i thread o i canali di ricezione, negando di fatto l'accesso a tutti gli altri consumatori (\emph{noisy neighbor problem}).

\subsection{Meccanismi Di Isolamento E Partizionamento Delle Risorse}

Il pattern Bulkhead implementa la compartimentazione stagna operando su due fronti:

\begin{enumerate}
    \item \textbf{Partizionamento delle istanze di servizio}: raggruppare le istanze fisiche o i container in cluster separati in base al profilo di carico o alla criticità del consumatore. Un crash nel gruppo assegnato ai carichi batch a bassa priorità non tocca minimamente il gruppo di istanze dedicate alle transazioni bancarie critiche.
    \item \textbf{Partizionamento delle risorse di comunicazione (Connection Pools dedicate)}: assegnare a ciascun servizio dipendente un pool di connessioni circoscritto e isolato. Se il servizio $S_1$ smette di rispondere, saturerà esclusivamente il proprio connection pool dedicato; il pool riservato alla comunicazione con $S_2$ rimarrà intonso, preservando intatta la funzionalità residua del sistema (\emph{graceful degradation}).
\end{enumerate}

\subsection{Approfondimento Tecnologico: Gestione E Ruolo Dei Connection Pool}

Nelle architetture distribuite, stabilire una connessione di rete non è un'operazione neutra. Il diagramma concettuale del funzionamento di un pool di connessioni evidenzia la sequenza di gestione del ciclo di vita:

\begin{enumerate}
    \item \textbf{Connection Management (Overhead di creazione e distruzione)}: stabilire una connessione a livello di trasporto verso un database o un'API remota impone l'handshake a tre vie TCP (\emph{Three-Way Handshake}), la negoziazione crittografica dei certificati TLS/SSL e la successiva autenticazione di sessione. Creare e abbattere un socket per ogni singola richiesta applicativa comporterebbe un consumo proibitivo di CPU, un'elevata allocazione di descrittori di file (\emph{file descriptors}) a livello kernel e un'inutile latenza addizionale.
    \item \textbf{Pooling Delle Connessioni}: il client alloca e mantiene attiva una collezione di socket TCP persistenti pre-stabiliti (\emph{keep-alive}). Quando un'attività necessita di comunicare con un backend, richiede un socket al pool locale, lo impiega per la trasmissione e, al termine, lo restituisce immediatamente al pool.
    \item \textbf{Riutilizzo Ed Efficienza}: l'eliminazione dell'overhead di apertura e chiusura socket riduce l'RTT percepito e abbatte drasticamente il carico sui router intermedi e sui server di destinazione.
    \item \textbf{Limiti Di Concorrenza (\emph{Connection Limits})}: il pool stabilisce un tetto massimo di connessioni simultanee ammissibili verso una destinazione. Questa limitazione funge da duplice barriera:
    \begin{itemize}[noitemsep]
        \item Protegge il client dall'allocare un numero insostenibile di socket verso un servizio problematico.
        \item Tutela il servizio backend a valle, impedendogli di ricevere un volume di connessioni concorrenti tale da provocarne il crash per esaurimento di memoria o denial of service interno.
    \end{itemize}
\end{enumerate}

Grazie a questo isolamento compartimentato, il sistema ottiene i tre benefici capitali del Bulkhead:
\begin{itemize}[noitemsep]
    \item Immunità ai guasti a catena: le anomalie rimangono confinate nel rispettivo perimetro.
    \item Conservazione parziale dell'esperienza utente: l'applicazione continua a erogare le funzioni non compromesse.
    \item Differenziazione della Quality of Service (QoS): possibilità di riservare pool con quote garantite a consumatori prioritari o paganti.
\end{itemize}


% AA --------------------------------------------------------------------------------------------

\section{Gestione Avanzata Delle Connessioni E Pattern Di Resilienza}

Nelle architetture distribuite e orientate ai microservizi, la comunicazione di rete rappresenta uno degli elementi più critici per determinare la scalabilità, le prestazioni e l'affidabilità complessiva dell'intero sistema. Quando una molteplicità di istanze client interagisce in modo concorrente con svariati servizi di \emph{backend}, la gestione efficiente del ciclo di vita delle connessioni a livello di trasporto diviene un requisito architetturale primario. L'apertura e la chiusura continua di socket di rete e sessioni di autenticazione comporta infatti un \emph{overhead} computazionale considerevole; per ovviare a tali inefficienze, si ricorre strutturalmente a meccanismi di riutilizzo controllato delle risorse.

\subsection{Principi Di Connection Pooling Nei Sistemi Distribuiti}

Il \textbf{Connection Pooling} è un pattern fondamentale che prevede il mantenimento di un insieme (\emph{pool}) di connessioni pre-allocate, persistenti e condivise, pronte a essere assegnate alle operazioni che ne fanno richiesta per poi essere restituite al termine dell'elaborazione. Questo approccio evita il costo sistematico del \emph{three-way handshake} del protocollo TCP, della negoziazione TLS/SSL e dell'autenticazione applicativa a ogni singola transazione.

Nelle implementazioni avanzate di \emph{connection pooling}, oltre alla logica basilare di assegnazione e rilascio, entrano in gioco tre proprietà cardine:

\begin{itemize}[noitemsep]
    \item \textbf{Connection Recycling (Riciclo delle Connessioni):} oltre a consentire il riutilizzo intensivo dei canali aperti, il gestore del pool implementa politiche di riciclo o rinfresco periodico delle connessioni. Questa operazione previene il degrado dei descrittori di socket legato a problemi di connessioni obsolete (\emph{stale connections}), chiusure silenti da parte dei firewall o dei \emph{load balancer} intermedi lungo il percorso di rete, nonché potenziali perdite di risorse di memoria (\emph{resource leaks}).
    \item \textbf{Configurabilità (Tuning dei Parametri):} un pool di connessioni industriale espone una gamma di parametri operativi configurabili per adattarsi ai carichi di lavoro. Tra questi rientrano il numero massimo di connessioni parallele ammissibili (\emph{maximum pool size}), i tempi limite di inattività (\emph{idle timeout}) oltre i quali le connessioni in eccesso vengono rilasciate, nonché i meccanismi di accodamento temporale e le strategie di acquisizione (\emph{acquisition timeout}) ed espulsione forzata (\emph{eviction policies}).
    \item \textbf{Fault Tolerance (Tolleranza ai Guasti):} per prevenire che l'applicazione tenti di utilizzare connessioni compromesse, i pool moderni integrano meccanismi attivi di verifica dello stato di salute (\emph{health checks} preliminari o \emph{test-on-borrow}), oltre a strategie di riconnessione automatica trasparente capaci di gestire disconnessioni transitorie o fluttuazioni di rete senza propagare l'errore al codice applicativo.
\end{itemize}

L'efficienza ottenuta mediante il pooling di connessione costituisce il fondamento per ottimizzare l'utilizzo delle risorse, innalzare il \emph{throughput} applicativo e garantire la stabilità di sistemi distribuiti eterogenei.

\subsection{Pattern Di Isolamento Delle Risorse: Il Bulkhead Pattern}

Nel dominio dell'ingegneria del software resiliente, il \textbf{Bulkhead Pattern} mutua il proprio nome dalle paratie stagne impiegate nell'ingegneria navale. Nei natanti, lo scafo è suddiviso internamente in compartimenti ermetici indipendenti affinché una falla in una singola sezione confini l'allagamento al comparto leso, impedendo l'affondamento dell'intera nave.

Trasposto nell'architettura dei sistemi cloud-native, il pattern impone la compartimentazione strutturale delle risorse computazionali (quali thread, connessioni di rete, memoria e CPU) allocate a specifici carichi di lavoro o servizi downstream. In questo modo, l'esaurimento delle risorse causato da un servizio degradato non si propaga a valanga (\emph{cascading failure}) all'intero ecosistema.

\subsubsection{Isolamento Mediante Pool Di Connessioni E Thread Pool}

Una tipica declinazione del pattern consiste nel disaccoppiare i pool di connessione assegnati a ciascun servizio dipendente. Se un'applicazione client condividesse un unico pool di connessioni o un unico pool di thread per interrogare molteplici servizi remoti, il blocco di uno solo di essi assorbirebbe rapidamente tutte le risorse libere, paralizzando anche le chiamate indirizzate a servizi perfettamente operativi.

% [IMMAGINE DEL FILE DI SLIDE CON NOME "CloudComputing.pdf"] Diagramma architetturale del Bulkhead pattern con carichi di lavoro (Workload 1 e 2) e connection pool isolati per i servizi Service A, Service B e Service C — Slide 2 / Pagina 2 del PDF

Come schematizzato nella topologia di isolamento, le richieste provenienti dai diversi carichi di lavoro applicativi (\emph{Workload 1} e \emph{Workload 2}) accedono ai servizi a valle tramite canali rigorosamente distinti:
\begin{itemize}[noitemsep]
    \item \emph{Service A} dispone di un proprio \textbf{Connection Pool} dedicato e di un relativo pool di thread;
    \item \emph{Service B} e \emph{Service C} operano con i propri pool dedicati e indipendenti.
\end{itemize}
Qualora \emph{Service A} incorra in guasti, latenze anomale o blocchi completi, l'esaurimento dei descrittori di rete e dei thread associati resta confinato all'interno del compartimento dedicato a \emph{Service A}. I flussi di elaborazione che necessitano di interagire con \emph{Service B} e \emph{Service C} continuano a essere serviti regolarmente senza subire interruzioni o contese di risorse.

A piccola scala e con volumi di traffico ridotti, il costo computazionale legato all'apertura e chiusura sequenziale delle connessioni di rete e dei socket verso i database può apparire trascurabile. Tuttavia, al crescere delle dimensioni applicative e del livello di concorrenza, la continua creazione e distruzione di connessioni genera un sovraccarico insostenibile su kernel, stack di rete e motori database. Il \textbf{Database Connection Pooling} risponde a questa problematica mantenendo canali di comunicazione aperti e pronti all'uso; combinato con la paratizzazione (\emph{bulkheading}), esso garantisce che picchi anomali di query su un archivio dati non privino di risorse altri archivi o servizi concorrenti.

\subsubsection{Isolamento A Livello Di Istanze Di Servizio}

Un'altra modalità fondamentale di applicazione del Bulkhead Pattern riguarda il confinamento dei consumatori (\emph{clients}) su istanze applicative dedicate o partizionate, impedendo il cosiddetto fenomeno del \emph{noisy neighbor} (il "vicino rumoroso").

% [IMMAGINE DEL FILE DI SLIDE CON NOME "CloudComputing.pdf"] Schema di partizionamento delle istanze di servizio: Client 1, Client 2 e Client 3 assegnati ciascuno a una Service Instance dedicata per il medesimo servizio Service A — Slide 3 / Pagina 3 del PDF

Nello scenario descritto:
\begin{itemize}[noitemsep]
    \item Una molteplicità di entità client (\emph{Client 1}, \emph{Client 2}, \emph{Client 3}) deve consumare le funzionalità offerte da un medesimo servizio logico (\emph{Service A}).
    \item L'infrastruttura assegna a ciascun client un'istanza fisica o logica separata (\emph{Service Instance}).
    \item Se \emph{Client 1} invia un volume anomalo di traffico saturando l'istanza assegnata, l'effetto di sovraccarico rimane circoscritto. Poiché ciascuna istanza è isolata a livello di risorse operative, le richieste di \emph{Client 2} e \emph{Client 3} continuano a essere elaborate regolarmente dalle rispettive istanze operative.
\end{itemize}

\subsubsection{Considerazioni Architetturali E Granularità Del Partizionamento}

La corretta implementazione delle paratie stagne richiede un attento bilanciamento ingegneristico tra il grado di isolamento garantito e il costo computazionale introdotto:

\begin{itemize}
    \item \textbf{Requisiti di Business e Tecnici:} la definizione delle partizioni deve riflettere la criticità dei processi aziendali e i vincoli tecnologici dei carichi di lavoro.
    \item \textbf{Costo di Isolamento e Overhead:} all'aumentare dell'isolamento cresce l'onere in termini di licenze, memoria consumata, complessità di manutenzione e latenza. Il confinamento può avvenire a livello di processi di sistema operativo, thread pool, semafori applicativi, fino alla segregazione hardware-assisted tramite macchine virtuali (VM) o container leggeri. L'utilizzo di container (come Docker/Kubernetes) rappresenta un ottimo compromesso tra isolamento dei vincoli di risorsa e ridotto overhead di virtualizzazione.
    \item \textbf{Sinergia tra Pattern di Resilienza:} il Bulkhead Pattern non opera in isolamento, ma si integra in modo sinergico con il \textbf{Retry Pattern}, il \textbf{Circuit Breaker Pattern} e i meccanismi di limitazione del tasso di ingresso (\textbf{Throttling/Rate Limiting}).
    \item \textbf{Isolamento Asincrono tramite Code:} nelle comunicazioni asincrone, l'isolamento si ottiene assegnando code di messaggi distinte per ciascun flusso o tipologia di servizio. In alternativa, è possibile impiegare un gruppo condiviso di nodi di elaborazione che governano il prelievo (\emph{dequeue}) e la distribuzione dei task attraverso algoritmi di schedulazione equi.
    \item \textbf{Livello di Granularità e Multi-Tenancy:} nella gestione di ambienti multi-tenant, l'architetto può optare per un isolamento rigido (\emph{single-tenant bulkhead}, dove ogni utente/tenant dispone di un pool o un set di nodi dedicato) oppure per raggruppare sottoinsiemi di tenant all'interno di partizioni condivise, mitigando i costi infrastrutturali a fronte di un isolamento parziale.
    \item \textbf{Framework di Mercato:} nei contesti di sviluppo software, la compartimentazione a livello consumer viene realizzata tramite librerie specializzate per l'ingegneria della resilienza come \textbf{Resilience4j} (in ambito Java/JVM) o \textbf{Polly} (in ambito .NET), capaci di gestire thread pool isolati e semafori di concorrenza.
\end{itemize}

\subsubsection{Esempio Applicativo: Isolamento Di Risorse Tramite Kubernetes}

In ambienti cloud-native moderni, una paratia solida a livello di infrastruttura viene istanziata definendo vincoli rigidi di allocazione di risorse (\emph{requests} e \emph{limits}) nei manifest di orchestrazione dei carichi di lavoro containerizzati.

\begin{lstlisting}[language=yaml, caption={Definizione di un Pod Kubernetes con quote vincolate di risorse}]
apiVersion: v1
kind: Pod
metadata:
  name: drone-management
spec:
  containers:
  - name: drone-management-container
    image: drone-service
    resources:
      requests:
        memory: "64Mi"
        cpu: "250m"
      limits:
        memory: "128Mi"
        cpu: "1"
\end{lstlisting}

Nel manifest sopra riportato, il container operante all'interno del Pod `drone-management` viene confinato entro margini operativi precisi. Il costrutto `requests` specifica la dotazione minima garantita dallo scheduler del cluster ($64\,\text{MiB}$ di RAM e $250\,\text{m}$, ovvero $0{,}25$ di core CPU). Il costrutto `limits` fissa il tetto massimo invalicabile ($128\,\text{MiB}$ di memoria e $1$ core CPU intero). Qualora il container subisse un sovraccarico computazionale anomalo o un consumo incontrollato di memoria, il sottosistema del sistema operativo (\emph{cgroups}) e il demone \textbf{Kubelet} interverranno (ad esempio terminando il processo tramite l'\emph{Out-Of-Memory Killer} o limitando i cicli di clock con il \emph{CPU throttling}), impedendo la propagazione dell'esaurimento di risorse agli altri Pod co-allocati sul medesimo nodo worker.

\newpage

\subsection{Pattern Di Interruzione Del Flusso Di Chiamate: Il Circuit Breaker Pattern}

Nei sistemi distribuiti, i guasti possono essere categorizzati in due macro-famiglie:
\begin{enumerate}[noitemsep]
    \item \textbf{Guasti Transitori:} causati da micro-interruzioni di rete, picchi improvvisi di traffico temporaneo o lentezze microscopiche degli apparati. Tali fallimenti tendono ad autorisolversi rapidamente e vengono affrontati con successo tramite il \textbf{Retry Pattern} (reiterazione della richiesta a intervalli cadenzati).
    \item \textbf{Guasti Non Transitori o Catastrofici:} dovuti al crash definitivo di un database, errori fatali di configurazione, indisponibilità totale di un'infrastruttura cloud o lunghi periodi di manutenzione.
\end{enumerate}

In presenza di guasti non transitori, insistere ciecamente nel ritentare l'invocazione di un servizio compromesso è fortemente controproducente. In primo luogo, satura il chiamante mantenendo thread bloccati in attesa di timeout lunghi; in secondo luogo, inonda il servizio di destinazione (\emph{thundering herd problem}), impedendogli qualunque tentativo di recupero termico e trascinando verso il collasso altre parti indipendenti del sistema.

Il \textbf{Circuit Breaker Pattern} risolve esattamente questa criticità. Agendo come l'interruttore magnetotermico o salvavita negli impianti elettrici domestici, il circuito "scatta" (\emph{trips}) nel momento in cui rileva un tasso anomalo di fallimenti, interrompendo immediatamente il flusso delle chiamate successive verso il microservizio in crisi e facendole fallire all'istante (\emph{fail-fast}). In tal modo si preserva la capacità computazionale del client e si offre alla risorsa remota il tempo necessario per ristabilirsi.

\subsubsection{Funzionamento Della Macchina A Stati Finiti}

Il meccanismo logico del pattern agisce come un \textbf{proxy} frapposto tra il client e l'operazione remota potenzialmente instabile. Tale proxy monitora l'esito delle invocazioni recenti ed è modellabile come una \textbf{Macchina a Stati Finiti (FSM)} caratterizzata da tre stati operativi: \textbf{Closed}, \textbf{Open} e \textbf{Half-Open}.

% [IMMAGINE DEL FILE DI SLIDE CON NOME "CloudComputing.pdf"] Diagramma di transizione degli stati della macchina a stati finiti del Circuit Breaker: Closed, Open e Half-Open con relativi contatori e timer — Slide 9 / Pagina 9 del PDF

\begin{theorem}[Stati Operativi del Circuit Breaker]
Il ciclo di vita del Circuit Breaker è descritto formalmente dalle seguenti configurazioni e regole di transizione:
\begin{itemize}
    \item \textbf{Closed (Chiuso):} Il circuito è a riposo e consente il normale passaggio delle richieste applicative verso il servizio di destinazione. Il proxy tiene traccia del conteggio dei fallimenti recenti avvenuti entro una finestra temporale prefissata. All'ingresso nello stato, il contatore dei fallimenti viene azzerato (`reset failure counter`). Se una chiamata fallisce, il contatore viene incrementato. Fintanto che il tasso o numero di errori rimane al di sotto di una determinata soglia critica, il circuito resta \emph{Closed}. Se la soglia di fallimento viene raggiunta o superata (\emph{Failure threshold reached}), il circuito scatta e transita istantaneamente nello stato \textbf{Open}.
    \item \textbf{Open (Aperto):} In questo stato, il circuito intercetta tutte le invocazioni in arrivo senza inoltrarle alla destinazione: la richiesta fallisce immediatamente (`return failure`) sollevando un'eccezione o restituendo una risposta di cortesia predefinita (\emph{fallback}). Contestualmente all'ingresso nello stato \emph{Open}, viene avviato un contaminuti di sicurezza (`start timeout timer`). Lo scopo di questo timer è accordare al servizio bersaglio una finestra di quiete protetta per consentirne il ripristino. Fintanto che il timer scorre, nessun pacchetto applicativo raggiunge la risorsa degradata. Alla scadenza del timer (\emph{Timeout timer expired}), il circuito transita nello stato intermedio \textbf{Half-Open}.
    \item \textbf{Half-Open (Semi-Aperto):} Il circuito opera in modalità di collaudo e autorizza il passaggio di un contingente strettamente limitato di richieste applicative esplorative verso il servizio downstream. All'ingresso, il contatore dei successi viene azzerato (`reset success counter`). Se le richieste pilota vengono elaborate con successo, il contatore di successo viene incrementato; al raggiungimento di una determinata soglia favorevole (\emph{Success count threshold reached}), il sistema deduce che la causa di malfunzionamento è stata risolta e transita nuovamente nello stato ordinario \textbf{Closed}, resettando integralmente i contatori. Se viceversa anche una singola operazione esplorativa dovesse fallire (\emph{Operation failed}), il circuito deduce che l'anomalia persiste: transita immediatamente a ritroso verso lo stato \textbf{Open} e fa ripartire il contaminuti di attesa.
\end{itemize}
\end{theorem}

L'esistenza dello stato \emph{Half-Open} risulta cruciale per salvaguardare il servizio convalescente: se il circuito passasse direttamente da \emph{Open} a \emph{Closed}, l'improvviso arrivo a valanga di tutti i carichi concorrenti accumulati rischierebbe di mandare nuovamente in crash istantaneo la risorsa appena riavviata.

\subsubsection{Distinzione Concettuale Rispetto Al Retry Pattern}

È essenziale non confondere gli scopi ingegneristici dei pattern di resilienza:
\begin{itemize}[noitemsep]
    \item Il \textbf{Retry Pattern} si propone di \emph{reiterare} un'operazione nell'aspettativa fiduciosa che l'errore sia un'anomalia microscopica passeggera.
    \item Il \textbf{Circuit Breaker Pattern} agisce in senso diametralmente opposto, mirando a \emph{impedire} l'esecuzione di un'operazione la cui probabilità di insuccesso è stimata pressoché certa.
\end{itemize}

I due pattern possono essere composti architetturalmente: un'applicazione può utilizzare il Retry per invocare una procedura protetta internamente da un Circuit Breaker. Tuttavia, la logica del retry deve essere istruita per riconoscere le eccezioni originate dall'interruzione rapida del circuito aperto, arrestando immediatamente i tentativi di reiterazione per non invalidare l'obiettivo del salvavita.

\subsubsection{Dinamiche Operative E Strategie Di Adattamento}

L'adozione del Circuit Breaker conferisce stabilità sistemica durante i guasti e minimizza il degrado prestazionale, abbattendo la latenza percepita dal client grazie al \emph{fail-fast}. Se il circuito genera notifiche o eventi a ogni transizione di stato, tali segnali possono essere inviati alle piattaforme di monitoraggio e \emph{observability} per generare allarmi operativi in tempo reale.

Il pattern supporta inoltre strategie avanzate di adattamento parametrico:
\begin{itemize}
    \item \textbf{Timeout Dinamici Crescenti:} il tempo di permanenza nello stato \emph{Open} può essere incrementato in modo esponenziale o progressivo (ad esempio da pochi secondi iniziali a svariati minuti se i tentativi in \emph{Half-Open} falliscono ripetutamente).
    \item \textbf{Gestione del Fallback:} anziché limitarsi a sollevare un'eccezione, il circuito nello stato \emph{Open} può attivare logiche di ripiego restituendo risposte di default statiche o recuperando dati parziali dalla memoria cache locale, offrendo all'utente finale un degrado controllato della funzionalità (\emph{graceful degradation}).
\end{itemize}

\subsubsection{Problematiche Realizzative E Considerazioni Pratiche}

L'ingegnerizzazione di un Circuit Breaker in ambienti di produzione su larga scala comporta la risoluzione di complessi problemi operativi:

\begin{itemize}
    \item \textbf{Concorrenza Elevata:} Il medesimo circuito può essere attraversato da decine di migliaia di richieste parallele. La gestione atomica dei contatori di stato e delle transizioni della FSM non deve introdurre sezioni critiche bloccanti o degradare la reattività del sistema attraverso \emph{lock contention}.
    \item \textbf{Accelerated Circuit Breaking (Scatto Immediato):} Se il servizio remoto risponde esplicitamente con codici HTTP indicativi di saturazione severa, quali \textbf{HTTP 429} (\emph{Too Many Requests} generato da logiche di rate-limiting) o \textbf{HTTP 503} (\emph{Service Unavailable}), il circuito può essere configurato per scattare immediatamente a \emph{Open} senza attendere il superamento della soglia numerica di errori. Tali risposte possono veicolare nell'header un'indicazione temporale (es. `Retry-After`), che il circuito può utilizzare per settare dinamicamente la durata del proprio timer di \emph{Open}.
    \item \textbf{Differenziazione delle Eccezioni:} I fallimenti non possiedono tutti la medesima gravità sistemica. Un'eccezione causata da timeout per rallentamento temporaneo può richiedere una tolleranza più elevata (più tentativi ammissibili prima di scattare), mentre un'eccezione di rete irreversibile o di rifiuto totale della connessione (\emph{Connection Refused}) deve forzare lo scatto con una soglia minima.
    \item \textbf{Differenziazione delle Risorse Sottostanti (Resource Differentiation):} È rischioso utilizzare un unico circuito monolitico per una risorsa che maschera partizioni o nodi indipendenti. In un database frammentato su più nodi (\emph{sharded database}), se un singolo shard è irraggiungibile mentre tutti gli altri sono sani, un circuito globale unificato finirebbe per inibire l'accesso all'intero archivio, bloccando indebitamente query destinate a nodi operativi. In tali contesti, il Circuit Breaker deve essere istanziato con granularità per-nodo o per-shard.
    \item \textbf{Testing e Probing Attivo (Health Check):} Oltre all'attesa passiva del timer per passare a \emph{Half-Open}, il circuito può effettuare periodicamente un controllo attivo (\emph{ping}) invocando uno specifico endpoint di diagnosi esposto dal servizio remoto, conformemente ai dettami dell'\textbf{Health Endpoint Monitoring Pattern}.
    \item \textbf{Accodamento e Rigiocabilità delle Richieste (Replaying):} Nei casi in cui le operazioni abbiano carattere asincrono o transazionale differibile, il circuito aperto può serializzare le richieste fallite su un giornale di log persistente (\emph{journaling}), predisponendone l'esecuzione differita non appena il servizio torna operativo.
    \item \textbf{Intervento Manuale (Manual Override):} A fronte di anomalie complesse con tempi di recupero impredicibili, deve essere consentito all'amministratore di sistema di forzare manualmente lo stato del circuito (apertura preventiva durante un disservizio pianificato, o chiusura forzata con azzeramento dei contatori a manutenzione avvenuta).
    \item \textbf{Inappropriate Timeouts su Servizi Esterni:} Se un servizio esterno possiede impostazioni di timeout intrinseche eccessivamente permissive o lunghe, i thread del client possono rimanere congelati in attesa della risposta remota prima che il circuito possa conteggiare il fallimento. Se il fenomeno interessa molteplici client concorrenti, si produce una progressiva saturazione del thread pool complessivo prima che l'interruttore intervenga efficacemente.
\end{itemize}

\newpage

\section{Modelli Di Coordinamento E Disaccoppiamento Dei Messaggi}

Nelle architetture a microservizi distribuite, le transazioni di business end-to-end richiedono la cooperazione coordinata di svariati moduli funzionali indipendenti. Il governo del flusso di elaborazione e la comunicazione tra servizi sollevano questioni di scalabilità e accoppiamento architetturale fondamentali.

I tre pattern messaggistici primari analizzati sono:
\begin{itemize}[noitemsep]
    \item \textbf{Choreography Pattern:} delega alle singole entità applicative le decisioni operative sul flusso di lavoro, eliminando coordinatori monolitici.
    \item \textbf{Priority Queue Pattern:} ottimizza l'ordine di elaborazione garantendo priorità alle richieste computazionalmente o funzionalmente urgenti.
    \item \textbf{Publisher-Subscriber (Pub/Sub) Pattern:} consente la notifica asincrona di eventi a una molteplicità di destinatari eterogenei senza accoppiare i mittenti ai destinatari.
\end{itemize}

\subsection{Pattern Di Coordinamento Distribuito: Il Choreography Pattern}

In contrapposizione al modello di **Orchestrazione** (in cui un'entità centralizzata funge da "direttore d'orchestra" invocando i singoli servizi mediante chiamate sincrone punto-a-punto e tracciando rigidamente l'avanzamento dei task), il \textbf{Choreography Pattern} adotta un paradigma di collaborazione decentralizzato ed emergente. I servizi operano in piena autonomia decisionale, reagendo e scambiandosi eventi mediante un'infrastruttura di intermediazione messaggistica (\emph{message broker}) come Apache Kafka o RabbitMQ.

\subsubsection{Orchestrazione vs Coreografia: Analisi Dei Trade-Off}

Sebbene l'orchestrazione consenta una rapida visibilità sull'intero processo, essa introduce punti di debolezza significativi:
\begin{itemize}[noitemsep]
    \item \textbf{Accoppiamento Forte:} L'orchestratore deve conoscere la semantica e le responsabilità specifiche di ciascun servizio; qualunque variazione o aggiunta di un componente richiede la ricompilazione o riconfigurazione del flusso dell'orchestratore centrale.
    \item \textbf{Bottleneck Prestazionale e SPOF:} Il nodo orchestratore rischia di tramutarsi nel collo di bottiglia dell'intera architettura e in un \emph{Single Point of Failure}.
\end{itemize}

Nel modello coreografato, invece, i servizi interagiscono conformemente a un principio asincrono \emph{event-driven}:

% [IMMAGINE DEL FILE DI SLIDE CON NOME "CloudComputing.pdf"] Confronto visivo tra Orchestrator centralizzato (chiamate a Service A, B, C) e Choreography decentralizzata tramite Message Broker — Slide 17 / Pagina 17 del PDF

\begin{enumerate}
    \item Una richiesta client iniziale deposita un messaggio o evento primario su una coda o topic del broker.
    \item I servizi sottoscritti a quella tipologia di evento ne ricevono copia, avviano l'elaborazione locale e producono un nuovo evento attestante l'esito (positivo o negativo) dell'operazione.
    \item Altri componenti a valle registrati su tali eventi secondari recepiscono la notifica e intraprendono il passo logico successivo della transazione, perpetuando il flusso senza che alcuna entità mantenga uno stato centralizzato globale.
    \item Qualora si verifichino guasti temporanei, la logica del broker di messaggi o dei servizi può tentare la riesecuzione (\emph{retry}) del task fallito.
\end{enumerate}

\subsubsection{Casi D'Uso Indicati}

Il Choreography Pattern è particolarmente indicato nei seguenti scenari:
\begin{itemize}[noitemsep]
    \item Sistemi soggetti a frequente evoluzione funzionale, in cui è necessario aggiungere, aggiornare o rimuovere microservizi con minimo impatto e senza riprogrammare processi centralizzati.
    \item Architetture in cui l'orchestratore centrale genera saturazioni critiche della larghezza di banda o ritardi di accodamento.
    \item Ambienti basati su \textbf{Serverless Computing} e architetture FaaS (\emph{Function-as-a-Service}), dove i servizi sono istanze effimere e prive di stato avviate unicamente all'occorrenza di un evento scatenante e terminate non appena elaborato il payload.
\end{itemize}

\subsubsection{Complessità Operative E Considerazioni Architetturali}

La decentralizzazione del controllo introduce problematiche non banali nella gestione dei flussi applicativi:
\begin{itemize}
    \item \textbf{Gestione Sequenziale Rigida:} L'orchestrazione parallela di processi indipendenti è semplice, ma coordinare transazioni che esigono un ordine strettamente sequenziale (ad esempio: \emph{Service C} può elaborare solo dopo che \emph{Service A} e \emph{Service B} hanno terminato entrambi con successo) impone l'articolazione di molteplici topic dedicati o logiche di sincronizzazione tra broker.
    \item \textbf{Tracciamento Distribuito (Distributed Tracing):} All'aumentare esponenziale del numero di servizi indipendenti, ricostruire l'albero di esecuzione e identificare la causa scatenante di un'anomalia richiede l'adozione sistematica di identificatori di correlazione globali iniettati negli header dei messaggi.
    \item \textbf{Gestione del Fallimento e Transazioni Compensative:} In assenza di un coordinatore che orchestra il \emph{rollback}, la gestione dell'errore è interamente delegata alla cooperazione distribuita. Quando un task fallisce, il microservizio deve emettere un evento di errore; i componenti che avevano preventivamente eseguito passi con successo devono intercettare tale evento ed eseguire \textbf{transazioni compensative} (\emph{compensating transactions}) per annullare gli effetti collaterali già consolidati (conformemente al pattern \textbf{Saga}). Se il componente guasto fallisce prima di poter emettere l'evento di errore, il sistema deve affidarsi a meccanismi di rilevazione per timeout.
    \item \textbf{Carico di Resilienza Distribuito:} La responsabilità di governare retry, guasti transitori e scadenze ricade per intero su ogni singolo servizio, aumentando la complessità del codice di sviluppo rispetto all'incapsulamento in un orchestratore centrale.
\end{itemize}

\subsubsection{Esempio Realistico: Sistema Di Drone Delivery E Dead-Letter Queue}

A titolo esemplificativo, si consideri la gestione logistica di un sistema di consegna merci tramite droni (\emph{Drone Delivery app}).

% [IMMAGINE DEL FILE DI SLIDE CON NOME "CloudComputing.pdf"] Diagramma applicativo di Choreography per Drone Delivery con microservizi Package, Drone Scheduler, Delivery e relative Dead-Letter Queues (DLQ) — Slide 21 / Pagina 21 del PDF

La transazione di accettazione di una consegna richiede l'interazione tra tre microservizi disaccoppiati:
\begin{itemize}[noitemsep]
    \item \textbf{Package Service:} responsabile della creazione fisica o dell'aggiornamento dell'ordine del pacco;
    \item \textbf{Drone Scheduler Service:} incaricato di allocare un drone idoneo alla missione e riservarne la capacità di volo;
    \item \textbf{Delivery Service:} deputato alla verifica e al monitoraggio in tempo reale dello stato di avanzamento della spedizione.
\end{itemize}

In luogo di un orchestratore, le richieste di prenotazione immesse dal client transitano per un modulo di ingresso (\emph{Ingestion}) che notifica i messaggi su un bus condiviso. Ciascun microservizio reagisce in modo autonomo producendo a sua volta messaggi per il servizio successivo.

In tale architettura gioca un ruolo fondamentale il meccanismo della \textbf{Dead-Letter Queue (DLQ)}, detta anche coda dei messaggi non recapitabili o non elaborabili:

\begin{theorem}[Dead-Letter Queue - DLQ]
Una \textbf{Dead-Letter Queue} è una struttura di stoccaggio ausiliaria mantenuta dai gestori di messaggistica (\emph{queue manager}) in cui vengono instradati e confinati tutti i messaggi che non possono essere recapitati alla destinazione ordinaria (ad esempio per assenza della coda di destinazione, saturazione dello spazio disponibile o errori di conversione e corruzione del tracciato dati), oppure che non possono essere elaborati dal consumatore per il verificarsi reiterato di eccezioni fatali.
\end{theorem}

L'instradamento in DLQ impedisce che un messaggio velenoso (\emph{poison message}) blocchi all'infinito la testa della coda applicativa principale; il payload corrotto o non gestibile viene isolato per consentire agli operatori l'analisi diagnostica e il reinoltro a posteriori.

\newpage

\subsection{Pattern Di Diffusione Dei Messaggi: Il Publisher-Subscriber Pattern}

Nelle infrastrutture cloud su larga scala, la propagazione unidirezionale di notifiche o dati applicativi verso una molteplicità di componenti indipendenti non può essere scalata allocando code dedicate punto-a-punto per ogni singolo consumatore. Il \textbf{Publisher-Subscriber Pattern} (abbreviato in \textbf{Pub/Sub}) realizza un disaccoppiamento completo a livello spaziale, temporale e sincronico tra i produttori dell'informazione (\emph{Publisher}) e i fruitori della stessa (\emph{Subscriber}).

\subsubsection{Architettura Di Riferimento}

% [IMMAGINE DEL FILE DI SLIDE CON NOME "CloudComputing.pdf"] Schema architetturale Pub/Sub: Publisher, Input Channel, Message Broker, Output Channels e molteplici Subscriber — Slide 23 / Pagina 23 del PDF

L'infrastruttura di un sottosistema Pub/Sub si compone di:
\begin{enumerate}
    \item \textbf{Canale di Ingresso (Input Channel):} il mittente (publisher) impacchetta l'informazione in un messaggio strutturato secondo un formato condiviso e lo immette sul canale di input senza dover conoscere l'identità, il numero o la locazione fisica dei consumatori.
    \item \textbf{Intermediario di Messaggistica (Message Broker / Event Bus):} entità infrastrutturale dotata di logiche di instradamento ad altissime prestazioni che riceve il messaggio dal canale d'ingresso e ne duplica il riferimento indirizzandolo ai canali di uscita corrispondenti a tutti e soli i sottoscrittori interessati.
    \item \textbf{Canali di Uscita (Output Channels):} una collezione di code o canali logici associati a ciascun sottoscrittore (subscriber) da cui il ricevente preleva i messaggi asincronamente.
\end{enumerate}

\subsubsection{Benefici Architetturali Del Pub/Sub}

L'adozione del modello a pubblicazione/sottoscrizione offre rilevanti vantaggi sistemici:
\begin{itemize}
    \item \textbf{Disaccoppiamento dei Sottosistemi:} I moduli possono essere gestiti, riavviati e riscalati in modo del tutto indipendente. Se uno o più subscriber risultano temporaneamente offline, i messaggi rimangono memorizzati nel broker in attesa della riconnessione del nodo.
    \item \textbf{Incremento della Scalabilità e Reattività del Mittente:} Il publisher deposita il messaggio sul canale d'ingresso in una frazione di millisecondo e torna immediatamente alle proprie elaborazioni primarie, delegando al broker l'onere dell'inoltro concorrente.
    \item \textbf{Elaborazione Differita o Programmata (Deferred Processing):} I consumatori non devono elaborare i dati in tempo reale, ma possono drenare le code durante le ore di minor carico (\emph{off-peak}), attuando politiche di livellamento dei carichi (\emph{load leveling}).
    \item \textbf{Interoperabilità Eterogenea:} Sistemi sviluppati su linguaggi, runtime e piattaforme operative eterogenee (nonché scenari ibridi on-premises e cloud pubblico) possono interoperare scambiandosi payload neutri (es. JSON, Avro, Protobuf).
    \item \textbf{Collaudabilità e Manutenibilità:} Le interfacce a canale semplificano il monitoraggio di rete e il tracciamento dei payload ai fini di audit, ispezione e test di integrazione senza interferire con la logica produttiva.
\end{itemize}

\subsubsection{Considerazioni Realizzative E Vincoli Di Progettazione}

La corretta implementazione del pattern Pub/Sub in ambienti mission-critical richiede l'analisi di vincoli rigorosi:

\begin{itemize}
    \item \textbf{Utilizzo di Tecnologie Consolidate:} È sconsigliato sviluppare sottosistemi di messaggistica personalizzati \emph{in-house}. È opportuno adottare soluzioni gestite erogate dai \emph{cloud provider} (ad esempio **Azure Service Bus**, **Azure Event Hubs**, **Azure Event Grid**) oppure prodotti di consolidata maturità architetturale (quali **Apache Kafka**, **RabbitMQ**, **Apache ActiveMQ**, **NATS** o **Redis Pub/Sub**).
    \item \textbf{Filtraggio e Selezione dei Messaggi (Message Subsetting):} Raramente un subscriber necessita dell'intero volume di messaggi prodotto. Il broker deve consentire il restringimento dei dati recapitati mediante:
    \begin{itemize}[noitemsep]
        \item \emph{Filtraggio per Argomenti (Topics):} ogni topic definisce un canale di trasmissione tematico dedicato, a cui i consumatori possono iscriversi selettivamente;
        \item \emph{Filtraggio per Contenuto (Content Filtering):} il broker ispeziona i metadati o il payload del messaggio, instradandolo verso il sottoscrittore solo se soddisfa determinati predicati logici definiti dal consumer;
        \item \emph{Sottoscrizioni con Caratteri Jolly (Wildcard Subscribers):} possibilità per i sottoscrittori di legarsi a insiemi gerarchici di topic tramite wildcard (es. `ordini.europa.*`).
    \end{itemize}
    \item \textbf{Sicurezza degli Accessi:} Le connessioni ai canali devono essere governate da stringenti criteri di autenticazione e autorizzazione (es. TLS mutuo, SAS token, RBAC) per scongiurare intercettazioni illecite (\emph{eavesdropping}) o iniezioni abusive di messaggi.
    \item \textbf{Unidirezionalità e Comunicazione Bidirezionale:} I canali Pub/Sub operano strutturalmente a senso unico (\emph{unidirectional}). Se un subscriber deve riscontrare l'elaborazione o inviare un esito di stato al publisher originale, l'interazione deve essere integrata con il **Request/Reply Pattern**, il quale alloca un canale dedicato aggiuntivo di risposta (\emph{reply channel}) tracciato tramite correlazione di ID.
    \item \textbf{Ordinamento Dei Messaggi E Idempotenza:} Nelle reti distribuite su larga scala, \textbf{l'ordine di ricezione dei messaggi da parte delle istanze consumer non è garantito}, potendo differire dall'ordine temporale esatto di emissione. L'elaborazione applicativa deve pertanto essere rigorosamente progettata per essere \textbf{idempotente}:
\end{itemize}

\begin{theorem}[Idempotenza dell'Elaborazione]
Un'operazione $f$ si definisce \textbf{idempotente} se l'effetto prodotto dall'applicazione ripetuta della medesima operazione su uno stato di sistema $x$ equivale a quello prodotto dalla sua prima e singola esecuzione:
\[
\boxed{f(f(x)) = f(x)}
\]
\end{theorem}

L'introduzione dell'idempotenza (ad esempio scartando chiavi transazionali duplicate archiviate in cache distribuita o impostando aggiornamenti di stato condizionali) garantisce l'integrità e la coerenza logica dei dati a fronte di recapiti non sequenziali, consegne multiple (\emph{at-least-once delivery}) o ripetizioni forzate da meccanismi di retry di rete.

% AA -----------------------------------------------------------------------------------------------------------

\section{Approfondimenti Sui Modelli Di Messaggistica: Pub/Sub E Code Di Priorità}

Nei sistemi distribuiti complessi e nelle architetture cloud-native a microservizi, il disaccoppiamento tra componenti mediante l'intermediazione di messaggi asincroni introduce sfide operative rilevanti. Se da un lato l'infrastruttura di trasmissione protegge i singoli servizi dal blocco sincronico e assorbe i picchi di carico, dall'altro impone la gestione rigorosa di anomalie quali messaggi malevoli, duplicazioni da guasti transitori, scadenze temporali e differenziazione dei livelli di servizio.

\subsection{Problematiche Avanzate Nel Publisher-Subscriber Pattern}

L'adozione del pattern \textbf{Publisher-Subscriber (Pub/Sub)} comporta una serie di criticità intrinseche legate alla natura distribuita dei canali di comunicazione, che gli sviluppatori e gli architetti del software devono affrontare mediante adeguate contromisure infrastrutturali e logiche.

\subsubsection{Priorità Dei Messaggi E Messaggi Veliniferi}

All'interno di un flusso uniforme di eventi, non tutte le informazioni possiedono la medesima urgenza o la stessa validità sintattica e semantica:
\begin{itemize}
    \item \textbf{Message Priority (Priorità del Messaggio):} determinate soluzioni aziendali impongono che specifici messaggi vengano ricevuti ed elaborati prima di altri, alterando il canonico ordine di arrivo. In tali circostanze, l'infrastruttura Pub/Sub deve cooperare con il \textbf{Priority Queue Pattern}, il quale fornisce il substrato architetturale per garantire la consegna privilegiata delle comunicazioni ad alta rilevanza.
    \item \textbf{Poison Messages (Messaggi Veliniferi o "Avvelenati"):} un messaggio strutturato in modo non conforme (\emph{malformed message}) oppure un'operazione che richieda l'accesso a risorse non disponibili può causare il fallimento sistematico o il crash dell'istanza di servizio consumatrice. Se il sistema rimettesse ciecamente il messaggio fallito in cima alla coda ordinaria, si innescherebbe un ciclo infinito di fallimento e riavvio (\emph{crash loop}). L'architettura deve impedire il reinserimento di tali payload tossici: i dettagli del messaggio devono essere catturati e archiviati separatamente per consentire l'ispezione analitica da parte degli amministratori. Piattaforme di livello enterprise, come \textbf{Azure Service Bus}, forniscono nativamente questa salvaguardia attraverso la funzionalità di \textbf{Dead-Letter Queue (DLQ)}.
\end{itemize}

\subsubsection{Duplicazione, Scadenza E Schedulazione Temporale}

La semantica di trasporto su canali inaffidabili impone precise strategie di gestione temporale e di deduplicazione:
\begin{itemize}
    \item \textbf{Repeated Messages (Messaggi Ripetuti):} in un ambiente cloud, lo stesso messaggio può essere pubblicato più volte. Ad esempio, il mittente potrebbe depositare con successo l'evento sul broker ma subire un crash o un'interruzione di rete un istante prima di registrarne l'avvenuto invio; al riavvio, una nuova istanza del mittente re-inoltrerà il medesimo messaggio. Per conseguire una semantica di consegna \emph{at-most-once} o prevenire elaborazioni multiple, l'infrastruttura di messaggistica deve implementare il rilevamento e la rimozione dei duplicati (\emph{de-duplication}) confrontando gli identificatori univoci dei messaggi (\emph{Message ID}). Qualora il broker sottostante non supporti nativamente il \emph{de-duping}, è vincolo assoluto che la logica di consumo applicativa sia rigorosamente \textbf{idempotente}.
    \item \textbf{Message Expiration (Scadenza dei Messaggi / Time-To-Live):} un'informazione ha spesso una validità temporale finita. Se un messaggio non viene elaborato entro una finestra temporale utile (a causa di code congestionate o rallentamenti a valle), perde di significato e la sua esecuzione risulterebbe dannosa o inutile. Il mittente specifica quindi un tempo di scadenza nei metadati del payload; l'entità ricevente (o il broker stesso) verifica tale campo scartando il messaggio qualora sia scaduto.
    \item \textbf{Message Scheduling (Schedulazione dei Messaggi):} talvolta un messaggio deve essere temporaneamente posto sotto vincolo di attesa (\emph{embargoed}), impedendone la lavorazione fino al raggiungimento di una specifica data e ora programmata. Il middleware trattiene il messaggio e lo rende visibile ai consumatori solo a partire dall'istante prestabilito.
\end{itemize}

\subsubsection{Contesti Di Applicabilità Del Pub/Sub}

La scelta di adottare il pattern Pub/Sub deve scaturire da un'attenta valutazione dei requisiti di integrazione. Il pattern risulta vantaggioso quando:
\begin{itemize}[noitemsep]
    \item L'applicazione deve diffondere broadcast la medesima informazione a una popolazione cospicua di consumatori;
    \item È richiesta l'interconnessione tra sistemi eterogenei, sviluppati in modo indipendente con stack tecnologici, linguaggi di programmazione o protocolli dissimili;
    \item Il mittente può disinteressarsi della risposta immediata da parte dei ricevitori, operando in modalità disaccoppiata;
    \item Il dominio applicativo tollera modelli di consistenza eventuale (\emph{eventual consistency}) per i propri archivi dati;
    \item I consumatori presentano profili di disponibilità, finestre di manutenzione e orari di attività differenti rispetto al mittente.
\end{itemize}

Al contrario, il pattern è sconsigliato quando:
\begin{itemize}[noitemsep]
    \item L'applicazione si interfaccia con un numero molto ristretto di consumatori che richiedono ciascuno dati strutturalmente e completamente diversi;
    \item L'interazione applicativa esige requisiti di risposta sincrona in tempo reale (\emph{real-time} o \emph{near real-time}).
\end{itemize}

\newpage

\subsection{Pattern Di Prioritizzazione Delle Richieste: Il Priority Queue Pattern}

Nelle architetture cloud, le applicazioni delegano tipicamente i compiti computazionalmente intensivi a processi in background mediante l'interposizione di code. Sebbene nella maggioranza dei casi l'ordine di arrivo dei task sia ininfluente, in molteplici contesti di business è mandatorio garantire che determinate richieste ricevano un trattamento preferenziale rispetto a compiti a bassa priorità già accodati in precedenza.

Il \textbf{Priority Queue Pattern} si prefigge di ordinare e distribuire i task in modo tale che le richieste caratterizzate da un livello di priorità superiore vengano ricevute e processate più rapidamente rispetto a quelle con priorità inferiore. Tale schema è essenziale per implementare garanzie differenziate nei \textbf{Service Level Agreement (SLA)} stipulati con i clienti.

\subsubsection{Coda Singola Con Riordinamento Interno}

Le code di messaggistica ordinarie operano secondo la classica semantica FIFO (\emph{First-In, First-Out}), in cui il consumatore riceve gli elementi esattamente nella sequenza cronologica di immissione. Alcuni sistemi di messaggistica evoluti integrano tuttavia il supporto nativo per la prioritizzazione all'interno della stessa struttura dati.

% [IMMAGINE DEL FILE DI SLIDE CON NOME "CloudComputing"] Schema architetturale del Priority Queue pattern con coda singola che riordina i messaggi da 1 (priorità bassa) a 3 (priorità alta) verso molteplici consumatori — Slide 5 / Pagina 5 del PDF

Nel modello a coda singola:
\begin{enumerate}
    \item L'applicazione client, all'atto dell'inserimento (\emph{enqueue}), assegna a ogni messaggio un indicatore numerico o logico di priorità (ad esempio, con scala da $1$ per la priorità minima a $3$ per la priorità massima).
    \item Il middleware provvede automaticamente a riorganizzare i puntatori interni della coda, spostando i messaggi a priorità più elevata verso la testa del buffer di lettura.
    \item I nodi consumatori (\emph{Consumer}), prelevando dalla testa, ricevono ed elaborano le richieste urgenti prima di quelle a bassa priorità, indipendentemente dall'istante di arrivo originario.
\end{enumerate}

Tale meccanismo apporta benefici evidenti:
\begin{itemize}[noitemsep]
    \item Soddisfa i vincoli di business differenziando la disponibilità e le prestazioni per gruppi specifici di utenti;
    \item Minimizza i costi di esercizio: in condizioni di budget o risorse limitate, è possibile ridurre il numero complessivo di consumatori. I messaggi ad alta priorità continueranno a essere evasi prima, mentre quelli a bassa priorità sconteranno ritardi maggiori;
    \item Ottimizza le risorse partizionando l'esecuzione: i task critici vengono serviti tempestivamente, mentre le lavorazioni accessorie possono essere smaltite nelle fasce orarie a minor traffico.
\end{itemize}

\subsubsection{Topologia A Code Multiple E Assegnazione Dei Pool}

Molti broker di messaggistica non implementano internamente code con riordinamento prioritario. L'alternativa architetturale consiste nell'istanziare code distinte per ciascun livello di priorità.

% [IMMAGINE DEL FILE DI SLIDE CON NOME "CloudComputing"] Diagramma a code multiple separate per priorità con pool di consumatori indipendenti dimensionati asimmetricamente — Slide 6 / Pagina 6 del PDF

In questa configurazione, l'onere di instradamento ricade sull'applicazione mittente, che indirizza il messaggio al canale dedicato in base alla sua classe di appartenenza (es. \emph{Priority 1}, \emph{Priority 2}, \emph{Priority 3}). Esistono due varianti principali per l'organizzazione dei consumatori:

\begin{itemize}
    \item \textbf{Pool di Consumatori Dedicati e Asimmetrici:} ogni coda dispone di un gruppo indipendente di istanze consumatrici. Per le code ad alta priorità si alloca un pool di consumatori più cospicuo ed eseguito su istanze di calcolo dotate di hardware più performante (più vCPU e memoria); per le code a bassa priorità si impiega un pool ristretto con risorse minime, o se ne sospende del tutto l'ascolto nei periodi di picco.
    \item \textbf{Pool Unico con Polling Gerarchico:} un singolo insieme di consumatori controlla sequenzialmente le code, interrogando prima la coda a priorità massima e procedendo a svuotare le code secondarie unicamente quando i canali più urgenti sono privi di messaggi.
\end{itemize}

\begin{table}[H]
\centering
\caption{Confronto semantico tra Pool Singolo e Pool Multipli nel Priority Queue Pattern}
\label{tab:priority_queue_comparison}
\begin{tabularx}{\textwidth}{p{3.5cm} X X}
\toprule
\textbf{Caratteristica} & \textbf{Approccio a Pool Singolo} & \textbf{Approccio a Pool Multipli} \\
\midrule
\textbf{Precedenza di Servizio} & I messaggi ad alta priorità scavalcano sempre in modo assoluto i messaggi secondari. & I messaggi a priorità più alta sono evasi più celermente, ma le code minori avanzano parallelamente. \\
\addlinespace
\textbf{Rischio di Starvation} & Elevato: se il flusso di richieste ad alta priorità è ininterrotto, i messaggi a bassa priorità rischiano di non essere mai elaborati. & Nullo o controllato: purché vi sia almeno un consumatore attivo assegnato alla coda secondaria, i task avanzano con continuità. \\
\addlinespace
\textbf{Complessità di Scheduling} & Richiede logiche complesse di preemption o algoritmi di scansione gerarchica tra le code. & Semplice e disaccoppiata: ciascun pool consuma indipendentemente con logica FIFO locale. \\
\addlinespace
\textbf{Dimensionamento Hardware} & Hardware omogeneo condiviso per tutti i tipi di carico. & Possibilità di allocare macchine con profili computazionali mirati al tipo di coda. \\
\bottomrule
\end{tabularx}
\end{table}

\subsubsection{Considerazioni Di Progettazione E Autoscaling}

La formalizzazione del pattern richiede la definizione di parametri quantitativi stringenti:
\begin{itemize}
    \item \textbf{Definizione Contestuale della Priorità:} la priorità deve riflettere metriche temporali oggettive. Ad esempio, assegnare "alta priorità" deve tradursi formalmente nel vincolo ingegneristico che il messaggio venga completato entro $10\,\text{secondi}$ dall'immissione.
    \item \textbf{Preemption:} se si adotta un pool singolo e sopraggiunge una richiesta prioritaria mentre un consumatore è impegnato su un task ordinario lungo, occorre prevedere se sia necessario sospendere o interrompere (\emph{preempt}) il task in esecuzione per liberare la risorsa.
    \item \textbf{Mitigazione della Starvation (Invecchiamento Dinamico):} nei sistemi che gestiscono internamente la prioritizzazione, si può mitigare la \emph{starvation} dei compiti a bassa priorità introducendo una politica di invecchiamento (\emph{aging}): la priorità del messaggio viene incrementata dinamicamente in funzione del tempo trascorso in coda.
    \item \textbf{Criteri Logici di Segmentazione:} la priorità può discendere da logiche di business non esplicitate dal mittente ma dedotte dal profilo utente (ad esempio differenziando tra clienti paganti \emph{"fee-paying"} e utenti gratuiti \emph{"non-fee paying"}).
    \item \textbf{Costi Economici delle Query:} molti servizi cloud commerciali (come AWS SQS o Azure Queue Storage) addebitano costi operativi per ogni singola operazione di invocazione o interrogazione (\emph{polling}). La scansione ripetuta di molteplici code vuote moltiplica il traffico di rete e i costi di fatturazione mensile.
    \item \textbf{Autoscaling Dinamico:} il dimensionamento dei pool di consumatori deve adattarsi reattivamente. Monitorando la velocità di svuotamento e la lunghezza della coda (\emph{queue depth}), i motori di orchestrazione (come l'Horizontal Pod Autoscaler di Kubernetes basato su metriche custom) possono scalare orizzontalmente le repliche dei consumatori all'aumentare dei task accumulati.
\end{itemize}

\newpage

\subsection{Esempio Realizzativo: Code Di Priorità Con Redis Ed RQ}

Un'implementazione pratica del pattern a code multiple può essere realizzata sfruttando un'infrastruttura di memorizzazione in memoria con **Redis** e il framework di gestione code **RQ (Redis Queue)** in ambiente Python.

Il codice seguente illustra la suddivisione tra le code logiche e la delega condizionale in base alla classe di servizio del cliente:

\begin{lstlisting}[language=python, caption={Implementazione del Priority Queue Pattern con Redis ed RQ}]
import time
from redis import Redis
from rq import Queue

# 1. Connessione all'infrastruttura di messaggistica Redis
redis_conn = Redis(host='localhost', port=6379)

# 2. Definizione delle code logiche separate per livello di priorita'
#    high_priority_queue gestisce il canale urgente
#    low_priority_queue gestisce i task ordinari di background
high_priority_queue = Queue('high', connection=redis_conn)
low_priority_queue = Queue('low', connection=redis_conn)

def process_transaction(user_type, amount):
    """
    Simulazione del compito di backend eseguito dal worker
    """
    print(f"[worker] Processing {user_type} transaction of ${amount}...")
    time.sleep(2)  # Simula l'onere computazionale dell'elaborazione
    print(f"[worker] Finished {user_type} transaction.")

def incoming_request_handler(user_type, amount):
    """
    Produttore (Gateway): instrada il task verso la coda opportuna
    in funzione della sottoscrizione dell'utente
    """
    if user_type == "PREMIUM":
        # Inserimento nella coda ad alta priorita'
        high_priority_queue.enqueue(process_transaction, user_type, amount)
        print(f"[Gateway] Sent PREMIUM request to HIGH priority queue.")
    else:
        # Inserimento nella coda a bassa priorita'
        low_priority_queue.enqueue(process_transaction, user_type, amount)
        print(f"[Gateway] Sent FREE request to LOW priority queue.")

if __name__ == '__main__':
    # Simulazione della ricezione asincrona
    # Anche se la richiesta FREE perviene prima cronologicamente...
    incoming_request_handler("FREE", 10)

    # ...la richiesta PREMIUM inviata successivamente dovra' essere privilegiata
    incoming_request_handler("PREMIUM", 500)
\end{lstlisting}

\subsubsection{Dinamica Operativa Dei Worker}

La separazione logica in due code distinte (`high` e `low`) non è sufficiente se il demone di esecuzione dei task non viene istruito circa l'ordine gerarchico di prelievo. Il comportamento differenziale si determina avviando il processo worker con la sintassi:

\begin{lstlisting}[language=bash, caption={Avvio del processo worker con priorita' gerarchica}]
rq worker high low
\end{lstlisting}

Il comando impone al worker una specifica direttiva di consumo:
\begin{enumerate}
    \item Il processo effettua preliminarmente il polling e lo scaricamento di tutti i task presenti sulla coda indicata per prima (`high`).
    \item Solo nel momento in cui la coda `high` risulta completamente svuotata, il worker preleva i compiti accumulati nella coda secondaria (`low`).
    \item Se durante l'elaborazione di un job secondario o al termine dello stesso sopraggiungono nuovi messaggi nella coda `high`, il worker ritorna immediatamente a processare la coda prioritaria.
\end{enumerate}

In questo modo, la transazione `PREMIUM` simulata nello script Python verrà sistematicamente acquisita ed evasa dal sistema prima della richiesta `FREE`, realizzando concretamente l'isolamento prestazionale e il rispetto delle gerarchie operative prefissate.

% BB -----------------------------------------------------------------------------------------------------------

% ==============================================================================
% CAPITOLO: VIRTUALIZZAZIONE E TECNOLOGIE DI CONTAINERIZZAZIONE
% ==============================================================================

\section{Fondamenti Di Virtualizzazione E Requisiti Architetturali}

\subsection{Dal Modello Monolitico Alla Virtualizzazione Delle Risorse}

Nel paradigma computazionale classico antecedente l'avvento della virtualizzazione moderna, ciascun calcolatore fisico ospitava una singola istanza di sistema operativo monolitico. In questo assetto, l'intero spettro delle risorse hardware --- l'unità centrale di elaborazione (\emph{Central Processing Unit}, CPU), la memoria primaria ad accesso casuale (\emph{Random Access Memory}, RAM), i canali di ingresso/uscita (I/O) e i dispositivi di memorizzazione secondaria --- era posto sotto l'autorità esclusiva e diretta del \emph{kernel} di quell'unico sistema operativo. Tale configurazione comportava forti inefficienze strutturali:
\begin{itemize}
    \item \textbf{Sotto-utilizzazione cronica}: gran parte dei carichi applicativi aziendali presentava requisiti di picco concentrati in intervalli ridotti, lasciando la capacità computazionale media inutilizzata per oltre l'80\% del tempo operativo.
    \item \textbf{Rigidità di confinamento}: per prevenire conflitti a livello di librerie condivise, dipendenze binarie o arresti anomali a catena causati da singoli applicativi, gli amministratori di sistema adottavano la prassi \emph{"one-app-per-server"}, moltiplicando i costi di approvvigionamento hardware, alimentazione elettrica e manutenzione dei centri elaborazione dati.
\end{itemize}

% [IMMAGINE DEL FILE DI SLIDE CON NOME "CloudComputing_2"] Confronto architetturale tra un sistema non virtualizzato con OS monolitico e un sistema virtualizzato con Virtualization Layer e molteplici container/VM — Slide 2 / Pagina 2 del PDF

La \textbf{virtualizzazione} supera questa limitazione introducendo uno strato software intermediario posizionato direttamente tra l'hardware sottostante e i sistemi operativi guest. Questo strato disaccoppia la risorsa fisica dalla sua rappresentazione logica, realizzando una suddivisione coordinata della macchina fisica in molteplici ambienti di esecuzione isolati.
La virtualizzazione dell'elaborazione (\emph{Compute Virtualization}) si articola in tre macro-aree inscindibili:
\begin{enumerate}
    \item \textbf{Virtualizzazione della CPU}: permette a molteplici flussi di istruzioni appartenenti a differenti sistemi operativi di condividere i core fisici mediante tecniche di multiplexing temporale e partizionamento trasparente dei registri.
    \item \textbf{Virtualizzazione della memoria}: sintetizza spazi di indirizzamento contigui e isolati per ciascun ambiente guest, mappando la memoria virtuale del guest sulla memoria fisica reale della macchina ospitante.
    \item \textbf{Virtualizzazione dell'I/O, dello storage e della rete}: converte i controller fisici, le interfacce di rete (NIC) e i blocchi disco in astrazioni logiche indipendenti, instradando gli interrupt e i trasferimenti DMA (\emph{Direct Memory Access}) attraverso opportuni filtri di controllo.
\end{enumerate}

Grazie a questo strato di astrazione, ciascun sistema operativo ospitato (\emph{Guest OS}) opera sotto la perfetta illusione di esercitare il controllo totale ed esclusivo su un calcolatore fisico indipendente. In virtù di questo principio, una singola piattaforma hardware fisica può gestire carichi di lavoro eterogenei (ad esempio istanze multiple di distribuzioni Linux affiancate a istanze Microsoft Windows) oppure carichi omogenei con requisiti operativi distinti.

---

\subsection{Architettura Del Calcolatore E Set Di Istruzioni (ISA)}

Per comprendere i requisiti necessari a virtualizzare un processore, è indispensabile formalizzare il concetto di \textbf{Instruction Set Architecture (ISA)}. L'ISA costituisce il contratto fondamentale stipulato a livello di interfaccia tra il software (inclusi i compilatori e il sistema operativo) e la microarchitettura hardware della CPU. L'ISA definisce non solo l'insieme di istruzioni binarie che il processore è in grado di decodificare ed eseguire, ma anche i meccanismi con cui la CPU interagisce con la gerarchia di memoria, le periferiche di I/O e la gestione degli stati interni.

% [IMMAGINE DEL FILE DI SLIDE CON NOME "CloudComputing_2"] Diagramma a livelli dell'architettura di calcolo (Applications, Libraries, OS, ISA, Hardware) e schema strutturale del Virtual Machine Monitor — Slide 3 / Pagina 3 del PDF

I componenti strutturali primari definiti da un'ISA comprendono:
\begin{itemize}
    \item \textbf{Istruzioni in linguaggio macchina e Codici Operativi (Opcode)}: l'insieme delle operazioni primitive binarie direttamente comprensibili dall'unità di decodifica della CPU. Ogni istruzione è strutturata con un campo operativo (\emph{opcode}), che dichiara l'operazione elementare da compiere (es. manipolazioni aritmetiche, trasferimenti di dati, istruzioni di salto condizionato, operazioni di confronto logico), e da campi operando, che identificano i registri di origine, di destinazione o gli indirizzi di memoria coinvolti.
    \item \textbf{Insieme dei registri (\emph{Register File})}: la specifica dell'architettura dei registri interni del processore, includendo i registri generici (\emph{General Purpose Registers}), il puntatore alle istruzioni (\emph{Program Counter} o Instruction Pointer), il registro di stato e di flag (\emph{Condition Codes}) e i registri speciali deputati al controllo del processore (es. puntatori alle tabelle di paginazione, registri di segmento).
    \item \textbf{Modalità di indirizzamento (\emph{Addressing Modes})}: le regole sintattiche e circuitali che stabiliscono come un operando debba essere reperito in memoria:
    \begin{itemize}
        \item \emph{Indirizzamento immediato}: il valore dell'operando è incorporato direttamente all'interno della parola dell'istruzione macchina.
        \item \emph{Indirizzamento diretto}: l'istruzione codifica in modo esplicito l'indirizzo della locazione di memoria in cui risiede il dato.
        \item \emph{Indirizzamento indiretto}: l'istruzione specifica un registro che contiene il puntatore all'indirizzo di memoria effettivo (\emph{Effective Address}).
    \end{itemize}
    \item \textbf{Tipi di dati e operazioni supportate}: il catalogo dei formati elementari manipolabili nativamente dall'unità aritmetico-logica (ALU) e dall'unità di calcolo in virgola mobile (FPU), quali interi con o senza segno a 32/64 bit, numeri floating-point a singola o doppia precisione conforme allo standard IEEE 754, e istruzioni SIMD vettoriali.
    \item \textbf{Supporto alla concorrenza e al parallelismo}: costrutti a livello di istruzione progettati per la sincronizzazione hardware atomica (come le istruzioni \texttt{test-and-set}, \texttt{compare-and-swap} o le primitive \emph{load-linked/store-conditional}), indispensabili per implementare lock e semafori in ambienti multiprocessore simmetrici (SMP) e multi-threading.
    \item \textbf{Gestione di interruzioni ed eccezioni}: le procedure hardware con cui la CPU interrompe il normale flusso di esecuzione sequenziale a fronte di anomalie interne sincrone (eccezioni software, es. divisione per zero, \emph{page fault}) o segnali elettrici asincroni provenienti da dispositivi periferici (interruzioni hardware, es. arrivo di un frame di rete o completamento di una transazione disco), definendo il salvataggio automatico del contesto e la chiamata ai relativi vettori di gestione (\emph{Interrupt Service Routine}, ISR).
    \item \textbf{Modalità di esecuzione (\emph{Execution Modes})}: la distinzione gerarchica tra livelli di privilegio nell'esecuzione delle istruzioni, comunemente strutturata nella separazione tra \textbf{modalità utente} (\emph{user mode}) e \textbf{modalità privilegiata} (\emph{privileged mode} o \emph{kernel mode}).
\end{itemize}

---

\subsection{Separazione Dei Privilegi: API, Chiamate Di Sistema E User ISA}

In un sistema operativo moderno, la sicurezza, la stabilità e l'isolamento dei processi si basano su una rigida stratificazione gerarchica delle interfacce. Tale disaccoppiamento garantisce che i programmi applicativi non possano manomettere lo stato del sistema né accedere arbitrariamente all'hardware:

\begin{figure}[H]
    \centering
    \begin{tikzpicture}[>=Stealth, node distance=1.2cm, auto]
        \node [draw, rectangle, fill=backcolour, minimum width=9cm, minimum height=0.7cm] (apps) {\textbf{Applicazioni Utente}};
        \node [draw, rectangle, fill=theoryblue!20, minimum width=9cm, minimum height=0.7cm, below of=apps] (libs) {\textbf{Librerie di Sistema} (es. \texttt{glibc})};
        \node [draw, rectangle, fill=orange!20, minimum width=9cm, minimum height=0.7cm, below of=libs] (os) {\textbf{Sistema Operativo} (\emph{Kernel})};
        \node [draw, rectangle, fill=codegray!30, minimum width=9cm, minimum height=0.7cm, below of=os] (hw) {\textbf{Hardware Fisico} (CPU, RAM, Dispositivi)};

        \draw [->, thick] (apps) -- node[right, font=\footnotesize] {Chiamate API (\emph{Application Programming Interface})} (libs);
        \draw [->, thick] (libs) -- node[right, font=\footnotesize] {Chiamate di Sistema (\emph{System Calls} / Trap)} (os);
        \draw [->, thick] (os) -- node[right, font=\footnotesize] {Istruzioni Privilegiate dell'ISA} (hw);

        \draw [->, thick, dashed] (apps.west) to [bend right=45] node[left, font=\footnotesize] {User ISA (istruzioni unprivileged dirette)} (hw.west);
    \end{tikzpicture}
    \caption{Stratificazione delle interfacce dal software applicativo all'hardware.}
    \label{fig:layer_interfaces}
\end{figure}

Analizzando i differenti livelli evidenziati nello schema:
\begin{enumerate}
    \item \textbf{User ISA}: rappresenta il sottoinsieme dell'ISA accessibile quando il processore esegue codice in modalità utente. Questo insieme include operazioni aritmetico-logiche, manipolazione dei dati nei registri generici e salti condizionati locali. Le istruzioni che compongono lo User ISA non hanno facoltà di modificare lo stato globale del calcolatore, di riconfigurare la tabella delle pagine o di disabilitare gli interrupt hardware.
    \item \textbf{API (\emph{Application Programming Interface})}: costituiscono l'interfaccia ad alto livello esposta agli sviluppatori. Progettate per isolare il codice sorgente dai dettagli di basso livello del sistema operativo sottostante, le API incapsulano operazioni complesse (es. manipolazione di strutture grafiche, gestione di file system, invio di dati su socket di rete). Esse risiedono primariamente nello spazio utente all'interno di librerie condivise (ad es. la \texttt{glibc} in ambiente GNU/Linux o le Win32 API in Windows).
    \item \textbf{Chiamate di Sistema (\emph{System Calls})}: costituiscono il canale controllato attraverso cui un processo che opera in modalità utente richiede un servizio con privilegi elevati al kernel. La chiamata di sistema provoca una transizione controllata dello stato della CPU da modalità utente a modalità supervisore mediante un'istruzione di eccezione software dedicata (come \texttt{sysenter}/\texttt{sysexit} o \texttt{syscall}/\texttt{sysret} su x86). Il kernel intercetta la richiesta, verifica l'identità e i permessi del processo chiamante, esegue l'operazione (allocazione di pagine fisiche, invio di un pacchetto di rete, lettura dal controller di archiviazione) e restituisce il risultato ritornando alla modalità utente.
    \item \textbf{ISA Del Sistema Operativo}: non si riferisce strettamente a una specifica sintattica di istruzioni binarie, bensì all'architettura interna con cui il kernel governa l'accesso ai componenti del sistema. Definisce la logica algoritmica attraverso cui il sistema operativo organizza il file system, gestisce i descrittori dei processi, programma il controllore degli interrupt e implementa la memoria virtuale.
\end{enumerate}

---

\subsection{Architettura Del Virtual Machine Monitor: Dispatcher, Allocator E Interpreti}

Il perno tecnologico della virtualizzazione dell'hardware è il \textbf{Virtual Machine Monitor (VMM)}, comunemente denominato \textbf{hypervisor}. Sebbene nella terminologia odierna i due vocaboli vengano frequentemente impiegati come sinonimi, formalmente il termine \emph{VMM} pone l'accento sul modulo logico incaricato di supervisionare e isolare l'esecuzione delle macchine virtuali in esecuzione, mentre il termine \emph{hypervisor} definisce l'architettura complessiva di sistema che gestisce l'intero ciclo di vita delle risorse.

Un VMM rispetta il classico teorema di virtualizzazione formalizzato da Gerald J. Popek e Robert P. Goldberg (1974). Secondo tale formulazione, un'architettura di calcolo è puramente virtualizzabile se tutte le istruzioni \emph{sensibili} (ovvero quelle che modificano la configurazione delle risorse o il cui comportamento dipende dalla modalità di esecuzione) sono un sottoinsieme rigoroso delle istruzioni \emph{privilegiate} (ossia quelle che sollevano un trap hardware se eseguite in modalità utente).

Un Virtual Machine Monitor è strutturato internamente in tre macro-componenti coordinati:
\begin{enumerate}
    \item \textbf{Dispatcher}: costituisce il router di controllo del VMM. Ha il compito primario di decidere quale flusso di codice debba essere mandato in esecuzione e di determinare il cambio di contesto (\emph{context switch}) tra le diverse macchine virtuali o tra una macchina virtuale e l'hypervisor. Il dispatcher gestisce la sospensione e il ripristino delle VM, comportandosi analogamente a un piano di schedulazione ad alta granularità. Inoltre, quando una VM esegue un'istruzione che scatena un trap, è il dispatcher a intercettare l'evento e a smistare il controllo al modulo preposto per la corretta gestione.
    \item \textbf{Allocator}: rappresenta l'amministratore delle risorse fisiche di calcolo. Ha la responsabilità di ripartire la CPU, la memoria principale e i dispositivi di I/O tra le molteplici VM attive, garantendo assenza di conflitti e isolamento reciproco:
    \begin{itemize}
        \item \emph{Memoria fisica}: l'allocator garantisce che due macchine virtuali non possano sovrascrivere i rispettivi spazi di memoria, gestendo tabelle di traduzione tra la memoria fittizia del guest e la memoria fisica dell'host.
        \item \emph{Tempo di CPU}: applica politiche di accodamento e ripartizione equa (\emph{fair-share}) affinché nessuna VM possa monopolizzare i core di calcolo a detrimento delle altre.
        \item \emph{Canali di I/O}: serializza gli accessi al disco e smista i frame di rete in arrivo verso i rispettivi sottosistemi guest.
    \end{itemize}
    \item \textbf{Insieme degli Interpreti (\emph{Interpreter Routines})}: costituisce una collezione di routine software dedicate, associate ciascuna a una specifica istruzione privilegiata o sensibile dell'ISA. Dal momento che alle VM non è consentito eseguire direttamente istruzioni privilegiate sull'hardware nativo (per ovvi requisiti di integrità del sistema), il tentativo di esecuzione scatena un'eccezione che viene intercettata dal dispatcher e passata all'interprete competente. L'interprete simula gli effetti dell'istruzione privilegiata all'interno dello stato fittizio della VM (\emph{virtual register state}) e fa avanzare il puntatore delle istruzioni, restituendo il controllo alla macchina guest.
\end{enumerate}

\newpage

% ==============================================================================
% SEZIONE 2: TASSONOMIA DEGLI HYPERVISOR ED EVOLUZIONE DELLE TECNOLOGIE
% ==============================================================================

\section{Tassonomia Degli Hypervisor Ed Evoluzione Delle Tecnologie}

\subsection{Hypervisor Di Tipo 1 (Bare-Metal) E Tipo 2 (Hosted)}

L'implementazione del livello di virtualizzazione può avvenire a differenti altezze dello stack software, dando luogo a due modelli architetturali:

\begin{figure}[H]
    \centering
    \begin{subfigure}[b]{0.48\textwidth}
        \centering
        \begin{tikzpicture}[node distance=0.8cm, every node/.style={rectangle, draw, minimum width=5.5cm, minimum height=0.6cm, font=\small}]
            \node[fill=orange!30] (app1) {App A \quad | \quad App B};
            \node[fill=yellow!20, below of=app1] (gos1) {Guest OS};
            \node[fill=theoryblue!30, below of=gos1, yshift=-0.3cm] (hyp1) {\textbf{Hypervisor Tipo 1 (Bare-Metal)}};
            \node[fill=codegray!30, below of=hyp1] (hw1) {\textbf{Hardware Fisico}};

            \draw[<->, thick] (app1) -- (gos1);
            \draw[<->, thick] (gos1) -- (hyp1);
            \draw[<->, thick] (hyp1) -- (hw1);
        \end{tikzpicture}
        \caption{Hypervisor di Tipo 1 (Bare-Metal)}
        \label{fig:type1_hyp}
    \end{subfigure}
    \hfill
    \begin{subfigure}[b]{0.48\textwidth}
        \centering
        \begin{tikzpicture}[node distance=0.8cm, every node/.style={rectangle, draw, minimum width=5.5cm, minimum height=0.6cm, font=\small}]
            \node[fill=orange!30] (app2) {App A \quad | \quad App B};
            \node[fill=yellow!20, below of=app2] (gos2) {Guest OS};
            \node[fill=theoryblue!20, below of=gos2] (hyp2) {\textbf{Hypervisor Tipo 2 (Hosted)}};
            \node[fill=green!20, below of=hyp2] (hostos) {Host Operating System};
            \node[fill=codegray!30, below of=hostos] (hw2) {\textbf{Hardware Fisico}};

            \draw[<->, thick] (app2) -- (gos2);
            \draw[<->, thick] (gos2) -- (hyp2);
            \draw[<->, thick] (hyp2) -- (hostos);
            \draw[<->, thick] (hostos) -- (hw2);
        \end{tikzpicture}
        \caption{Hypervisor di Tipo 2 (Hosted)}
        \label{fig:type2_hyp}
    \end{subfigure}
    \caption{Confronto strutturale tra architetture hypervisor di Tipo 1 e Tipo 2.}
    \label{fig:hyp_comparison}
\end{figure}

\subsubsection{Hypervisor Di Tipo 1 (Bare-Metal)}
Un hypervisor di Tipo 1 opera a diretto contatto con il silicio, senza alcun sistema operativo host interposto. Esso agisce come un kernel specializzato e minimale, integrando al proprio interno i driver dei dispositivi fisici, lo scheduler dei task e la gestione dell'allocazione di memoria primordiale.
\begin{itemize}
    \item \textbf{Funzionamento}: intercetta direttamente le istruzioni privilegiate e di protezione emesse dai Guest OS, assegna le code temporali dei core fisici alle macchine virtuali e invia direttamente i risultati dell'elaborazione hardware ai contesti guest.
    \item \textbf{Vantaggi}: assicura prestazioni prossime a quelle native, latenze ridottissime nell'inoltro delle interruzioni e massima densità computazionale. Garantisce inoltre l'esecuzione parallela di sistemi operativi radicalmente eterogenei.
    \item \textbf{Svantaggi}: il kernel dello strato di virtualizzazione presenta elevata complessità di progettazione e sviluppo; inoltre, la compatibilità hardware è ristretta alla lista di dispositivi esplicitamente supportata dai driver integrati nel VMM.
    \item \textbf{Esempi notevoli}: VMware ESXi, Microsoft Hyper-V e Xen.
\end{itemize}

\subsubsection{Hypervisor Di Tipo 2 (Hosted)}
Un hypervisor di Tipo 2 viene installato ed eseguito come un normale processo dello spazio utente al di sopra di un sistema operativo convenzionale (\emph{Host OS}).
\begin{itemize}
    \item \textbf{Funzionamento}: la macchina virtuale viene mappata come un thread o sottoprocesso del processo hypervisor. L'accesso alle risorse fisiche (dischi, memoria, schede di rete) non avviene in maniera diretta, bensì delegando le operazioni alle chiamate di sistema del sistema operativo host.
    \item \textbf{Vantaggi}: estrema semplicità di installazione e portabilità. L'hypervisor beneficia dell'ampio ecosistema di driver già supportati dal sistema operativo host.
    \item \textbf{Svantaggi}: introduce un significativo degrado delle prestazioni dovuto all'elevato sovraccarico (\emph{overhead}) causato da molteplici cambi di contesto: un'istruzione privilegiata del guest deve essere prima intercettata dall'hypervisor, poi convertita in una system call per l'host OS e infine eseguita sull'hardware, raddoppiando le transizioni di ring della CPU.
    \item \textbf{Esempi notevoli}: Oracle VM VirtualBox e VMware Workstation.
\end{itemize}

---

\subsection{Le Tre Generazioni Di Soluzioni Software: Dalla Binary Translation All'Assistenza Hardware}

L'architettura originaria dei processori x86 a 32 bit (IA-32) presentava una violazione formale del teorema di Popek-Goldberg: esistevano circa 17 istruzioni sensibili (tra cui \texttt{POPF}, \texttt{PUSHF}, \texttt{SMSW}) che, se eseguite in modalità utente (Ring 1 o Ring 3), non scatenavano alcuna eccezione o trap hardware, ma fallivano silenziosamente oppure producevano comportamenti imprevisti. Per aggirare questo ostacolo, l'industria ha attraversato tre distinte generazioni tecnologiche:

% [IMMAGINE DEL FILE DI SLIDE CON NOME "CloudComputing_2"] Evoluzione temporale delle tre generazioni di virtualizzazione: Full Virtualization (Binary Translation), Paravirtualization e Hardware-assisted Virtualization — Slide 11 / Pagina 11 del PDF

\subsubsection{1ª Generazione: Virtualizzazione Completa Con Riscrizione Binaria (\emph{Binary Translation})}
Sviluppata alla fine degli anni '90 (introdotta da VMware e implementata anche da tool di emulazione come QEMU e soluzioni software Microsoft), la \textbf{Full Virtualization} consente l'esecuzione di sistemi operativi guest interamente non modificati.
\begin{itemize}
    \item \textbf{Meccanismo}: il VMM esamina dinamicamente a tempo di esecuzione i blocchi di codice macchina emessi dal guest. Le istruzioni non privilegiate vengono lasciate passare immutate per essere eseguite direttamente sul processore fisico (\emph{Direct Execution}). Al contrario, le istruzioni privilegiate o sensibili prive di trap vengono intercettate, tradotte o sostituite dinamicamente (\emph{Dynamic Binary Translation}) con sequenze equivalenti di istruzioni sicure capaci di richiamare le strutture del VMM.
    \item \textbf{Impatto computazionale}: comporta un sensibile decadimento prestazionale dovuto al costo di interpretazione e riscrittura del codice a runtime. Rappresenta tuttavia la soluzione d'elezione per carichi di tipo \emph{legacy} in cui il codice sorgente del kernel guest non è modificabile.
\end{itemize}

\subsubsection{2ª Generazione: Paravirtualizzazione}
Sviluppata dai ricercatori dell'Università di Cambridge con il progetto Xen all'inizio degli anni 2000, la \textbf{Paravirtualizzazione} adotta un approccio cooperativo tra il sistema operativo guest e l'hypervisor.
\begin{itemize}
    \item \textbf{Meccanismo}: il kernel del Guest OS viene modificato a livello di codice sorgente, rendendolo consapevole di trovarsi all'interno di un ambiente virtualizzato. Tutte le istruzioni privilegiate o le interazioni critiche con l'hardware vengono sostituite a tempo di compilazione con chiamate esplicite dirette verso il VMM, denominate \textbf{hypercall} (il cui funzionamento è concettualmente analogo a quello delle system call verso il kernel).
    \item \textbf{Impatto computazionale}: abbatte drasticamente il sovraccarico rispetto alla traduzione binaria, eliminando la necessità di intercettare istruzioni per tentativi ed errori. Il principale vincolo risiede nell'obbligo di poter ricompilare il kernel guest, rendendo impraticabile la virtualizzazione di sistemi operativi proprietari a codice chiuso (es. vecchie versioni di Windows) su hypervisor puramente paravirtualizzati.
\end{itemize}

\subsubsection{3ª Generazione: Virtualizzazione Assistita Da Hardware (\emph{Hardware-Assisted Virtualization})}
A partire dal 2005--2006, i produttori di silicio hanno introdotto estensioni architetturali native all'ISA (Intel VT-x con istruzioni VMX e AMD-V con istruzioni SVM) per colmare le lacune architetturali della piattaforma x86.
\begin{itemize}
    \item \textbf{Meccanismo}: il processore definisce due nuove modalità ortogonali di funzionamento hardware:
    \begin{itemize}
        \item \emph{VMX Root Operation}: la modalità pienamente privilegiata in cui esegue l'hypervisor.
        \item \emph{VMX Non-Root Operation}: la modalità riservata all'esecuzione del Guest OS e dei suoi applicativi.
    \end{itemize}
    Quando il Guest OS in modalità non-root esegue un'istruzione sensibile o si verifica un evento hardware critico, la CPU esegue una transizione controllata dall'hardware denominata \textbf{VM-Exit}, trasferendo l'esecuzione alla modalità root dell'hypervisor. Terminato l'intervento dell'allocator o dell'interprete, l'hypervisor esegue l'istruzione speciale \texttt{VMLAUNCH} o \texttt{VMRESUME}, provocando una \textbf{VM-Entry} che ripristina all'istante il contesto hardware del guest.
    \item \textbf{Impatto computazionale}: garantisce prestazioni vicine a quelle del bare-metal per l'esecuzione del codice, supporta sistemi operativi guest completamente non modificati ed elimina l'onere della traduzione binaria dinamica.
\end{itemize}

---

\subsection{Cronologia Delle Tecnologie Di Virtualizzazione (1960–2020)}

L'evoluzione della virtualizzazione ha attraversato oltre sessant'anni di scoperte architetturali, muovendo dai mainframe fino all'era cloud-native.

% [IMMAGINE DEL FILE DI SLIDE CON NOME "CloudComputing_2"] Timeline storico-tecnologica della virtualizzazione dal 1960 al 2020 (da IBM CP-40 a Docker, Kata Containers e Podman) — Slide 12 / Pagina 12 del PDF

\begin{itemize}
    \item \textbf{Anni '60 (1960--1967)}: pionieri come General Electric (GE), Bell Labs e International Business Machines (IBM) affrontano l'esigenza di far convivere carichi legacy con le nuove generazioni di mainframe. Nel 1967 nasce l'architettura \textbf{IBM CP-40 / CP-67}, il primo sistema commerciale a supportare la virtualizzazione completa. A differenza del classico \emph{time-sharing} (in cui la memoria e la CPU sono divise staticamente tra gli utenti loggati), il sistema CP fornisce a ciascun operatore un'istanza virtuale separata del sistema operativo, aumentando l'affidabilità (il crash del guest di un utente non compromette il mainframe) e la sicurezza.
    \item \textbf{1982}: introduzione della chiamata di sistema \texttt{chroot} all'interno di Unix versione 7, segnando la prima forma rudimentale di isolamento dello spazio di memoria secondaria tramite la reimpostazione della directory radice di un processo.
    \item \textbf{Fine anni '80}: comparsa dei primi emulatori software per workstation Unix e Macintosh finalizzati a consentire l'esecuzione di applicativi sviluppati per ambiente DOS e Microsoft Windows.
    \item \textbf{1995}: lancio da parte di Sun Microsystems del linguaggio \textbf{Java} e della relativa \emph{Java Virtual Machine} (JVM), che astrae l'esecuzione del codice intermedio (\emph{bytecode}) rispetto all'architettura fisica sottostante.
    \item \textbf{1999}: fondazione e rilascio da parte di \textbf{VMware} della prima tecnologia di virtualizzazione completa per processori x86 basata su traduzione binaria dinamica (\textbf{VMware Workstation}).
    \item \textbf{2001}: rilascio di \textbf{VMware ESX Server}, il primo hypervisor bare-metal di Tipo 1 orientato al comparto enterprise per x86, affiancato da \textbf{VMware GSX Server}, soluzione hosted di Tipo 2 concepita per operare al di sopra di un kernel Windows o Linux.
    \item \textbf{2003}: rilascio di \textbf{Xen} da parte dei ricercatori dell'Università di Cambridge, primo hypervisor open-source per x86 strutturato sulla paravirtualizzazione.
    \item \textbf{2004}: rilascio di \textbf{Virtual Server} da parte di Microsoft.
    \item \textbf{2005}: Sun Microsystems introduce \textbf{Solaris Containers} (basati su \textbf{Solaris Zones}) all'interno del sistema operativo Solaris 10. Questa tecnologia sancisce l'affermazione della virtualizzazione a livello di sistema operativo (\emph{OS-level virtualization}), permettendo la coesistenza di molteplici ambienti virtuali completamente isolati ma governati da un unico kernel condiviso.
    \item \textbf{2005}: comparsa del progetto open-source \textbf{OpenVZ}, che introduce la virtualizzazione a container per il kernel Linux.
    \item \textbf{2007}: nel mese di febbraio, l'infrastruttura \textbf{KVM} (\emph{Kernel-based Virtual Machine}) viene formalmente integrata all'interno del ramo principale del kernel Linux (\emph{mainline}). KVM trasforma il kernel Linux in un hypervisor integrando il modulo di controllo con le estensioni di virtualizzazione hardware (Intel VT-x / AMD-V), delegando a una variante di QEMU l'emulazione delle periferiche di I/O.
    \item \textbf{2008}: debutto di \textbf{LXC} (\emph{Linux Containers}), la prima soluzione integrata che fonde nativamente due pilastri del kernel Linux: i \emph{Control Groups} (cgroups) e i \emph{Namespaces}.
    \item \textbf{2013}: nascita del progetto open-source \textbf{Docker}, che automatizza la distribuzione di container tramite immagini immutabili e introduce un formato standard universale per il packaging del software.
    \item \textbf{2020 ed evoluzione recente}: consolidamento di strumenti modulari e runtime orientati alla sicurezza come \textbf{Podman} (architettura senza daemon centrale) e \textbf{Kata Containers} (container eseguiti entro micro-VM per unire la velocità dei container con l'isolamento hardware degli hypervisor).
\end{itemize}

\newpage

% ==============================================================================
% SEZIONE 3: LA RIVOLUZIONE DEI CONTAINER E L'ECOSISTEMA DOCKER
% ==============================================================================

\section{La Rivoluzione Dei Container E L'Ecosistema Docker}

\subsection{La Crisi Della Complessità: La "Matrice Infernale" E L'Analogia Intermodale}

Con la progressiva transizione dalle architetture monolitiche ai moderni sistemi distribuiti e a microservizi, l'ingegneria del software si è scontrata con un problema strutturale di complessità combinatoria. Da un lato si assiste a una \textbf{molteplicità di stack applicativi}: i sistemi eterogenei integrano contemporaneamente linguaggi di programmazione eterogenei, runtime specifici, server web e database con requisiti e vincoli contrastanti (ad es. un sito web statico servito da Nginx con OpenSSL, un frontend web in Ruby on Rails, un backend di calcolo asincrono in Python con librerie grafiche OpenCV e broker Celery, code di messaggistica Redis, basi di dati relazionali PostgreSQL e cluster analitici Apache Hadoop). Dall'altro lato sussiste una \textbf{molteplicità di ambienti hardware di esecuzione}: le macchine virtuali di sviluppo sui computer portatili dei programmatori, i server fisici dedicati ai test di qualità (QA), i nodi di calcolo locali (\emph{on-premises}), le infrastrutture private di disaster recovery e i molteplici fornitori di Cloud pubblico.

% [IMMAGINE DEL FILE DI SLIDE CON NOME "CloudComputing_2"] Rappresentazione concettuale della sfida architetturale tra molteplicità di stack software e molteplicità di infrastrutture di distribuzione — Slide 14 / Pagina 14 del PDF

In assenza di standardizzazione, configurare $N$ stack software distinti per funzionare in modo prevedibile su $M$ differenti ambienti operativi dà origine alla cosiddetta \textbf{Matrice Infernale} (\emph{The Matrix From Hell}), in cui il numero di combinazioni di compatibilità da testare, manutenere e risolvere cresce proporzionalmente al prodotto $N \times M$.

% [IMMAGINE DEL FILE DI SLIDE CON NOME "CloudComputing_2"] La "Matrix from Hell": tabella bidimensionale con punti interrogativi che illustra la complessità NxM prima dell'avvento dei container — Slide 15 / Pagina 15 del PDF

Il problema generato da questa matrice riproduce esattamente la crisi logistica che ha afflitto il trasporto mercantile prima degli anni '60. Storicamente, il trasporto merci prevedeva una moltitudine eterogenea di beni (casse di legno, sacchi di juta, botti, veicoli pesanti, pianoforti) trasportati su una molteplicità di vettori logistici (carrelli elevatori, vagoni ferroviari, camion, navi cargo, gru portuali). Il trasbordo di merci deperibili o incompatibili richiedeva continue manipolazioni manuali, ritardi nelle banchine e rischi di danneggiamento.

% [IMMAGINE DEL FILE DI SLIDE CON NOME "CloudComputing_2"] Il trasporto merci pre-1960 e la matrice di incompatibilità tra merci eterogenee e vettori di trasporto — Slide 16-17 / Pagine 16-17 del PDF

La svolta epocale nel commercio marittimo coincise con l'introduzione del \textbf{container marittimo intermodale} standardizzato (standard ISO). Il container astrae il contenuto: una volta sigillato, all'operatore portuale non interessa la natura del bene trasportato; l'infrastruttura (gru, motrici, vagoni ferroviari) deve interfacciarsi unicamente con i punti di ancoraggio standardizzati della scatola.

% [IMMAGINE DEL FILE DI SLIDE CON NOME "CloudComputing_2"] La soluzione: il container intermodale standardizzato per il trasporto merci — Slide 18 / Pagina 18 del PDF

\textbf{Docker} ha trasposto la medesima intuizione ingegneristica nell'industria del software. Il container applicativo diventa l'unità intermodale del codice.

% [IMMAGINE DEL FILE DI SLIDE CON NOME "CloudComputing_2"] Docker come container intermodale per il codice software: risoluzione della matrice infernale — Slide 19-20 / Pagine 19-20 del PDF

---

\subsection{Concetto Di Container E Separazione Delle Competenze (Dev vs Ops)}

Un container costituisce un'unità software standardizzata che incapsula al proprio interno il codice sorgente compilato unitamente a tutte le sue dipendenze operative: file binari, librerie di sistema, interpreti di runtime, variabili d'ambiente e file di configurazione statici.
L'adozione di Docker ridefinisce in modo rigoroso la \textbf{separazione delle responsabilità} (\emph{Separation of Concerns}) tra lo sviluppo software e l'esercizio infrastrutturale:

% [IMMAGINE DEL FILE DI SLIDE CON NOME "CloudComputing_2"] Schema della separazione delle responsabilità (Separation of Concerns): Dan lo Sviluppatore vs Oscar l'Operatore di Sistema — Slide 21 / Pagina 21 del PDF

\begin{itemize}
    \item \textbf{Lo Sviluppatore (\emph{Dan the Developer})}: si focalizza su quanto risiede \emph{all'interno} del container. Cura la scrittura del codice, la selezione delle librerie, la gestione dei pacchetti specifici della propria applicazione e la consistenza dei dati. In virtù del formato container, qualsiasi server Linux di destinazione presenta la medesima interfaccia di runtime, scongiurando il classico paradosso operativo: \emph{"sul mio computer funzionava"}.
    \item \textbf{L'Ingegnere Operativo (\emph{Oscar the Ops Guy})}: si concentra su quanto risiede \emph{all'esterno} del container. Governa il monitoraggio dello stato di salute, l'aggregazione centralizzata dei log, i criteri di routing e configurazione di rete, le politiche di accesso remoto e l'allocazione delle risorse fisiche. Dal punto di vista dell'infrastruttura, tutti i container espongono la medesima interfaccia uniforme per il ciclo di vita: operazioni come \texttt{start}, \texttt{stop}, \texttt{restart}, \texttt{attach} e la migrazione tra host si eseguono sempre allo stesso modo, a prescindere dal linguaggio con cui l'applicativo è stato compilato.
\end{itemize}

---

\subsection{Confronto Architetturale Diretto: Macchine Virtuali vs Container}

Il confronto strutturale tra macchine virtuali e container riflette due filosofie di astrazione profondamente distinte:

% [IMMAGINE DEL FILE DI SLIDE CON NOME "CloudComputing_2"] Confronto architetturale diretto e dettagliato tra Macchine Virtuali (Hypervisor Type 2) e Container Docker — Slide 22 / Pagina 22 del PDF

\begin{table}[H]
    \centering
    \small
    \begin{tabularx}{\textwidth}{p{3.5cm} X X}
        \toprule
        \textbf{Caratteristica} & \textbf{Macchine Virtuali (VM)} & \textbf{Container (Docker)} \\
        \midrule
        \textbf{Livello di Astrazione} & Virtualizzazione dell'hardware fisico (CPU, RAM, NIC virtuali). & Virtualizzazione a livello di sistema operativo (\emph{OS-level virtualization}). \\
        \midrule
        \textbf{Kernel del S.O.} & Ciascuna VM esegue un proprio kernel guest completo e indipendente. & Tutti i container condividono lo stesso kernel del sistema operativo host. \\
        \midrule
        \textbf{Isolamento} & Isolamento hardware forte (fornito dal VMM e dai privilegi della CPU). & Isolamento logico dei processi basato su costrutti del kernel (\emph{namespaces}). \\
        \midrule
        \textbf{Impronta di Memoria} & Pesante (centinaia di MB o GB per ciascuna istanza solo per il kernel guest). & Molto leggera (pochi megabyte; include solo i binari e le librerie necessarie). \\
        \midrule
        \textbf{Tempo di Avvio} & Lento (da decine di secondi a svariati minuti per il boot completo del guest). & Quasi istantaneo (millisecondi, equivalente all'avvio di un normale processo Linux). \\
        \midrule
        \textbf{Prestazioni di I/O} & Minori a causa dell'emulazione dell'hardware e dei frequenti cambi di contesto. & Prestazioni prossime a quelle native (accesso diretto alle chiamate di sistema dell'host). \\
        \midrule
        \textbf{Portabilità} & Limitata dai formati di immagine disco proprietari e dalla pesantezza dei file. & Elevata grazie a formati di immagine a strati leggeri e universali. \\
        \bottomrule
    \end{tabularx}
    \caption{Quadro comparativo esaustivo tra Macchine Virtuali e Container.}
    \label{tab:vm_vs_docker}
\end{table}

\newpage

% ==============================================================================
% SEZIONE 4: PRIMITIVI DEL KERNEL LINUX PER LA CONTAINERIZZAZIONE
% ==============================================================================

\section{Primitivi Del Kernel Linux Per La Containerizzazione}

Docker non costituisce un ambiente di esecuzione a sé stante dotato di un proprio kernel monolitico, bensì opera sfruttando un insieme di primitive integrate nel kernel Linux:

% [IMMAGINE DEL FILE DI SLIDE CON NOME "CloudComputing_2"] Stack architetturale dei componenti del kernel Linux alla base di Docker (libcontainer, cgroups, namespaces, netlink, capabilities, selinux, apparmor) — Slide 23 / Pagina 23 del PDF

---

\subsection{Namespaces E Partizionamento Delle Risorse}

I \textbf{namespaces} rappresentano la tecnologia del kernel Linux concepita per partizionare e isolare le risorse di sistema in modo tale che gruppi distinti di processi visualizzino insiemi mutualmente esclusivi di risorse globali. A partire dalla versione di kernel 5.6, sono formalmente presenti 8 tipologie di namespaces. I principali per il funzionamento dei container sono:
\begin{itemize}
    \item \textbf{Mount Namespace (\texttt{mnt})}: isola la gerarchia dei punti di montaggio del file system. Al momento della creazione di un nuovo namespace di mount, la tabella corrente viene clonata, ma qualsiasi punto di montaggio creato, spostato o smontato successivamente rimane confinato localmente, consentendo al container di disporre di una propria visione esclusiva della root (\texttt{/}).
    \item \textbf{Network Namespace (\texttt{net})}: virtualizza l'intero stack di rete. Ciascun network namespace possiede una tabella di instradamento (\emph{routing table}), catene di firewalling indipendenti (\texttt{iptables}/\texttt{nftables}), porte TCP/UDP isolate e un insieme esclusivo di interfacce di rete. Alla creazione, include unicamente l'interfaccia di loopback (\texttt{lo}); interfacce virtuali (come coppie di interfacce di tipo \emph{veth}) possono essere istanziate e spostate all'interno per collegare il container a bridge software.
    \item \textbf{PID Namespace (\texttt{pid})}: isola lo spazio degli identificatori numerici di processo (\emph{Process ID}). I namespace PID sono intrinsecamente gerarchici: il primo processo eseguito nel container assume il valore $\text{PID} = 1$ (diventando il processo genitore di default e gestendo la raccolta degli stati dei processi orfani), mentre all'esterno, visto dallo spazio del sistema operativo host, il medesimo processo assume un valore di PID standard (ad es. $\text{PID} = 15420$).
    \item \textbf{IPC Namespace (\texttt{ipc})}: isola i canali di comunicazione inter-processo di tipo SysV e le code di messaggi POSIX, impedendo che processi in container separati possano accedere alla memoria condivisa altrui.
    \item \textbf{UTS Namespace (\texttt{uts})}: permette al container di definire un proprio nome host (\emph{hostname}) e dominio di rete, disaccoppiati da quelli del nodo host fisico.
    \item \textbf{User Namespace (\texttt{user})}: rimappa gli identificativi utente (\texttt{UID}) e gruppo (\texttt{GID}) del container su intervalli differenti dell'host. Un processo può operare come utente root ($\text{UID} = 0$) all'interno del container pur corrispondendo a un normale utente non privilegiato ($\text{UID} > 10000$) sul sistema host, riducendo i rischi di compromissione del sistema.
    \item \textbf{Cgroup Namespace}: isola la visibilità della radice dei controller cgroup.
    \item \textbf{Time Namespace}: isola l'orologio interno di sistema (\texttt{CLOCK\_MONOTONIC} e \texttt{CLOCK\_BOOTTIME}).
\end{itemize}

---

\subsection{Control Groups (Cgroups): Allocazione E Limiti Di Risorsa}

Se i namespaces determinano \emph{cosa un processo può vedere}, i \textbf{Control Groups} (comunemente noti come \textbf{cgroups}) stabiliscono \emph{quante risorse un processo può consumare}. Introdotti dagli ingegneri di Google nel 2006 (originariamente con il nome di \emph{Process Containers}) e poi inglobati nel kernel Linux nel 2008, i cgroups forniscono quattro funzionalità critiche:
\begin{enumerate}
    \item \textbf{Limitazione delle risorse (\emph{Resource Limits})}: consentono di impostare tetti massimi invalicabili all'uso di CPU, memoria RAM di lavoro, swap su disco o banda I/O. Se un processo tenta di eccedere il quantitativo di memoria allocato dal cgroup, interviene il sottosistema del kernel sollevando un evento di \emph{Out-of-Memory (OOM) Killer} mirato al solo cgroup colpevole, senza arrestare il sistema host.
    \item \textbf{Prioritizzazione (\emph{Prioritization})}: in condizioni di saturazione o contesa concorrente delle risorse (es. contesa della CPU tra processi affamati di calcolo), i cgroups distribuiscono quote di utilizzo tramite pesi percentuali (\emph{CPU shares}) e algoritmi di tipo \emph{Completely Fair Scheduler} (CFS).
    \item \textbf{Contabilità (\emph{Accounting})}: effettuano il tracciamento dettagliato e costante del consumo effettivo di risorse a livello di gruppo, rendendo accessibili metriche di carico a tool di telemetria e orchestratori.
    \item \textbf{Controllo (\emph{Control})}: consentono di congelare lo stato operativo (\emph{freeze}), arrestare o risvegliare all'istante l'intero albero di processi associati a un cgroup mediante un unico comando atomico inviato al file system virtuale di gestione.
\end{enumerate}

---

\subsection{Sicurezza E Strumenti Di Controllo: Capabilities, AppArmor, SELinux E Libvirt}

A completamento dell'infrastruttura di confinamento, il kernel Linux integra moduli dedicati alla sicurezza e all'interfacciamento hardware:
\begin{itemize}
    \item \textbf{Netlink}: protocollo socket-based per la comunicazione bidirezionale tra lo spazio utente e i moduli interni del kernel. Docker lo impiega per configurare dinamicamente le rotte di rete, le interfacce virtuali e le tabelle di indirizzi.
    \item \textbf{Linux Capabilities}: nei sistemi UNIX tradizionali, il modello dei permessi è binario: un processo è gestito come non privilegiato ($\text{UID} \neq 0$), oppure possiede poteri assoluti come superutente root ($\text{UID} = 0$), eludendo qualsiasi controllo del kernel. Il meccanismo delle \emph{Capabilities} scompone i privilegi di root in oltre 40 privilegi atomici e granulari (es. \texttt{CAP\_NET\_ADMIN} per modificare le interfacce di rete, \texttt{CAP\_SYS\_TIME} per impostare l'orologio, \texttt{CAP\_NET\_BIND\_SERVICE} per aprire porte inferiori alla 1024). Docker avvia i container rimuovendo per impostazione predefinita le capacità rischiose (come \texttt{CAP\_SYS\_ADMIN}), riducendo la superficie di attacco anche qualora un attaccante ottenga privilegi di root dentro il container.
    \item \textbf{AppArmor e Security-Enhanced Linux (SELinux)}: implementazioni di sistemi a controllo d'accesso obbligatorio (\emph{Mandatory Access Control}, MAC) all'interno del sottosistema LSM (\emph{Linux Security Modules}).
    \begin{itemize}
        \item \emph{AppArmor}: definisce profili di sicurezza basati sui percorsi del file system (\emph{name-based MAC}), vincolando rigorosamente a quali file e capacità applicative il binario containerizzato abbia diritto di accedere.
        \item \emph{SELinux}: associa etichette di sicurezza a ciascun oggetto del sistema (processi, file, descrittori), applicando politiche di transizione basate su modelli a minimi privilegi.
    \end{itemize}
    \item \textbf{Libvirt}: toolkit e libreria API di livello sistemistico aperta (compatibile con C, Python, Go) che standardizza la gestione di molteplici piattaforme di virtualizzazione, tra cui KVM, QEMU, Xen, VMware ESXi, BHyve ed estensioni per container.
\end{itemize}

\newpage

% ==============================================================================
% SEZIONE 5: ARCHITETTURA, OGGETTI E CICLO DI VITA DI DOCKER
% ==============================================================================

\section{Architettura, Oggetti E Ciclo Di Vita Di Docker}

\subsection{Architettura Client-Server E Componenti Principali}

Docker adotta un'architettura rigorosamente disaccoppiata di tipo \textbf{Client-Server}:

% [IMMAGINE DEL FILE DI SLIDE CON NOME "CloudComputing_2"] Diagramma dell'architettura Client-Server di Docker: Client, Docker Host con dockerd, Registri e interazioni API — Slide 24 / Pagina 24 del PDF

L'architettura si articola su tre componenti principali:
\begin{enumerate}
    \item \textbf{Docker Client (\texttt{docker})}: costituisce l'interfaccia a riga di comando (CLI) mediante cui l'utente interagisce con il sistema. Quando viene digitato un comando (come \texttt{docker build}, \texttt{docker run} o \texttt{docker pull}), il client trasforma la richiesta in una chiamata HTTP conforme alle specifiche \textbf{Docker Engine REST API} e la inoltra al daemon. Il client può comunicare con un daemon locale (tramite socket UNIX di dominio, tipicamente collocato in \texttt{/var/run/docker.sock}) oppure con un daemon remoto attraverso una connessione di rete protetta da cifratura TLS.
    \item \textbf{Docker Daemon (\texttt{dockerd})}: il processo in background in esecuzione continua sul nodo host. Ascolta costantemente le richieste provenienti dall'API REST e ha il compito operativo di creare, avviare, monitorare e distruggere gli oggetti Docker (immagini, container, reti virtuali e volumi di archiviazione). Inoltre, può comunicare con altri daemon per coordinare cluster di nodi.
    \item \textbf{Docker Registry}: il registro remoto o locale preposto all'archiviazione e alla catalogazione delle immagini dei container. Il registro predefinito di riferimento è il registro pubblico \textbf{Docker Hub}, ma all'interno di reti aziendali possono essere istanziati registri privati per ragioni di sicurezza e conformità.
\end{enumerate}

Per agevolare l'integrazione con sistemi esterni, Docker distribuisce librerie di sviluppo ufficiali (\textbf{SDK}) per linguaggi quali Go e Python, che consentono di programmare e automatizzare le operazioni di gestione dei container senza passare per la shell di sistema.

---

\subsection{Oggetti Docker: Immagini A Livelli (CoW), Container E Registri}

Il funzionamento dell'ecosistema poggia sulla manipolazione di oggetti software ben definiti:

\begin{figure}[H]
    \centering
    \begin{tikzpicture}[node distance=1.8cm, auto]
        \node [draw, rectangle, fill=orange!20, text width=3.5cm, align=center, minimum height=1.2cm] (dockerfile) {\textbf{Dockerfile}\\\footnotesize Ricetta testuale a istruzioni};
        \node [draw, rectangle, fill=theoryblue!20, text width=3.5cm, align=center, minimum height=1.2cm, right of=dockerfile, xshift=2.5cm] (image) {\textbf{Docker Image}\\\footnotesize Template immutabile a strati (Read-Only)};
        \node [draw, rectangle, fill=darkgreen!20, text width=3.5cm, align=center, minimum height=1.2cm, right of=image, xshift=2.5cm] (container) {\textbf{Docker Container}\\\footnotesize Istanza di runtime eseguibile con layer R/W};

        \draw [->, thick] (dockerfile) -- node[above, font=\footnotesize] {\texttt{docker build}} (image);
        \draw [->, thick] (image) -- node[above, font=\footnotesize] {\texttt{docker run}} (container);
    \end{tikzpicture}
    \caption{Flusso di transizione da Dockerfile a Immagine e a Container in esecuzione.}
    \label{fig:dockerfile_to_container}
\end{figure}

\begin{itemize}
    \item \textbf{Immagine (\emph{Docker Image})}: rappresenta un pacchetto immutabile e di sola lettura contenente il software e i metadati di configurazione necessari per istanziare un ambiente di esecuzione. Le immagini vengono generate a partire da un file dichiarativo denominato \textbf{Dockerfile}. Ogni istruzione contenuta nel Dockerfile (ad es. \texttt{FROM}, \texttt{RUN}, \texttt{COPY}) genera un nuovo livello (\emph{layer}) differenziale sovrapposto ai precedenti mediante un file system unificato di tipo \textbf{Union File System} (come Overlay2). Grazie a questo modello, i layer sottostanti sono immutabili e possono essere condivisi contemporaneamente da decine di immagini diverse; in caso di modifica del Dockerfile, vengono ricompilati solo i layer modificati, minimizzando il consumo di banda e velocizzando le build.
    \item \textbf{Container}: è l'istanza viva ed eseguibile di un'immagine. Quando un'immagine viene avviata, il Docker Engine alloca sopra la pila di layer immutabili in sola lettura un sottile livello modificabile in lettura e scrittura (\textbf{Read-Write Container Layer}), utilizzando un meccanismo a \textbf{Copy-on-Write (CoW)}:
    \begin{itemize}
        \item Se il processo in esecuzione legge un file appartenente all'immagine, esso viene letto direttamente dai layer inferiori in sola lettura.
        \item Se il processo modifica un file preesistente, il file viene copiato al volo nel layer superiore di scrittura e lì modificato, lasciando intatta l'immagine base.
        \item Alla rimozione del container, tutti i dati contenuti nel layer di scrittura vengono eliminati, a meno che non siano stati esplicitamente mappati su \textbf{volumi} persistenti gestiti sull'host.
    \end{itemize}
\end{itemize}

---

\subsection{Anatomia Dettagliata Di Un Comando: Il Flusso Operativo Di \texttt{docker run}}

Per comprendere la sequenza di interazioni hardware e software orchestrate dal daemon a fronte di una richiesta utente, analizziamo l'esecuzione del comando di avvio per una sessione interattiva con una distribuzione Ubuntu:

\begin{lstlisting}[language=bash]
$ docker run -i -t ubuntu /bin/bash
\end{lstlisting}

I passaggi esecutivi interni che si susseguono comprendono:
\begin{enumerate}
    \item \textbf{Ispezione del repository locale}: il client trasmette la richiesta a \texttt{dockerd}. Il daemon effettua una ricerca nella memoria locale per verificare se l'immagine con tag \texttt{ubuntu:latest} sia già presente nella cache dei layer del nodo host.
    \item \textbf{Download dell'immagine (\emph{Image Pulling})}: qualora l'immagine non sia presente localmente, il daemon contatta automaticamente il registro configurato di default (Docker Hub) ed esegue il download sequenziale di tutti i layer costitutivi mancanti, assemblandoli nel motore locale di storage.
    \item \textbf{Creazione strutturale del container}: il daemon invoca le primitive di creazione allocando l'identificatore univoco del container, configurando i puntatori dei descrittori e inizializzando lo stato a riposo.
    \item \textbf{Allocazione del layer Read-Write}: sopra lo stack di layer immutabili dell'immagine Ubuntu viene istanziato un layer effimero di lettura e scrittura basato sul driver Overlay2.
    \item \textbf{Inizializzazione delle interfacce di rete}: il daemon predispone un'interfaccia di rete virtuale (\emph{veth pair}) assegnandone un'estremità all'interno del nuovo network namespace del container e l'altra al bridge software dell'host (tipicamente \texttt{docker0}). Assegna contestualmente un indirizzo IP privato dinamico (estratto dalla subnet del bridge), imposta le regole di routing e abilita la risoluzione DNS.
    \item \textbf{Allocazione dello pseudo-terminale e aggancio dello standard input}:
    \begin{itemize}
        \item L'opzione \texttt{-i} (\texttt{--interactive}) mantiene aperto lo stream dello standard input (\texttt{stdin}) anche se il processo non è in foreground, permettendo di immettere dati da tastiera.
        \item L'opzione \texttt{-t} (\texttt{--tty}) alloca una pseudo-terminale (TTY virtuale) associata al processo, formattando correttamente i caratteri di escape e il prompt dei comandi della shell all'interno della finestra del terminale dell'utente.
    \end{itemize}
    \item \textbf{Esecuzione del processo}: il daemon manda in esecuzione il binario specificato (nel nostro caso \texttt{/bin/bash}) come processo $\text{PID} = 1$ all'interno del nuovo PID namespace.
    \item \textbf{Cessazione e arresto}: quando l'utente digita il comando \texttt{exit} all'interno della shell bash, il processo con $\text{PID} = 1$ termina regolarmente con codice di ritorno. Il daemon rileva la terminazione del processo principale e passa il container nello stato di arresto (\emph{Stopped}). È fondamentale sottolineare che \emph{il container non viene automaticamente rimosso}: il suo stato e i file modificati nel layer R/W restano salvati su disco e possono essere riavviati con \texttt{docker start} oppure eliminati definitivamente con \texttt{docker rm}.
\end{enumerate}

---

\subsection{Distribuzione E Sicurezza Delle Immagini: Docker Hub E Docker Content Trust}

La rapida adozione dei container ha reso la sicurezza della supply chain del software un elemento critico:
\begin{itemize}
    \item \textbf{Docker Desktop}: suite software progettata per ambienti workstation (macOS e Windows) che incorpora al proprio interno un'istanza virtualizzata minimale di Linux utile a far girare il daemon \texttt{dockerd}, il client \texttt{docker}, \textbf{Docker Compose} per la gestione multi-container, un cluster locale \textbf{Kubernetes} a singolo nodo, la \textbf{Docker Dashboard} grafica e moduli per il salvataggio sicuro delle credenziali utente nel portachiavi di sistema (\texttt{docker-credential-helpers}).
    \item \textbf{Docker Hub}: costituisce il più vasto registro pubblico mondiale per container, ospitando sia immagini ufficiali certificate di fornitori terzi sia repository di progetti open source e cataloghi privati a pagamento.
    \item \textbf{Docker Content Trust (DCT)}: modulo basato sullo standard The Update Framework (TUF) e sul componente Notary. DCT integra firme digitali crittografiche per consentire la verifica dell'autenticità e dell'integrità dei dati scambiati con i registri remoti. Quando DCT è attivo, il client impedisce l'esecuzione del comando \texttt{docker pull} o \texttt{docker run} qualora l'immagine remota non sia stata esplicitamente firmata con una chiave crittografica valida dal relativo autore certificato, proteggendo il sistema da attacchi di tipo man-in-the-middle, corruzione accidentale o iniezione di malware nei tag delle immagini.
\end{itemize}

\newpage

% ==============================================================================
% SEZIONE 6: ANALISI COMPARATIVA DELLE PRESTAZIONI: BARE-METAL, DOCKER E KVM
% ==============================================================================

\section{Analisi Comparativa Delle Prestazioni: Bare-Metal, Docker E KVM}

L'impatto architetturale introdotto dallo strato di virtualizzazione è stato oggetto di rigorosa validazione sperimentale da parte dei laboratori di ricerca IBM all'interno dello studio comparativo del luglio 2014, che ha messo a confronto le prestazioni di tre configurazioni operative equivalenti: \textbf{Bare-Metal Nativo}, \textbf{Container Docker} e macchine virtuali basate sull'hypervisor \textbf{KVM}.

I test sono stati condotti su piattaforme server di classe enterprise \textbf{IBM System x3650 M4}, equipaggiate con processori Intel Sandy Bridge-EP Xeon E5-2665 (frequenza base di 2.4 GHz, turbo boost fino a 3.0 GHz, 16 core fisici con Hyper-Threading abilitato per un totale di 32 core logici) e una dotazione di memoria primaria pari a 256 GB di RAM.

---

\subsection{Prestazioni Di I/O E Sovraccarico Di Emulazione}

Nei centri elaborazione dati e nelle architetture cloud, l'accesso a storage a blocchi su reti SAN (\emph{Storage Area Network}) è fondamentale per garantire elevato throughput e rigida consistenza dei dati transazionali.

% [IMMAGINE DEL FILE DI SLIDE CON NOME "CloudComputing_2"] Grafici e tabelle sulle prestazioni di I/O (Throughput IOPS e netperf) tra Bare-Metal, Docker e KVM (IBM Research Report 2014) — Slide 29 / Pagina 29 del PDF

I benchmark di throughput di I/O casuale (\emph{Random Read}, \emph{Random Write} e \emph{Random Mixed}) espressi in operazioni al secondo (\textbf{IOPS}) mostrano che:
\begin{itemize}
    \item \textbf{Docker} ottiene valori di IOPS quasi sovrapponibili a quelli registrati sul sistema bare-metal nativo (circa 85.000 IOPS in lettura, 110.000 in scrittura e 105.000 nel misto), poiché i processi all'interno dei container utilizzano i medesimi driver di blocco e il medesimo page cache manager del kernel ospitante, azzerando i livelli di traduzione intermedi.
    \item \textbf{KVM} mostra un decadimento marcato del throughput, registrando valori dimezzati o inferiori (circa 40.000 IOPS). La causa strutturale risiede nel fatto che ciascuna operazione di I/O emessa dal guest deve attraversare il layer di emulazione gestito da \textbf{QEMU}.
\end{itemize}

Anche adottando interfacce di rete paravirtualizzate ad alte prestazioni come i driver \texttt{virtio} (che evitano l'emulazione a basso livello della scheda di rete fisica), i risultati di rete ottenuti tramite il benchmark \texttt{netperf} (flusso TCP\_STREAM con segmenti da 16 KB su interfaccia da \SI{1}{\giga\bit/\second}) evidenziano un limite intrinseco della virtualizzazione basata su hypervisor:

\begin{table}[H]
    \centering
    \small
    \begin{tabularx}{\textwidth}{X c c}
        \toprule
        \textbf{Metodo di Virtualizzazione I/O} & \textbf{Throughput (\si{\mega\bit/\second})} & \textbf{Utilizzo CPU (\%)} \\
        \midrule
        Bare-Metal Nativo & 950 & 20\% \\
        Assegnazione Dispositivo Diretto (\emph{Device Assignment} / SR-IOV) & 950 & 25\% \\
        Rete Paravirtualizzata (\emph{Virtio}) & 950 & 50\% \\
        Emulazione Completa (\emph{Full Emulation}) & 250 & 100\% \\
        \bottomrule
    \end{tabularx}
    \caption{Prestazioni di rete e consumo di CPU nei diversi paradigmi di virtualizzazione I/O (Dati IBM Research).}
    \label{tab:io_performance}
\end{table}

Come si evince dai dati:
\begin{itemize}
    \item Sia il Bare-Metal sia il driver paravirtualizzato saturano il link fisico raggiungendo il massimo teorico utile di \SI{950}{\mega\bit/\second}.
    \item L'approccio paravirtualizzato richiede tuttavia un assorbimento di risorse di CPU del \textbf{50\%}, a fronte del solo \textbf{20\%} registrato dal Bare-Metal.
\end{itemize}
Questo divario è riconducibile all'elevata frequenza di \textbf{cambi di contesto} e alla generazione di eventi di trap (VM-Exit) richiesti per ogni pacchetto trasferito: la gestione delle interruzioni virtuali deve comunque passare attraverso la mediazione dell'hypervisor, saturando i core di calcolo in attività di gestione del contesto anziché nell'esecuzione di calcoli utili.

---

\subsection{Latenza Real-Time E Test Con Cyclictest}

In contesti computazionali critici (telecomunicazioni 5G, controllo robotico in tempo reale, finanza ad alta frequenza), il parametro chiave non è solo il throughput medio, bensì la \textbf{latenza di risposta deterministica}.

\begin{theorem}[\textbf{Latenza}]
La \textbf{latenza deterministica} è formalmente definita come l'intervallo temporale che intercorre tra il verificarsi di un evento di stimolo fisico o logico $t_0$ e l'istante effettivo $t_1$ in cui il thread preposto a gestire tale evento viene messo in esecuzione dalla CPU:
\[
    \Delta t_{\text{latency}} = t_1 - t_0
\]
Esempi applicativi includono l'intervallo tra la generazione di un'interruzione hardware e l'avvio della relativa routine di servizio (ISR), il tempo che separa lo scadere di un timer ad alta precisione dalla schedulazione del thread dormiente, o l'intervallo tra l'arrivo di un pacchetto di rete al controller e l'inizio della sua elaborazione da parte dello stack di protocollo.
\end{theorem}

% [IMMAGINE DEL FILE DI SLIDE CON NOME "CloudComputing_2"] Grafici della distribuzione temporale della latenza misurata con Cyclictest su Host nativo, Docker e KVM — Slide 30 / Pagina 30 del PDF

La quantificazione sperimentale della latenza è stata condotta impiegando \textbf{Cyclictest}, lo strumento standard per il benchmarking di sistemi a bassa latenza e kernel real-time. Il test è stato eseguito su processori Intel Ivy Bridge 4 core con Hyper-Threading (8 core logici) a 2.2 GHz con 8 GB di RAM:
\begin{itemize}
    \item \textbf{Host Nativo}: presenta una latenza contenuta rigorosamente entro i \SI{25}{\micro\second}, con una dispersione limitata attorno alla media.
    \item \textbf{Container Docker}: esibisce un comportamento identico a quello dell'host fisico, mantenendosi costantemente al di sotto dei \SI{25}{\micro\second}, in virtù del fatto che il processo del container condivide direttamente il medesimo scheduler CFS del kernel host senza intermediari.
    \item \textbf{KVM}: registra una degradazione considerevole, con picchi di latenza dispersi che salgono fino a \textbf{\SI{9000}{\micro\second}} (9 millisecondi). Questo comportamento irregolare è determinato dalle latenze introdotte dal meccanismo di \emph{vCPU preemption}, dai ritardi di iniezione degli interrupt virtualizzati da parte dell'hypervisor e dal tempo speso per transitare avanti e indietro tra la modalità di root e non-root (costo di salvataggio e ripristino dei registri della VM durante le VM-Exit).
\end{itemize}

---

\subsection{Calcolo Scientifico E Vincoli Di Ottimizzazione Hardware (Linpack)}

Un'ulteriore evidenza empirica emerge dall'analisi delle prestazioni di calcolo scientifico e intensivo, misurate mediante il benchmark \textbf{Linpack} (risoluzione di sistemi densi di equazioni lineari ad alte prestazioni, espresso in GigaFLOPS, ovvero miliardi di operazioni in virgola mobile per secondo):

% [IMMAGINE DEL FILE DI SLIDE CON NOME "CloudComputing_2"] Grafico a barre del benchmark Linpack GFLOPS per Native, Docker e KVM (IBM Research Report 2014) — Slide 31 / Pagina 31 del PDF

\begin{itemize}
    \item \textbf{Configurazione Bare-Metal}: raggiunge circa \SI{290}{\giga\text{FLOPS}}.
    \item \textbf{Container Docker}: ottiene lo stesso identico valore di picco, attestandosi a \SI{290}{\giga\text{FLOPS}}.
    \item \textbf{Macchina Virtuale KVM}: registra un crollo verticale delle prestazioni, arrestandosi a circa \SI{140}{\giga\text{FLOPS}} (un calo superiore al 50\%).
\end{itemize}

La spiegazione teorica di questo fenomeno risiede nei meccanismi interni di autotuning delle librerie matematiche avanzate (es. Intel MKL, BLAS, LAPACK). Tali librerie analizzano all'avvio la topologia dell'hardware mediante l'istruzione assembler \texttt{CPUID}, rilevando il modello esatto di core, la presenza di set di istruzioni vettoriali avanzate (AVX, AVX2, FMA) e la dimensione precisa della gerarchia delle memorie cache (L1, L2, L3) per suddividere le matrici in blocchi dimensionalmente allineati alle cache del silicio (\emph{cache tiling}).

All'interno di un container Docker, l'applicazione vede direttamente la CPU fisica reale dell'host e ne sfrutta tutte le estensioni e le caratteristiche architetturali. In un ambiente virtualizzato KVM, invece, l'hypervisor astrae e uniforma la topologia del processore per consentire la migrazione a caldo delle VM tra nodi eterogenei: di conseguenza, la libreria matematica non riceve i parametri precisi della gerarchia di memoria e ripiega su algoritmi di calcolo conservativi e generici, rinunciando alle ottimizzazioni e dimezzando il rendimento computazionale.

% BB ------------------------------------------------------------------------------------------------------------

% ==============================================================================
% SEZIONE 1: CREAZIONE DI IMMAGINI PERSONALIZZATE MEDIANTE DOCKERFILE
% ==============================================================================

\section{Creazione Di Immagini Personalizzate Mediante Dockerfile}

\subsection{Il Ciclo Di Vita Del Container E La Filosofia Del Singolo Task}

Per comprendere a fondo la logica con cui vengono assemblate le immagini Docker, è necessario rinfrescare una distinzione concettuale fondamentale rispetto all'ambiente operativo tradizionale. All'interno di un sistema operativo convenzionale o di una macchina virtuale, il kernel inizializza il processo capostipite (come \texttt{systemd} o \texttt{init}, caratterizzato da $\text{PID} = 1$), il quale procede ad attivare una moltitudine di demoni di sottofondo: server SSH, logger di sistema, pianificatori cron e sottosistemi di monitoraggio.

Al contrario, la tecnologia dei container Linux nasce per incapsulare un \textbf{singolo compito computazionale} (\emph{single unit task}) o un unico processo applicativo principale. Ne consegue una proprietà architetturale cardine:
\begin{theorem}[\textbf{Accoppiamento del Ciclo di Vita}]
Il ciclo di vita di un container coincide esattamente con la durata operativa del processo primario avente $\text{PID} = 1$ all'interno del proprio PID namespace. Nel momento esatto in cui tale processo termina la propria esecuzione (sia essa un'uscita con successo con codice di stato zero, un'eccezione non gestita o la ricezione di un segnale di terminazione standard come \texttt{SIGTERM} o \texttt{SIGKILL}), il container cessa immediatamente di esistere come entità attiva e passa nello stato di arresto (\emph{Exited}).
\end{theorem}

Un container non è progettato per fungere da contenitore statico di un intero sistema operativo general-purpose, bensì come un involucro isolato e riproducibile per l'esecuzione deterministica di un applicativo. Lo strumento formale impiegato per dichiarare la composizione di tale ambiente è il \textbf{Dockerfile}.

---

\subsection{Anatomia Di Un Dockerfile E Meccanismo Di Layering}

Il \textbf{Dockerfile} è un documento testuale contenente un elenco sequenziale e ordinato di istruzioni che definiscono come assemblare un'immagine Docker a partire da un contesto di build (\emph{build context}). Ciascuna istruzione viene interpretata ed eseguita automaticamente tramite il comando di compilazione dell'engine:
\begin{lstlisting}[language=bash]
# Sintassi base per l'avvio della compilazione nel contesto corrente
$ cd /path/to/app
$ docker build -t nome-immagine:tag .

# Varianti con percorsi espliciti o tramite il sottosistema avanzato Buildx
$ docker build ./
$ docker buildx build -t nome-immagine:tag ./
\end{lstlisting}

% [IMMAGINE DEL FILE DI SLIDE CON NOME "CloudComputing_2"] Esempio di Dockerfile per un'applicazione Indico su CentOS 7 affiancato dalle istruzioni di build CLI — Slide 1 / Pagina 28 del PDF

L'aspetto cardine dell'architettura di compilazione risiede nella struttura a \textbf{livelli sovrapposti} (\emph{layer stacking}):
\begin{itemize}
    \item Ogni istruzione modificativa presente nel Dockerfile (in particolar modo \texttt{RUN}, \texttt{COPY} e \texttt{ADD}) istanzia un container temporaneo intermedio, vi esegue il comando prescritto, congela le sole differenze binarie apportate al file system (\emph{delta changeset}) e le consolida (\emph{commit}) come un nuovo \textbf{layer di sola lettura} munito di un proprio hash crittografico univoco (SHA-256).
    \item Grazie alla sovrapposizione mediata da un file system ad unione (\emph{Union File System}, come il driver Linux \textbf{Overlay2}), i livelli sottostanti non modificati vengono preservati e riutilizzati tramite meccanismi di memorizzazione nella cache locale (\emph{build cache}), riducendo sia i tempi di ricompilazione che l'occupazione di banda e memoria secondaria.
\end{itemize}

Consideriamo il seguente esempio iniziale di Dockerfile orientato alla configurazione dell'applicativo web di conferenze scientifiche \emph{Indico} su base CentOS:
\begin{lstlisting}[language=dockerfile]
FROM centos:7
RUN yum install -y python-devel python-virtualenv
RUN virtualenv /opt/indico/venv
RUN pip install indico
COPY entrypoint.sh /opt/indico/entrypoint.sh
EXPOSE 8000
ENTRYPOINT /opt/indico/entrypoint.sh
\end{lstlisting}

---

\subsection{Sintassi E Direttive Fondamentali: FROM, ARG, WORKDIR Ed ENV}

La compilazione di una nuova immagine richiede la definizione chiara dell'albero di dipendenze:
\begin{itemize}
    \item \textbf{\texttt{FROM}}: costituisce l'istruzione obbligatoria con cui deve aprirsi qualsiasi Dockerfile funzionante. La clausola specifica l'\textbf{Immagine Genitore} (\emph{Parent Image} o \emph{Base Image}) da cui deriva il nuovo artefatto:
\begin{lstlisting}[language=dockerfile]
# Esempio su container Linux Microsoft Windows Server Core
FROM mcr.microsoft.com/windows/servercore
\end{lstlisting}
    Se l'immagine richiesta non è presente nel deposito locale dell'host, il motore Docker procede al recupero automatico contattando il registro remoto configurato. L'istruzione \texttt{FROM} può essere preceduta esclusivamente da commenti, da direttive di parsing del compilatore (es. \texttt{\# syntax=...}) o dalla direttiva \textbf{\texttt{ARG}} a portata globale, impiegata per parametrizzare la versione dell'immagine base:
\begin{lstlisting}[language=dockerfile]
ARG VERSION=3.9
FROM python:${VERSION}
\end{lstlisting}
    \item \textbf{\texttt{WORKDIR}}: imposta la directory di lavoro corrente all'interno dell'immagine per tutte le successive istruzioni \texttt{RUN}, \texttt{CMD}, \texttt{ENTRYPOINT}, \texttt{COPY} e \texttt{ADD}. Se il percorso specificato non esiste nell'albero delle cartelle del container, il compilatore provvede a crearlo automaticamente. L'uso di \texttt{WORKDIR} è considerata una buona norma architetturale poiché previene l'impiego dispersivo di comandi concatenati come \texttt{RUN cd /dir \&\& ...}.
    \item \textbf{\texttt{ENV}}: definisce variabili d'ambiente accessibili sia durante la fase di build dell'immagine sia a runtime all'interno del container. Tali variabili vengono interpretate ed espanse dinamicamente all'interno del Dockerfile da un sottoinsieme tassativo di istruzioni:
    \begin{center}
        \texttt{ADD}, \texttt{COPY}, \texttt{ENV}, \texttt{EXPOSE}, \texttt{FROM}, \texttt{LABEL}, \texttt{STOPSIGNAL}, \texttt{USER}, \texttt{VOLUME}, \texttt{WORKDIR}, \texttt{ONBUILD}
    \end{center}
\end{itemize}

Si consideri il seguente frammento esplicativo sull'interpretazione delle variabili d'ambiente:
\begin{lstlisting}[language=dockerfile]
FROM busybox
ENV FOO=/bar
WORKDIR ${FOO}          # Il compilatore espande in: WORKDIR /bar
ADD . $FOO              # Il compilatore espande in: ADD . /bar
COPY \$FOO /quux        # Carattere di escape: copia il file denominato "$FOO" in /quux
\end{lstlisting}

---

\subsection{Gestione Dei File Nel File System Del Container: COPY vs ADD}

Il trasferimento di file e cartelle dall'host di compilazione al file system dell'immagine può essere implementato tramite due istruzioni funzionalmente simili ma contraddistinte da requisiti di trasparenza operativi distinti:

\begin{figure}[H]
    \centering
    \begin{tikzpicture}[node distance=2cm, auto]
        \node [draw, rectangle, fill=theoryblue!20, minimum width=3.5cm, minimum height=1cm, align=center] (host) {\textbf{Host di Build}\\(Directory locale/Context)};
        \node [draw, rectangle, fill=darkgreen!20, minimum width=3.5cm, minimum height=1cm, align=center, right of=host, xshift=4cm] (image) {\textbf{Container Image}\\(File System a livelli)};

        \draw [->, thick] (host.north east) -- node[above, font=\footnotesize] {\texttt{COPY <src> <dest>}} (image.north west);
        \draw [->, thick, dashed] (host.south east) -- node[below, font=\footnotesize] {\texttt{ADD <tar.gz|url> <dest>}} (image.south west);
    \end{tikzpicture}
    \caption{Canali di inserimento file nel file system dell'immagine: \texttt{COPY} trasparente vs \texttt{ADD} con scompattazione/download implicito.}
    \label{fig:copy_vs_add}
\end{figure}

\begin{enumerate}
    \item \textbf{L'istruzione \texttt{COPY}}: è la direttiva di elezione per il trasferimento standard di file e cartelle. Presenta una sintassi esplicita e controllata:
\begin{lstlisting}[language=dockerfile]
COPY ["<sorgente>", "<destinazione>"]

# Sintassi per ambienti Linux e percorsi Windows
COPY test1.txt /tmp/
COPY test1.txt c:/temp/
\end{lstlisting}
    La sorgente deve risiedere rigorosamente all'interno del contesto di build (la directory passata come argomento a \texttt{docker build}), impedendo qualsiasi tentativo di accedere a percorsi assoluti esterni come \texttt{COPY ../../../etc/shadow .} per ragioni di sicurezza e riproducibilità.
    \item \textbf{L'istruzione \texttt{ADD}}: include due comportamenti impliciti di secondo livello che ne rendono sconsigliato l'uso generico:
    \begin{itemize}
        \item Se la sorgente è un archivio compresso locale riconosciuto in un formato standard (\texttt{tar}, \texttt{tar.gz}, \texttt{tgz}, \texttt{tar.bz2}, \texttt{tar.xz}), l'istruzione provvede ad \textbf{estrarlo automaticamente} all'interno della directory di destinazione nel container.
        \item Consente l'inserimento di URL remoti come sorgente, scaricando il payload direttamente all'interno dell'immagine (sebbene tale pratica sia considerata un anti-pattern, in quanto genera layer non ottimizzati che non puliscono i file temporanei).
    \end{itemize}
\end{enumerate}

La linea guida ingegneristica standard stabilisce pertanto che l'uso di \texttt{COPY} sia sempre preferibile in virtù della sua natura trasparente; l'utilizzo di \texttt{ADD} rimane giustificato unicamente per l'iniezione con estrazione atomica di archivi compressi locali alla radice del file system, come nel classico pattern di assemblaggio di una rootfs minimale:
\begin{lstlisting}[language=dockerfile]
ADD rootfs.tar.xz /
\end{lstlisting}

---

\subsection{Esecuzione Di Comandi: RUN Nella Forma Shell Ed Exec}

L'istruzione \textbf{\texttt{RUN}} è il costrutto fondamentale impiegato per configurare il file system dell'immagine, installando pacchetti applicativi, compilando librerie e instradando configurazioni software. Il risultato di ciascuna istruzione \texttt{RUN} viene consolidato in un nuovo layer dell'immagine.

L'istruzione può essere formulata in due varianti sintattiche profonde:

\subsubsection{1. Forma Shell (\emph{Shell Form})}
\begin{lstlisting}[language=dockerfile]
RUN <comando>
\end{lstlisting}
In questa modalità, il comando specificato viene invocato come parametro di una shell di sistema sottostante, che per impostazione predefinita corrisponde a \texttt{/bin/sh -c} sulle distribuzioni Linux e a \texttt{cmd /S /C} sui sistemi Windows. Di conseguenza, il processo avviato dal kernel Linux non sarà il comando finale bensì la shell stessa:
\begin{center}
    \texttt{/bin/sh -c "apt-get update \&\& apt-get install -y curl"}
\end{center}
La presenza della shell garantisce l'espansione automatica delle variabili d'ambiente (es. \texttt{\$HOME}, \texttt{\$PATH}), l'interpretazione dei caratteri jolly (\emph{wildcard}) e l'uso di operatori di concatenazione (\texttt{\&\&}, \texttt{||}, \texttt{|}).

\subsubsection{2. Forma Exec (\emph{Exec Form})}
\begin{lstlisting}[language=dockerfile]
RUN ["<eseguibile>", "<parametro 1>", "<parametro 2>"]
\end{lstlisting}
Questa formulazione, analizzata come un array JSON rigoroso racchiuso tra parentesi quadre e virgolette doppie, invoca l'eseguibile sul kernel mediante la chiamata di sistema primordiale di esecuzione (\texttt{execve}), \textbf{senza richiamare alcuna shell di sistema intermediaria}.
Ne discende che:
\begin{itemize}
    \item Non si verifica alcuna sostituzione o espansione automatica delle variabili d'ambiente. Ad esempio, la sequenza:
\begin{lstlisting}[language=dockerfile]
RUN ["echo", "$HOME"]
\end{lstlisting}
    non stamperà il percorso della cartella home dell'utente, bensì produrrà letteralmente la stringa testuale \texttt{"\$HOME"}.
    \item Qualora sia necessaria l'elaborazione della shell mantenendo la sintassi exec, è obbligatorio invocare la shell in modo esplicito come binario primario:
\begin{lstlisting}[language=dockerfile]
RUN ["/bin/sh", "-c", "echo $HOME"]
\end{lstlisting}
    \item Tale meccanismo consente inoltre di scavalcare la shell minimale predefinita specificando interpreti alternativi dotati di feature avanzate, come ad esempio la shell Bash:
\begin{lstlisting}[language=dockerfile]
RUN ["/bin/bash", "-c", "echo hello"]
\end{lstlisting}
\end{itemize}

---

\subsection{Configurazione Dell'Esecuzione: Il Dualismo Tra ENTRYPOINT E CMD}

Mentre l'istruzione \texttt{RUN} opera esclusivamente a tempo di compilazione (\emph{build time}), definendo i layer statici del file system, le istruzioni \textbf{\texttt{ENTRYPOINT}} e \textbf{\texttt{CMD}} operano a tempo di esecuzione (\emph{runtime}), stabilendo quale processo debba essere avviato quando viene creato un container a partire dall'immagine. Sebbene condividano la possibilità di essere espresse sia in \emph{exec form} (raccomandata) che in \emph{shell form}, esse svolgono ruoli architetturali distinti:

\begin{table}[H]
    \centering
    \small
    \begin{tabularx}{\textwidth}{p{3cm} X X}
        \toprule
        \textbf{Istruzione} & \textbf{Ruolo Architetturale} & \textbf{Comportamento all'avvio con \texttt{docker run}} \\
        \midrule
        \textbf{\texttt{CMD}} & Definisce il comando predefinito e/o i parametri opzionali di fallback per il container. & Se l'utente specifica qualsiasi argomento a riga di comando (es. \texttt{docker run img bash}), l'intera istruzione \texttt{CMD} viene sovrascritta e ignorata. \\
        \midrule
        \textbf{\texttt{ENTRYPOINT}} & Fissa l'eseguibile immutabile del container, trasformando l'immagine in un binario eseguibile a riga di comando. & Non viene sovrascritto dagli argomenti passati da riga di comando; al contrario, i parametri immessi dall'utente vengono accodati (\emph{appended}) come argomenti aggiuntivi all'eseguibile. \\
        \bottomrule
    \end{tabularx}
    \caption{Confronto operativo tra le direttive runtime \texttt{CMD} ed \texttt{ENTRYPOINT}.}
    \label{tab:cmd_vs_entrypoint}
\end{table}

Entrambe le istruzioni sottostanno alla regola secondo cui, qualora nel Dockerfile siano presenti molteplici dichiarazioni di \texttt{CMD} o di \texttt{ENTRYPOINT}, \textbf{unicamente l'ultima occorrenza dichiarata avrà effetto}, mentre tutte le precedenti verranno scartate dall'interprete.

La potenza di questi strumenti si esprime nell'uso combinato:
\begin{lstlisting}[language=dockerfile]
ENTRYPOINT ["/opt/indico/entrypoint.sh"]
CMD ["run", "--host=0.0.0.0"]
\end{lstlisting}
In questo assetto:
\begin{itemize}
    \item L'invocazione minimale standard \texttt{docker run indico-app} comporterà l'esecuzione del comando:
\begin{center}
    \texttt{/opt/indico/entrypoint.sh run --host=0.0.0.0}
\end{center}
    \item L'invocazione esplicita con parametri aggiuntivi \texttt{docker run indico-app worker} manterrà inalterato l'eseguibile prioritario definito in \texttt{ENTRYPOINT}, sostituendo la clausola \texttt{CMD} e generando il comando finale:
\begin{center}
    \texttt{/opt/indico/entrypoint.sh worker}
\end{center}
\end{itemize}

È fondamentale evidenziare che l'impiego della \emph{shell form} per \texttt{ENTRYPOINT} (es. \texttt{ENTRYPOINT /path/bin}) costituisce una prassi sconsigliata: avviando l'eseguibile come figlio di \texttt{/bin/sh}, quest'ultima non inoltrerà per default i segnali POSIX (\texttt{SIGTERM}) al processo applicativo, impedendo un corretto arresto controllato (\emph{graceful shutdown}) del container. La \emph{exec form} assicura invece che il processo applicativo acquisisca direttamente il ruolo di $\text{PID} = 1$.

---

\subsection{Esposizione Di Porte Di Rete: La Direttiva EXPOSE}

L'istruzione \textbf{\texttt{EXPOSE}} documenta su quali porte di rete e con quali protocolli di livello trasporto il container accetta connessioni in ingresso a runtime.
La sintassi prevede la dichiarazione della porta associata facoltativamente al protocollo, con TCP impostato come scelta predefinita:
\begin{lstlisting}[language=dockerfile]
EXPOSE 80/tcp
EXPOSE 80/udp
\end{lstlisting}

Dal punto di vista architetturale e sistemistico, è opportuno chiarire che l'istruzione \texttt{EXPOSE} funge essenzialmente da \textbf{metadato documentativo}: essa non effettua alcuna apertura o binding effettivo di porte sull'interfaccia di rete dell'host fisico, né configura regole nel firewall di sistema. Affinché una porta del container risulti accessibile dall'esterno della macchina host, l'operatore deve associare esplicitamente la porta a tempo di avvio tramite l'opzione di inoltro (\emph{port forwarding}) \texttt{-p} del comando \texttt{docker run} (ad es. \texttt{-p 8080:80}).

\newpage

% ==============================================================================
% SEZIONE 2: PERSISTENZA DEI DATI IN DOCKER: VOLUMI, BIND MOUNTS E TMPFS
% ==============================================================================

\section{Persistenza Dei Dati In Docker: Volumi, Bind Mounts E Tmpfs}

\subsection{Limiti Del File System Di Scrittura Effimero (Copy-on-Write)}

Come analizzato precedentemente, un container in esecuzione sovrappone ai livelli statici di sola lettura dell'immagine un sottile strato scrivibile (\textbf{Read-Write Container Layer}). La scrittura su tale strato è governata dal paradigma \textbf{Copy-on-Write (CoW)}: se un processo tenta di modificare un blocco preesistente appartenente a un layer inferiore, il file system intercetta l'operazione di I/O, clona il file originale e ne memorizza le variazioni esclusivamente nel layer superiore modificabile.

Sebbene tale architettura garantisca l'isolamento e la rapida creazione di nuove istanze, essa introduce severe criticità ingegneristiche quando applicata a carichi con persistenza di stato (\emph{stateful workloads}):
\begin{itemize}
    \item \textbf{Natura effimera}: nel momento in cui il container viene distrutto (mediante \texttt{docker rm}), l'intero layer scrivibile associato viene deallocato e rimosso dalla memoria secondaria, causando la perdita irreparabile di tutti i dati applicativi ivi contenuti.
    \item \textbf{Penalità di I/O}: la mediazione dell'algoritmo di Copy-on-Write e dei driver di unione (come Overlay2) comporta un sovraccarico non trascurabile in scrittura, riducendo le prestazioni delle basi di dati ad alta frequenza di transazione.
    \item \textbf{Crescita delle dimensioni dell'immagine}: i file temporanei o i log accumulati nel layer di scrittura incrementano la dimensione complessiva del container su disco, ostacolandone la migrazione e il monitoraggio.
\end{itemize}

---

\subsection{Tassonomia Dello Storage: Named Volumes, Bind Mounts E Tmpfs Mounts}

Per svincolare la persistenza dei dati dal ciclo di vita del container, Docker mette a disposizione tre distinti meccanismi di montaggio:

% [IMMAGINE DEL FILE DI SLIDE CON NOME "CloudComputing_2"] Diagramma strutturale dei meccanismi di montaggio dello storage in Docker: Bind Mount, Named Volume e Tmpfs Mount — Slide 13 / Pagina 40 del PDF

\begin{figure}[H]
    \centering
    \begin{tikzpicture}[>=Stealth, node distance=1.5cm]
        % Disegno Host
        \draw[thick, fill=backcolour] (0,0) rectangle (12, 5.5);
        \node[anchor=north west, font=\bfseries] at (0.2, 5.3) {Macchina Host};

        % File system host
        \draw[thick, fill=white] (0.5, 0.5) rectangle (5.5, 4.5);
        \node[anchor=north west, font=\small\bfseries] at (0.7, 4.3) {File System Host};

        % Docker Storage Area
        \draw[thick, dashed, fill=orange!15] (1.0, 0.8) rectangle (5.0, 2.5);
        \node[font=\footnotesize\bfseries, align=center] at (3.0, 1.65) {Area Riservata a Docker\\(\texttt{/var/lib/docker/volumes/})};

        % Percorso arbitrario
        \draw[thick, fill=gray!20] (1.0, 2.8) rectangle (5.0, 3.8);
        \node[font=\footnotesize\bfseries, align=center] at (3.0, 3.3) {Percorso Arbitrario Host\\(es. \texttt{/opt/data} o \texttt{C:\textbackslash dev})};

        % Memoria Host
        \draw[thick, fill=white] (6.5, 0.5) rectangle (11.5, 2.5);
        \node[font=\small\bfseries] at (9.0, 1.5) {Memoria RAM di Sistema};

        % Container
        \draw[thick, fill=theoryblue!20] (6.5, 3.0) rectangle (11.5, 4.8);
        \node[font=\bfseries] at (9.0, 3.9) {Container};
        \node[font=\footnotesize] at (9.0, 3.4) {Punto di Montaggio Target (\texttt{/data})};

        % Frecce di montaggio
        \draw[<->, thick, theoryblue] (3.0, 2.5) to [bend right=20] node[left, font=\scriptsize\bfseries] {Volume} (7.5, 3.0);
        \draw[<->, thick, darkgreen] (3.0, 3.8) to [bend left=15] node[above, font=\scriptsize\bfseries] {Bind Mount} (6.5, 4.2);
        \draw[<->, thick, red] (9.0, 2.5) -- node[right, font=\scriptsize\bfseries] {tmpfs Mount} (9.0, 3.0);
    \end{tikzpicture}
    \caption{Architettura e allocazione delle aree di persistenza dello storage in ambiente Docker.}
    \label{fig:storage_architectures}
\end{figure}

\begin{enumerate}
    \item \textbf{I Volumi Gestiti (\emph{Named Volumes})}: costituiscono il meccanismo standard raccomandato per la persistenza. I dati vengono allocati all'interno di una porzione dedicata e segregata del file system dell'host gestita esclusivamente dall'engine Docker (su Linux posizionata in \texttt{/var/lib/docker/volumes/}).
    I vantaggi strutturali dei volumi comprendono:
    \begin{itemize}
        \item \emph{Indipendenza e ciclo di vita autonomo}: il contenuto di un volume sopravvive intatto all'arresto e all'eliminazione del container che lo impiega, e non incrementa le dimensioni del layer scrivibile del container.
        \item \emph{Gestione centralizzata tramite CLI/API}: i volumi costituiscono entità autonome gestibili nativamente tramite comandi dedicati (\texttt{docker volume create}, \texttt{docker volume ls}, \texttt{docker volume rm}).
        \item \emph{Popolamento automatico iniziale}: se un container monta un nuovo volume vuoto su un percorso interno contenente file preesistenti (ad es. file di configurazione predefiniti forniti dall'immagine), l'engine provvede a clonare tali dati all'interno del volume prima del rilascio.
        \item \emph{Portabilità multi-piattaforma}: operano con medesima semantica sia su container Linux che su ambienti Windows, ed esibiscono prestazioni di gran lunga superiori rispetto ai bind mounts in ambienti virtualizzati di sviluppo come Docker Desktop per macOS o Windows.
        \item \emph{Estendibilità mediante Volume Drivers}: supportano plugin di terze parti per instradare lo storage su file system distribuiti di rete (NFS, Ceph) o piattaforme di cloud storage con cifratura integrata trasparente.
    \end{itemize}
    \item \textbf{I Montaggi Diretti (\emph{Bind Mounts})}: permettono di mappare direttamente un file o una directory arbitraria dell'host all'interno del file system del container, facendovi riferimento mediante il percorso assoluto della macchina host.
    \begin{itemize}
        \item \emph{Caratteristiche e limiti}: sebbene storicamente molto performanti, i bind mount dipendono rigidamente dalla struttura delle cartelle e dal sistema operativo dell'host ospitante, minando la riproducibilità architetturale. Inoltre, non possono essere monitorati o manipolati direttamente tramite la CLI di Docker. Risultano tuttavia ideali in fase di sviluppo per iniettare codice sorgente aggiornato in tempo reale all'interno del container senza dover rieseguire la build.
    \end{itemize}
    \item \textbf{I Montaggi in Memoria (\emph{tmpfs Mounts})}: riservati a scritture temporanee volatili, mappano una porzione della memoria RAM dell'host all'interno del container. I dati non transitano né sul disco host né sul layer scrivibile Copy-on-Write, garantendo le massime prestazioni possibili per l'elaborazione di chiavi di cifratura effimere, file temporanei ad alta frequenza o segreti di processo che non devono lasciare traccia persistente.
\end{enumerate}

---

\subsection{Analisi Delle Strategie Di Persistenza: Il Caso Di Studio Di Redis}

L'applicazione pratica delle strategie di gestione dello stato emerge nell'analisi di un database in-memory ad alte prestazioni come **Redis** (\emph{Remote Dictionary Server}). In assenza di configurazioni correttive, Redis conserva le strutture dati esclusivamente nella memoria RAM allocata dal processo. Per garantire tolleranza ai guasti, il motore offre diverse modalità operative di salvataggio su storage secondario persistente:

\begin{enumerate}
    \item \textbf{RDB (\emph{Redis Database Backup})}: effettua istantanee puntuali a intervalli temporali definiti (\emph{point-in-time snapshots}) salvando l'intero dataset compresso su un file binario su disco (\texttt{dump.rdb}). Presenta un'eccellente compattezza e rapidità di ripristino all'avvio, ma introduce il rischio di perdita dei dati elaborati tra un intervallo di snapshotting e l'altro in caso di arresto anomalo.
    \item \textbf{AOF (\emph{Append Only File})}: registra in modo incrementale ogni singola operazione di scrittura ricevuta dal server all'interno di un file di log testuale, formattato secondo il protocollo nativo Redis. All'avvio, il motore esegue nuovamente l'intera sequenza di comandi, ricostruendo lo stato esatto del dataset. Redis esegue algoritmi in background (\emph{AOF Rewrite}) per compattare il file qualora le dimensioni crescano oltre soglia. L'approccio assicura massima durabilità ma comporta un footprint di I/O maggiore.
    \item \textbf{Modalità a Persistenza Disabilitata (\emph{No Persistence})}: elimina completamente la scrittura su memoria secondaria; le strutture dati sopravvivono unicamente per la durata del ciclo di vita del server.
    \item \textbf{Assetto Combinato (\emph{RDB + AOF})}: integra entrambe le strategie nella medesima istanza. In fase di riavvio, il motore caricherà obbligatoriamente il file AOF, poiché esso garantisce intrinsecamente una ricostruzione più aggiornata e completa del dataset rispetto allo snapshot RDB.
\end{enumerate}

Nel caso in cui Redis venga impiegato in qualità di pura cache transitoria a bassa latenza, è comune disabilitare la persistenza configurando esplicitamente i flag di avvio all'interno di un file di orchestrazione:
\begin{lstlisting}[language=bash]
command: redis-server --save "" --appendonly no
\end{lstlisting}
La stringa vuota passata al parametro \texttt{--save ""} inibisce l'allocazione automatica di snapshot periodici RDB, mentre \texttt{--appendonly no} disattiva la generazione del registro cumulativo AOF, garantendo che l'istanza operi in uno stato puramente volatile senza gravare sul sottosistema di storage del nodo host.

\newpage

% ==============================================================================
% SEZIONE 3: ARCHITETTURA DI RETE IN DOCKER
% ==============================================================================

\section{Architettura Di Rete In Docker}

\subsection{Il Sottosistema Di Rete Pluggabile E Il Modello CNM/Drivers}

Uno dei fattori abilitanti dell'architettura Docker è la capacità di interconnettere carichi applicativi distribuiti garantendo isolamento o integrazione trasparente con le reti fisiche preesistenti. L'infrastruttura di rete di Docker è completamente agnostica rispetto al sistema operativo dell'host ed è concepita secondo un modello modulare e innestabile (\emph{pluggable}), denominato **Container Network Model (CNM)**.

All'interno del CNM, l'engine astrae la connettività per mezzo di driver intercambiabili specializzati, sollevando l'applicazione dall'onere di conoscere la topologia sottostante: due container possono cooperare tramite socket di rete indipendentemente dal fatto che si trovino sullo stesso nodo o su cluster eterogenei distribuiti.

---

\subsection{Tassonomia Dei Driver Di Rete}

Docker mette a disposizione un catalogo nativo di driver di rete specializzati:

\begin{itemize}
    \item \textbf{\texttt{bridge}}: è il driver di rete predefinito. Quando un container viene avviato senza parametri di rete espliciti, esso viene automaticamente collegato a una rete di tipo bridge. I container appartenenti alla medesima rete bridge possono scambiarsi pacchetti IP in un dominio di broadcast isolato, pur rimanendo protetti rispetto all'esterno.
    \item \textbf{\texttt{host}}: rimuove la barriera di isolamento di rete tra il container e il sistema operativo host. Il container non acquisisce un proprio stack IP o una propria tabella di routing separata all'interno di un network namespace, bensì effettua il binding dei propri servizi direttamente sulle porte e sulle interfacce fisiche dell'host. Se un'applicazione containerizzata con driver host apre la porta 80, essa risulterà accessibile direttamente all'indirizzo IP fisico del server senza passare da alcun livello di port mapping.
    \item \textbf{\texttt{overlay}}: consente di creare una rete virtuale cifrata a livello 2/3 sovrapposta (\emph{overlay network}) che estende il piano di connettività attraverso molteplici daemon Docker residenti su host fisici distinti. Costituisce la spina dorsale per la comunicazione scalabile tra servizi distribuiti in cluster **Docker Swarm**, eliminando la necessità di configurare instradamenti statici a livello di router di sistema operativo.
    \item \textbf{\texttt{macvlan}}: permette di assegnare a ciascun container un indirizzo fisico MAC univoco dedicato, facendolo apparire alle apparecchiature di rete (switch, router fisici) come un calcolatore fisico reale direttamente collegato alla LAN. Il driver instrada il traffico basandosi sulla commutazione di livello 2; rappresenta la scelta primaria per migrare applicazioni legacy o macchine virtuali vincolate alla presenza di un indirizzo IP appartenente alla medesima subnet aziendale non instradata.
    \item \textbf{\texttt{ipvlan}}: conferisce all'amministratore un controllo granulare sia sugli schemi di indirizzamento IPv4 che IPv6. Rispetto a macvlan, consente di operare sia a livello 2 (gestendo i tag 802.1Q delle VLAN di trunking) sia a livello 3 (IPvlan L3), demandando l'instradamento direttamente all'infrastruttura fisica sottostante (\emph{underlay network}) ed evitando i limiti imposti dal filtraggio di indirizzi MAC multipli sulle porte degli switch fisici.
    \item \textbf{\texttt{none}}: disabilita completamente lo stack di rete per il container. All'interno del rispettivo network namespace viene inizializzata unicamente l'interfaccia di loopback (\texttt{lo}), escludendo qualsiasi scheda virtuale. Risulta idoneo per l'esecuzione di calcoli sensibili o batch di compressione batch completamente segregati.
    \item \textbf{Plugin di rete di terze parti}: interfacce esterne fornite da produttori o dalla community (es. Calico, Weave Net, Cisco ACI) per integrare l'ecosistema Docker all'interno di complesse architetture Software-Defined Networking (SDN).
\end{itemize}

---

\subsection{Rete Bridge Utente E Meccanismi Di Port Mapping (Iptables E NAT)}

L'architettura basata su bridge implementa la connettività locale tramite componenti virtuali del kernel Linux:

% [IMMAGINE DEL FILE DI SLIDE CON NOME "CloudComputing_2"] Topologia di rete Bridge: molteplici reti bridge isolate (bridgenet1, bridgenet2, bridgenet3) su Docker Host differenti — Slide 17 / Pagina 44 del PDF

L'amministratore può creare una o più reti bridge isolate tramite CLI:
\begin{lstlisting}[language=bash]
$ docker network create -d bridge --subnet 10.0.0.0/24 --gateway 10.0.0.1 bridgenet1
\end{lstlisting}
L'impiego di una rete **bridge definita dall'utente** (\emph{user-defined bridge}) offre un vantaggio chiave rispetto alla rete bridge predefinita (\texttt{docker0}): integra un servizio di \textbf{risoluzione DNS automatico}. Due o più container collegati alla medesima rete utente possono comunicare semplicemente invocando i reciproci nomi di container (es. \texttt{ping redis-db}), senza richiedere l'assegnazione o la manipolazione manuale di indirizzi IP statici.

% [IMMAGINE DEL FILE DI SLIDE CON NOME "CloudComputing_2"] Schema del Port Mapping e della traslazione NAT tramite bridge Linux tra porta host (8080) e porta container (80) — Slide 18 / Pagina 45 del PDF

Il meccanismo di inoltro del traffico tra la rete fisica e l'ambiente isolato del container è illustrato nella configurazione di **Port Mapping**:
\begin{lstlisting}[language=bash]
$ docker container run -d -p 8080:80 --name webserver nginx
\end{lstlisting}

La dinamica esecutiva è regolata dai seguenti passaggi:
\begin{enumerate}
    \item Il container riceve un indirizzo IP privato non instradabile all'esterno (es. \texttt{10.0.0.8}) assegnato all'interfaccia virtuale collegata al bridge interno.
    \item La macchina host espone una propria interfaccia fisica dotata di un indirizzo instradabile sulla LAN o su Internet (es. \texttt{172.14.3.55}).
    \item L'engine Docker modifica automaticamente le tabelle di filtraggio del kernel dell'host (\textbf{\texttt{iptables}} in ambiente Linux) aggiungendo una regola di traslazione di indirizzo a destinazione (\textbf{Destination NAT}, DNAT) nella catena di routing:
\begin{center}
    Pacchetto in arrivo: \texttt{172.14.3.55:8080} $\longrightarrow$ Riscrittura DNAT $\longrightarrow$ Inoltro a: \texttt{10.0.0.8:80}
\end{center}
    \item I pacchetti di risposta emessi dal container subiscono la traslazione simmetrica di indirizzo sorgente (\textbf{Source NAT}, SNAT/Masquerade), consentendo la corretta ricezione da parte del client mittente originario.
\end{enumerate}

La flessibilità della clausola \texttt{-p} consente configurazioni mirate:
\begin{itemize}
    \item \texttt{-p 8080:80}: associa la porta TCP 80 del container alla porta TCP 8080 di tutte le interfacce disponibili dell'host (\texttt{0.0.0.0}).
    \item \texttt{-p 192.168.1.100:8080:80}: vincola l'apertura della porta 8080 esclusivamente all'indirizzo IP host specificato (\texttt{192.168.1.100}), impedendo l'accesso da altre interfacce di rete del server.
    \item \texttt{-p 8080:80/udp}: definisce esplicitamente il protocollo di trasporto UDP al posto del default TCP.
\end{itemize}

\newpage

% ==============================================================================
% SEZIONE 4: STRUMENTI DI GESTIONE E ORCHESTRAZIONE: DA DOCKER CLI A DOCKER COMPOSE E KUBERNETES
% ==============================================================================

\section{Strumenti Di Gestione E Orchestrazione: Da Docker CLI A Docker Compose E Kubernetes}

\subsection{Comandi Fondamentali Della Docker CLI Per Immagini E Container}

L'interazione ordinaria con l'infrastruttura Docker avviene tramite la sintassi strutturata della CLI, che raggruppa i comandi in oggetti logici:

\begin{lstlisting}[language=bash]
# 1. Recupero di un'immagine da registro remoto
$ docker image pull node:latest

# 2. Elenco delle immagini disponibili nella memoria locale
$ docker image ls

# 3. Avvio di un container in background (-d) con port mapping (-p) e nome assegnato
$ docker container run -d -p 5000:5000 --name node node:latest

# 4. Monitoraggio dei container in esecuzione attiva
$ docker container ps

# 5. Invio del segnale di arresto al processo principale (SIGTERM / SIGKILL)
$ docker container stop node

# 6. Rimozione definitiva del container dallo storage dell'host
$ docker container rm node

# 7. Rimozione dell'immagine locale
$ docker image rmi <image id>

# 8. Compilazione di una nuova release dell'immagine
$ docker build -t node:2.0 .
# In alternativa con il sottosistema Buildx multi-architettura:
$ docker buildx build -t node:2.0 .

# 9. Pubblicazione dell'immagine compilata su registro remoto
$ docker image push node:2.0

# 10. Manuale in linea dei comandi e sintassi
$ docker --help
\end{lstlisting}

---

\subsection{Limiti Dell'Approccio Manuale E Nascita Di Docker Compose}

Nella progettazione di un sistema moderno, l'applicazione non è quasi mai un blocco monolitico, bensì un insieme cooperante di servizi distribuiti (frontend web, API REST, code di messaggi, database relazionali, cache in-memory).
Gestire una simile architettura mediante comandi elementari da riga di comando (approccio \emph{naive}) introduce severe criticità operative:
\begin{itemize}
    \item \textbf{Mancanza di orchestrazione delle dipendenze}: ciascun container deve essere compilato, configurato e avviato singolarmente con una complessa sequenza di comandi manuali, rispettando rigidamente l'ordine temporale dei servizi (ad es. un backend non può partire finché il relativo database non è pronto all'ascolto).
    \item \textbf{Cablaggio manuale delle reti}: le reti virtuali e i volumi devono essere pre-allocati separatamente, e i parametri di connessione devono essere passati a mano tramite flag lunghi e soggetti a errori di digitazione.
\end{itemize}

% [IMMAGINE DEL FILE DI SLIDE CON NOME "CloudComputing_2"] Confronto concettuale tra l'approccio Naive (gestione manuale container per container) e Docker Compose (orchestrazione dichiarativa unificata tramite compose.yml) — Slide 8 / Pagina 35 del PDF

**Docker Compose** risolve queste inefficienze offrendo uno strumento dichiarativo per descrivere e governare applicazioni multi-container all'interno di un unico file di configurazione formattato in sintassi YAML (\texttt{docker-compose.yml}).

Il ciclo di vita con Docker Compose si standardizza in un flusso a tre passi:
\begin{enumerate}
    \item Definizione dell'ambiente applicativo di ciascun componente mediante il rispettivo \textbf{Dockerfile}, garantendo la replicabilità delle dipendenze binarie.
    \item Descrizione organica di tutti i servizi, delle porte da mappare, delle reti e dei volumi all'interno del file \textbf{\texttt{docker-compose.yml}}.
    \item Avvio atomico dell'intero ecosistema mediante un singolo comando:
\begin{lstlisting}[language=bash]
$ docker compose up      # Versione moderna integrata nel client Docker
$ docker-compose up      # Eseguibile binario Python storico della v1
$ docker compose down    # Arresto ordinato e rimozione di container e reti
\end{lstlisting}
\end{enumerate}

---

\subsection{Struttura E Sintassi Di Un File docker-compose.yml}

Un file Compose specifica la composizione architetturale dei servizi cooperanti. Si consideri il seguente esempio basato sulla specifica versione 2 che coordina un'applicazione web Angular, un backend Node/Express e un database documentale MongoDB:

\begin{lstlisting}[language=yaml]
version: '2'
services:
  angular:
    build: client                # Directory contenente il Dockerfile del frontend
    ports:
      - "4200:4200"              # Esposizione dell'interfaccia utente

  express:
    build: api                   # Directory contenente il Dockerfile del backend
    ports:
      - "3977:3977"              # Esposizione dell'endpoint REST API

  database:
    image: mongo                 # Recupero dell'immagine ufficiale da Docker Hub
    ports:
      - "27017:27017"            # Esposizione della porta nativa di MongoDB
\end{lstlisting}

% [IMMAGINE DEL FILE DI SLIDE CON NOME "CloudComputing_2"] Struttura dichiarativa di un file compose.yml con definizione dei servizi Angular, Express e Database Mongo — Slide 9 / Pagina 36 del PDF

In un secondo scenario architetturale, l'applicazione richiede il supporto congiunto di un database relazionale PostgreSQL e di un'istanza Redis configurata con persistenza disabilitata:

\begin{lstlisting}[language=yaml]
version: "2"
services:
  my-application:
    build: ./                    # Compilazione del Dockerfile nella directory corrente
    ports:
      - "8000:8000"
    environment:
      - CONFIG_FILE              # Iniezione di variabili d'ambiente di configurazione

  db:
    image: postgres              # Istanza del database PostgreSQL

  redis:
    image: redis
    command: redis-server --save "" --appendonly no  # Disabilitazione persistenza
    ports:
      - "6379"                   # Apertura della porta di comunicazione interna
\end{lstlisting}

---

\subsection{Cenni Di Orchestrazione Su Scala: Docker Swarm vs Kubernetes}

Nel momento in cui l'architettura applicativa richiede di scalare orizzontalmente su decine o centinaia di nodi fisici per gestire picchi di carico e garantire tolleranza ai guasti, l'impiego di un singolo motore Docker o di Docker Compose non è più sufficiente. Emerge quindi la necessità di piattaforme di **Orchestrazione di Container**:

% [IMMAGINE DEL FILE DI SLIDE CON NOME "CloudComputing_2"] Schema architetturale a microservizi di riferimento eShopOnContainers di Microsoft (Client apps, Docker host, Microservizi con database autonomi, Event Bus e API Gateway) — Slide 21 / Pagina 48 del PDF

Il panorama dell'orchestrazione su cluster è dominato da due soluzioni:
\begin{enumerate}
    \item \textbf{Docker Swarm}: la soluzione di clustering nativa realizzata e manutenuta direttamente da Docker.
    \begin{itemize}
        \item \emph{Vantaggi}: estrema semplicità di installazione e configurazione iniziale; riutilizza la medesima sintassi dichiarativa della CLI Docker e dei file Compose; integra funzionalità native di bilanciamento del carico interno sui servizi.
        \item \emph{Limiti}: scarsa flessibilità per flussi di deployment complessi e politiche avanzate di schedulazione.
    \end{itemize}
    \item \textbf{Kubernetes (K8s)}: la piattaforma open source di riferimento industriale, sviluppata inizialmente da Google per orchestrare le proprie infrastrutture globali e poi donata alla \emph{Cloud Native Computing Foundation} (CNCF).
    \begin{itemize}
        \item \emph{Architettura}: adotta un modello a nodi gerarchici diviso tra piano di controllo (\emph{Control Plane} o \emph{Master}) e nodi operativi (\emph{Worker Nodes}), articolando i carichi di lavoro in costrutti modulari quali **Pod** (la più piccola unità di esecuzione, composta da uno o più container strettamente accoppiati), **Namespaces**, **Services**, **ConfigMaps** e controller di stato.
        \item \emph{Vantaggi competitivi}: vastissima comunità open-source; supporto per qualsiasi sistema operativo e hardware; capacità intrinseche di scalabilità automatica (\emph{Horizontal Pod Autoscaling}) e auto-riparazione (\emph{Self-Healing}, tramite riavvio o rischedulazione dei Pod guasti); monitoraggio e telemetria integrabili con ampi ecosistemi; disponibilità di offerte gestite native su tutti i principali provider cloud (\textbf{Google GKE}, \textbf{Microsoft AKS}, \textbf{Amazon EKS}).
    \end{itemize}
\end{enumerate}

\newpage

% ==============================================================================
% SEZIONE 5: CASO DI STUDIO APPLICATIVO: SISTEMA MULTI-CONTAINER FLASK, NEO4J E JUPYTER
% ==============================================================================

\section{Caso Di Studio Applicativo: Sistema Multi-Container Flask, Neo4j E Jupyter}

\subsection{Requisiti Funzionali E Architettura Del Sistema Distribuito}

Come applicazione pratica delle metodologie di orchestrazione esaminate, si analizza un'architettura distribuita eterogenea basata su un modello client-server che integra l'analisi di dati grafici e la consultazione web:

% [IMMAGINE DEL FILE DI SLIDE CON NOME "CloudComputing_2"] Architettura di rete distribuita dell'applicazione di studio: Flask (porta 5000), Neo4j (porte 7474 e 7687 Bolt) e Jupyter Notebook (porta 8888) — Slide 23 / Pagina 50 del PDF

Il sistema comprende tre componenti distinti racchiusi in container indipendenti:
\begin{enumerate}
    \item \textbf{Container Web Frontend/Backend (Flask)}: espone l'interfaccia utente e gestisce le rotte HTTP per la consultazione e la ricerca di film e attori.
    \item \textbf{Container Database (Neo4j)}: gestisce l'archiviazione persistente delle entità e delle loro relazioni topologiche.
    \item \textbf{Container di Calcolo Interattivo (Jupyter Notebook)}: fornisce un ambiente di esecuzione scientifico impiegato per l'inserimento iniziale dei dati, l'analisi statistica e l'amministrazione del grafo mediante script Python.
\end{enumerate}

---

\subsection{Componenti Software E Protocolli Di Comunicazione}

L'architettura impiega componenti e protocolli eterogenei:
\begin{itemize}
    \item \textbf{Flask}: micro-framework per il linguaggio Python concepito per la realizzazione rapida di servizi web e API REST. La sua architettura è conforme alle specifiche dello standard **Web Server Gateway Interface (WSGI)**, il quale definisce una convenzione di chiamata universale tra i server web (es. Gunicorn, uWSGI) e le applicazioni scritte in linguaggio Python.
    \item \textbf{Neo4j}: motore per database orientato ai grafi (\emph{Graph Database}). A differenza dei database relazionali (basati su tabelle e chiavi esterne) o documentali, Neo4j memorizza le informazioni sotto forma di **nodi** (entità), **relazioni** (collegamenti orientati con valenza semantica) e **proprietà** (metadati chiave-valore associati a nodi e archi). L'interrogazione del grafo avviene mediante il linguaggio dichiarativo **Cypher**.
    \item \textbf{Protocollo Bolt}: è il protocollo applicativo a livello di trasporto proprietario ottimizzato per Neo4j. Progettato per essere fortemente tipizzato e binario, Bolt viene instradato su una normale connessione TCP (sulla porta standard \texttt{7687}) o via WebSocket, consentendo l'esecuzione remota ad alta velocità di query Cypher e la serializzazione binaria dei risultati verso i driver applicativi ufficiali Python.
    \item \textbf{Driver di Base Dati}: modulo software applicativo incaricato di implementare la conformità con uno standard di connessione verso il DBMS (quali ODBC, JDBC o il driver nativo Bolt).
    \item \textbf{Jupyter Notebook}: ambiente di calcolo interattivo basato su interfaccia web che combina codice eseguibile, formule matematiche, testo formattato in Markdown e grafici visuali in una struttura a celle sequenziali.
\end{itemize}

---

\subsection{Implementazione Del Backend E Frontend In Flask}

L'interfaccia grafica dell'applicazione adotta il paradigma di rendering server-side:

% [IMMAGINE DEL FILE DI SLIDE CON NOME "CloudComputing_2"] Albero dei file del progetto visualizzato nell'editor Visual Studio Code: la cartella templates e gli script di backend — Slide 24 / Pagina 51 del PDF

La cartella di sviluppo presenta la seguente organizzazione:
\begin{itemize}
    \item La directory \texttt{src/templates/} contiene i file sorgente HTML dinamici: \texttt{index.html} (pagina iniziale dell'interfaccia utente), \texttt{actors.html} (visualizzazione dell'elenco degli attori), \texttt{movies.html} (catalogo delle pellicole), \texttt{movies\_actor.html} (relazioni film-attori) e \texttt{actors\_movie.html}.
    \item La generazione dinamica delle pagine è affidata alla funzione \texttt{render\_template} del modulo \texttt{flask.templating}, che compila le pagine interpolando le strutture dati restituite dalle query con i modelli HTML basandosi sul motore di templating **Jinja2**.
    \item La raccolta dell'input utente fa largo uso del tag HTML \texttt{<input>}. La flessibilità del tag è governata dall'attributo fondamentale \texttt{type}, che supporta oltre venti definizioni semantiche standard:
    \begin{quote}
        \texttt{button}, \texttt{checkbox}, \texttt{color}, \texttt{date}, \texttt{datetime-local}, \texttt{email}, \texttt{file}, \texttt{hidden}, \texttt{image}, \texttt{month}, \texttt{number}, \texttt{password}, \texttt{radio}, \texttt{range}, \texttt{reset}, \texttt{search}, \texttt{submit}, \texttt{tel}, \texttt{text} (valore di default), \texttt{time}, \texttt{url}, \texttt{week}.
    \end{quote}
\end{itemize}

L'instradamento delle richieste (\textbf{App Routing}) associa i percorsi degli indirizzi URL alle rispettive funzioni Python mediante il decoratore di metodo \texttt{@app.route()}:
\begin{lstlisting}[language=python]
@app.route("/")
def index():
    return "Hello world"
\end{lstlisting}

\begin{theorem}[\textbf{Gestione dei Metodi HTTP in Flask}]
Per impostazione predefinita, se una funzione decorata con \texttt{@app.route()} non dichiara esplicitamente il parametro \texttt{methods}, essa accetta unicamente richieste di tipo \textbf{GET} e \textbf{HEAD}.
La ricezione di un metodo non dichiarato (come una sottomissione di form tramite \textbf{POST}) non invoca la funzione e produce automaticamente sul client la risposta di errore \textbf{HTTP 405 Method Not Allowed}.
\end{theorem}

Affinché la rotta possa elaborare i dati inviati dagli utenti tramite un form HTML, è indispensabile dichiarare esplicitamente i verbi autorizzati:
\begin{lstlisting}[language=python]
@app.route("/search", methods=['GET', 'POST'])
def search_handler():
    if request.method == 'POST':
        # Elaborazione logica della richiesta di ricerca nel grafo
        pass
    return render_template('index.html')
\end{lstlisting}
Il metodo \texttt{HEAD} viene agganciato automaticamente ogni volta che è consentito il metodo \texttt{GET}; in tal caso Flask esegue la funzione restituendo al client esclusivamente le intestazioni HTTP e scartando il corpo della risposta. Le richieste di pre-flight \texttt{OPTIONS} sono invece gestite in autonomia dal server per negoziare le intestazioni di sicurezza senza attivare il codice della rotta.

---

\subsection{Il Dockerfile Dell'Applicazione Web}

Il container preposto all'esecuzione dell'interfaccia utente e della logica server-side di Flask viene compilato a partire dal seguente \textbf{Dockerfile}:

\begin{lstlisting}[language=dockerfile]
FROM python:3.9
WORKDIR /GraphApp

# Copia del punto di ingresso server-side
COPY src/main.py .

# Copia dei template HTML per il motore Jinja2
COPY src/templates/index.html ./templates/index.html
COPY src/templates/actors.html ./templates/actors.html
COPY src/templates/movies.html ./templates/movies.html
COPY src/templates/movies_actor.html ./templates/movies_actor.html
COPY src/templates/actors_movie.html ./templates/actors_movie.html

# Risoluzione e installazione delle dipendenze di terze parti
COPY requirements.txt requirements.txt
RUN pip install -r requirements.txt
\end{lstlisting}

La dinamica di compilazione prevede:
\begin{enumerate}
    \item Il recupero dell'immagine genitore certificata \texttt{python:3.9} dai registri di Docker Hub.
    \item La creazione e impostazione della directory di lavoro \texttt{/GraphApp} nel container.
    \item L'iniezione differenziale tramite \texttt{COPY} dello script di controllo \texttt{main.py} e dell'albero dei template.
    \item L'importazione del file di tracciamento delle librerie (\texttt{requirements.txt}) contenente le dipendenze software (in particolare le versioni compatibili dei pacchetti \texttt{flask} e del driver \texttt{neo4j}).
    \item L'esecuzione del comando \texttt{pip install -r requirements.txt}, che consolida un layer statico contenente tutti i binari delle librerie esterne compilate.
\end{enumerate}

---

\subsection{Orchestrazione Con Docker Compose: Reti, Volumi E Dipendenze Di Avvio}

I tre servizi applicativi vengono coordinati dal descrittore dichiarativo \textbf{\texttt{docker-compose.yml}} basato sulle specifiche Compose v3.8:

\begin{lstlisting}[language=yaml]
version: '3.8'

services:
  app:
    build: .
    container_name: graphapp
    restart: always
    depends_on:
      - neo4j
    ports:
      - "5000:5000"
    command: python main.py
    networks:
      - graph_network

  neo4j:
    image: neo4j:4.3.7-enterprise
    container_name: neo4j
    restart: always
    ports:
      - "7474:7474"              # Porta per l'interfaccia grafica Web HTTP
      - "7687:7687"              # Porta di comunicazione per il protocollo binario Bolt
    volumes:
      - ./data:/data             # Montaggio di persistenza del grafo su directory host
    environment:
      - NEO4J_ACCEPT_LICENSE_AGREEMENT=yes
      - NEO4J_AUTH=none          # Disabilitazione autenticazione per finalita di test
    networks:
      - graph_network

  jupyter:
    image: jupyter/minimal-notebook:latest
    container_name: jupyter
    restart: always
    depends_on:
      - neo4j
    ports:
      - "8888:8888"              # Porta per la dashboard di sviluppo web
    networks:
      - graph_network

networks:
  graph_network:
    name: graph_network          # Rete bridge isolata creata ad hoc per l'applicazione
\end{lstlisting}

L'architettura configurata all'interno del file manifest evidenzia specifiche scelte di design:
\begin{itemize}
    \item \textbf{Segregazione di Rete}: tutti i container dichiarano l'appartenenza a una rete utente dedicata denominata \texttt{graph\_network}. Questo garantisce la totale protezione rispetto ad altri carichi di lavoro presenti sull'host e abilita la risoluzione DNS automatica tra i nodi (ad es. l'applicazione Flask e Jupyter raggiungono il database riferendosi direttamente all'endpoint con nome \texttt{bolt://neo4j:7687}).
    \item \textbf{Persistenza dei Dati}: il servizio \texttt{neo4j} dichiara un bind mount mappando la cartella locale dell'host \texttt{./data} sul percorso interno \texttt{/data} del container, assicurando che la struttura del grafo sopravviva a qualsiasi riavvio o distruzione del container.
    \item \textbf{Vincoli di Ordine Temporale (\texttt{depends\_on})}: la direttiva \texttt{depends\_on} formalizza una relazione d'ordine tassativa per il motore di orchestrazione:
    \begin{itemize}
        \item \emph{Fase di creazione e avvio}: i servizi vengono avviati secondo l'ordine di dipendenza. Nel manifest d'esempio, il motore garantisce che il servizio \texttt{neo4j} venga instanziato \textbf{prima} che i container \texttt{app} e \texttt{jupyter} vengano avviati.
        \item \emph{Fase di arresto (\texttt{docker compose down})}: il motore inverte specularmente l'ordine di terminazione, arrestando prima i servizi dipendenti (\texttt{app} e \texttt{jupyter}) e soltanto per ultimo il database di supporto (\texttt{neo4j}).
    \end{itemize}
\end{itemize}

---

\subsection{Operazioni Nel Jupyter Notebook: Ingestione Dati E Query Cypher}

L'alimentazione dei dati e la sperimentazione scientifica sul database avvengono sfruttando l'istanza Jupyter avviata sulla porta \texttt{8888}:

\begin{enumerate}
    \item All'atto dell'avvio con il comando:
\begin{lstlisting}[language=bash]
$ docker compose up -d
\end{lstlisting}
    l'operatore consulta i log di sistema generati dal container Jupyter per estrarre l'URL di autenticazione provvisto del token crittografico effimero generato a runtime:
\begin{lstlisting}[language=bash]
$ docker logs jupyter
# Output: http://127.0.0.1:8888/?token=3b8a1c9...
\end{lstlisting}
    \item All'interno dell'interfaccia web del notebook, vengono installate le dipendenze di connessione nell'ambiente Python del container:
\begin{lstlisting}[language=python]
import sys
!{sys.executable} -m pip install neo4j
!{sys.executable} -m pip install flask
\end{lstlisting}
    \item Viene mandato in esecuzione lo script \texttt{jupyter\_notebook.txt}, che si connette via protocollo Bolt all'endpoint \texttt{bolt://neo4j:7687} e inserisce l'insieme di nodi rappresentanti attori e pellicole cinematografiche, collegandoli con relazioni orientate (es. \texttt{(:Actor)-[:ACTED\_IN]->(:Movie)}).
\end{enumerate}

Le operazioni di interrogazione e pulizia dell'ambiente database vengono formulate mediante la sintassi dichiarativa del linguaggio **Cypher**:
\begin{itemize}
    \item \textbf{Visualizzazione dell'intero grafo}: la seguente clausola individua qualsiasi entità nodale $n$ appartenente allo spazio vettoriale del database e ne restituisce la collezione completa:
\begin{lstlisting}[language=sql]
MATCH (n) RETURN n;
\end{lstlisting}
    \item \textbf{Eliminazione atomica del database}: per cancellare l'intero contenuto del grafo senza violare i vincoli di integrità referenziale, è necessario scindere preventivamente tutti gli archi associati ai nodi. La clausola \texttt{DETACH DELETE} elimina in un'unica transazione atomica sia le relazioni incidenti che i nodi associati:
\begin{lstlisting}[language=sql]
MATCH (n) DETACH DELETE n;
\end{lstlisting}
\end{itemize}

% BB ---------------------------------------------------------------------------------------------------------------------------

% ==============================================================================
% CAPITOLO: ARCHITETTURA E COMPONENTI DI OPENSTACK
% ==============================================================================

\section{Introduzione Ad OpenStack E Modello Architetturale}

Nel paradigma del \emph{Cloud Computing}, l'infrastruttura come servizio (\textbf{Infrastructure as a Service}, IaaS) rappresenta il livello fondamentale in cui le risorse fisiche di calcolo, memorizzazione e connettività di rete vengono astratte, aggregate e messe a disposizione sotto forma di entità virtuali configurabili programmaticamente. \textbf{OpenStack} si posiziona in questo contesto come un vero e proprio sistema operativo per il cloud (\emph{Cloud Operating System}), progettato per orchestrare e gestire grandi pool di risorse eterogenee distribuite all'interno di uno o più centri di elaborazione dati (\emph{datacenter}).

\subsection{Principi Ispiratori E Astrazione Delle Risorse}

L'obiettivo cardine di OpenStack consiste nel superare i vincoli proprietari delle infrastrutture tradizionali mediante l'impiego di standard aperti e architetture modulari. Il sistema consente agli operatori e agli amministratori di sistema di mantenere un controllo capillare sulle risorse fisiche sottostanti (\emph{Standard Hardware}), abilitando al contempo un modello \emph{self-service} per gli utenti finali. Questi ultimi possono richiedere, allocare e deallocare istanze virtuali, volumi di memorizzazione e topologie di rete arbitrarie in totale autonomia, interagendo sia attraverso un'interfaccia grafica web (\textbf{Horizon Dashboard}) sia mediante \textbf{API RESTful} ed interfacce a riga di comando (\textbf{CLI}).

% [IMMAGINE DEL FILE DI SLIDE CON NOME "CloudComputing_3.pdf"] Diagramma architetturale a livelli di OpenStack: interfaccia Dashboard Horizon, servizi di Compute, Networking e Storage posizionati sopra i servizi condivisi e l'hardware standard — Slide 2 / Pagina 2 del PDF

Come schematizzato nel livello concettuale fondamentale, OpenStack si colloca tra l'hardware grezzo e le applicazioni utente:
\begin{itemize}[noitemsep]
    \item \textbf{Hardware Standard}: server con architetture convenzionali x86\_64 (CPU, RAM), dischi locali o collegati tramite Storage Area Network (SAN), e interfacce di rete fisiche (NIC).
    \item \textbf{OpenStack Shared Services}: servizi trasversali condivisi (autenticazione, catalogo delle risorse, message broker e database relazionali).
    \item \textbf{Core IaaS Services}: i motori operativi dedicati alla computazione (\emph{Compute}), alla rete virtuale (\emph{Networking}) e all'archiviazione (\emph{Storage}).
    \item \textbf{Interfacce Di Controllo E Accesso}: il cruscotto web unificato (\emph{Dashboard}) e gli strati di esportazione delle API pubbliche accessibili dall'esterno.
\end{itemize}

\subsection{Ecosistema Esteso: Strumenti Di Automazione E Servizi Terzi}

La natura aperta di OpenStack ne permette l'integrazione fluida con piattaforme e strumenti di gestione superiori, abilitando scenari che spaziano dalla gestione del metallo nudo (\emph{Bare Metal}) alla virtualizzazione tramite hypervisor, fino all'orchestrazione avanzata di \emph{container}.

% [IMMAGINE DEL FILE DI SLIDE CON NOME "CloudComputing_3.pdf"] Schema di integrazione di OpenStack con tool nativi (SDK, Horizon) e servizi terzi (Kubernetes, Cloud Foundry, Terraform) sopra risorse Bare Metal, VM e Container — Slide 3 / Pagina 3 del PDF

Gli operatori possono scegliere se impiegare gli strumenti integrati nativi forniti dal progetto, quali l'\textbf{OpenStack SDK} e l'interfaccia \textbf{Horizon}, oppure disaccoppiare la gestione operativa appoggiandosi a tool di automazione e piattaforme esterne:
\begin{itemize}
    \item \textbf{Terraform}: è uno strumento \emph{open-source} di \emph{Infrastructure as Code} (IaC) sviluppato da HashiCorp. Consente di definire, creare e aggiornare l'infrastruttura di rete, calcolo e memorizzazione tramite file di configurazione dichiarativi scritti in \textbf{HCL} (\emph{HashiCorp Configuration Language}). Invece di invocare manualmente le API di OpenStack, l'ingegnere specifica lo stato desiderato dell'infrastruttura (macchine virtuali, indirizzi IP fluttuanti, regole di sicurezza) e il motore di Terraform pianifica ed esegue le chiamate necessarie per convergere a tale topologia in maniera riproducibile e deterministica.
    \item \textbf{Cloud Foundry}: rappresenta una piattaforma aperta di tipo \emph{Platform as a Service} (PaaS). Astrae completamente l'infrastruttura computazionale sottostante agli sviluppatori software, automatizzando il ciclo di vita applicativo, il bilanciamento del carico, la scalabilità orizzontale automatica (\emph{auto-scaling}) e la registrazione dei microservizi, basandosi sull'incapsulamento del codice all'interno di container applicativi standardizzati.
    \item \textbf{Kubernetes}: orchestratore di carichi di lavoro containerizzati, che può utilizzare OpenStack sia come fornitore IaaS di macchine virtuali o nodi bare-metal su cui poggiare i nodi \emph{Control Plane} e \emph{Worker}, sia per richiedere dinamicamente volumi a blocchi e bilanciatori di carico esterni tramite i rispettivi plugin \emph{Cloud Provider Interface} (CPI) e \emph{Container Storage Interface} (CSI).
\end{itemize}

\subsection{Tipologie Di Moduli E Componenti Operativi}

Da un punto di vista sistemistico, i vari servizi che compongono OpenStack non costituiscono un monolite software, bensì una costellazione di micro-moduli distribuiti che rientrano in tre precise categorie esecutive:
\begin{enumerate}
    \item \textbf{Demoni (\emph{Daemons})}: processi eseguiti costantemente in \emph{background} sui server host. Nei sistemi operativi Linux, questi demoni vengono tipicamente gestiti come unità di servizio da \texttt{systemd} (ad esempio, \texttt{nova-api}, \texttt{neutron-server}, \texttt{glance-api}). Hanno il compito di rimanere in ascolto di richieste di rete o di consumare code di messaggi.
    \item \textbf{Script E Strumenti Ausiliari}: file eseguibili e automazioni utilizzati per inizializzare ambienti virtuali, gestire migrazioni di schemi di database relazionali (es. tramite \texttt{alembic}), verificare la conformità funzionale dell'infrastruttura ed eseguire test di integrazione continui.
    \item \textbf{Interfacce A Riga Di Comando (\emph{CLI})}: strumenti client (unificati sotto il comando \texttt{openstack}) che consentono a utenti e amministratori di formattare parametri, generare richieste HTTP strutturate e sottomettere chiamate verso gli endpoint delle API di ciascun servizio.
\end{enumerate}

Tutti i servizi OpenStack autenticano obbligatoriamente ogni singola transazione per mezzo di un componente centrale comune di gestione dell'identità (\textbf{Keystone}). Le comunicazioni tra i diversi servizi avvengono tramite API pubbliche RESTful, fatto salvo per le procedure amministrative interne che richiedono privilegi specifici di sistema.

% [IMMAGINE DEL FILE DI SLIDE CON NOME "CloudComputing_3.pdf"] Architettura logica dettagliata di OpenStack con interconnessioni tra componenti, bus AMQP, DB condiviso e API gateway — Slide 4 / Pagina 4 del PDF

\subsection{Architettura Concettuale Delle Interdipendenze}

Il disaccoppiamento architetturale impone che ogni progetto gestisca il proprio dominio di competenza e cooperi con gli altri mediante contratti di interfaccia ben definiti. Internamente, ciascun servizio OpenStack si compone di processi distinti:
\begin{itemize}[noitemsep]
    \item Almeno un processo \textbf{API Process}, che riceve le chiamate in ingresso (over HTTP/HTTPS), valida la sintassi della richiesta, verifica l'autenticità del token crittografico rilasciato da Keystone e la traduce in un comando interno.
    \item Uno o più processi di elaborazione e coordinamento (\textbf{Worker Process}, \textbf{Schedulers}, \textbf{Conductors}), incaricati di eseguire l'effettiva logica di business, interagire con le code asincrone e impartire istruzioni ai driver dell'hypervisor o dei dispositivi di rete fisici.
\end{itemize}

% [IMMAGINE DEL FILE DI SLIDE CON NOME "CloudComputing_3.pdf"] Diagramma concettuale delle relazioni tra servizi (Keystone, Horizon, Nova, Neutron, Cinder, Swift, Glance, Trove, Ironic, Sahara, Heat) — Slide 5 / Pagina 5 del PDF

Tra i moduli ausiliari integrati si distinguono:
\begin{itemize}
    \item \textbf{Trove}: servizio di \emph{Database as a Service} (DBaaS), ideato per orchestrare istanze di database relazionali (es. MySQL, PostgreSQL) o non relazionali (NoSQL come MongoDB, Cassandra), sollevando l'utente dalla complessità di backup, ripristino, configurazione e manutenzione dei cluster dati.
    \item \textbf{Ironic}: modulo specializzato nel \emph{bare-metal provisioning}. Nato originariamente come driver sperimentale scorporato da Nova (\emph{nova-baremetal-driver}), Ironic consente di distribuire immagini del sistema operativo direttamente su server fisici privi di hypervisor intermedio, sfruttando protocolli di accensione e controllo fuori banda (\textbf{IPMI}, Redfish) e avvii di rete tramite \textbf{PXE} (\emph{Preboot Execution Environment}).
    \item \textbf{Sahara}: piattaforma progettata per fornire agli utenti un meccanismo dichiarativo e automatizzato per il dispiegamento e la configurazione rapida di cluster dedicati al trattamento di grandi volumi di dati (\emph{Big Data Processing Frameworks}), quali Apache Hadoop, Apache Spark e Apache Storm.
\end{itemize}

\newpage

% ==============================================================================
% SEZIONE 2: COMUNICAZIONE INTER-SERVIZIO E MESSAGE ORIENTED MIDDLEWARE
% ==============================================================================

\section{Comunicazione Inter-Servizio E Message-Oriented Middleware (MOM)}

Un'infrastruttura cloud come OpenStack costituisce a tutti gli effetti un \textbf{sistema distribuito su larga scala}. I processi che la compongono possono risiedere su nodi distinti all'interno della stessa \emph{server farm} o persino su datacenter dislocati in diverse aree geografiche. Tali processi devono cooperare per presentarsi all'esterno come un'applicazione monolitica coerente e resiliente, garantendo un'elevata tolleranza ai guasti transitori di rete o di elaborazione.

\subsection{Limiti Delle Architetture Tradizionali: Database E REST Sincrono}

Nella progettazione dei canali di interconnessione tra componenti interni ad alta frequenza di scambio, la condivisione dello stato tramite un database centralizzato o per mezzo di chiamate sincrone HTTP/REST presenta criticità strutturali non trascurabili:
\begin{itemize}
    \item \textbf{Condivisione Tramite Database Relazionale}: sebbene una base di dati conforme ai criteri ACID garantisca consistenza e persistenza durevole, essa è ottimizzata per memorizzare dati strutturati e recuperare informazioni di stato, non per fare da tramite attivo al flusso di messaggi (\emph{inter-process communication}). L'adozione del database come intermediario costringe a ricorrere a costosi meccanismi di polling ciclico, determinando contese sui blocchi di lettura/scrittura (\emph{locking overhead}) e colli di bottiglia prestazionali difficilmente scalabili.
    \item \textbf{Chiamate Sincrone HTTP/REST}: il paradigma REST basato su HTTP assicura un'eccellente interoperabilità verso l'esterno ed è leggero, ma impone un accoppiamento temporale stretto di tipo sincrono (richiesta-risposta bloccante). Se il destinatario di una chiamata è momentaneamente irraggiungibile, congestionato o in fase di riavvio, il tentativo di connessione fallisce o va in \emph{timeout}, provocando la perdita irreversibile della trasmissione a meno di non implementare complesse logiche applicative di riprova (\emph{retry}).
\end{itemize}

\subsection{Il Modello Store-And-Forward E I Pattern A Code}

Per ovviare a tali problemi, OpenStack adotta un'architettura lassamente accoppiata basata su un \textbf{Message-Oriented Middleware} (MOM), impiegando un sistema di messaggistica asincrono con paradigma \emph{Store-and-Forward}.

\begin{theorem}[Principio Dello Store-And-Forward]
Nel modello \emph{Store-and-Forward}, il mittente non recapita il messaggio direttamente al processo destinatario finale, bensì lo affida a un intermediario software specializzato, denominato \textbf{Message Broker}. Il broker assume la custodia formale del messaggio, salvandolo temporaneamente in una coda di memoria o su disco non volatile (\emph{store}) e provvedendo a inoltrarlo (\emph{forward}) non appena il processo consumatore risulta disponibile, registrato e pronto alla lavorazione.
\end{theorem}

Questo meccanismo disaccoppia i processi secondo tre assi fondamentali:
\begin{itemize}[noitemsep]
    \item \textbf{Disaccoppiamento Spaziale}: il produttore non ha necessità di conoscere la locazione fisica (indirizzo IP e porta) del consumatore.
    \item \textbf{Disaccoppiamento Temporale}: produttore e consumatore non devono essere attivi contemporaneamente durante la transazione.
    \item \textbf{Disaccoppiamento Delle Prestazioni}: un picco improvviso di richieste viene assorbito dalla capacità di bufferizzazione della coda, evitando il collasso per saturazione dei nodi di calcolo.
\end{itemize}

I pattern architetturali implementati sono prevalentemente due: il modello punto-a-punto (\textbf{Queue-based Point-to-Point}), utilizzato per il bilanciamento del carico tra worker omogenei (\emph{competing consumers}), e il modello pubblicazione-sottoscrizione (\textbf{Publisher-Subscriber}), in cui un evento o una notifica di stato viene distribuita in \emph{broadcast} selettivo a molteplici agenti interessati.

\subsection{Il Protocollo AMQP E L'Integrazione Con RabbitMQ}

Lo standard aperto adottato da OpenStack a livello applicativo è l'\textbf{Advanced Message Queuing Protocol} (\textbf{AMQP}).

AMQP è un protocollo binario operante al livello 7 dello stack OSI, progettato specificamente per ambienti enterprise che esigono garanzie formali sull'affidabilità e l'integrità dei flussi informativi. Esso presuppone un protocollo di trasporto affidabile sottostante, identificato nello standard \textbf{TCP} (\emph{Transmission Control Protocol}), e integra meccanismi nativi di autenticazione e cifratura fondati su standard industriali quali \textbf{SASL} (\emph{Simple Authentication and Security Layer}) e \textbf{TLS} (\emph{Transport Layer Security}).

La specifica formale di AMQP si articola in quattro layer cooperanti:
\begin{enumerate}
    \item \textbf{Type System}: un sistema formale di codifica unificata dei tipi di dato primitivi e composti, indipendente dal linguaggio di programmazione dei singoli nodi.
    \item \textbf{Protocollo Asincrono Simmetrico}: un insieme di convenzioni di framing per il trasferimento asincrono bidirezionale di frame di controllo e di carico utile tra due processi (\emph{peer-to-peer}).
    \item \textbf{Formato Standard Del Messaggio}: una struttura dati uniforme ed estendibile che separa le intestazioni di instradamento (\emph{headers}), le proprietà di consegna e il corpo del messaggio (\emph{payload}).
    \item \textbf{Funzionalità Di Messaggistica Standardizzate}: le entità logiche che governano l'elaborazione del traffico: \emph{Exchanges} (che ricevono i messaggi e ne valutano i criteri di routing), \emph{Queues} (strutture FIFO che accumulano i pacchetti) e \emph{Bindings} (le regole formali di associazione tra exchange e code).
\end{enumerate}

\subsubsection{Garanzie Di Consegna Del Protocollo}

Il protocollo consente di formalizzare rigorosi contratti di affidabilità, parametrizzabili in base alla criticità dell'operazione:
\begin{itemize}
    \item \textbf{At-most-once}: il messaggio viene recapitato al più una volta, o non recapitato affatto. Nessun riscontro (\emph{acknowledgment}) viene atteso; le perdite dovute a crash sono tollerate a fronte della massima rapidità trasmissiva.
    \item \textbf{At-least-once}: il messaggio raggiunge con certezza la destinazione, ma possono verificarsi duplicazioni qualora l'ACK di conferma del consumatore vada smarrito sulla rete.
    \item \textbf{Exactly-once}: semantica ideale in cui la consegna del messaggio è garantita esattamente una singola volta, resa possibile dalla combinazione di canali idempotenti, transazioni distribuite o filtri di deduplicazione a livello broker.
\end{itemize}

All'interno dell'architettura di riferimento OpenStack, l'implementazione predefinita del message broker è affidata a \textbf{RabbitMQ}, il quale risiede sul nodo di controllo (\textbf{Controller Node}).

Ad esempio, all'interno del sottosistema di calcolo (\textbf{Nova}), componenti come \texttt{nova-api}, \texttt{nova-scheduler} e \texttt{nova-conductor} non si scambiano mai messaggi di controllo diretti su socket TCP privati, ma pubblicano e prelevano messaggi attraverso code gestite da RabbitMQ. Parallelamente, lo stato persistente dell'infrastruttura (metadati dei server, mappature degli indirizzi, allocazioni di quote) viene salvato su una base di dati relazionale ACID, tipicamente implementata con cluster \textbf{MariaDB} o \textbf{MySQL} (oppure SQLite per installazioni minime di testing).

\newpage

% ==============================================================================
% SEZIONE 3: TASSONOMIA DEI SERVIZI OPENSTACK
% ==============================================================================

\section{Tassonomia Dei Servizi Fondamentali E Ausiliari}

L'ecosistema OpenStack adotta una suddivisione rigorosa delle responsabilità architetturali, assegnando a ciascun macro-ambito un progetto specifico contraddistinto da un nome convenzionale.

\subsection{Servizi Di Infrastruttura Primari (Core Services)}

\begin{itemize}
    \item \textbf{Dashboard (Horizon)}: interfaccia utente su protocollo web (HTTPS). Funge da punto di accesso unificato grafico (\emph{self-service portal}) per interagire con i vari sottosistemi OpenStack. Consente di avviare macchine virtuali, associare indirizzi IP pubblici/fluttuanti, definire router virtuali e configurare le regole di isolamento e controllo degli accessi.
    \item \textbf{Compute (Nova)}: il nucleo computazionale dell'intera piattaforma. Ha la responsabilità esclusiva di orchestrare il ciclo di vita completo delle istanze di calcolo (macchine virtuali o server fisici). Nova si occupa di:
    \begin{itemize}[noitemsep]
        \item Ricevere richieste di creazione di istanze tramite la propria API.
        \item Selezionare il nodo host ottimale su cui distribuire il carico (\emph{scheduling}).
        \item Interfacciarsi con l'hypervisor locale per avviare, arrestare, mettere in pausa, ridimensionare o deallocare le risorse.
    \end{itemize}
    \item \textbf{Networking (Neutron)}: realizza il paradigma di connettività di rete distribuita come servizio (\textbf{Network-Connectivity-as-a-Service}). Fornisce una completa astrazione logica su switch, router, subnet, indirizzi IP e firewall. Presenta un'architettura altamente modulare basata su plugin e driver, consentendo l'integrazione di SDN avanzati, apparati hardware multi-vendor e tecnologie di virtualizzazione di rete tradizionali.
\end{itemize}

\subsection{Servizi Di Supporto E Condivisi (Shared Services)}

I servizi condivisi garantiscono il supporto orizzontale a tutti i moduli del cloud:
\begin{itemize}
    \item \textbf{Identity Service (Keystone)}: fornisce l'infrastruttura centralizzata di autenticazione (\emph{AuthN}) e autorizzazione (\emph{AuthZ}). Conserva il catalogo globale degli endpoint (\textbf{Service Catalog}), fornendo agli utenti e ai servizi gli indirizzi URL univoci a cui inoltrare le chiamate API.
    \item \textbf{Image Service (Glance)}: agisce da registro e archivio centrale per le immagini disco dei sistemi operativi delle macchine virtuali (immagini ISO, QCOW2, RAW). Quando Nova avvia un'istanza, invoca Glance per reperire l'immagine di partenza e copiarla sul disco di avvio locale o su un volume storage.
    \item \textbf{Telemetry (Ceilometer)}: monitora, cattura e raccoglie metriche d'uso relative a CPU, memoria, storage e traffico di rete generato da ciascun tenant o singola istanza. I dati archiviati da Ceilometer sono impiegati per finalità di fatturazione analitica (\emph{billing}), auditing, benchmarking prestazionale, e come segnale di attivazione per le politiche di scalabilità orizzontale.
\end{itemize}

\subsection{Servizi Di Memorizzazione (Storage Services)}

OpenStack disaccoppia la memorizzazione in due paradigmi architetturali profondamente eterogenei:

\begin{itemize}
    \item \textbf{Object Storage (Swift)}: sistema di archiviazione per dati non strutturati (file multimediali, backup, archivi statici) accessibile esclusivamente tramite API HTTP RESTful. A differenza di un comune file system, non supporta il montaggio di directory logiche né la gestione di puntatori POSIX. I file vengono trattati come oggetti immutabili corredati di metadati. L'architettura è nativamente distribuita e scalabile orizzontalmente (\emph{scale-out}): per massimizzare la tolleranza ai guasti, ogni oggetto viene duplicato su molteplici dischi e server differenti dislocati nel cluster, garantendo consistenza e durabilità anche a fronte di rotture hardware simultanee.
    \item \textbf{Block Storage (Cinder)}: fornisce dispositivi di memorizzazione persistente a blocchi (\emph{persistent block storage}) alle istanze virtuali in esecuzione. Un volume Cinder può essere formattato con qualsiasi file system nativo ed essere montato come disco rigido secondario (o volume di boot primario) da parte della macchina virtuale. Grazie a un'architettura flessibile a driver aperti, Cinder astrae controller di storage aziendali eterogenei (dischi locali LVM, Ceph, iSCSI, Fibre Channel, storage proprietari SAN/NAS).
\end{itemize}

\subsection{Servizi Di Orchestrazione Superiore}

\begin{itemize}
    \item \textbf{Orchestration (Heat)}: motore dichiarativo di orchestrazione di applicazioni cloud complesse. Heat consente di formalizzare l'intera composizione di un'architettura (server, reti, gruppi di sicurezza, volumi di archiviazione e regole di auto-scaling) all'interno di modelli testuali leggibili. Supporta il formato nativo \textbf{HOT} (\emph{Heat Orchestration Template}), strutturato in sintassi YAML, e garantisce piena compatibilità retroattiva con i template di \textbf{AWS CloudFormation}, esponendo sia un'API REST nativa OpenStack sia una Query API compatibile con CloudFormation.
\end{itemize}

\newpage

% ==============================================================================
% SEZIONE 4: ARCHITETTURA DI RIFERIMENTO DEI NODI E LAYOUT DI RETE
% ==============================================================================

\section{Architettura Di Riferimento Dei Nodi E Layout Di Rete}

Un'installazione reale di OpenStack distribuisce i servizi sopra una serie di nodi fisici specializzati, delineando un'architettura modulare in cui ciascuna classe di host adempie a una funzione computazionale o di controllo ben precisa.

% [IMMAGINE DEL FILE DI SLIDE CON NOME "CloudComputing_3.pdf"] Requisiti hardware per la configurazione di riferimento: Controller Node, Compute Node 1, Block Storage Node 1, Object Storage Node 1 e 2 — Slide 14 / Pagina 14 del PDF

\subsection{Tipologie Di Nodi E Requisiti Hardware}

La configurazione di riferimento comunemente raccomandata prevede la separazione in quattro tipologie di server fisici:

\begin{enumerate}
    \item \textbf{Controller Node (Nodo Di Controllo)}:
    Costituisce il punto nevralgico di gestione del cloud. Ospita i servizi di identità (Keystone), il catalogo immagini (Glance), le porzioni di coordinamento e gestione delle API di Compute e Networking (\texttt{nova-api}, \texttt{neutron-server}), i rispettivi agenti di monitoraggio, l'interfaccia Horizon, oltre ai servizi ausiliari critici: il database relazionale (SQL), il broker di messaggistica (RabbitMQ) e il demone di sincronizzazione temporale \textbf{NTP} (\emph{Network Time Protocol}), fondamentale per la coerenza dei certificati e dei log distribuiti.
    \begin{itemize}[noitemsep]
        \item \emph{Requisiti hardware minimi}: 1--2 CPU, \SI{8}{\giga\byte} di RAM, \SI{100}{\giga\byte} di storage su disco.
        \item \emph{Connettività}: richiede obbligatoriamente un minimo di **due interfacce di rete fisiche (NIC)**.
    \end{itemize}

    \item \textbf{Compute Node (Nodo Di Calcolo)}:
    È la macchina fisica deputata all'esecuzione effettiva delle istanze virtuali utente. Esegue la componente \emph{hypervisor} di Compute; nell'ecosistema OpenStack standard, l'ipervisore di riferimento predefinito è \textbf{KVM} (\emph{Kernel-based Virtual Machine}). Il nodo ospita inoltre l'agente locale di Neutron (\texttt{neutron-agent}) incaricato di collegare le interfacce delle VM ai virtual switch locali (es. Open vSwitch o Linux Bridge) e di applicare i meccanismi di filtraggio dei pacchetti tramite i \textbf{Security Groups}.
    \begin{itemize}[noitemsep]
        \item \emph{Requisiti hardware minimi}: 2--4+ core CPU (con supporto alla virtualizzazione hardware VT-x/AMD-V abilitata), \SI{8}{\giga\byte}+ di RAM (scalabile a seconda del numero di istanze ospitate), \SI{100}{\giga\byte}+ di archiviazione locale.
        \item \emph{Connettività}: richiede un minimo di **due interfacce di rete (NIC)**. L'infrastruttura è orizzontalmente scalabile mediante l'aggiunta di $N$ nodi di calcolo.
    \end{itemize}

    \item \textbf{Block Storage Node (Nodo Di Memorizzazione A Blocchi)}:
    Nodo opzionale deputato a ospitare i volumi persistenti allocati da Cinder e i servizi di file system condiviso (\emph{Manila}). Contiene dischi fisici addizionali o partizioni dedicate (es. \texttt{/dev/sdb}, \texttt{/dev/sdc}) che vengono aggregate in volumi logici tramite LVM o storage dedicati.
    \begin{itemize}[noitemsep]
        \item \emph{Requisiti hardware minimi}: 1--2 CPU, \SI{4}{\giga\byte} di RAM, \SI{100}{\giga\byte}+ di spazio disco per i volumi.
        \item \emph{Connettività}: richiede almeno **una interfaccia di rete (NIC)**. Nelle installazioni di produzione è fortemente raccomandato l'uso di una rete dedicata isolata per evitare interferenze con il traffico di gestione.
    \end{itemize}

    \item \textbf{Object Storage Nodes (Nodi Di Memorizzazione Ad Oggetti)}:
    Server fisici che contengono i dischi dedicati alla persistenza di account, container e oggetti gestiti da Swift. Si utilizzano solitamente cluster composti da nodi multipli (Node 1, Node 2, \dots) per soddisfare il fattore di replica dei blocchi.
    \begin{itemize}[noitemsep]
        \item \emph{Requisiti hardware minimi}: 1--2 CPU, \SI{4}{\giga\byte}+ di RAM, \SI{100}{\giga\byte}+ di storage grezzo (\texttt{/dev/sdb}, \texttt{/dev/sdc}).
        \item \emph{Connettività}: richiede almeno **una interfaccia di rete fisica (NIC)**.
    \end{itemize}
\end{enumerate}

\subsection{Topologia E Segmentazione Delle Reti Fisiche E Virtuali}

La corretta segregazione del traffico è un requisito primario per garantire sicurezza, determinismo prestazionale e isolamento tra tenant.

% [IMMAGINE DEL FILE DI SLIDE CON NOME "CloudComputing_3.pdf"] Layout topologico di rete: indirizzamento IP su Management Network (10.0.0.0/24), Provider Network (203.0.113.0/24), NAT e Internet — Slide 18 / Pagina 18 del PDF

L'architettura di rete si articola su domini logici distinti:
\begin{itemize}
    \item \textbf{Management Network} (\texttt{10.0.0.0/24}): una rete privata interna impiegata esclusivamente per la comunicazione amministrativa tra i nodi OpenStack. Su questo canale viaggiano le query tra le API e il database MariaDB, i flussi AMQP del broker RabbitMQ e le comunicazioni di stato tra i demoni.
    Ad esempio:
    \begin{itemize}[noitemsep]
        \item Controller Node: \texttt{10.0.0.11/24} (su \texttt{Interface 1})
        \item Compute Node 1: \texttt{10.0.0.31/24} (su \texttt{Interface 1})
        \item Block Storage Node 1: \texttt{10.0.0.41/24} (su \texttt{Interface 1})
        \item Object Storage Node 1: \texttt{10.0.0.51/24} (su \texttt{Interface 1})
        \item Object Storage Node 2: \texttt{10.0.0.52/24} (su \texttt{Interface 1})
    \end{itemize}
    \item \textbf{Provider Network} (\texttt{203.0.113.0/24}): piano di rete che connette l'infrastruttura virtuale alla rete fisica esterna (Internet o Intranet aziendale). Sui nodi Compute e Controller, la seconda interfaccia di rete (\texttt{Interface 2}) è direttamente collegata a questo segmento, permettendo l'esposizione degli indirizzi IP esterni e l'attraversamento di gateway con \textbf{NAT} (\emph{Network Address Translation}).
\end{itemize}

\subsection{Alternative Di Virtualizzazione Con Neutron: Provider Vs. Self-Service}

Nell'implementazione di Neutron si distinguono due modelli operativi:

\begin{figure}[H]
\centering
\begin{tabularx}{\textwidth}{p{4.2cm}X}
\toprule
\textbf{Opzione Di Rete} & \textbf{Caratteristiche E Dettagli Implementativi} \\
\midrule
\textbf{Opzione 1: Provider Networks} &
Rappresenta l'approccio più semplice ed essenziale. Il servizio di rete opera prevalentemente al \textbf{Layer 2} dello stack ISO/OSI (bridging e switching locale). Le reti virtuali create dalle VM vengono collegate direttamente alle interfacce di rete fisiche mediante \textbf{VLAN} (standard IEEE 802.1Q).
\vspace{0.15cm}
Il routing al \textbf{Layer 3} non viene effettuato da OpenStack, ma è demandato interamente all'infrastruttura dei router fisici aziendali. Neutron si limita a distribuire la configurazione IP alle macchine virtuali mediante un server interno di \textbf{DHCP}. \\
\addlinespace
\textbf{Opzione 2: Self-Service Networks} &
Estende l'infrastruttura introducendo il routing di \textbf{Layer 3} virtualizzato e distribuito. Consente ai singoli tenant di configurare reti private personalizzate e indipendenti, disaccoppiate dalla topologia fisica mediante protocolli di \textbf{incapsulamento overlay}, in particolare \textbf{VXLAN} (\emph{Virtual Extensible LAN}).
\vspace{0.15cm}
L'inoltro dei pacchetti verso le reti fisiche e verso Internet avviene mediante router software interni gestiti da Neutron con tecniche di \textbf{SNAT} (\emph{Source NAT}) per il traffico in uscita e \textbf{Floating IP} (\emph{DNAT}) per il traffico in ingresso. Questa opzione è indispensabile per erogare servizi a valore aggiunto quali il \textbf{Load Balancer as a Service} (\textbf{LBaaS}) e il \textbf{Firewall as a Service} (\textbf{FWaaS}). \\
\bottomrule
\end{tabularx}
\caption{Confronto tra i modelli di virtual networking in OpenStack Neutron.}
\end{figure}

\newpage

% ==============================================================================
% SEZIONE 5: GESTIONE DELLE ISTANZE E OVERCOMMIT DELLE RISORSE
% ==============================================================================

\section{Gestione Operativa Delle Istanze E Overcommit Delle Risorse}

L'utilizzo pratico di OpenStack ruota attorno alla gestione delle entità operative fondamentali: progetti, utenti, immagini di avvio e profili dimensionali di calcolo.

\subsection{Identità Di Base E Concetto Di Istanza}

Un'installazione di base definisce convenzionalmente due tenant (progetti) e account predefiniti:
\begin{itemize}[noitemsep]
    \item \textbf{admin}: utente privilegiato, titolare dei diritti amministrativi globali sull'intero cloud.
    \item \textbf{demo}: utente privo di privilegi infrastrutturali, impiegato per collaudi funzionali all'interno del progetto omonimo.
\end{itemize}

Nel lessico di OpenStack, una macchina virtuale in esecuzione viene formalmente definita **istanza** (\emph{instance}). L'avvio di un'istanza richiede la fornitura di una specifica **immagine** (\emph{image}), ovvero una copia statica contenente il disco di un sistema operativo preconfigurato. OpenStack garantisce la compatibilità con una vasta gamma di distribuzioni Linux (\emph{Ubuntu, RedHat, Fedora, SUSE, Debian}) e con sistemi operativi Microsoft Windows.

Per le attività di collaudo rapido e sviluppo si impiega frequentemente l'immagine minimale **CirrOS** (es. \texttt{cirros-0.3.2-x86\_64-uec}): una distribuzione Linux leggera, appositamente ottimizzata per richiedere risorse minime (pochi megabyte di RAM) e convalidare l'integrità del percorso di rete e dello storage.

\subsection{Gestione Delle Coppie Di Chiavi Crittografiche (Keypair)}

L'accesso sicuro alle istanze Linux appena istanziate non avviene tramite credenziali statiche memorizzate nell'immagine, bensì tramite crittografia asimmetrica:
\begin{itemize}
    \item L'utente genera o importa in OpenStack una **coppia di chiavi** (\textbf{Keypair}), costituita da una chiave pubblica e da una chiave privata (generabile tramite la dashboard Horizon o mediante l'utility Unix \texttt{ssh-keygen}).
    \item La **chiave pubblica** viene registrata su OpenStack e, in fase di bootstrap dell'istanza, iniettata all'interno del file \texttt{\textasciitilde/.ssh/authorized\_keys} del sistema operativo guest tramite il meccanismo di metadata service (\textbf{cloud-init}).
    \item La **chiave privata** deve essere scaricata e conservata gelosamente dall'utente; essa funge da credenziale crittografica per l'accesso remoto via shell sicura (\texttt{ssh -i chiave.pem utente@ip}).
\end{itemize}

\subsection{Definizione Dei Profili Di Risorsa (Flavors)}

La quantità di hardware virtuale assegnata a un'istanza è vincolata a un profilo rigido denominato **Flavor**. Un flavor quantifica e definisce:
\begin{enumerate}[noitemsep]
    \item Il numero di processori virtuali (**vCPU**).
    \item La memoria volatile riservata (**RAM** in MB o GB).
    \item La dimensione del disco di root principale (\textbf{Root Disk}).
    \item Eventuali dischi ausiliari effimeri (\textbf{Ephemeral Disk}) e aree di swap.
\end{enumerate}

Per l'immagine CirrOS, ad esempio, si associa tipicamente il profilo minimale \texttt{m1.nano} (generalmente caratterizzato da 1 vCPU, \SI{64}{\mega\byte} o \SI{128}{\mega\byte} di RAM e \SI{1}{\giga\byte} di storage).

% [IMMAGINE DEL FILE DI SLIDE CON NOME "CloudComputing_3.pdf"] Schermata della dashboard Horizon con la tabella riassuntiva delle istanze attive (Status, Task, Power State, Time since created) — Slide 25 / Pagina 25 del PDF

\subsection{Il Principio Di Overcommit Di CPU E Memoria}

I carichi di lavoro all'interno di un'infrastruttura cloud presentano tipicamente una natura stocastica: raramente tutte le macchine virtuali attive impegnano contemporaneamente il \SI{100}{\percent} dei cicli di clock della CPU o la totalità della RAM allocata. Per massimizzare il ritorno economico e l'efficienza di utilizzazione dei server fisici, OpenStack consente la sovra-allocazione controllata delle risorse fisiche, operazione nota come \textbf{Overcommitting}.

\begin{theorem}[Rapporto Di Allocazione (Overcommit Ratio)]
Dato un nodo di calcolo con $C_{\text{fisiche}}$ core fisici di calcolo e una quantità $M_{\text{fisica}}$ di memoria RAM installata, le risorse virtuali complessivamente allocabili dallo scheduler Nova sono determinate da:
\begin{align}
    \sum vCPU &\le C_{\text{fisiche}} \cdot R_{\text{CPU}} \\
    \sum RAM_{\text{allocata}} &\le M_{\text{fisica}} \cdot R_{\text{RAM}}
\end{align}
dove $R_{\text{CPU}}$ e $R_{\text{RAM}}$ rappresentano i rispettivi tassi di overcommit definiti nella configurazione di Nova.
\end{theorem}

I valori predefiniti utilizzati dal servizio Compute in OpenStack sono:
\begin{itemize}
    \item \textbf{CPU Allocation Ratio}: $\mathbf{16:1}$. Significa che su un host fisico dotato di 16 core fisici (o thread hardware con Hyper-Threading abilitato), lo scheduler accetta di allocare fino a:
    \[
        16 \times 16 = 256 \text{ vCPU}
    \]
    \item \textbf{RAM Allocation Ratio}: $\mathbf{1.5:1}$. Su un server equipaggiato con \SI{64}{\giga\byte} di memoria fisica, Nova consentirà la creazione di istanze per un monte complessivo nominale di:
    \[
        64 \times 1.5 = \SI{96}{\giga\byte} \text{ di RAM virtuale}
    \]
\end{itemize}

\subsubsection{Analisi Ingegneristica Dei Trade-Off}

L'adozione dell'overcommit introduce un chiaro compromesso progettuale:
\begin{itemize}
    \item \textbf{Vantaggi}: drastico incremento della densità di macchine virtuali ospitate per singolo nodo server (\emph{instance packing}), con conseguente abbattimento dei costi di capitale (\textbf{CapEx}) e operativi (\textbf{OpEx}) legati a consumo energetico e raffreddamento.
    \item \textbf{Svantaggi e Rischi}: qualora molteplici istanze concentrino picchi computazionali nello stesso istante, si genera contesa sulle cache del processore e sulla larghezza di banda del bus di memoria. Tale fenomeno, noto come problema del \emph{Noisy Neighbor}, causa una sensibile degradazione delle prestazioni e una variabilità della latenza di risposta. Nel caso della RAM, l'esaurimento della memoria fisica reale costringe il kernel Linux dell'host ad attivare lo swap su disco o a invocare l'\textbf{OOM Killer} (\emph{Out-Of-Memory Killer}), con il rischio di terminare bruscamente processi critici dell'ipervisore.
\end{itemize}

\newpage

% ==============================================================================
% SEZIONE 6: PARTIZIONAMENTO E GERARCHIA DI ISOLAMENTO
% ==============================================================================

\section{Partizionamento E Gerarchia Di Isolamento Nel Cloud}

La scalabilità e la resilienza di un'infrastruttura cloud dipendono dalla capacità di segmentare le risorse in domini operativi e domini di guasto indipendenti. OpenStack formalizza una suddivisione strutturata su tre livelli gerarchici: **Regions**, **Availability Zones (AZ)** e **Host Aggregates**.

% [IMMAGINE DEL FILE DI SLIDE CON NOME "CloudComputing_3.pdf"] Schema gerarchico del partizionamento: Region 1 e Region 2, suddivise in Availability Zone (AZ 1, AZ 2) e mappatura trasversale con Host Aggregates (SSD, High CPU, NIC 10GB) — Slide 29 / Pagina 29 del PDF

\subsection{Regioni (Regions)}

Una **Region** rappresenta un'installazione autonoma e completa di OpenStack, comprensiva dei propri endpoint API dedicati, del proprio database, del broker di messaggistica e del proprio pool indipendente di nodi di rete e calcolo.

Le caratteristiche essenziali di una Region sono:
\begin{itemize}
    \item \textbf{Condivisione Di Keystone Ed Horizon}: pur possedendo infrastrutture computazionali fisicamente separate (spesso collocate in siti geografici o città distinte, ad esempio \texttt{us-east} ed \texttt{eu-west}), le diverse Regioni condividono un'unica istanza centralizzata del servizio di identità (\textbf{Keystone}) e dell'interfaccia grafica (\textbf{Horizon}).
    \item L'Identity API v3 mappa le Region come entità strutturate ad albero, consentendo di istituire sotto-regioni gerarchiche. La prima regione creata assume di default la denominazione standard \texttt{RegionOne}.
    \item Il traffico di rete tra regioni diverse transita necessariamente su reti pubbliche WAN e non beneficia dell'astrazione diretta di Layer 2 di Neutron.
\end{itemize}

\subsection{Zone Di Disponibilità (Availability Zones)}

Un'\textbf{Availability Zone} (\textbf{AZ}) costituisce un raggruppamento logico di nodi di calcolo (\emph{compute nodes}) situati all'interno della medesima Region, isolati tra loro in termini di \textbf{domini di guasto} (\emph{fault domains}).

\begin{itemize}
    \item Tipicamente, zone di disponibilità differenti insistono su circuiti di alimentazione elettrica indipendenti, gruppi di continuità (UPS) distinti, generatori ausiliari dedicati, armadi rack separati e apparati di commutazione di rete ridondati.
    \item \textbf{Visibilità per l'utente finale}: a differenza degli Host Aggregates, le Availability Zone sono esplicitamente **esposte agli utenti**. In fase di avvio di una VM, l'utente può esplicitare la zona desiderata:
    \begin{lstlisting}[language=bash]
openstack server create --flavor m1.small --image cirros \
    --availability-zone AZ_1 web-server-01
    \end{lstlisting}
    Ciò consente la progettazione di architetture applicative ad alta affidabilità (\textbf{High Availability}): distribuendo repliche di un servizio su AZ differenti, il disservizio catastrofico di un intero rack o ramo elettrico non compromette la continuità operativa dell'applicazione.
\end{itemize}

\subsection{Aggregati Di Host (Host Aggregates)}

Gli \textbf{Host Aggregates} rappresentano un meccanismo di categorizzazione logica interna dei nodi di calcolo, riservato all'amministrazione dell'infrastruttura.
\begin{itemize}
    \item \textbf{Invisibilità Agli Utenti}: gli aggregati di host non sono direttamente visibili o selezionabili dai tenant nei propri comandi di avvio istanza. Solo l'amministratore può crearli, modificarli ed associarvi nodi fisici.
    \item \textbf{Uso Dei Metadati E Compute Scheduler}: ciascun aggregato viene arricchito con coppie chiave-valore di metadati (\emph{metadata}) descrittive delle caratteristiche hardware peculiari dei nodi associati (ad esempio: dischi a stato solido \texttt{SSD=true}, processori dedicati ad elevata frequenza di clock \texttt{High\_CPU=true}, o schede di rete ad alta velocità \texttt{NIC=10GB}).
    \item \textbf{Binding Con I Flavor}: l'amministratore può associare le medesime chiavi di metadati (\emph{extra specs}) a specifici \emph{flavor}. Quando il modulo \texttt{nova-scheduler} deve posizionare un'istanza, legge i metadati richiesti dal flavor e restringe la scelta dell'hypervisor esclusivamente ai server fisici che appartengono all'aggregato provvisto di tale tag hardware.
\end{itemize}

\begin{figure}[H]
\centering
\begin{tabularx}{\textwidth}{p{3.5cm}p{3.5cm}X}
\toprule
\textbf{Livello Gerarchico} & \textbf{Visibilità Utente} & \textbf{Obiettivo Architetturale} \\
\midrule
\textbf{Region} & Visibile (tramite selezione endpoint) & Massima separazione geografica e indipendenza dei servizi API; condivisione dell'infrastruttura di autenticazione Keystone. \\
\addlinespace
\textbf{Availability Zone (AZ)} & Visibile (esposta in fase di provisioning) & Tolleranza ai guasti infrastrutturali fisici (energia, condizionamento, connettività di rack) per alta affidabilità applicativa. \\
\addlinespace
\textbf{Host Aggregate} & Invisibile (riservato all'amministratore) & Segmentazione interna e pilotaggio dello scheduler di calcolo su nodi con specifiche caratteristiche hardware tramite metadati. \\
\bottomrule
\end{tabularx}
\caption{Confronto sinottico tra i livelli di partizionamento di OpenStack.}
\end{figure}

\newpage

% ==============================================================================
% SEZIONE 7: IL SERVIZIO DI IDENTITÀ KEYSTONE
% ==============================================================================

\section{Il Servizio Di Identità: Keystone}

In un sistema cloud distribuito composto da decine di servizi cooperanti, la gestione dell'autenticazione, dell'autorizzazione e la scoperta dinamica degli endpoint costituiscono il prerequisito fondamentale di sicurezza. Tale responsabilità è interamente demandata a \textbf{Keystone}, il servizio di identità di riferimento per l'intero ecosistema OpenStack.

% [IMMAGINE DEL FILE DI SLIDE CON NOME "CloudComputing_3.pdf"] Architettura funzionale di Keystone: interfaccia API e backend per Identity, Catalog, Policy, Token e Resource — Slide 30 / Pagina 30 del PDF

\subsection{Casi D'Uso Fondamentali E Competenze Chiave}

Il ruolo operativo di Keystone si articola su quattro pilastri concettuali:
\begin{enumerate}
    \item \textbf{Autenticazione E Servizio Di Identità (Identity Service)}: valida le credenziali fornite dall'utente (coppie username/password, certificati crittografici X.509, chiavi API) interfacciandosi con molteplici backend di memorizzazione. Oltre al database relazionale interno SQL, Keystone supporta l'integrazione con directory aziendali esterne mediante il protocollo \textbf{LDAP} o Active Directory, oltre a protocolli di federazione d'identità (\textbf{SAML2}, \textbf{OpenID Connect}).
    \item \textbf{Emissione E Convalida Dei Token (Token Service)}: all'atto di un'autenticazione andata a buon fine, il servizio genera un token crittografico effimero che rappresenta l'identità confermata del richiedente e l'insieme dei ruoli assegnati. Nelle successive interazioni verso gli altri moduli OpenStack (Nova, Neutron, Glance), il client non trasmette più le credenziali originarie, ma esibisce il token all'interno dell'header HTTP delle richieste.
    \item \textbf{Registro Degli Endpoint (Service Catalog)}: funge da catalogo centrale in cui ogni servizio OpenStack registra la propria presenza e i propri URL di ascolto (\emph{endpoints}). Gli endpoint sono categorizzati in tre tipologie di accesso:
    \begin{itemize}[noitemsep]
        \item \emph{Public}: esposti su Internet o alla rete aziendale per l'accesso da parte degli utenti e delle CLI.
        \item \emph{Internal}: utilizzati per la comunicazione ad alta velocità sulla rete di gestione riservata tra servizi interni.
        \item \emph{Admin}: dedicati a operazioni di configurazione privilegiata e manutenzione infrastrutturale.
    \end{itemize}
    \item \textbf{Autorizzazione Basata Su Ruoli E Politiche (Policy Service)}: implementa un motore decisionale basato su \textbf{RBAC} (\emph{Role-Based Access Control}). Ciascun servizio OpenStack interroga il motore di policy (descritto formalmente tramite file JSON o YAML, quali \texttt{policy.yaml}) per determinare se il set di ruoli incapsulato nel token dell'utente consenta l'esecuzione della specifica azione API richiesta (es. riavviare un server, creare una subnet, eliminare un volume).
\end{enumerate}

\subsection{Architettura Interna Di Keystone}

Per soddisfare tali requisiti in modo scalabile, Keystone è organizzato internamente come un insieme di micro-servizi modulari esposti tramite un'API gateway unificata:

% [IMMAGINE DEL FILE DI SLIDE CON NOME "CloudComputing_3.pdf"] Elenco dei sotto-servizi architetturali di Keystone: Identity, Resource, Assignment, Token, Catalog, Policy — Slide 31 / Pagina 31 del PDF

\begin{itemize}
    \item \textbf{Identity Service}: modulo preposto alla gestione degli account utente e dei gruppi di appartenenza, interfacciato con i backend SQL o LDAP.
    \item \textbf{Resource Service}: gestisce la gerarchia dei contenitori organizzativi all'interno del cloud, ossia i domini (\emph{Domains}) e i progetti o tenant (\emph{Projects/Tenants}).
    \item \textbf{Assignment Service}: mantiene e governa l'associazione formale che lega un'identità (utente o gruppo) a un ruolo specifico all'interno di un determinato dominio o progetto.
    \item \textbf{Token Service}: governa il ciclo di vita, il formato (es. token crittografici non persistiti di tipo \emph{Fernet}), la verifica e la revoca anticipata dei token emessi.
    \item \textbf{Catalog Service}: archivio dei servizi registrati e della matrice di endpoint associati, fondamentale per consentire ai client la risoluzione dinamica delle rotte API.
    \item \textbf{Policy Service}: motore formale di convalida delle regole RBAC, responsabile dell'applicazione uniforme dei vincoli di sicurezza su tutto l'ambiente cloud.
\end{itemize}

Grazie a questo impianto a servizi disaccoppiati, qualsiasi interazione nel datacenter OpenStack garantisce tracciabilità, isolamento multi-tenant e robustezza contro accessi non autorizzati.

% CC -----------------------------------------------------------------------------------------------------------

% ==============================================================================
% CAPITOLO: ARCHITETTURA DI OPENSTACK KEYSTONE E MECCANISMI DI AUTENTICAZIONE
% ==============================================================================

\section{Architettura Di Keystone: Il Resource Service E La Struttura Gerarchica}

All'interno dell'ecosistema OpenStack, il componente preposto alla gestione centralizzata dell'identità, dell'autenticazione, dell'autorizzazione e del catalogo dei servizi è \textbf{Keystone}. Prima di esaminare i flussi crittografici di autenticazione, è essenziale comprendere come le risorse organizzative e gli accessi siano strutturati concettualmente. La gestione di tali entità fa capo al \textbf{Resource Service}, il quale definisce e mantiene le relazioni tra domini e progetti.

\subsection{Domini E Progetti: I Confini Dell'Isolamento Multi-Tenant}

In un'infrastruttura di tipo \emph{Infrastructure as a Service} (IaaS), la segregazione delle risorse e delle responsabilità amministrative è un vincolo architetturale fondamentale. Keystone affronta tale esigenza introducendo una gerarchia esplicita:

\begin{itemize}[noitemsep]
    \item \textbf{Domini (Domains)}: rappresentano il contenitore di più alto livello per progetti, utenti e gruppi. Introdotti in modo strutturale con l'API Identity v3, i domini fungono da confini amministrativi completi (\emph{administrative boundaries}) e possono corrispondere a singole aziende, dipartimenti o tenant logici dell'operatore cloud. Ogni dominio definisce un proprio \emph{namespace} per i nomi delle entità in esso contenute. L'amministratore di sistema può delegare pieni privilegi di gestione a un \textbf{amministratore di dominio} (\emph{Domain Administrator}), il quale è autorizzato a creare utenti, gruppi e progetti all'interno del dominio assegnatogli, nonché ad attribuire i ruoli operativi. All'atto dell'installazione, Keystone predispone un dominio predefinito denominato formalmente \texttt{'Default'}.
    \item \textbf{Progetti (Projects o Tenants)}: costituiscono l'unità fondamentale di possesso (\emph{base unit of ownership}) all'interno del cloud. Qualsiasi risorsa computazionale, di memorizzazione o di rete creata in OpenStack (come un'istanza computazionale Nova, un volume a blocchi Cinder o una rete Neutron) deve obbligatoriamente appartenere a uno specifico progetto. Ogni progetto deve essere vincolato a un unico dominio genitore. Se durante la creazione del progetto il dominio di riferimento viene omesso, l'entità viene associata d'ufficio al dominio predefinito \texttt{'Default'}. A seconda del modello operativo del fornitore di servizi, il concetto di progetto mappa direttamente su un account cliente, una divisione aziendale o un gruppo di lavoro.
\end{itemize}

\subsection{Politiche Di Unicità Degli Attributi Nell'API Identity v3}

La transizione dall'API v2 alla v3 di Keystone ha trasformato il modello di visibilità e univocità degli identificatori. Nel modello v3, l'introduzione dei domini elimina il vincolo di unicità globale per la maggior parte delle entità, sostituendolo con una contestualizzazione basata sul \emph{namespace} del dominio di appartenenza.

\begin{table}[H]
\centering
\small
\begin{tabularx}{\textwidth}{l p{5cm} X}
\toprule
\textbf{Entità} & \textbf{Ambito di Unicità} & \textbf{Descrizione Operativa} \\
\midrule
\textbf{Domain Name} & Globale (\emph{globally unique}) & Identifica univocamente il dominio nell'intera istanza Keystone; non possono esistere due domini con il medesimo nome. \\
\textbf{Project Name} & Locale al Dominio & Univoco esclusivamente all'interno del dominio genitore; domini distinti possono contenere progetti omonimi. \\
\textbf{User Name} & Locale al Dominio & Identifica l'utente nel perimetro del suo dominio; utenti in domini diversi possono avere la stessa stringa identificativa. \\
\textbf{Group Name} & Locale al Dominio & Univoco all'interno del dominio proprietario per l'aggregazione di utenti. \\
\textbf{Role Name} & Locale al Dominio & Definisce l'insieme di privilegi ed è vincolato al dominio proprietario di definizione. \\
\bottomrule
\end{tabularx}
\caption{Regole di unicità degli identificatori nominali nell'API Identity v3 di Keystone.}
\end{table}

Questo disaccoppiamento garantisce un isolamento architetturale ottimale: due organizzazioni indipendenti ospitate sulla medesima cloud possono nominare i propri progetti \texttt{"production"} o i propri utenti amministrativi \texttt{"admin"} senza collisioni logiche a livello di database.

\newpage

\section{Identity Service, Componenti Architetturali E Service Catalog}

Il nucleo operativo di Keystone è costituito dall'\textbf{Identity Service}. Esso rappresenta il primo punto di contatto per qualsiasi entità (utente umano, script di automazione o microservizio di sistema) che debba operare sull'infrastruttura. L'Identity Service assolve a due compiti primari: la convalida delle credenziali di accesso e la gestione dell'anagrafica di utenti e gruppi.

\subsection{Utenti E Gruppi Come Consumatori Di API}

Nel contesto di Keystone:
\begin{itemize}[noitemsep]
    \item Un \textbf{Utente (User)} rappresenta il consumatore finale delle API, inteso come un individuo o un account di servizio programmabile. Ogni utente è posseduto da un singolo dominio specifico; pertanto, il suo nome è vincolato all'unicità locale e non globale.
    \item Un \textbf{Gruppo (Group)} è un contenitore logico che aggrega più utenti appartenenti al medesimo dominio. L'utilità principale del gruppo risiede nella semplificazione dell'amministrazione dei privilegi: concedere o revocare un ruolo operativo a un gruppo (assegnato su un intero dominio o su un singolo progetto) riflette istantaneamente le autorizzazioni su tutti i membri presenti o futuri di quel gruppo.
\end{itemize}

\subsection{Componenti Interni Dell'Identity Service}

L'architettura interna dell'Identity Service è modularizzata per garantire scalabilità e interoperabilità con infrastrutture enterprise preesistenti. Si suddivide in tre macro-componenti:

\begin{enumerate}
    \item \textbf{Server}: è il demone applicativo centrale che espone un'interfaccia di tipo RESTful su protocollo HTTP/HTTPS. Riceve le richieste di autenticazione e autorizzazione, coordina la validazione dei dati e restituisce i metadati di sicurezza.
    \item \textbf{Moduli (WSGI Middleware)}: pipeline di intercettazione collocate a monte o a valle del server. Questi moduli esaminano le richieste HTTP in transito, estraggono le credenziali o i token contenuti negli header (ad esempio il campo \texttt{X-Auth-Token}) e le inoltrano al server centrale per la verifica crittografica e dei permessi.
    \item \textbf{Driver}: componenti di astrazione dell'accesso ai dati integrati nel server centrale. Permettono all'Identity Service di interrogare sorgenti informative eterogenee, sia native interne (come un database relazionale SQL basato su MariaDB o PostgreSQL) sia esterne preesistenti nell'organizzazione.
\end{enumerate}

Un esempio primario di integrazione tramite driver è rappresentato dal protocollo \textbf{LDAP} (\emph{Lightweight Directory Access Protocol}). LDAP è uno standard applicativo aperto e vendor-neutral concepito per la gestione e la consultazione distribuita di servizi di directory su reti IP. Poiché un servizio di directory organizza le informazioni in modo gerarchico e ad alte prestazioni di lettura (analogamente a una rubrica aziendale strutturata), l'adozione di un driver LDAP consente a Keystone di validare gli account utente direttamente a partire da directory corporate come Microsoft Active Directory o OpenLDAP, evitando duplicazioni anagrafiche.

\subsection{Il Service Catalog E La Segregazione Degli Endpoint Di Rete}

Una volta che l'utente o il servizio ha autenticato la propria identità, deve essere in grado di localizzare e invocare i vari componenti di OpenStack (Nova per il calcolo, Glance per le immagini, Cinder per i volumi, Neutron per la rete). Questa funzionalità di \emph{service discovery} è implementata tramite il \textbf{Service Catalog} (catalogo dei servizi).

Il catalogo è una tabella gestita da Keystone che associa a ciascun servizio registrato una serie di \textbf{endpoint} (punti di terminazione API identificati da URI di rete). Per ciascun servizio, l'architettura supporta endpoint distribuiti su tre tipologie distinte:

\begin{itemize}[noitemsep]
    \item \textbf{Endpoint Pubblici (\emph{Public})}: esposti verso reti esterne o non fidate (Internet). Consentono ai clienti e tenant di inviare comandi di gestione per le proprie risorse cloud.
    \item \textbf{Endpoint Interni (\emph{Internal})}: accessibili esclusivamente all'interno della rete locale dell'infrastruttura. Sono riservati alla comunicazione est-ovest inter-servizio (ad esempio quando Nova deve richiedere un volume a Cinder o una rete a Neutron).
    \item \textbf{Endpoint Amministrativi (\emph{Admin})}: dedicati alle operazioni ad alto privilegio, quali la creazione e la rimozione di utenti, tenant e quote. Per impostazione predefinita, le API invocate su questo endpoint consentono modifiche strutturali precluse alle interfacce pubbliche e interne.
\end{itemize}

In un'infrastruttura di produzione conforme ai moderni requisiti di sicurezza, queste tre classi di endpoint non condividono la medesima interfaccia fisica o logica di rete, ma risiedono su piani di indirizzamento e VLAN separati. Ciò garantisce che i vettori di attacco provenienti dall'esterno non possano raggiungere le interfacce amministrative o il traffico di controllo dei nodi computazionali.

\newpage

\section{Assignment Service E Gestione Dei Privilegi}

Il completamento del modello di sicurezza in Keystone richiede il legame operativo tra l'identità (chi è l'utente) e le risorse accessibili (su cosa può operare). Tale associazione è governata dall'\textbf{Assignment Service}.

\subsection{Definizione E Assegnazione Dei Ruoli}

Un \textbf{Ruolo (Role)} incapsula una determinata personalità autorizzativa, a cui è associato un insieme strutturato di permessi e diritti di esecuzione all'interno dei file di policy (\texttt{policy.yaml} o \texttt{policy.json}) dei singoli microservizi OpenStack. Quando Keystone autentica con successo un utente, rilascia un token contenente la lista dei ruoli a lui assegnati. All'atto della ricezione del token, il servizio target interpreta tale vettore di ruoli e determina se consentire o vietare l'operazione richiesta.

I ruoli presentano le seguenti caratteristiche operative:
\begin{itemize}[noitemsep]
    \item Possono essere concessi con ambito a livello di **intero dominio** (\emph{domain level}) oppure a livello di **singolo progetto** (\emph{project level}).
    \item Possono essere associati direttamente a un **singolo utente** oppure a un **gruppo di utenti**.
    \item Il nome del ruolo (\emph{Role Name}) è univoco nell'ambito del dominio proprietario che lo ha istanziato.
\end{itemize}

Per associare un utente a un progetto conferendogli un ruolo operativo, l'operatore impiega la riga di comando OpenStack. Il costrutto richiede di specificare la terna utente-progetto-ruolo:

\begin{lstlisting}[language=bash]
openstack role add --user USER_NAME --project PROJECT_NAME ROLE_NAME
\end{lstlisting}

L'esecuzione di tale comando istruisce l'Assignment Service a persistere l'associazione, facendo sì che ogni successivo token generato per quell'utente con ambito (\emph{scoped}) sul progetto indicato contenga il set di privilegi associato a \texttt{ROLE\_NAME}.

% ==============================================================================
% SEZIONE: EVOLUZIONE DEI TOKEN
% ==============================================================================

\section{L'Evoluzione Dei Meccanismi Di Autenticazione: Dalle Credenziali Ai Token}

L'interazione con un'infrastruttura di controllo cloud distribuita avviene tramite scambi asincroni di messaggi basati sul paradigma REST sopra il protocollo HTTP. La gestione delle credenziali di sicurezza in tale ambiente impone stringenti trade-off di natura sistemistica.

\subsection{Inadeguatezza Del Modello Username/Password Diretto}

L'approccio più banale consisterebbe nel trasmettere le credenziali primitive (nome utente e password) a ogni singola invocazione di endpoint API. Questo pattern presenta criticità ingegneristiche inaccettabili:
\begin{itemize}[noitemsep]
    \item \textbf{Esposizione continua sul canale di rete}: ogni richiesta HTTP moltiplica la superficie di intercettazione; la compromissione di un singolo pacchetto espone permanentemente l'identità dell'utente.
    \item \textbf{Persistenza insicura lato client}: per non costringere l'operatore o lo script a digitare manualmente la password per centinaia di micro-operazioni (creazione di VM, allocazione di porte di rete, montaggio dischi), le credenziali dovrebbero essere salvate in chiaro o memorizzate localmente in memoria, rendendole vulnerabili ad accessi non autorizzati sull'endpoint dell'amministratore.
\end{itemize}

\subsection{Il Principio Del Token Temporaneo}

Per risolvere queste problematiche, Keystone implementa l'autenticazione basata su \textbf{Token}. Un token è un artefatto temporaneo rilasciato a fronte di una singola validazione iniziale delle credenziali. Esso possiede una finestra temporale di validità rigorosamente limitata (\emph{Time-to-Live}, TTL).

I vantaggi architetturali del token sono duplici:
\begin{enumerate}
    \item L'esposizione delle credenziali primarie è ridotta al solo scambio iniziale con Keystone.
    \item Qualora il token venga intercettato durante il transito di rete o sottratto da una cache locale compromessa, la finestra temporale a disposizione dell'attaccante per utilizzarlo prima della scadenza è estremamente ristretta.
\end{enumerate}

La generazione e la verifica di questi artefatti sono demandate al \textbf{Token Service}. Il formato crittografico del token generato da Keystone è definito globalmente all'interno del file di configurazione principale:

\begin{lstlisting}[language=ini]
# File di configurazione /etc/keystone/keystone.conf
[token]
provider = fernet
\end{lstlisting}

Nelle versioni storiche e moderne di OpenStack si sono succedute diverse generazioni di token: \textbf{UUID}, \textbf{PKI/PKIZ}, \textbf{Fernet} e infine \textbf{JWS} (introdotto nella release Stein). Attualmente, l'unico provider abilitato come predefinito e pienamente raccomandato in produzione è \textbf{Fernet}.

\newpage

\section{Token UUID: Architettura Stateful E Limiti Di Scalabilità}

Il primo formato implementato in OpenStack è basato su identificatori univoci universali (\textbf{UUID}, \emph{Universally Unique Identifier}).

\subsection{Formato E Flusso Operativo}

I token UUID sono stringhe casuali generate conformemente allo standard **UUID versione 4**, aventi una dimensione fissa di 128 bit (16 byte), comunemente rappresentate in formato esadecimale a blocchi separati da trattini (36 caratteri, es. \texttt{12ea1917-4641-49dc-ad68-d0dc1eb30842}).

Dal punto di vista architetturale, il token UUID è puramente referenziale: non contiene al proprio interno alcun metadato informativo o autorizzativo (\emph{opaque token}). Esso funge semplicemente da chiave primaria (\emph{handle}) puntatore a una riga registrata nel database relazionale di Keystone.

Il ciclo di vita si articola in due fasi operative:

\begin{enumerate}
    \item \textbf{Fase 1: Acquisizione del token}
    \begin{itemize}[noitemsep]
        \item L'utente invia le proprie credenziali (username e password) all'endpoint di Keystone.
        \item Keystone valida le credenziali interrogando il backend Identity (SQL o LDAP).
        \item A fronte dell'esito positivo, Keystone genera un UUIDv4 casuale e memorizza nel proprio database persistente (tabella \texttt{token}) il token stesso, associandovi l'identità dell'utente, i ruoli, il progetto target e il timestamp di scadenza.
        \item Il token viene restituito al client nel corpo o nell'header della risposta HTTP.
    \end{itemize}
    \item \textbf{Fase 2: Utilizzo e validazione}
    \begin{itemize}[noitemsep]
        \item L'utente formula una richiesta verso un servizio applicativo (ad esempio invia a Nova una richiesta di accensione VM), includendo il token nell'header HTTP: \texttt{X-Auth-Token: \$token}.
        \item Il middleware WSGI del servizio ricevente intercetta la richiesta. Non potendo estrarre informazioni dal token UUID, invia una richiesta REST a Keystone (operazione di \emph{token validation}).
        \item Keystone esegue una query sul proprio database SQL ricercando il record indicizzato dal valore UUID.
        \item Se il record esiste, non è scaduto e non compare tra quelli revocati, Keystone restituisce al servizio ricevente l'esito positivo unitamente a tutti i metadati dell'utente e dei suoi ruoli. In caso contrario, risponde con un codice di errore HTTP 404/401.
        \item Il servizio applicativo applica le sue regole di policy e processa la richiesta dell'utente.
    \end{itemize}
\end{enumerate}

% [IMMAGINE DEL FILE DI SLIDE CON NOME "CloudComputing_3.pdf"] Diagramma del flusso di autenticazione e validazione dei token UUID tra User, Keystone, SQL DB e OpenStack API — Slide 39 / Pagina 10 del PDF

\subsection{Criticità Architetturali Del Modello UUID}

Sebbene semplice da implementare, il modello basato su token UUID evidenzia colli di bottiglia prestazionali che ne hanno precluso l'utilizzo su cloud su larga scala:

\begin{itemize}
    \item \textbf{Traffico di rete e amplificazione dei messaggi}: poiché i servizi consumatori non possono ispezionare il contenuto del token, \emph{ogni singola operazione} all'interno del cloud richiede una chiamata sincrona di validazione verso Keystone. Una chiamata complessa (come l'avvio di un'istanza che coinvolge Nova, Glance, Neutron e Cinder) scatena una cascata di validazioni di rete, sovraccaricando il control plane di Keystone.
    \item \textbf{Degrado delle prestazioni del database relazionale (\emph{Database Bloat})}: ogni token creato deve essere salvato su disco. In installazioni con migliaia di utenti, nodi e microservizi che effettuano polling continuo, la tabella dei token cresce a dismisura. Le query di ricerca sul database per validare la stringa UUID richiedono scansioni di indici via via più lente, degradando i tempi di risposta dell'intera infrastruttura.
    \item \textbf{Manutenzione complessa e pulizia dei record (\emph{Pruning})}: la persistenza di tutti i token impone l'esecuzione frequente di job di manutenzione batch (tramite il comando \texttt{keystone-manage token\_flush}) per eliminare i record scaduti o revocati, operazione che introduce contese di lock sul database.
\end{itemize}

\newpage

\section{Token PKI E PKIZ: Firma Crittografica E Il Problema Del Payload}

Per ovviare all'obbligo di contattare Keystone a ogni singola richiesta API, la comunità di OpenStack ha sviluppato i formati basati su crittografia a chiave pubblica: i token \textbf{PKI} (\emph{Public Key Infrastructure}) e la loro variante compressa \textbf{PKIZ}.

\subsection{Architettura PKI E Convalida Distribuita}

Nel modello PKI, Keystone cessa di essere un mero database di lookup e assume il ruolo di una vera e propria **Certification Authority (CA)**. All'atto dell'inizializzazione del servizio, vengono generati quattro elementi crittografici conformi allo standard X.509:
\begin{itemize}[noitemsep]
    \item La chiave privata della Certification Authority (\emph{CA Private Key}).
    \item Il certificato pubblico della CA (\emph{CA Certificate}).
    \item La chiave privata di firma del servizio Identity (\emph{Signing Private Key}).
    \item Il certificato pubblico di firma (\emph{Signing Certificate}).
\end{itemize}

Questi file, gestiti tramite l'utility \texttt{keystone-manage}, risiedono nei percorsi specificati in \texttt{keystone.conf} e sono protetti da rigidi permessi di lettura riservati all'utente di sistema di Keystone.

\usetikzlibrary{positioning, shapes.geometric, arrows.meta}

\begin{figure}[H]
    \centering
    \begin{tikzpicture}[
        >=Stealth,
        node distance=2.2cm,
        auto,
        block/.style={
            rectangle,
            draw=theoryblue,
            fill=backcolour,
            text width=3.8cm,
            text centered,
            rounded corners,
            minimum height=1.2cm,
            font=\small
        },
        cloud/.style={
            shape=ellipse,
            draw=theoryblue,
            fill=white,
            text centered,
            minimum height=1.2cm,
            font=\small
        },
        arrow/.style={
            thick,
            ->,
            >=Stealth
        }
    ]

        \node[cloud] (user) {Client / Utente};

        \node[block, right=3.5cm of user] (keystone)
            {\textbf{Keystone}\\Firma con Signing Private Key};

        \node[block, above=1.5cm of user] (api)
            {\textbf{OpenStack Service (Nova)}\\
            Verifica con Signing Certificate locale};

        \draw[arrow]
            (user) --
            node[below, font=\footnotesize] {1. Login (User/Pass)}
            (keystone);

        \draw[arrow]
            (keystone) |-
            node[above, near start, font=\footnotesize]
            {2. Token PKI (documento CMS firmato)}
            (user);

        \draw[arrow]
            (user) --
            node[left, font=\footnotesize] {3. Richiesta API con Token}
            (api);

        \draw[arrow, dashed]
            (api) --
            node[right, font=\footnotesize]
            {Cache certificati e CRL (una tantum)}
            (keystone);

    \end{tikzpicture}

    \caption{Flusso architetturale della validazione distribuita nei token PKI.}
\end{figure}

% [IMMAGINE DEL FILE DI SLIDE CON NOME "CloudComputing_3.pdf"] Diagramma di interazione per l'autenticazione PKI/PKIZ con verifica locale tramite chiave pubblica presso l'API del servizio — Slide 44 / Pagina 15 del PDF

Il token PKI non è un identificatore astratto, bensì un **documento crittografico strutturato** conforme alla specifica **Cryptographic Message Syntax (CMS)**, standard IETF per messaggi protetti mediante firma o cifratura. Al suo interno è serializzata la risposta di validazione completa di Keystone, comprendente:
\begin{itemize}[noitemsep]
    \item L'identità dell'utente, l'ID di dominio e di progetto.
    \item L'intero set dei ruoli assegnati.
    \item La data e l'ora esatta di scadenza.
    \item L'intero **Service Catalog** (tutti i servizi del cloud, le regioni geografiche e gli URL degli endpoint pubblici, interni e di amministrazione).
\end{itemize}

Questo aggregato viene firmato digitalmente con la Signing Key di Keystone e codificato in formato testuale.

\subsection{Meccanismo Di Verifica Locale}

Il vantaggio cardine dei token PKI è l'eliminazione della richiesta sincrona di convalida verso Keystone:
\begin{enumerate}
    \item Ciascun servizio di OpenStack scarica localmente da Keystone, una tantum all'avvio, il certificato di firma, il certificato della CA e la lista di revoca dei certificati (\textbf{CRL}, \emph{Certificate Revocation List}).
    \item Quando riceve una richiesta dall'utente, il servizio estrae il token PKI dall'header HTTP, converte il documento CMS e verifica crittograficamente la firma digitale impiegando il certificato di Keystone presente nella propria cache locale.
    \item Se la firma è integra, il token non è scaduto e non compare nella CRL, il servizio considera la richiesta autorizzata senza inviare alcun byte di traffico di rete verso Keystone.
\end{enumerate}

\subsection{Il Limite Delle Dimensioni E L'Introduzione Di PKIZ}

Nonostante l'eccellente disaccoppiamento di rete, i token PKI hanno introdotto un problema collaterale critico: la **dimensione del payload**.

Poiché il token incorpora l'intero catalogo dei servizi:
\begin{itemize}[noitemsep]
    \item Un token PKI di base, configurato con un unico endpoint, ha un'impronta di circa \SI{1700}{\byte}.
    \item All'aumentare dei servizi installati e delle regioni configurate, la dimensione cresce proporzionalmente, superando frequentemente la soglia di \SI{8}{\kilo\byte}.
\end{itemize}

Questa crescita dimensionale comporta anomalie operative severe:
\begin{itemize}
    \item La maggior parte dei server web (Apache HTTPD, Nginx) e dei bilanciatori di carico ha un limite predefinito per la dimensione massima degli header HTTP (solitamente configurato a \SI{8}{\kilo\byte}). Quando il token PKI supera tale soglia, le chiamate API vengono respinte con errori di tipo \texttt{HTTP 400 Bad Request} o \texttt{HTTP 494 Request Header Too Large}.
    \item L'elevato volume di dati da trasmettere a ogni singola chiamata HTTP incrementa la latenza e consuma banda sul piano di controllo.
\end{itemize}

Per mitigare tale criticità, è stata introdotta la variante **PKIZ**. I token PKIZ mantengono inalterata la logica di firma e validazione del PKI tradizionale, ma applicano una compressione con l'algoritmo **zlib** al payload CMS prima della codifica finale. Tuttavia, anche con la compressione, all'espandersi dell'infrastruttura il payload tendeva rapidamente a superare i limiti di sicurezza degli header HTTP, rendendo la famiglia PKI/PKIZ insostenibile nei cluster di grandi dimensioni.

\newpage

\section{Token Fernet: Risoluzione Del Trade-Off Tra Persistenza E Dimensione}

L'analisi dei fallimenti dei modelli precedenti ha portato alla progettazione dei token **Fernet**, che rappresentano il connubio ottimale tra leggerezza, assenza di persistenza su database e sicurezza crittografica.

\subsection{Natura Effimera E Serializzazione MessagePack}

A differenza dei token UUID e PKI, i token Fernet sono intrinsecamente **effimeri** (\emph{ephemeral} o \emph{non-persistent}):
\begin{itemize}[noitemsep]
    \item Non vengono in alcun modo memorizzati o indicizzati nel database relazionale di Keystone.
    \item Di conseguenza, il database non subisce alcun fenomeno di espansione incontrollata (\emph{bloat}), eliminando radicalmente l'onere operativo di eseguire periodicamente script di pulizia (\emph{pruning}) dei token scaduti.
    \item Non richiedono la replica sincrona di tabelle del database su cluster eterogenei; la validità di un token è verificata puramente mediante l'applicazione di algoritmi crittografici su chiavi simmetriche condivise.
\end{itemize}

Al fine di minimizzare l'ingombro del pacchetto, Fernet scarta i formati testuali verbosi (XML o JSON) per la serializzazione interna, adottando **MessagePack**. MessagePack è un formato di serializzazione binaria altamente ottimizzato che garantisce l'interscambio di strutture dati eterogenee con prestazioni e compattezza superiori rispetto a JSON. Ad esempio, piccoli interi vengono codificati in un singolo byte e le stringhe brevi richiedono un solo byte di overhead descrittivo.

% [IMMAGINE DEL FILE DI SLIDE CON NOME "CloudComputing_3.pdf"] Confronto binario tra formato JSON (27 byte) e MessagePack (18 byte) — Slide 47 / Pagina 18 del PDF

Grazie all'impiego congiunto di MessagePack e della crittografia simmetrica, la dimensione totale di un token Fernet si mantiene rigorosamente **al di sotto della soglia dei 250 byte**, risolvendo alla radice le problematiche di saturazione degli header HTTP.

\subsection{Struttura Crittografica Del Token Fernet}

Un token Fernet applica una crittografia a chiave simmetrica basata sul costrutto crittografico standard Fernet. Ciascuna chiave (\textbf{Fernet Key}) è in realtà una composizione di due sotto-chiavi da 128 bit ciascuna:
\begin{itemize}[noitemsep]
    \item Una chiave da 128 bit per la cifratura/decifratura mediante l'algoritmo **AES-128** in modalità CBC (\emph{Cipher Block Chaining}).
    \item Una chiave da 128 bit per la firma e verifica di integrità tramite **HMAC-SHA256**.
\end{itemize}

Il token serializzato è composto dai seguenti campi sequenziali, successivamente codificati mediante l'algoritmo **Base64 URL-safe**:

\begin{table}[H]
\centering
\small
\begin{tabularx}{\textwidth}{l l l X}
\toprule
\textbf{Campo} & \textbf{Valore / Tipo} & \textbf{Dimensione} & \textbf{Descrizione Operativa} \\
\midrule
\texttt{Version} & \texttt{0x80} & 8 bit (1 byte) & Specifica la versione del formato token Fernet. \\
\texttt{Timestamp} & Intero 64 bit & 64 bit (8 byte) & Orario di generazione in formato UTC (tempo assoluto). \\
\texttt{IV} & Vettore pseudocasuale & 128 bit (16 byte) & Vettore di Inizializzazione per la cifratura AES. \\
\texttt{Ciphertext} & Dati cifrati con AES & Variabile & Payload compresso e imbottito (\emph{padding}) contenente i dati di autorizzazione. \\
\texttt{HMAC} & Codice di autenticazione & 256 bit (32 byte) & Firma SHA256 calcolata sull'intero blocco precedente. \\
\bottomrule
\end{tabularx}
\caption{Campi strutturali di un token crittografico Fernet.}
\end{table}

% [IMMAGINE DEL FILE DI SLIDE CON NOME "CloudComputing_3.pdf"] Diagramma di generazione del token Fernet con evidenza di Timestamp, IV, Payload cifrato con AES e HMAC — Slide 52 / Pagina 23 del PDF

Il \textbf{Payload cifrato} contiene esclusivamente le informazioni strettamente necessarie a definire il contesto di sicurezza, serializzate in MessagePack:
\begin{itemize}[noitemsep]
    \item \textbf{Version}: definisce l'ambito di autorizzazione del token (\texttt{0} = unscoped, \texttt{1} = domain-scoped, \texttt{2} = project-scoped, \texttt{3} = trust-scoped, \texttt{4} = federated unscoped, \texttt{5} = federated project-scoped, \texttt{6} = federated domain-scoped).
    \item \textbf{User ID}: identificatore univoco dell'utente autenticato.
    \item \textbf{Methods}: lista dei metodi usati per l'autenticazione (es. password, token, oauth1, external).
    \item \textbf{Project ID / Domain ID}: identificatore dell'entità su cui il token è contestualizzato.
    \item \textbf{Expiration Time}: orario esatto in cui il token cessa di avere validità.
    \item \textbf{Audit IDs}: lista contenente fino a due stringhe pseudocasuali univoche impiegate dai sistemi di tracciamento e audit log per correlare l'uso del token senza esporne il contenuto riservato.
\end{itemize}

Nel meccanismo dei \textbf{Trusts}, una delega esplicita consente a un utente (\emph{trustor}) di trasferire temporaneamente specifici privilegi a un altro utente (\emph{trustee}), abilitando l'impersonificazione controllata tramite token con scope \texttt{trust-scoped}.

\subsection{Gestione Del Key Repository E Ciclo Di Vita Delle Chiavi}

Poiché Keystone adotta crittografia simmetrica, la sicurezza dell'intero sistema si basa sulla corretta custodia e distribuzione delle chiavi. Le chiavi risiedono nel filesystem locale nel percorso definito dalla direttiva \texttt{key\_repository} (di norma \texttt{/etc/keystone/fernet-keys/}).

All'interno di tale cartella, le chiavi assumono nomi puramente numerici interi non negativi e si distinguono in tre tipologie funzionali:

\begin{enumerate}
    \item \textbf{Primary Key}: esiste una sola chiave primaria nel repository. È l'unica chiave autorizzata sia a \textbf{cifrare} (generare nuovi token) sia a \textbf{decifrare}. Corrisponde sempre al file avente il \textbf{valore numerico più alto} presente nella directory.
    \item \textbf{Secondary Keys}: chiavi storiche che in passato hanno ricoperto il ruolo di chiave primaria. Vengono mantenute esclusivamente per \textbf{decifrare} token emessi precedentemente e ancora in circolazione entro i termini di scadenza. Non sono mai impiegate per generare nuovi token.
    \item \textbf{Staged Key}: chiave speciale preparatoria contraddistinta tassativamente dal nome numerico \texttt{0}. Ha il potere di decifrare i token ma non viene mai usata per cifrare finché permane nello stato di \emph{staged}. È destinata a essere promossa a nuova chiave primaria nella successiva rotazione.
\end{enumerate}

\subsection{Meccanica Della Rotazione Delle Chiavi}

Il ciclo di vita naturale di ciascuna chiave Fernet segue rigorosamente il percorso:
\[
\text{Staged Key (0)} \xrightarrow{\text{rotazione}} \text{Primary Key (N+1)} \xrightarrow{\text{rotazioni successive}} \text{Secondary Key (demoted)}
\]

% [IMMAGINE DEL FILE DI SLIDE CON NOME "CloudComputing_3.pdf"] Diagramma dell'evoluzione degli indici delle chiavi Fernet durante due fasi consecutive di rotazione — Slide 51 / Pagina 22 del PDF

La procedura operativa di rotazione si esegue mediante il comando:
\begin{lstlisting}[language=bash]
keystone-manage fernet_rotate --keystone-user keystone --keystone-group keystone
\end{lstlisting}

L'algoritmo di rotazione esegue i seguenti passaggi in sequenza transazionale:
\begin{enumerate}[noitemsep]
    \item Identifica il valore intero massimo attuale $N$ (la Primary Key corrente).
    \item Promuove la Staged Key corrente (il file con nome \texttt{0}) rinominandola con indice $N+1$. Essa diviene la nuova \textbf{Primary Key} a cui compete la cifratura di tutti i nuovi token.
    \item La vecchia chiave primaria di indice $N$ degrada automaticamente a \textbf{Secondary Key}.
    \item Genera una nuova chiave crittografica casuale e la scrive sul filesystem assegnandole l'indice \texttt{0}, che assume il ruolo di nuova \textbf{Staged Key}.
    \item Se il numero complessivo di chiavi secondarie eccede il limite configurato per coprire il TTL massimo dei token, le chiavi secondarie con gli indici più vecchi vengono rimosse dal filesystem.
\end{enumerate}

Questo meccanismo risolve brillantemente la sincronizzazione nei cluster multi-nodo. L'operatore può distribuire la nuova chiave su tutti i nodi Keystone dell'infrastruttura (ad esempio tramite rsync o pipeline Ansible) quando essa è ancora nello stato di Staged Key (\texttt{0}). Poiché la chiave \texttt{0} valida i token ma non ne emette di nuovi, la sua distribuzione non richiede tempi di convergenza istantanei né interruzioni del servizio; non appena la rotazione scatta sul nodo coordinatore e l'indice viene incrementato, tutti i nodi sono già in grado di validare immediatamente i messaggi crittografati con tale materiale.

% [IMMAGINE DEL FILE DI SLIDE CON NOME "CloudComputing_3.pdf"] Diagramma di convalida dei token Fernet con l'ausilio di Memcached e repository di chiavi — Slide 54 / Pagina 26 del PDF

% [IMMAGINE DEL FILE DI SLIDE CON NOME "CloudComputing_3.pdf"] Diagramma di flusso decisionale per la validazione di un token Fernet presso Keystone — Slide 55 / Pagina 27 del PDF

\newpage


\end{document}


